%\documentclass[aps,prl,singlecolumn,superscriptaddress]{revtex4-2}


\documentclass[a4paper,11pt]{article}
\usepackage{jheppub}
\setcounter{secnumdepth}{3}
% Useful packages
\usepackage{tikz-cd}
\usepackage{graphicx}
\usepackage{amsmath}
\usepackage{amssymb}
\usepackage{array,booktabs}
\usepackage{hyperref}
%\usepackage{showlabels}
\usepackage{amsthm}
\newtheorem{theorem}{Theorem}
\newtheorem{lemma}[theorem]{Lemma}
\newtheorem{corollary}[theorem]{Corollary}
\newtheorem{proposition}[theorem]{Proposition}
\theoremstyle{definition}
\newtheorem{definition}{Definition}[section]
\newtheorem{prop}{Conjecture}
\newtheorem{cor}{Corollary}

\theoremstyle{remark}
\newtheorem{remark}[theorem]{Remark}
%-------------------packages-------------------------
\usepackage{enumerate}
\usepackage[utf8]{inputenc}
\usepackage{amssymb,amsbsy,mathtools,amsmath}
\usepackage{braket}
\usepackage{graphicx}
\usepackage{xcolor}
%\usepackage{multirow}
\usepackage{hyperref}
%\usepackage{caption}
%\usepackage{subcaption}
%\mathcal{A}ptionsetup{compatibility=false}
\usepackage{natbib}
%\usepackage{dcolumn}% Align table columns on decimal point
\usepackage{bm}% bold math
\usepackage{soul}
\usepackage{showlabels}
%\newtheorem{proposition}{Proposition}
\newcommand{\E}{\mathbb{E}}
\newcommand{\dd}{\,\mathrm{d}}
%\newcommand{\mathcal{A}}{\mathcal{A}}
\newcommand{\cH}{\mathcal{H}}
\newcommand{\KLL}{D_{\mathrm{KL}}}
\newcommand{\Mop}{\mathsf{M}}

%Includes "References" in the table of contents
\usepackage[nottoc]{tocbibind}

\usepackage{bbm}
\usepackage{amsthm}
%\newtheorem{theorem}{Theorem}
%\newtheorem{lemma}[theorem]{Lemma}
%\newtheorem{corollary}[theorem]{Corollary}
%\theoremstyle{definition}
%\newtheorem{definition}{Definition}[section]
\usepackage{amssymb}
\newcommand{\id}{\mathbb{I}}
\newcommand{\Tr}{\text{Tr}}
\newcommand{\Div}{\operatorname{div}}
\usepackage{amsmath}
\usepackage{tikz}
\usepackage{braket}
\usepackage{quantikz}
\usepackage{graphicx}
\usepackage{caption}
\newcommand{\mA}{\mathcal{A}}
\newcommand{\mB}{\mathcal{B}}
\newcommand{\nn}{\nonumber}
\newcommand{\mL}{\mathcal{L}}
\newcommand{\lb}{\left(}
\newcommand{\rb}{\right)}
\newcommand{\mH}{\mathcal{H}}
\usepackage{subcaption}
\newcommand{\tOO}{\tilde{\mathcal{O}}}
\newcommand{\mO}{\mathcal{O}}
\newcommand{\mK}{\mathcal{K}}
\newcommand{\ep}{\epsilon}
\newcommand{\mbR}{\mathbb{R}}
\newcommand{\mM}{\mathcal{M}}
\newcommand{\mT}{\mathcal{T}}
\newcommand{\NL}[1]{\textcolor{blue}{#1}}
\newcommand{\p}{\partial}
\newcommand{\tr}{\text{tr}}
\newcommand{\tT}{\tilde{T}}
\newcommand{\Ad}{\operatorname{Ad}}
\newcommand{\Dom}{\operatorname{Dom}}
\newcommand{\mE}{\mathcal{E}}
\newcommand{\PSL}{\mathrm{PSL}(2,\mathbb{R})}
\newcommand{\PSLt}{\widetilde{\mathrm{PSL}(2,\mathbb{R})}}
\newcommand{\psl}{\mathfrak{psl}(2,\mathbb{R})}
\newcommand{\SO}[1]{\textcolor{purple}{#1}}
%\newcommand{\MM}[1]{\textcolor{red}{#1}}
\newcommand{\KL}[1]{\textcolor{brown}{#1}}
\newcommand{\YC}[1]{\textcolor{teal}{#1}}
\newcommand{\cL}{\mathcal{L}}
\begin{document}

% Title
\title{Subsystem Flows and Many-body Ergodicity in QFT}

\author{Yidong Chen$^1$,}
\affiliation{$^1$Department of Mathematics, University of Illinois at Urbana-Champaign, Illinois, IL 61801, USA}
\author{Nima Lashkari$^2$,}
\author{Kwing Lam Leung$^2$}
\affiliation{$^2$Department of Physics and Astronomy, Purdue University, West Lafayette, IN 47907, USA}
\emailAdd{yidongc2@illinois.edu}
\emailAdd{nima@purdue.edu}
\emailAdd{leung60@purdue.edu}

% Date
%\date{\today}

% Abstract
\abstract{Recent developments in the algebraic approach to bulk reconstruction in holography have highlighted the role of half-sided modular inclusions and modular intersections in reconstructing spacetime symmetry transformations. A key ingredient in these algebraic reconstructions is the concept of subsystem flow, which composes modular evolutions associated with an algebra and that of its subsystem. In the first part of this work, we motivate the notion of coarse-grainings in the classical Hamiltonian ergodic hierarchy, distinguishing two cases: subspace and subsystem subalgebras. To identify a classical avatar of modular flow, we associate flows to states, using the Poisson bracket on phase space. We discuss the generalization to the quantum ergodic theory, and in particular to the net of subsystems associated to subregions in many-body quantum systems. We focus on quantum ergodic classes that correspond to {\it non-expected} and {\it split} subsystems that arise in QFT. In the second part of this work, we compute subsystem flows corresponding to causal developments of balls in vacuum CFT and QFT and study their analytic properties. For inclusion of wedges in QFT and inclusion of ball-shaped regions in CFTs, the subregion flow is, up to a reparameterization, the same as time evolution with the global Hamiltonian. We compute the subregion flows in conformal perturbation theory and prove bounds on the comparison of subregion flows of balls with global time evolution in vacuum QFT.}



\maketitle

\NL{To do: 1) Merge appendix E with the end of section 2; 3) Change the title and content of Appendix C to the mathematical background for Section 5.}

\section{Introduction}


The passage from a reversible microscopic description to an apparently irreversible macroscopic one is often attributed to information loss during coarse-graining. Both in the classical and quantum settings, we have two types of coarse-graining that we call subspace and subalgebra coarse-graining.


The intuition of subspaces is a ``one-particle" system, where the particle is either in a set of states or in another set of states. Consider two separate distinguishable particles, each coming with $d_i$ states. The two-particle system has $d_1d_2$ states, whereas if we consider a one-particle subspace of it, we either have particle one in one of the $d_1$ states, or particle two in one of the $d_2$ states; therefore, in the one-particle setup, the number of states add up.

This information loss can be quantified by associating an entropy to a coarse-graining. A relative phase constructed from the fine-grained and coarse-grained densities generates this entropy and provides a first classical avatar of a coarse-graining flow. 


In modular theory, given a state $\omega$ on a von Neumann algebra $\mA$, we define the modular flow $\Delta_\omega^{-it}$, which can be identified as the unique automorphism of the algebra of observables (dynamics) that satisfies the KMS condition. Given a state $\omega$ on $\mA$ and a subalgebra $\mathcal{B}\subset \mA$, we define the {\it subalgebra flow} $\Delta_{\omega;\mA}^{-it}\Delta_{\omega;\mathcal{B}}^{it}$ flow that evolves backward in time by $\Delta_{\omega;B}^{it}$, and then forward in time by $\Delta_{\psi;A}^{-it}$. QFT associates to any inclusion of causally complete regions $B\subset A$, an inclusion of von Neumann algebras $\mA_B\subset \mA_A$. In this case, we call $\Delta_{\omega;A}^{-it}\Delta_{\omega;B}^{it}$ a {\it subregion flow}.

In holography, at the $G_N=0$ limit (infinite $N$), the subregion/subalgebra duality associates to every causally-complete subregion of the bulk spacetime a von Neumann subalgebra on the boundary. The subregion flows play a prominent role in the algebraic reconstruction of the bulk spacetime from the net of boundary large $N$ algebras. 

In the first part of this work, we motivate subregion/subalgebra flows from the point of view of classical Hamiltonian ergodic hierarchy. The standard ergodic hierarchy is formulated in terms of 

In the second part of this work, we study the subregion/subalgebra flows in the vacuum of QFT for the inclusion of causal developments of balls.




\NL{Should we focus on standard inclusions on top of ergodic}

Standard+maximal chaos (HMSI)  = PSL(2,R)

Standard+ ergodic (cyclic-separating)

\NL{Modular Berry phases, $[P_+,P_-]$ commuting with $K$}


\section{Subalgebra flows in quantum ergodic theory}%
\label{section:basics}
\label{subsec:quantum-coarse}
 
% %% ---- NOTATION BOX: entirely new (issue: notation summary) ----
% \begingroup\color{blue}
% \begin{tcolorbox}[
%   enhanced,breakable,
%   colback=blue!5!gray!8,colframe=blue!40!black,
%   title={\bfseries Notation summary},
%   fonttitle=\small,
%   before skip=6pt, after skip=10pt,
%   left=4pt, right=4pt, top=4pt, bottom=4pt]
% \small
% \begin{tabular}{@{}ll@{\quad}ll@{}}
% $\mA,\mB$          & von Neumann algebras ($\mB\subset\mA$)  &
% $\mA'$             & commutant of $\mA$ \\
% $\mB^c$            & relative commutant $\mB'\cap\mA$        &
% $B(\mH)$           & bounded operators on $\mH$ \\
% $\varphi,\omega$   & faithful normal states on $\mA$         &
% $\rho_\varphi$     & density matrix of $\varphi$ \\
% $\ket{\Omega}$     & GNS cyclic/separating vector for $\varphi$  &
% $e_\mB$            & Jones projection $\mH\to\overline{\mB\ket{\Omega}}$\\
% $\Delta_{\omega;\mA}$ & modular operator of $\omega$ on $\mA$  &
% $J_{\omega;\mA}$   & modular conjugation \\
% $K_{\omega;\mA}$   & modular Hamiltonian $-\log\Delta_{\omega;\mA}$ &
% $\sigma_t^\varphi$ & modular flow $\Delta_\varphi^{it}(\cdot)\Delta_\varphi^{-it}$\\
% $\mE:\mA\to\mB$    & normal conditional expectation            &
% $D_{\mA\subset\mB;\omega}(t)$ & subsystem flow (Def.~\ref{def:subsystemflow})\\
% $\mB_{hist}$       & $\bigvee_{t}\alpha_t(\mB)$ (history algebra)  &
% $\mB_{mem}$        & $\bigcap_{t}\alpha_t(\mB)$ (memory algebra)\\
% $S(\varphi\|\varphi_0)$ & Araki relative entropy             &
% $\Delta_{\varphi|\varphi_0;\mA}$ & relative modular operator\\
% \end{tabular}
 
% \medskip
% \noindent
% \textbf{Classical counterparts} (Appendix~\ref{app:classical}):
% $L^\infty(X,\mu)$ replaces $\mA$;
% probability density $\rho$ replaces the state;
% $L^2(X,\mu)$ is the GNS space;
% Poisson flow $\alpha_t^\rho$ generated by $K_\rho=-\log\rho$ replaces modular flow.
% \end{tcolorbox}
% \endgroup
 
%% ---- CLASSICAL-QUANTUM CORRESPONDENCE TABLE: entirely new (issue 16) ----
\begingroup\color{blue}
\begin{table}[ht]
\centering
\renewcommand{\arraystretch}{1.35}
\small
\begin{tabular}{@{}p{6.2cm}p{6.8cm}@{}}
\toprule
\textbf{Classical (Appendix~\ref{app:classical})} & \textbf{Quantum (Section~\ref{section:basics})} \\
\midrule
Commutative algebra $L^\infty(X,\Sigma,\mu)$ on Poisson manifold $X$
  & Von Neumann algebra $\mA$  \\
   Poisson bracket $\{f,g\}$ & Commutator $[a,b]$\\
Commutative algebra $L^\infty(X,\Sigma,\mu)$ on Symplectic manifold $X$
  & Von Neumann factor $\mA$  \\
Probability density $\rho\in L^1_+(X,\mu)$
  & Faithful normal state $\varphi$ on $\mA$ \\
$L^2(X,\mu)$ with cyclic vector $\ket{1}$
  & GNS Hilbert space $\mH_\varphi$ with $\ket{\Omega}$ \\
State Hamiltonian $K_\rho=-\log\rho$
  & Modular Hamiltonian $K_{\omega;\mA}=-\log\Delta_{\omega;\mA}$ \\
Poisson modular flow $\alpha_t^\rho(f)=f\circ\Phi_{-t}^\rho$
  & Modular flow $\sigma_t^\varphi(a)=\Delta_\varphi^{it}a\Delta_\varphi^{-it}$ \\
Classical KMS identity $\omega_\rho(\{f,g\})=\omega_\rho(f\delta_\rho(g))$
  & KMS condition at $\beta=-1$ (Eq.~\ref{eq:KMS-function}) \\
Conditional expectation $E_\mB f=\mathbb{E}_\mu[f|\Sigma_0]$
  & Normal conditional expectation $\mE_\mB:\mA\to\mB$ \\
KL-divergence $S(\rho\|\rho_0)=\int\rho\log(\rho/\rho_0)d\mu$
  & Araki relative entropy $S(\varphi\|\varphi_0)$ \\
Poisson centralizer $\mB_P^c=\{h:\{h,b\}=0\;\forall b\in\mB\}$
  & Relative commutant $\mB^c=\mB'\cap\mA$ \\
Product state $\rho=\rho_Y\rho_Z$ (no correlations)
  & Uncorrelated state $\omega=\omega_\mB\otimes\omega_R$ (Takesaki's condition) \\
K-property: $\bigcap_{t\ge0}\mB_t=\mathbb{C}1$
  & Quantum K-system: $\mB_{mem}=\mathbb{C}\id$ \\
\midrule
\multicolumn{2}{@{}l}{\itshape Key quantum departure: Takesaki's theorem
  (§\ref{subsec:quantum-coarse-2}) links modular invariance of $\mB$ to existence} \\
\multicolumn{2}{@{}l}{\itshape of $\varphi$-preserving $\mE:\mA\to\mB$;
  no classical analogue—every classical subalgebra always admits one.} \\
\bottomrule
\end{tabular}
\caption{Dictionary between the classical ergodic constructions of
  Appendix~\ref{app:classical} and their quantum counterparts in
  Section~\ref{section:basics}.}
\label{tab:classical-quantum-dictionary}
\end{table}
\endgroup
 
%% ---- PREAMBLE PARAGRAPH: substantially rewritten (issue 13, preamble fix) ----
%% Original had a dangling sentence and no road-map; new text is entirely new.
The observables of a physical system form a von Neumann algebra. Classical observables are in the center of this algebra. A quantum system with no non-trivial classical observables is a von Neumann \emph{factor}. For a finite quantum system, such factors are
finite matrix algebras\footnote{The factor associated to a $d$-level quantum system is the algebra of $d\times d$ complex matrices, $M(d,\mathbb{C})$}. 
In this work, we are primarily interested in the infinite-dimensional factors describing subsystems of many-body quantum systems in the thermodynamic and continuum limit.

Quantum systems come naturally equipped with a notion of modular dynamics that associates to every state a unitary flow. The set of operators that are invariant under modular flow (conserved modular charges) forms an algebra of modular charges \footnote{The subalgebra of modular charges is called {the centralizer} in mathematical jargon.}. Using the ergodic properties of modular flows, one classifies infinite-dimensional factors. To state the results easily, we consider connected modular two-point functions:
\begin{align}
    f_{ab,con}^{(\omega)}(t)=\omega(a_1 \sigma_\omega^t(a_2))-\omega(a_1)\omega(a_2)\ .
\end{align}
If there are classical observables, they are inside the intersection of modular charges of all states. Since our focus is modular ergodicity, we focus on von Neumann factors. They can be classified as
\begin{enumerate}
    \item {\bf Type I$_\infty$:} The modular flow of any state is quasi-periodic, never ergodic, and there are Poincaré recurrences. In other words, in any state, all modular two-point functions are quasi-periodic.
    \item {\bf Type II} There exist states where certain modular two-point functions can decay, but the generator of modular flow (modular Hamiltonian) is inside the algebra, which means that their correlators are time-independent and can never decay. Further, if there exists a state (finite trace) with trivial modular flow, we call it type II$_1$, otherwise, it is type II$_\infty$. 
    \item {\bf Type III:} The modular Hamiltonian is never inside the algebra. However, it could still be possible that for select $t$, there exist operators in the algebra such that $\sigma_\omega^t(a)=u_t a u_t^\dagger$. If this never occurs, we say the algebra is type III$_1$. Type III algebras admit no traces or density matrices.
\end{enumerate}
In a theory of ergodic theory specialized to modular flow of states (modular ergodic theory), type III$_1$ factors are the most mixing and chaotic. In physics, the local algebras of causally complete regions in QFT are type III$_1$ factors. 


In many-body finite quantum systems, the algebra of observables is the tensor product of single-body algebras. For a two-body finite system $A=BR$, if the algebra of $B$ is a factor, we say $B$ is a subsystem ({\it subfactor}) of $A$ which is described by an {\it inclusion} of von Neumann algebras: $\mB\subset \mA$. The algebra of $R$ can be viewed as the {\it relative commutant} of this inclusion: $\mB^c=\mB'\cap \mA$. Maps such as partial trace over $R$ are called {\it conditional expectations}: $\mE:\mA\to \mB$. In the more general case, where the algebras of $\mA$ and $\mB$ are not factors, these conditional expectations are interpreted as a generalization of partial trace to {\it coarse-graining}. 

Subfactors play a crucial role in modular ergodic theory. The intuition is that the modular flow of $A$ preserves $A$, but can mix $B$ and $R$. From the point of view of subsystem $B$, the relative commutant $\mB^c$ plays the role of environment degrees of freedom. The evolution of $B$ with the modular flow of $A$ sends it to 
\begin{align}
    \mB_t=\sigma_{\omega;A}^{it}(\mB)\ .
\end{align}
Since, in general, $\mB_t$ leaves $\mB$, the effective open quantum dynamics seen by in $B$ is described by the unital CP map
\begin{align}
    \mT_t(B\subset A;\omega)\equiv \mE\circ \sigma_{\omega;\mA}^t\circ \mE:\mB\to \mB
\end{align}
Note that this dynamics depends crucially on the state $\omega_A$. We say the dynamics can be {\it reduced} to $\mB$ if $\mB_t=\mB$ at all times. This means that information in $B$ never leaks outside. A key result of the theory of algebraic quantum error correction is that this occurs if and only if 
\begin{align}
    \forall b\in \mB: \qquad \sigma_{\omega;A}^t(b)=\sigma_{\omega;B}^t(b)\ .
\end{align}
Therefore, the map 
\begin{align}
    \mathcal{D}_t=\mE\circ \sigma_{\omega;A}^{it}\circ \sigma_{\omega;B}^{-t}\circ \mE
\end{align}
is a measure of the residual dynamics left by the
mismatch between the modular flow of the full algebra $\mA$ and that of a subalgebra
$\mB$. We call this the \emph{subsystem flow} 
(see Definition~\ref{def:subsystemflow}), and argue it plays a central role in quantum ergodic hierarchy (Section~\ref{subsec:qeh}), which classifies inclusions from ergodic through
mixing to quantum K-systems.  A key result, developed in
Section~\ref{subsec:geometry-K-systems}, is that for the \emph{modular} K-systems that
arise in QFT, this hierarchy collapses: a single nesting condition forces generating,
forgetfulness, and maximal chaos at the universal rate $\lambda=2\pi$.

In this section we develop the operator-algebraic framework that will be needed to
study subregion flows in QFT.  For completeness of presentation, and to further the intuition about the role of
subsystem flows in classical ergodic theory, we include a review of the relevant
classical concepts in Appendix~\ref{app:classical}. Throughout, the reader is invited to compare each
quantum construction with its classical counterpart collected in
Table~\ref{tab:classical-quantum-dictionary} and in Appendix~\ref{app:classical};
this comparison both motivates the definitions and highlights the genuinely new
structure that non-commutativity introduces.
 
%\endgroup
 

 
% -------------------------------------------------------------------

 
In this work, we consider dynamics that include the modular flow of arbitrary states.
This is the quantum analog of specializing to arbitrary Hamiltonian dynamics.
 
\subsection{Ergodicity in modular dynamics}

When the algebra $\mA$ has a non-trivial center $Z\subset \mA$, then this center lies
inside the centralizer of all modular flows:
\begin{align}
    \forall z\in Z,\;\forall t,\omega:\qquad \Delta_{\omega}^{it}z \Delta_\omega^{-it}=z.
\end{align}
These observables are conserved charges and obstruct ergodicity for any modular
dynamics.  That is why, in quantum ergodic theory, it is natural to restrict to
algebras with trivial center, i.e., factors.
 
%% "Given a fixed modular flow" — original had "Given a fixed modular," (missing word)
\textcolor{blue}{Given a fixed modular flow}, the condition of ergodicity translates to a trivial
centralizer:
\begin{align}
    \mA^\omega\equiv \{a\in \mA:\Delta_{\omega;A}^{it}a\Delta_{\omega;A}^{-it}=a\;\forall t\}=\mathbb{C} 1.
\end{align}
This can only occur in type $III_1$ algebras, as required by \emph{strong mixing} in
modular dynamics~\cite{}.  However, not every state of a type $III_1$ algebra has this
property.  As an example, consider the \emph{Araki-Powers states}: a countably infinite
set of 3-level quantum \textcolor{blue}{systems} in the state $\otimes_{n\in\mathbb{N}}\rho_n$.
If the density \textcolor{blue}{matrix} of each site
$\rho=p_1|1\rangle\langle1|+p_2|2\rangle\langle2|+(1-p_1-p_2)|3\rangle\langle3|$
with $p_1/p_2$ irrational, \textcolor{blue}{the resulting algebra} is type $III_1$; however, the
eigenstates $|i\rangle\langle i|_n$ lie in the centralizer of $\omega_A$.
 
% -------------------------------------------------------------------

\paragraph{Dynamics reducible to a subalgebra.}
Given a general dynamical evolution $\alpha:G\to\mA$, we say the inclusion $\mB\to\mA$
is \emph{reducible} if
\begin{align}
    \forall g\in G,\;b\in\mB:\qquad \alpha_g(b)\in \mB.
\end{align}
Then the dynamics can be \emph{reduced} to $\mB$ (also called \emph{autonomous
dynamics}).  We say the system is \emph{irreducible} if there is no non-trivial proper
subalgebra $\mB\subset \mA$ with reducible dynamics.
 
\paragraph{Modular dynamics and quantum coarse-graining.}
Since we are interested in ergodic modular dynamics, we restrict to subfactor
coarse-grainings. In the GNS Hilbert space, every observable $a\in \mA$ is associated
with a state that it creates as it acts on the vacuum $a|\omega^{1/2}\rangle$.  Modular
theory associates to this state a flow generated by the modular Hamiltonian
$K_{a\omega^{1/2};B}$ on subalgebra $\mB$.  \textcolor{blue}{This is the quantum analog of the
Poisson bracket closure condition for classical Poisson subalgebras: a von Neumann
subalgebra $\mB$ is stable under commutator flow just as a Poisson subalgebra is
stable under the Poisson bracket.}
 
For the modular flow of a state $\omega$ of $\mA$ to be reducible with respect to an
inclusion $\mB\subset \mA$ means that
\begin{align}
    \forall t\in\mathbb{R}, \qquad
    \mB_t=\Delta_{\omega;A}^{it}\mB\Delta_{\omega;A}^{-it}\subset \mB.
\end{align}
According to Takesaki, this holds if and only if there exists a $\omega$-preserving
normal conditional expectation $\mE_\omega: \mathcal{A} \to \mB$.  When this holds,
\textcolor{blue}{the following two properties are also satisfied:}
\begin{enumerate}
    \item The modular flows of $\omega_A$ and $\omega_B$ are consistent:
    \begin{align}
    \forall t\in\mathbb{R}:\qquad
    \Delta_{\omega;A}^{it}b\Delta_{\omega;A}^{-it}=\Delta_{\omega;B}^{it}b\Delta_{\omega;B}^{-it},
    \qquad b\in\mB.
    \end{align}
    \item For the Jones projection $e_B$:
    \begin{align}
        e_B\Delta_{\omega;A}^{it}e_B=\Delta_{\omega;B}^{-it}e_B.
    \end{align}
\end{enumerate}
This means that the modular flow of $\mA$ is autonomous with respect to the subfactor
$\mB\subset \mA$ if and only if the state $\omega$ has no correlations between $\mA$
and the relative commutant $R=\mB^c$.  In this case, there is a one-to-one
correspondence between states of $R$ and the conditional expectations $\mE:\mA\to \mB$.
 
The Jones projection represents the conditional expectation as an operator on the GNS
space.  For an uncorrelated state $\omega_A=\omega_R\otimes\omega_B$,
\begin{align}
    \mathrm{tr}_B\circ \sigma_{\omega;A}^{(t)}\circ \mathrm{tr}_B
    =\sigma_{\omega;B}^{(t)}\circ \mathrm{tr}_B, \qquad
    \sigma^{(t)}_{\omega;A}=\sigma^{(t)}_{\omega;R}\otimes \sigma^{(t)}_{\omega;B}.
\end{align}
 
\begin{definition}[Subsystem flow]\label{def:subsystemflow}
   For an inclusion of von Neumann algebras $\mB\subset \mA$ in state $\omega$, the
   \emph{subsystem flow} is
\begin{align}
    D_{A\subset B;\omega}(t)\equiv e_B\Delta_{\omega;A}^{it}\Delta_{\omega;B}^{-it}e_B.
\end{align}
\end{definition}
 
The subsystem flow is insensitive to any information in the orthogonal subspace
$(1-e_B)\mH$.  It captures the residual dynamics due to the failure of the
coarse-grained dynamics to capture the full microscopic dynamics.  When the Takesaki
condition holds, the coarse-grained dynamics captures the full microscopic dynamics and
the subsystem flow acts on $\mB$ unitarily; otherwise, it is a unital CP map.
 
%% Issue 9: unified S(·‖·) notation. Original mixed | and ‖.
The subsystem flow is the generating function for relative entropies.
\textcolor{blue}{Using $K_{\omega;A}-K_{\omega;B}=\log\Delta_{\omega;B}-\log\Delta_{\omega;A}$
and denoting Araki relative entropy as $S(\cdot\|\cdot)$ throughout:}
\begin{align}\label{eq:subsystem-flow-entropy}
    &i\frac{d}{dt}\langle e_B\Psi|\Delta_{\omega;A}^{it}\Delta_{\omega;B}^{-it}|e_B\Psi\rangle\big|_{t=0}
    =\langle e_B\Psi|K_{\omega;A}-K_{\omega;B}|e_B\Psi\rangle\nonumber\\
    &=\textcolor{blue}{S(e_B\Psi\|\Omega;A)-S(e_B\Psi\|\Omega;B)+S(e_B\Psi\|\Omega;B')-S(e_B\Psi\|\Omega;A')}
    \geq 0,
\end{align}
where we have used $B\subset A \Rightarrow A'\subset B'$ and monotonicity of relative
entropy.  This expression vanishes for all $e_B|\Psi\rangle$ if and only if the Takesaki
condition is satisfied.
 
%% QEC introductory sentence (readability fix)
\textcolor{blue}{The language of quantum error correction provides an intuitive picture: the
subspace $e_B\mH$ is the \emph{code subspace}, and the conditional expectation down to
$\mB$ is the \emph{noise map} (erasure of the relative commutant).}
Takesaki's condition is the statement that this noise map is exactly correctable for
all states in the code subspace if and only if, for all states in the code subspace,
there are no correlations between $B$ and the relative commutant.
 
The subsystem flow has the following properties:
\begin{theorem}\label{thm:subsystem-flow-properties}
The subsystem flow $\mathcal{D}_{A\subset B;\omega}(t) = e_B \Delta_{\omega;A}^{it}
\Delta_{\omega;B}^{-it} e_B$ captures the residual open-system dynamics of the
subsystem and satisfies:
\begin{enumerate}
    \item \textbf{Analytic extension:} extends to the strip
    $\mathcal{S}(0,-\tfrac{1}{2}) := \{z : 0 \ge \operatorname{Im}(z)
    \ge -\tfrac{1}{2}\}$ as a bounded analytic function.
    \item \textbf{Strong continuity:} bounded, with strongly continuous boundary values.
    \item \textbf{Vacuum invariance:} $\mathcal{D}_{A\subset B;\omega}(z)|\Omega\rangle
    = |\Omega\rangle$.
    \item \textbf{Subspace preservation:} $e_B \mathcal{D}_{A\subset B;\omega}(z)
    e_B = \mathcal{D}_{A\subset B;\omega}(z)$.
    \item \textbf{Generalized cocycle identity:}
    \begin{equation}
        \mathcal{D}_{A\subset B;\omega}(t+s)
        = \mathcal{D}_{A\subset B;\omega}(t)\,
          \sigma^{(t)}_{\omega;B}\!\left(\mathcal{D}_{A\subset B;\omega}(s)\right)
          + \mathcal{M}(t,s),
    \end{equation}
    where $\mathcal{M}(t,s)=e_B u_t e_B^\perp \sigma^{(t)}_{\omega;B}(u_s)e_B$ is the
    leakage term through $e_B^\perp=1-e_B$.
    \item \textbf{Dissipative generator:}
    $i \frac{d}{dt} \mathcal{D}_{A\subset B;\omega}(t)|_{t=0}
        = e_B (K_{\omega;A} - K_{\omega;B}) e_B$.
    \item \textbf{Takesaki equivalence:} $\mathcal{D}_{A\subset B;\omega}(t)=e_B$
    for all $t$ (and $\mathcal{M}(t,s)=0$) iff a state-preserving conditional
    expectation $\mathcal{E}: \mA \to \mB$ exists.
\end{enumerate}
\end{theorem}
 
% -------------------------------------------------------------------
\subsection{Quantum ergodic hierarchy}%
\label{subsec:qeh}
 
The modular flow of a state $\omega$ of $\mA$ might not be reducible on subalgebra
$\mB$, meaning that $\mB_t\subset \mA$ need not equal $\mB$.  As $\mB_t$ explores
$\mA$, it might not reach everywhere.  In that case, we have a proper \emph{history
subalgebra}:
$\bigvee_{t\in\mathbb{R}}\alpha_t(\mB)=\mB_{hist}$.
It is also possible that all $\mB_t$ retain a permanent memory of a set of operators,
giving a proper \emph{memory subalgebra}:
$\bigcap_{t\in\mathbb{R}}\alpha_t(\mB)=\mB_{mem}$.
 
\begin{definition}[Generating inclusion]\label{def:generating}
If the dynamics reaches everywhere in $\mA$, the inclusion is \emph{generating}:
\begin{align}\label{geneating}
    \bigvee_{t\in\mathbb{R}}\mB_t=\mA.
\end{align}
\end{definition}
 
\begin{definition}[Forgetful inclusion]
If the dynamics keeps no permanent memory, the inclusion is \emph{forgetful}:
$\bigcap_{t\in\mathbb{R}}\mB_t=\mathbb{C}\id$.
\end{definition}
 
The relative commutant of the total history of $\mathcal{B}$ is exactly the permanent
memory of its relative commutant:
$({\mathcal{B}_{hist}})^c = (\mathcal{B}^c)_{mem}$.
Similarly, $(\mathcal{B}^c)_{hist} \subseteq (\mathcal{B}_{mem})^c$.
 
\paragraph{The ergodic hierarchy and subsystem flows.}
For a $W^*$-dynamical system $(\mathcal{A}, \omega, \alpha_t)$:
\begin{enumerate}
    \item \textbf{Ergodicity:} fixed-point algebra trivial:
    $\{A \in \mathcal{A} \mid \alpha_t(A) = A\; \forall t\} = \mathbb{C}\mathbf{1}$.
 
    \item \textbf{Mixing:} \textcolor{blue}{A system is \emph{strongly mixing in the future} if all
    dynamical correlations decay asymptotically:
    $\lim_{t \to +\infty} \omega(A \alpha_t(B)) = \omega(A)\omega(B)$ for all
    $A, B \in \mathcal{A}$.  (Strong mixing in the past replaces $t\to+\infty$ by
    $t\to-\infty$.)}  Mixing guarantees that any localized excitation eventually spreads
    across the system but does not mandate a strictly monotonic flow of information.
 
    \item \textbf{Quantum $K$-systems:} nested
    ($\alpha_t(\mathcal{B}) \subset \mathcal{B}$ for $t \ge 0$),
    generating ($\mathcal{B}_{hist} = \mathcal{A}$), and forgetful
    ($\mathcal{B}_{mem} = \mathbb{C}\mathbf{1}$).
\end{enumerate}
 
\textcolor{blue}{The quantum error correction interpretation of the K-system condition is developed in
Section~\ref{subsec:QEC-ergodic}, where we will see that the inclusion $\mB\subset\mA$
in a $K$-system perfectly realizes a QEC duality: the nesting condition makes
$T_t = \mathcal{E}\circ\alpha_t|_\mathcal{B}$ a contractive semigroup; the generating
condition guarantees eventual exploration of the full algebra; and the trivial tail
condition forces the complementary environment to undergo complete asymptotic
thermalization.}
 
% -------------------------------------------------------------------
\subsection{Characteristic flows}
\label{subsec:char-flows}
 
\paragraph{Common cyclic and separating vector.}
Clearly, the action of the subsystem flow on general operators in $\mA$ is non-unitary.
It becomes unitary when $|\Omega\rangle$ is a \emph{common cyclic and separating
vector} for $\mA$ and $\mB$.  In that case the code subspace is the entire Hilbert
space, $e_B=\id$, and the standard cocycle identity holds ($\mM(t,s)=0$).  In the
AQFT literature, such subsystem flows are called \emph{characteristic functions}.  In
fact, the inclusion has a common cyclic and separating vector if and only if the
subsystem flow $D(t)$ satisfies properties 1--5 of Theorem~\ref{thm:subsystem-flow-properties}
with $e_B=\id$ and $\mM(t,s)=0$ (see \cite{} for a proof).
 
\begin{definition}[Characteristic function]\label{def:char-fn}
Given a pair of von Neumann algebras $\mathcal{B}\subset\mathcal{A}$ and a common
cyclic and separating vector $|\Omega\rangle\in\mathcal{H}$, the
\emph{characteristic function} is
\begin{equation}
    \mathcal{D}_{\mathcal{A}, \mathcal{B}}(t) := \Delta_\mathcal{A}^{-it}\Delta_\mathcal{B}^{it}.
\end{equation}
\end{definition}
 
We have already seen that $\mathcal{D}_{\mathcal{A}, \mathcal{B}}(t)$ extends to
imaginary values $\mathcal{D}_{\mathcal{A}, \mathcal{B}}(-i\alpha)$ for
$0\leq \alpha\leq 1/2$.  The crucial fact is that it extends to the entire strip
$\mathcal{S}(0, -1/2)$~\cite{revolutionizing}:
 
\begin{lemma}\label{lem:char-fn-analytic}
Given $\mathcal{B}\subset\mathcal{A}$ with common cyclic and separating vector
$|\Omega\rangle$, the characteristic function extends to a bounded analytic function on
$\mathcal{S}(0,-1/2)$, strongly continuous on the boundary.
\end{lemma}
 
We list the properties of the characteristic function~\cite{revolutionizing}:
\begin{enumerate}
    \item $\mathcal{D}_{\mathcal{A},\mathcal{B}}(t)$ and
    $\mathcal{D}_{\mathcal{A}, \mathcal{B}}(t - \frac{i}{2})$ are unitary and strongly
    continuous in $t$;
    \item $\mathcal{D}_{\mathcal{A},\mathcal{B}}(z)$ bounded analytic on $\mathcal{S}(0,-\frac{1}{2})$;
    \item $\mathcal{D}_{\mathcal{A},\mathcal{B}}(t)|\Omega\rangle = |\Omega\rangle$;
    \item Cocycle condition:
    $\mathcal{D}_{\mathcal{A},\mathcal{B}}(t+s)
    = \sigma_t^{|\Omega\rangle}(\mathcal{D}_{\mathcal{A},\mathcal{B}}(s))
      \mathcal{D}_{\mathcal{A},\mathcal{B}}(t)$;
    \item $\mathcal{D}_{\mathcal{A},\mathcal{B}}(t - \frac{i}{2})
    = J_\mathcal{A}\mathcal{D}_{\mathcal{A},\mathcal{B}}(t)J_\mathcal{B}$; in particular
    $\mathcal{D}(-\frac{i}{2}) = J_\mathcal{A}J_\mathcal{B}$, generating the Jones tower:
    $\mathcal{A}\supset\mathcal{B} \supset D^*\mathcal{A}D \supset D^*\mathcal{B}D\supset\cdots$;
    \item For all $t\in\mathbb{R}$: $\mathcal{D}(t)\,D^*\mathcal{A}D\,\mathcal{D}(t)^*\subset\mathcal{A}$.
\end{enumerate}
 
\begin{theorem}[\cite{revolutionizing}]\label{thm:char-fn-unique}
Each subalgebra $\mathcal{B}\subset\mathcal{A}$ for which $|\Omega\rangle$ is cyclic
is uniquely associated with its characteristic function.
\end{theorem}
 
An inclusion with a common cyclic and separating vector is always generating.  We
focus on type $III_1$ subfactors with a common cyclic and separating vector, the
generic setting in QFT subregion inclusions.
 
%% Issue 7 fixed and examples restructured as a ladder
In infinite dimensions, the structure of inclusions can be more exotic. \textcolor{blue}{The
three cases below form a ladder of increasing ``strangeness,'' distinguished by whether
and how the notion of an uncorrelated (product) state breaks down:}
 
\textcolor{blue}{\medskip
\noindent\textbf{Level 1 — Normal conditional expectations exist (finite index):}}
 
\paragraph{Araki-Powers inclusions (non-ergodic inclusion of type $III_1$ algebras).}
\textcolor{blue}{Let $\mB\subset\mA$ be a \emph{finite-index} inclusion of type $III_1$ factors},
meaning that the relative commutant \textcolor{blue}{$\mathcal{R}=\mB'\cap\mA$}
is a finite-dimensional matrix algebra $M_n(\mathbb{C})$.  By the Jones--Kosaki index
theory, \textcolor{blue}{$[\mathcal{A}:\mathcal{B}]=n^2<\infty$}.  Since $\mathcal{R}$
is finite-dimensional, it is type $I$, so the inclusion is normal (conormal) in the
sense that \textcolor{blue}{$\mathcal{B}\vee\mathcal{R}=\mathcal{A}$}.
Normal conditional expectations $\mE:\mA\to \mB$ exist and are parametrised by states
on $M_n(\mathbb{C})$.
 
The coarse-graining flow $U_t=\sigma_t^{\omega_\mathcal{A}}\circ(\sigma_t^{\omega_\mathcal{B}})^{-1}$
acts non-trivially only on $\mathcal{R}=M_n(\mathbb{C})$ and preserves it; it does
\emph{not} satisfy generating property~\eqref{geneating} and is \emph{not mixing}.
\textcolor{blue}{This inclusion is the quantum analog of an ergodic but non-mixing classical system: a
conditional expectation exists (the QEC recovery map is well-defined), but the
inclusion is non-ergodic because operators in $R$ that are eigenstates of $\omega_R$
lie in the centralizer of $\omega_A$.}
 
\textcolor{blue}{\medskip
\noindent\textbf{Level 2 — No conditional expectations, but non-trivial relative commutant:}}
 
\paragraph{Local algebras of QFT (no normal conditional expectations).}
In a local conformal net $\{\mathcal{A}(I)\}$ on the circle, let $I_1,I_2$ be disjoint
open intervals with closures sharing exactly one endpoint, and set
$I=\mathrm{int}(I_1\cup I_2)$.  Haag duality and strong additivity give
$\mathcal{A}(I_1)'\cap\mathcal{A}(I)=\mathcal{A}(I_2)$ and
$\mathcal{A}(I_1)\vee\mathcal{A}(I_2)=\mathcal{A}(I)$, so the inclusion
$\mathcal{B}:=\mathcal{A}(I_1)\subset\mathcal{M}:=\mathcal{A}(I)$ is conormal with
$\mathcal{R}=\mathcal{A}(I_2)$.  By Theorem~\ref{thm:split-no-expectation} and the fact
that adjacent touching algebras of a non-trivial conformal net never form a split pair
(Wiesbrock~\cite{Wiesbrock1993}), there is \emph{no} normal conditional expectation
$\mE:\mathcal{A}(I)\to\mathcal{A}(I_1)$.  If $\mE$ existed,
Proposition~\ref{prop:conormal-factorization} would give
$\mathcal{A}(I)\cong\mathcal{A}(I_1)\otimes\mathcal{A}(I_2)$, contradicting the
splitting obstruction.
 
\textcolor{blue}{This inclusion occupies a more extreme position than the Araki-Powers case: the
relative commutant is still non-trivial, but every normal conditional expectation is
obstructed.  Every state on $\mathcal{A}(I)$ is universally entangled between
$\mathcal{A}(I_1)$ and $\mathcal{A}(I_2)$; no QEC recovery map exists.}
 
\textcolor{blue}{\medskip
\noindent\textbf{Level 3 — No conditional expectations, no split, full universality:}}
 
\paragraph{General von Neumann algebras: rigorous framework.}%
\label{sec:conormality-expectations}
\textcolor{blue}{We now provide rigorous results covering all three cases.}
 
Let $\mB\subset\mA\subset B(\mathcal{H})$ be an inclusion of $\sigma$-finite type
$III_1$ factors.  The inclusion is \emph{standard} with respect to $\Omega\in\mathcal{H}$
if $\Omega$ is cyclic and separating for $\mA$, $\mB$, and the relative commutant
$\mB^c=\mB'\cap\mA$.  It is \emph{split} if a type $I$ factor $\mathcal{F}$ exists with
$\mB\subset\mathcal{F}\subset\mA$, and \emph{conormal} if $\mB\vee\mB^c=\mA$.
Throughout, $\mE:\mA\to\mB$ is assumed faithful and normal unless stated otherwise.
Tomiyama's theorem~\cite{Tomiyama1957,Tomiyama1959,Takesaki1972} gives
$\mE(b_1ab_2)=b_1\mE(a)b_2$ for $b_1,b_2\in\mB$, $a\in\mA$.
 
\begin{theorem}[Sufficient condition for no conditional expectation]%
\label{thm:split-no-expectation}
If $\mB$ is a type $III$ factor and $\mA\supset\mB$ is split, then no normal
conditional expectation $\mE:\mA\to\mB$ exists.
\end{theorem}
 
\begin{proof}
Let $\mB\subset\mathcal{F}\subset\mA$ be an interpolating type $I$ factor and suppose
$\mE:\mA\to\mB$ is normal.  Its restriction
$\mE_\mathcal{F}:=\mE|_\mathcal{F}:\mathcal{F}\to\mB$ is a normal norm-one projection
with range $\mB$.  Tomiyama's type theorem~\cite{Tomiyama1959,Takesaki1972} states
that a normally expected subalgebra of a type $I$ algebra is type $I$, contradicting
$\mB$ being type $III$.
\end{proof}
 
\begin{lemma}\label{lem:expectation-slice}
If $\mathcal{B}$ is a factor and $\mE:\mathcal{A}\to\mathcal{B}$ is a faithful normal
conditional expectation, there is a faithful normal state $\chi_R$ on $\mB^c$ such that
$\mE(c)=\chi_R(c)\mathbf{1}$, $\mE(bc)=\chi_R(c)b$.
\end{lemma}
 
\begin{proposition}[Conormality and tensor product]\label{prop:conormal-factorization}
Under the hypotheses of Lemma~\ref{lem:expectation-slice}, multiplication extends to
$m:\mathcal{B}\,\overline\otimes\,\mathcal{R}\overset{\cong}{\to}\mathcal{B}\vee\mathcal{R}$.
If the inclusion is conormal, $\mathcal{A}\cong \mathcal{B}\,\overline\otimes\, \mB^c$~\cite{Faulkner2020,FaulknerSperanza2024}.
\end{proposition}
 
The \emph{Pimsner--Popa index} is defined, and the \emph{index of the inclusion} is
$[\mA:\mB]_0=\inf_\mE \mathrm{Ind}(\mE)$.  When the relative commutant is
infinite-dimensional or no conditional expectations exist, the index is infinite.
 
% -------------------------------------------------------------------
\subsection{Geometry from modular K-systems}
\label{subsec:geometry-K-systems}
 
In this section, we have introduced the quantum ergodic hierarchy for a general
$W^*$-dynamical system $(\mA,\omega,\alpha_t)$.  The classical precursor, recalled in
Appendix~\ref{subsec:classical-ergodic-hierarchy}, already anticipated the quantum
theory: the one-sided inclusion $\alpha_t(\mB_+)\subseteq\mB_+$, $t\ge0$, was built
from an \emph{externally specified} measure-preserving flow.  In the quantum setting,
there is a distinguished choice of dynamics attached to every state: the modular flow
itself.  The resulting \emph{modular} K-systems are far more rigid than generic
K-systems: the single nesting condition already forces the generating and forgetful
conditions, forces maximal chaos with a universal rate, and, when strengthened to a
standard inclusion, forces an emergent continuous conformal symmetry.  These results
are due to Borchers~\cite{Borchers2000} and Wiesbrock~\cite{Wiesbrock1993}; their
formulation in the language of quantum ergodic theory follows~\cite{OFLLM2024}, with
the sharpened geometric statement due to~\cite{LLMO2025}.
 
\subsubsection*{Half-sided modular inclusions}
 
\begin{definition}[Half-sided modular inclusion]\label{def:HSMI}
Let $\mB\subset\mA$ be a proper inclusion of von Neumann factors with a common cyclic
and separating vector $|\Omega\rangle$.  The inclusion is a \emph{future half-sided
modular inclusion} (\textsc{hsmi}$+$), or $\mB$ is a \emph{modular future subalgebra}
of $\mA$, if
\begin{equation}
  \forall t\ge0:\qquad
  \Delta_{\omega;\mA}^{-it}\,\mB\,\Delta_{\omega;\mA}^{it}\subset\mB.
  \label{eq:HSMI}
\end{equation}
\textsc{hsmi}$-$ (a \emph{modular past subalgebra}) is the same condition with $t\le0$.
\end{definition}
 
\textcolor{blue}{The sign in~\eqref{eq:HSMI} is conventional but not arbitrary: setting
$\alpha_t:=\sigma_{-t}^{\omega;\mA}$ identifies~\eqref{eq:HSMI} exactly with the
abstract nesting condition $\alpha_t(\mB)\subset\mB$, $t\ge0$, required of a quantum
K-system.  A modular K-system is then the special case of
Section~\ref{subsec:qeh}'s general construction in which the external dynamics
$\alpha_t$ coincides with (minus) the algebra's own modular flow.  This is the precise
quantum sharpening of the classical anticipation in~\eqref{eq:classical-one-sided-inclusion}:
there the one-sided inclusion was imposed on an arbitrary $\alpha_t$; here it is a
rigid algebraic condition on the state $\omega$ and the subalgebra $\mB$ alone.}
 
\begin{theorem}[HSMI theorem; Borchers, Wiesbrock~{\cite{Borchers2000,Wiesbrock1993}},
see also~{\cite[Thm.~16]{OFLLM2024}}]
\label{thm:HSMI}
Suppose $\mB\subset\mA$ is \textsc{hsmi}$\pm$.  Then the operator
\begin{equation}
  G := \frac{1}{2\pi}\bigl(K_{\omega;\mA} - K_{\omega;\mB}\bigr)
  \label{eq:HSMI-generator}
\end{equation}
is positive (resp.\ negative) semi-definite for \textsc{hsmi}$+$ (resp.\ \textsc{hsmi}$-$),
and generates a unitary one-parameter group $U(s):=e^{isG}$ satisfying, for all
$s,t\in\mathbb{R}$,
\begin{align}
  \frac{1}{2\pi}\,[\,K_{\omega;\mA},\,G\,] &= -iG, \label{eq:HSMI-commutator}\\
  \Delta_{\omega;\mA}^{-it}\,U(s)\,\Delta_{\omega;\mA}^{it}
  = \Delta_{\omega;\mB}^{-it}\,U(s)\,\Delta_{\omega;\mB}^{it} &= U(e^{2\pi t}s),
  \label{eq:HSMI-chaos}\\
  J_{\omega;\mA}\,U(s)\,J_{\omega;\mA}
  = J_{\omega;\mB}\,U(s)\,J_{\omega;\mB} &= U(-s),
  \label{eq:HSMI-detailed-balance}\\
  \mB &= U(1)\,\mA\,U(-1). \label{eq:HSMI-translation}
\end{align}
\end{theorem}
 
Equation~\eqref{eq:HSMI-chaos} is the \emph{maximal modular chaos} relation.  We will
see in Section~3.3 that this is precisely the rate saturating the modular chaos bound.
Equation~\eqref{eq:HSMI-detailed-balance} is the quantum detailed balance condition.
 
\textcolor{blue}{The following proposition shows that $G$ is the positive-definite restriction of the
subsystem-flow generator already defined in Section~\ref{subsec:quantum-coarse-2}.}
 
\begin{proposition}\label{prop:G-is-subsystem-generator}
If $\mB\subset\mA$ is \textsc{hsmi}$+$, the common cyclic and separating vector gives
$e_\mB=\id$, and
\begin{equation}
  2\pi G
  = i\,\frac{d}{dt}D_{\mA\subset\mB;\omega}(t)\Big|_{t=0}
  = K_{\omega;\mA}-K_{\omega;\mB}.
\end{equation}
\end{proposition}
\begin{proof}
Since $|\Omega\rangle$ is cyclic for $\mB$, $e_\mB=\id$.  The claim follows from
comparing~\eqref{eq:HSMI-generator} with property~6 of Theorem~\ref{thm:subsystem-flow-properties}.
\end{proof}
 
\begin{theorem}[Half-sided translations; Borchers~{\cite{Borchers2000}},
see also~{\cite[Thm.~18]{OFLLM2024}}]
\label{thm:half-sided-translations}
Let $U(t)=e^{iGt}$ with $U(t)|\Omega\rangle=|\Omega\rangle$, and suppose
$\mB_t:=U(t)\mB U(t)^\dagger\subset\mB$ for all $t>0$.  The following are equivalent:
(1)~$G\ge0$;\quad
(2)~$\Delta_{\omega;\mA}^{1/2}U(t)=U(-t)\Delta_{\omega;\mA}^{1/2}$;\quad
(3)~$\Delta_{\omega;\mA}^{-is}U(t)\Delta_{\omega;\mA}^{is}=U(e^{2\pi s}t)$;\quad
(4)~$\mB_t$ is a modular future subalgebra of $\mB$ for every $t$.
\end{theorem}
 
\subsubsection*{Modular K-systems are automatically maximally chaotic}
 
\begin{corollary}[Maximal modular chaos $\equiv$ modular K-system;
{\cite[Cor.~17]{OFLLM2024}}]
\label{cor:maximal-chaos-K-system}
Let $\mB\subset\mA$ be \textsc{hsmi}$+$ with $\omega$ faithful normal on the type
$III_1$ factor $\mA$.  With $\alpha_t:=\sigma_{-t}^{\omega;\mA}$, the inclusion is a
quantum K-system: $\mB_{hist}=\mA$ and $\mB_{mem}=\mathbb{C}\id$.
\end{corollary}
 
\subsubsection*{Standard inclusions and the emergent conformal group}
 
\begin{theorem}[Standard inclusion implies $\PSL$;
Wiesbrock~{\cite{Wiesbrock1993}}, see also~{\cite[Thm.~29, $U=\id$]{LLMO2025}}]
\label{thm:wiesbrock-PSL}
Let $\mB\subset\mA$ be a standard \textsc{hsmi}$+$ inclusion with common cyclic and
separating vector $|\Omega\rangle$ for $\mA$, $\mB$, $\mB^c$.  Define
$\widehat{\mB} := \mB \vee J_{\omega;\mA}\,\mB\,J_{\omega;\mA}$,
and $G_0,G_1,G_2$ as in~\eqref{eq:three-generators} (with $G_1=G$ the HSMI generator).
Then $[G_0,G_1]=iG_1$, $[G_0,G_2]=-iG_2$, $[G_1,G_2]=-2iG_0$,
and $G_0,G_1,G_2$ exponentiate to a unitary representation of $\widetilde{\PSL}$,
with inversion map $I=J_{\omega;\mA}J_{\omega;\widehat{\mB}}$.
\end{theorem}
 
Theorem~\ref{thm:wiesbrock-PSL} is the precise sense in which geometry emerges from
ergodicity.  In Section~3 we will see this made concrete: the inclusion of wedges and
balls furnishes standard HSMI inclusions whose $\PSL$ representation is literally the
geometric action of boosts, null translations, and special conformal transformations.
The exponential dependence on $e^{2\pi z}$ throughout Sections~3.1 and~3.2 is not a
coincidence of the free-field examples: it is the universal signature of any standard
half-sided modular inclusion.
 
% -------------------------------------------------------------------
%% Issue 17: QEC material moved to its own subsection here
\begingroup\color{blue}\subsection{Quantum error correction interpretation of the ergodic hierarchy}
\label{subsec:QEC-ergodic}
 
The algebraic structures introduced in Sections~\ref{subsec:qeh} and~\ref{subsec:char-flows}
admit a natural interpretation in the language of quantum error correction (QEC).  This
dictionary illuminates why the K-system condition is the maximally chaotic regime in
which \emph{no} classical information can be permanently stored in the subalgebra.
 
At the most basic level, the quantum channel picture provides an intuitive way to think
about the subsystem flow: the conditional expectation $\mE:\mA\to\mB$ acts as the
\emph{erasure channel}, discarding the degrees of freedom in the relative commutant
$\mB^c=\mB'\cap\mA$, while the subsystem flow $D_{\mA\subset\mB;\omega}(t)$ measures
how much information the coarse-grained evolution fails to capture.  The
Takesaki-equivalence property~7 of Theorem~\ref{thm:subsystem-flow-properties} says
precisely that this mismatch vanishes if and only if the erasure is perfectly correctable.
 
More precisely: in the operator-algebraic formulation of QEC, the relative commutant
$\mB^c = \mB' \cap \mA$ serves as the complementary (or conjugate) subsystem to the
primary inclusion $\mB \subset \mA$.  Given a state-preserving normal conditional
expectation $\mathcal{E}: \mathcal{A} \to \mathcal{B}$, this map acts as the primary
\emph{quantum erasure channel}.  The complementary noise channel is given by
$\mathcal{E}^c: \mathcal{A} \to \mathcal{B}^c$.
 
The true open-system dynamics is described by CPU maps
$T_t = \mathcal{E} \circ \alpha_t|_\mathcal{B}$.  The memory algebra
$\mathcal{B}_{mem} = \bigcap_{t} \alpha_t(\mathcal{B})$ is the
\emph{decoherence-free subalgebra} (noiseless subsystem): the maximal subalgebra on
which $T_t$ acts purely reversibly.  The history algebra
$\mathcal{B}_{hist} = \bigvee_{t} \alpha_t(\mathcal{B})$ is the dynamical support or
\emph{information hull} of the code.  The history/memory duality
$({\mathcal{B}_{hist}})^c = ({\mathcal{B}^c})_{mem}$ dictates that operators permanently
hidden from an observer with access to the total history of $\mathcal{B}$ are precisely
the absolute invariants of the conjugate environment $\mathcal{B}^c$.
 
\paragraph{Decoupling theorem.}
The decoupling theorem states that a quantum code is exactly correctable against a
noise channel if and only if the complementary channel completely erases the logical
information.  Algebraically: if $\mathcal{B}_{hist} = \mathcal{A}$, the relative
commutant duality forces $(\mathcal{B}^c)_{mem} = \mathbb{C}\mathbf{1}$.
 
\paragraph{Relating QEC to the ergodic hierarchy.}
In a quantum K-system, the nesting condition ensures that the maps $T_t$ form a
strictly contractive semigroup.  The generating condition guarantees that the subsystem
flow eventually explores the entire global algebra.  The trivial tail condition forces
$(\mathcal{B}^c)_{mem}=\mathbb{C}\mathbf{1}$, so the complementary environment
undergoes complete asymptotic thermalization.  In the modular K-system case
(Corollary~\ref{cor:maximal-chaos-K-system}), the additional structure of
Theorem~\ref{thm:HSMI} makes these statements sharper: the HSMI generator $G$ has a
spectral interpretation as the exponentially amplified ``leakage Hamiltonian,'' and no
QEC code subspace can be protected against the modular erasure.
\endgroup
 

% \section{Subalgebra flows in quantum ergodic theory}\label{section:basics}
% \label{subsec:quantum-coarse}

% The observables of a physical system form a von Neumann algebra. If there are no classical observables (those that commute with all observables) this algebra is a factor; it has a trivial center. For a finite quantum system, such factors are $d\times d$ complex matrices. Infinite-dimensional factors are classified as type I$_\infty$, type II$_1$, type II$_\infty$, and type III factors. Type III algebras admit no density matrices and have no trace.

% Modular theory associates to every state a dynamical flow called modular flow $\sigma_{\omega;t}(a)=\Delta_{\omega,A}^{it}a\Delta_{\omega,A}^{-it}$ that in the case of type I and II algebras is simply $\alpha_{\omega,t}(a)=\omega^{it} a\omega^{-it}$, where $\omega$ is a density matrix. 

% The modular flow of a faithful normal state on a type $I_\infty$ algebra is always inner and hence non-ergodic on the algebra; ergodic modular dynamics requires a type II or III algebra \cite{}. In this work, we focus primarily on the ergodic properties of modular dynamics. In particular, we consider a von Neumann factor with the exotic property that the modular flow of every state is ergodic and strongly mixing. Such an algebra is type III$_1$ \cite{}.

% A composite system is then the tensor product of two such matrix algebras $\mA_B\otimes \mA_R=\mA_A$. Then, a subsystem is a von Neumann subfactor $\mA_B\subset \mA_A$. In this section, motivated by quantum ergodic theory for modular dynamics, we define the {\it subsystem flow} (see Definition \ref{def:subsystemflow}).
% We comment on its interpretation in modular ergodic theory, and it is natural to require chaotic dynamics to be ergodic as in (\ref{}) and the generating property in (\ref{}). This naturally lands us in the world of inclusion of type III$_1$ algebras with common cyclic and separating states, as realized in subregion algebras of QFT. 



% For completeness of presentation, and to further the intuition about the role of subsystem flows in classical ergodic theory, we have included a review of the relevant concepts in Appendix \ref{app:classical}.

% % The structures of Section~2 have direct quantum counterparts obtained by replacing the
% % commutative algebra $L^{\infty}(X,\Sigma,\mu)$ with a non-commutative von Neumann
% % algebra and the classical probability density with a normal faithful state.
% % We introduce these counterparts in the same order---observables, states, Hilbert space,
% % dynamics, coarse-graining, entropy---before specialising to the subspace and subsystem
% % constructions that parallel Sections~2.1 and~2.4.

% % -------------------------------------------------------------------

% \subsection{Review of quantum dynamical systems}

% In this subsection, we review and fix our notation. Readers familiar with standard definitions can skip this subsection.

% \paragraph{Algebra of quantum observables:}
% Quantum observables are described by a 
% \emph{von Neumann algebra} $\mathcal{A}$ a unital $*$-subalgebra of the bounded operators
% $B(\mathcal{H})$ on some separable Hilbert space $\mathcal{H}$, closed in the
% weak operator topology.
% % The commutativity of the classical algebra is replaced by the requirement that $\mathcal{A}$
% % be closed under the adjoint; it need not be abelian.  Indeed the classical theory is recovered
% % by taking $\mathcal{A}=L^\infty(X,\Sigma,\mu)$, acting by multiplication on
% % $\mathcal{H}=L^2(X,\Sigma,\mu)$.
% A von Neumann algebra $\mathcal{A}$ is a \emph{factor} if its center
% $\mathcal{Z}(\mathcal{A}) = \mathcal{A}\cap\mathcal{A}'$ consists only of scalar multiples
% of the identity, where $\mathcal{A}' = \{b\in B(\mathcal{H}) : [b,a]=0\text{ for all }a\in\mathcal{A}\}$
% is the commutant.  Physically, a factor algebra describes a system with no classical
% superselection sectors \footnote{As we review in Appendix \ref{app:classical}, in classical physics, the commutator is replaced by the Poisson bracket and the commutant $\mathcal{A}'$ is the quantum analog of the Poisson
% centralizer introduced in Appendix \ref{app:classical}. It consists of all observables invisible to $\mathcal{A}$. The classical analog of a factor is a non-degenerate Poisson bracket that restricts to a symplectic leaf.}.

% % -------------------------------------------------------------------
% \paragraph{Quantum states:}
% A \emph{normal state} on $\mathcal{A}$ is a positive, normalized, $\sigma$-weakly continuous
% linear functional $\varphi:\mathcal{A}\to\mathbb{C}$, with $\varphi(\id)=1$.  Normal states
% are in one-to-one correspondence with density matrices $\rho_\varphi\in\mathcal{A}_*^+$ via
% $\varphi(a)=\Tr(\rho_\varphi\, a)$ when $\mathcal{A}$ is represented on a Hilbert space with
% a trace.  In the finite-dimensional setting, these are precisely the positive semidefinite
% matrices of unit trace.

% A state $\varphi$ is \emph{faithful} if $\varphi(a^*a)=0$ implies $a=0$; equivalently,
% its density matrix has no zero eigenvalues.  Faithfulness is the quantum counterpart of the
% condition that the probability density $\rho_\nu$ is strictly positive almost everywhere,
% and it is required for the modular theory to be non-trivial.

% A state $\varphi$ on $\mathcal{A}$ is \emph{tracial} if $\varphi(ab)=\varphi(ba)$ for
% all $a,b\in\mathcal{A}$.  Tracial states play the role of the reference Liouville measure in classical dynamics:
% every normal state is absolutely continuous with respect to a faithful tracial state, in
% direct analogy with $\nu\ll\mu$ in the classical setting.  Not every von Neumann algebra
% admits a faithful tracial state; those that do are called \emph{finite} von Neumann algebras.
% In the infinite-dimensional algebras relevant to quantum field theory, one works instead with
% faithful normal states.

% % -------------------------------------------------------------------
% \paragraph{GNS Hilbert space and modular operators.}
% Given a von Neumann algebra $\mathcal{A}$ and a faithful normal state $\varphi$, the
% Gelfand--Naimark--Segal (GNS) construction produces a Hilbert space $\mathcal{H}_\varphi$,
% a cyclic unit vector $|\Omega\rangle\in\mathcal{H}_\varphi$ with
% $\varphi(a)=\langle\Omega|a|\Omega\rangle$, and a $*$-representation
% $\pi:\mathcal{A}\to B(\mathcal{H}_\varphi)$ \footnote{The classical analog is $\mathcal{H}=L^2(X,\mu)$
% with the constant function $|\id\rangle$ as the reference vector.}.

% The Tomita--Takesaki Modular theory assigns to the pair $(\mathcal{A},|\Omega\rangle)$ an
% antilinear operator $S_\varphi$ defined on the dense domain
% $\mathcal{A}|\Omega\rangle$ by
% \begin{equation}
%   S_\varphi\, a|\Omega\rangle = a^*|\Omega\rangle, \qquad a\in\mathcal{A}.
%   \label{eq:Tomita-S}
% \end{equation}
% Its polar decomposition $S_\varphi = J_\varphi\,\Delta_\varphi^{1/2}$ defines the
% \emph{modular conjugation} $J_\varphi$, an antiunitary involution, and the
% \emph{modular operator} $\Delta_\varphi$, a positive self-adjoint operator. 

% % -------------------------------------------------------------------
% \paragraph{Modular dynamics:}
% The Tomita--Takesaki theorem~\cite{Takesaki1970} guarantees that the one-parameter group
% of unitaries $\Delta_\varphi^{it}$, $t\in\mathbb{R}$, implements an automorphism group of
% $\mathcal{A}$:
% \begin{equation}
%   \sigma_t^\varphi(a) = \Delta_\varphi^{it}\, a\, \Delta_\varphi^{-it},
%   \qquad a\in\mathcal{A},\quad t\in\mathbb{R}.
%   \label{eq:modular-flow}
% \end{equation}
% This is the \emph{modular flow} (or \emph{modular automorphism group}) of $\varphi$. 

% The modular operator $\Delta_\varphi$ and $J_\varphi$ have no direct counterparts in classical dynamics. They are genuinely quantum objects encoding the non-commutativity of $\mathcal{A}$. However, in Appendix \ref{app:classical}, we argue that modular flow can be interpreted as a quantum generalization of {\it Poisson flow}. The intuition is that the GNS Hilbert space associates observables with states, and the modular flow associates a flow with a state generated by the modular Hamiltonian $K_\rho=-\log\rho$. In classical physics, the Radon-Nikodym derivative plays the role of $\rho$, and the Poisson bracket with $-\log\rho$ generates a flow $\alpha_t^\rho$ associated to every state; see Appendix \ref{app:classical}. The state $\varphi$ is invariant under its own modular flow:
% $\varphi\circ\sigma_t^\varphi = \varphi$ for all $t\in\mathbb{R}$, in direct analogy
% with the classical fact that $\omega_\rho$ is preserved by $\alpha_t^\rho$.

% % The modular operator $\Delta_\varphi$ and $J_\varphi$ have no direct counterparts in classical dynamics. In Appendix \ref{}, we argue that 

% % a
% % genuinely quantum object encoding the non-commutativity of $\mathcal{A}$.


% The modular flow satisfies the \emph{KMS condition} at inverse temperature $\beta=-1$:
% for all $a,b\in\mathcal{A}$ the function
% \begin{equation}
%   F_{a,b}(t) = \varphi\bigl(a\,\sigma_t^\varphi(b)\bigr)
%   \label{eq:KMS-function}
% \end{equation}
% extends to a function analytic in the strip $0<\operatorname{Im}(t)<1$ and continuous on its
% closure, satisfying the boundary condition $F_{a,b}(t+i)=\varphi(\sigma_t^\varphi(b)\,a)$.
% This is the precise quantum analog of the classical KMS identity~\eqref{eq:KMSclassical},
% which in the classical case reduced to an integration-by-parts identity for the Poisson
% bracket.  Indeed, differentiating~\eqref{eq:KMS-function} and setting $t=0$ recovers
% \begin{equation}
%   \varphi([a,b]) = \varphi(a\,\delta_\varphi(b)),
%   \label{eq:KMS-infinitesimal}
% \end{equation}
% where $\delta_\varphi(b)=[\log\Delta_\varphi,b]$ is the derivation generating $\sigma^\varphi$.
% In the abelian case, $[a,b]=0$ and this becomes $0=\varphi(a\{b,\log\rho\})$, matching~\eqref{eq:KMSclassical}.

% In this work, we consider dynamics that include the modular flow of arbitrary states. This is the quantum analog of specializing to arbitrary Hamiltonian dynamics.

% % Crucially, Takesaki's theorem characterises the modular flow as the \emph{unique} $\sigma$-weakly
% % continuous one-parameter automorphism group of $\mathcal{A}$ for which $\varphi$ is KMS at
% % $\beta=-1$.  This uniqueness has no classical analog: in the abelian case any Hamiltonian flow
% % that preserves $\varphi$ is KMS, so the Poisson structure is needed to single out one.


% \paragraph{Ergodicity in modular dynamics:}

% When the algebra $\mA$ has a non-trivial center $Z\subset \mA$, then this center is inside the centralizer of all modular flows:
% \begin{align}
%     \forall z\in Z,\forall t,\omega:\qquad \Delta_{\omega}^{it}z \Delta_\omega^{-it}=z
% \end{align}
% This means that these observables are conserved charges, and an obstruction to ergodicity in modular dynamics of arbitrary states. That is why, in quantum ergodic theory, it is natural to restrict to algebras with trivial center, i.e., factors.

% Given a fixed modular, the condition of ergodicity translates to a trivial centralizer:
% \begin{align}
%     \mA^\omega\equiv \{a\in \mA:\Delta_{\omega;A}^{it}a\Delta_{\omega;A}^{-it}=a\qquad \forall t\}=\mathbb{C} 1
% \end{align}
% This will only occur in type III$_1$ algebras, as required by {\it strong mixing} in modular dynamics \cite{}. However, not every state of type a III$_1$ algebra has this property. As an example, consider the states of Araki-Powers: a countably infinite set of 3-level quantum systmes in the state $\otimes_{n\in\mathbb{N}}\rho_n$. If the density martix of each site $\rho=p_1\ket{1}\bra{1}+p_2\ket{2}\bra{2}+(1-p_1-p_2)\ket{3}\bra{3}$ with $p_1/p_2$ irrational, the algebra the resulting algebra is type III$_1$; however, the eigenstates $\ket{i}\bra{i}_n$ are in the centralizer of the algebra.


% \paragraph{}

% % -------------------------------------------------------------------

% \subsection{Subsystem quantum coarse-grainings}
% \label{subsec:quantum-coarse}


% \paragraph{Quantum coarse-grainings.}
% A \emph{quantum coarse-graining} is a von Neumann subalgebra $\mathcal{B}\subseteq\mathcal{A}$
% together with a \emph{normal conditional expectation}
% $\mE_{\mathcal{B}}:\mathcal{A}\to\mathcal{B}$, a completely positive, unital,
% $\mathcal{B}$-bimodular map satisfying
% \begin{equation}
%   \mE_{\mathcal{B}}(b_1 a b_2) = b_1\,\mE_{\mathcal{B}}(a)\,b_2, \qquad a\in\mathcal{A},\;b_1,b_2\in\mathcal{B}.
%   \label{eq:CE-bimodular}
% \end{equation}
% This is the quantum analog of the classical conditional expectation
% $\mE_{\mathcal{B}}(f) = \mathbb{E}_\mu[f|\Sigma_0]$ we review  in Appendix \ref{app:classical}.
% On the GNS Hilbert space, $\mE_{\mathcal{B}}$ extends to an orthogonal projection
% $e_{\mathcal{B}}:\mathcal{H}_\varphi\to\overline{\mathcal{B}|\Omega\rangle}$, the {\it Jones
% projection} associated to the inclusion $\mathcal{B}\subset\mathcal{A}$. It is the operator representative of the superoperator $\mE_\mB$ in the GNS Hilbert space. 

% \paragraph{Entropy of a coarse-graining.}
% Given a faithful normal state $\varphi$ on $\mathcal{A}$, a von Neumann subalgebra
% $\mathcal{B}\subseteq\mathcal{A}$, and a $\varphi$-preserving conditional expectation
% $E:\mathcal{A}\to\mathcal{B}$, the coarse-grained state is
% $\varphi_0 = \varphi\circ \mE$, and the \emph{quantum relative entropy} of $\varphi$ with
% respect to $\varphi_0$ is
% \begin{equation}
%   S(\varphi\|\varphi_0) = \langle\Omega|\log\Delta_{\varphi|\varphi_0;\mathcal{A}} |\Omega\rangle=i\frac{d}{dt}\braket{\Omega|\Delta_{\varphi|\varphi_0}^{it}\Omega}\Big|_{t=0},
%   \label{eq:relative-entropy}
% \end{equation}
% where $\Delta_{\varphi|\varphi_0;\mathcal{A}}$ is the relative modular operator of $\varphi$ and $\varphi_0$ on $\mathcal{A}$:
% \begin{align}
%   &S_{\varphi|\varphi_0;A}a\ket{\Omega}=a^*\ket{\Omega_0},\qquad \Delta_{\varphi|\varphi_0;A}=S^\dagger_{\phi|\phi_0;A}S_{\phi|\phi_0;A}\nn\\
%   &\langle a\Omega|\Delta_{\varphi|\varphi_0}|b\Omega\rangle = \varphi_0(b^*a),
%   \qquad a,b\in\mathcal{A},
%   \label{eq:relative-modular}
% \end{align}
% which generalises the classical ratio (Radon-Nikodym derivative) $\rho/\rho_0$. Quantum relative entropy is the quantum analog of the KL-divergence~\eqref{eq:KL-classical}:
% \begin{equation}
%   S(\rho\|\rho_0) = \int_X \rho\log\!\left(\frac{\rho}{\rho_0}\right)d\mu = H(\rho_0)-H(\rho).
%   \label{eq:classical-KL}
% \end{equation}
% The non-negativity $S(\varphi\|\varphi_0)\geq 0$, with equality iff $\varphi=\varphi_0$ on
% $\mathcal{A}$ (i.e.\ iff $\varphi$ is already supported on $\mathcal{B}$), is the
% quantum analog of the monotonicity of entropy under coarse-graining: discarding information
% can only increase entropy.


% % This means that we have an effective description of the dynamics on $\mB$. We say the inclusion is dynamically invariant 
% % \begin{align}
% %     \mE_B\circ \alpha
% % \end{align}

% \paragraph{Different types of coarse-grainings:}

% Just as in the classical theory, there are two natural classes of subalgebras that we call subspaces and subsystems. To build intuition, we first motivate them intuitively using finite quantum systems (matrix algebras):
% \begin{enumerate}
%     \item {\bf Subspace inclusions:} The Hilbert space splits into the subspace and its orthogonal complement: $\mH=P\mH\oplus(1-P)\mH$. The subalgebra of {\it coarse} consists of operators $\mA_P=\pi(\mA)=P\mA P$. We have the inclusion $\pi_*(\mA_P)=\mA_P\oplus \mathbb{C}(1-P)\subset \mA$ \footnote{The inclusion is the pull-back of the projection map.}.
    
%     The inclusion has a non-trivial center $\mathbb{C}P\oplus \mathbb{C}(1-P)=Z_P\subset \pi_*(\mA_P)$.
    
%     Coarse-grainings are described by unital CP maps that project down to the {\it coarse subalgebra}. In mathematical terminology, they are called {\it conditional expectations}. There is a one-to-one correspondence between conditional expectation $\mE_{\omega_{1-P}}:\mA\to \mA_P$ and states $\omega_P$ on the complementary subspace $(1-P)\mH$:
%     \begin{align}
%     \forall a\in \mA:\qquad \mE_P(a)=Pa P\oplus \omega_{1-P}((1-P)a(1-P))\   
%     \end{align}
% A state $\omega$ (density matrix) on $\mA$ is preserved by this conditional expectation if and only if it has no off-diagonal elements between $P\mH$ and $(1-P)\mH$. These are density matrices that are decohered in the binary measurement basis: $\{P,1-P\}$. 
    
% The subspace inclusions above are particularly relevant for single-particle systems, or Hilbert spaces that do not decompose into a tensor product of smaller ones. For an increasing family of projections $P_i\leq P_{i+1}$, we have
% $\mathcal{A}_{P_i}\subseteq\mathcal{A}_{P_i}$, giving a net of coarse observable subspaces that parallels the classical net $\mathcal{A}_{O_i}\subseteq\mathcal{A}_{O_{i+1}}$ for regions
% $O_i\subseteq O_{i+1}$ on the phase space of a single particle (see Appendix \ref{app:classical}). Since the focus of our work is many-body quantum systems and QFT, in this work, we do not consider such inclusions.

% \item {\bf Subsystem inclusions:} The Hilbert space decomposes as a tensor product $\mH_A=\mH_B\otimes \mH_R$, the {\it coarse operators} are $\mA_B:\mH_B\to \mH_B$, and the operators $\mA_R:\mH_R\to \mH_R$ are in the {\it relative commutant} of the inclusion $\mA_B\otimes 1_R\subset \mA_A$. In this case, the inclusion has a trivial center, and we say we have a {\it subfactor}. Most of this work will focus on subfactors. As we review in Appendix \ref{app:classical}, this is a quantum generalization of the classical setup where the phase space splits as $X=Y\times Z$. 

% There is a one-to-one correspondence between conditional expectations $\mE_{\omega_R}:\mA_A\to \mA_B$, and states of the relative commutant $\omega_R$:
%     \begin{align}\label{condexp}
%        \forall b\in \mA_B, r\in \mA_R:\qquad  \mE_{\omega_R}(b\otimes r)=b\: \omega_R(r)\ .
%     \end{align}
% A state is invariant under this coarse-graining if and only if it has no correlations between $B$ and $R$; i.e., $\omega_A=\omega_B\otimes \omega_R$. This is equivalent to saying $I_\omega(A:R)=0$, or 
% \begin{align}
%     \forall t\in\mathbb{R}, b\in\mA_B:\qquad \omega_A^{it} (b\otimes 1_R) \omega_A^{-it}= \omega_B^{it}b\omega_B^{-it}\otimes 1_R
% \end{align}
    
% This type of inclusion is natural in the multi-particle quantum systems we study here. In a many-body quantum system, we can have an increasing family of subsystems, e.g., half-lines in an infinite 1d spin chain or wedges in QFT, $W_i\leq W_{i+1}$; we have
% $\mathcal{A}_{W_i}\subseteq\mathcal{A}_{W_{i+1}}$, giving a net of coarse observable subsystems. See Appendix (\ref{}) for a classical analog of subsystem nets in classical field theory. In this work, we focus in such inclusions.



% \item {\bf Mixed inclusions:} A more general setup occurs when $\mH_A=\oplus_a \mH_B^{(a)}\otimes \mH_R^{(a)}$. This is a mixture of subspaces and subsystems, and is, in fact, the most general inclusion for a finite quantum system:
% \begin{align}
% \pi_*(\mA_B)=\oplus_a \mA_B^{(a)}\otimes 1_R^{(a)}  
% \end{align}
% This inclusion has a non-trivial center $\oplus_a \mathbb{C}1_B^{(a)}\otimes \mathbb{C}1_R^{(a)}$. Once again, the states of the relative commutant $\omega_R$ provide conditional expectations from $\mE:\mA_A\to \mA_B$ given by (\ref{condexp}); however, this correspondence is no longer one-to-one. These conditional expectations leave states of the following form invariant: 
% \begin{align}
%     \omega=\oplus_a p_a \omega_B^{(a)}\otimes \omega_R^{(a)}\ .
% \end{align}
% This is equivalent to the condition (\ref{condexp}). Another way of 
% saying that the state on $BR$ has $I(B:R|Z)=0$, i.e., it is quantum Markov. These are special separate states on $AR$, namely those for which the relative entropy of coarse-graining is zero: 
% \begin{align}
%     S(\omega_A\|\omega_A\circ \mE)=0\ .
% \end{align}
% In the language of quantum error correction, these are states of $A$ that are exactly correctable under the erasure of $R$. Alternatively, one can think of these states as maximum entropy states among those that reduce to $\omega_B$. 
% \end{enumerate}

% In Appendix \ref{app:classical}, we argue that modular flow is a quantum analog of Hamiltonian flow. This motivates restricting allowed dynamics to the modular flow of arbitrary states. We are interested in the study of ergodicity in modular dynamics. If the algebra has a center, that will be left untouched by the modular flow of any state, and an obstruction to ergodicity. Therefore, we are going to focus on case 2 above: subsystems or subfactors.





% %\medskip
% % \noindent\textit{Subspace coarse-grainings.}
% % Let $p\in\mathcal{A}$ be a projection.  The \emph{corner} subalgebra
% % \begin{equation}
% %   \mathcal{A}_p = p\mathcal{A}p \oplus \mathbb{C}(1-p)
% %   \label{eq:corner-subalgebra}
% % \end{equation}
% % is the quantum analog of the region algebra $\mathcal{A}_O$ defined in~\eqref{eq:AO}.
% % Observables in $\mathcal{A}_p$ are completely resolved within the range of $p$ and retain
% % only the event ``outside $p$'' on its complement, exactly as $\mathcal{A}_O$ resolved the
% % microstate within $O$ and identified the complement.  The corresponding conditional
% % expectation is
% % \begin{equation}
% %   E_p(a) = pap + \varphi\bigl((1-p)a(1-p)\bigr)\cdot(1-p),
% %   \qquad a\in\mathcal{A},
% %   \label{eq:CE-subspace}
% % \end{equation}
% % which preserves the state $\varphi$ restricted to the range of $p$ and discards information
% % in the complementary sector.  For an increasing family of projections $p_1\leq p_2$, we have
% % $\mathcal{A}_{p_1}\subseteq\mathcal{A}_{p_2}$, giving a net of coarse observable algebras
% % that parallels the classical net $\mathcal{A}_{O_1}\subseteq\mathcal{A}_{O_2}$ for
% % $O_1\subseteq O_2$.


% \paragraph{Dynamics reducible to a subalgebra:} Given a general dynamical evolution, $\alpha:G\to \mA$, we say an inclusion $\mB\to \mA$ is reducible if 
% \begin{align}
%     \forall g\in G,b\in\mB:\qquad \alpha_g(b)\in \mB\ . 
% \end{align}
% Then, the dynamics can be {\it reduced} to $\mB$. In dynamical systems, this is sometimes called {\it autonomous dynamics}, as well.

% \paragraph{Irreducible dynamics:} Given a general dynamical system, we say it is irreducible if there is no non-trivial proper subalgebra $\mB\subset \mA$ with reducible dynamics.


% \paragraph{Modular dynamics and Quantum Coarse-graining:} Since we are interested in studying modular dynamics that are ergodic, we are going to restrict to subfactor coarse-grainings. In the GNS Hilbert space, every observable $a\in \mA$ is associated with a state that it creates as it acts on the vacuum $\pi_\omega:a\to \ket{a}_\omega=a\ket{\omega^{1/2}}$. Modular theory associates to this state a flow $\Delta_{a\omega^{1/2};A}^{it}$. In other words, we have
% \begin{align}
%     \forall b_1,b_2\in \mB, \partial_t b_2=[K_{b_1\omega^{1/2};B},b_2]\in \mB
% \end{align}
% where $K_{b_1\omega^{1/2};B}$ is the modular Hamiltonian associated to the state $b_1\ket{\omega_B^{1/2}}$ on subalgebra $\mB$. This is the quantum analog of the Poisson bracket $\{b_1,b_2\}$, which means that a von Neumann subalgebra is a quantum analog of a Poisson subalgebra (the subalgebra is closed under the Poisson bracket). 

% For the modular flow of a state $\omega$ of $\mA$ to be reducible with respect to an inclusion $\mB\subset \mA$ means that we should have
% \begin{align}
%     \forall t\in\mathbb{R}, \qquad \mB_t=\Delta_{\omega;A}^{it}\mB\Delta_{\omega;A}^{-it}\subset \mB
% \end{align}
% According to Takesaki, this happens if and only if there exists a $\omega$-preserving normal conditional expectation $\mE_\omega: \mathcal{A} \to B$. When this holds, we further have the following two properties
% \begin{enumerate}
%     \item  The modular flows of $\omega_A$ and $\omega_B$ are consistent
% \begin{align}
% \forall t\in\mathbb{R}:\qquad    \Delta_{\omega;A}^{it}b\Delta_{\omega;A}^{-it}=\Delta_{\omega;B}^{it}b\Delta_{\omega;B}^{-it}
% \end{align}   
% \item 

% \begin{align}
%     e_B\Delta_{\omega;A}^{it}e_B=\Delta_{\omega;B}^{-it}e_B
% \end{align}
% where $e_B$ is the Jones projection associated to $\mE_\omega$.
% \end{enumerate}

% This means that the modular flow of $\mA$ is autonomous with respect to the subfactor $\mB\subset \mA$ if and only if the state $\omega$ has no correlations between $\mA$ and the relative commutant $R$. In this case, there is a one-to-one correspondence between states of $R$ and the conditional expectations from $\mE:\mA\to \mB$. 

% The Jones projection is the GNS operator representation of the superoperator that is the conditional expectation; i.e., partial trace over the relative commutant. Then, we have
% \begin{align}
%     \tr_B\circ \sigma_{\omega;A}^{(t)}\circ \tr_B=\sigma_{\omega;B}^{(t)}\circ \tr_B
% \end{align}
% which can be arranged from an uncorrelated state $\omega$ where
% \begin{align}
% \sigma^{(t)}_{\omega;A}=\sigma^{(t)}_{\omega;R}\otimes \sigma^{(t)}_{\omega;B}
% \end{align}
% because 
% \begin{align}
%     \omega_A=\omega_R\otimes \omega_B\ .
% \end{align}

% The discussion above motivates the following definition:
% \begin{definition}[Subsystem flow]\label{def:subsystemflow}
%    For an inclusion of von Neumann algebras $\mB\subset \mA$ in state $\omega$ we define the {\it subsystem flow}
% \begin{align}
%     D_{A\subset B;\omega}(t)\equiv e_B\Delta_{\omega;A}^{it}\Delta_{\omega;B}^{-it}e_B
% \end{align}
% \end{definition}
% The subsystem flow is insensitive to any information in the orthogonal subspace to the Jones projection $(1-e_B)\mH$. It captures the residual dynamics due to the failure of the reduced dynamics (coarse-grained dynamics) to capture the full (microscopic) dynamics.

% When the Takesaki condition above holds, the coarse-grained dynamics captures the full microscopic dynamics, and the subsystem flow acts on the subalgebra $\mB$ unitarily; however, for a general operator, or in the absence of the Takesaki condition, it describes an open quantum system dynamics; i.e., it is a unital CP map. 

% The subsystem flow can be viewed as the generating function for relative entropies:
% \begin{align}
%     &i\frac{d}{dt}\braket{e_B\Psi|\Delta_{\omega;A}^{it}\Delta_{\omega;B}^{-it}|e_B\Psi}\big|_{t=0}=\braket{e_B\Psi|K_{\omega;A}-K_{\omega;B}|e_B\Psi}\nn\\
%     &=S(e_B\Psi|\Omega;A)-S(e_B\Psi\|\Omega;B)+S(e_B\Psi|\Omega;B')-S(e_B\Psi\|\Omega;A') \geq 0
% \end{align}
% where we have used if $B\subset A$ then $A'\subset B'$, and the monotonicity of relative entropy. The expression above vanishes for all vectors $e_B\ket{\Psi}$ if and only if the Takesaki condition is satisfied. 

% In the language of quantum error correction, the subspace $e_B\mH$ is called the code subspace, and 
% the conditional expectation down to $\mB$ can be interpreted as the noise map, which is erasure (partial trace) of the relative commutant. Takesaki's condition is the statement that this noise map is exactly correctable for all states in the code subspace if and only if, for all states in the code subspace, there are no correlations between $B$ and the relative commutant.

% The subsystem flow we defined has the following properties:
% \begin{theorem}
% The subsystem flow $\mathcal{D}_{A\subset B;\omega}(t) = e_B \Delta_{\omega;A}^{it} \Delta_{\omega;B}^{-it} e_B$ captures the residual open-system dynamics of the subsystem and satisfies the following properties:
% \begin{enumerate}
%     \item \textbf{Analytic Extension:} The flow extends from the real line $t$ to a bounded operator-valued function $\mathcal{D}_{A\subset B;\omega}(z) = e_B \Delta_{\omega;A}^{z} \Delta_{\omega;B}^{-z} e_B$ analytic on the strip $\mathcal{S}(0,- \tfrac{1}{2}) := \{z \in \mathbb{C} : 0 \ge \operatorname{Im}(z) \ge - \tfrac{1}{2}\}$.
%     \item \textbf{Strong Continuity:} $\mathcal{D}_{A\subset B;\omega}(t)$ is bounded and its boundary values on the strip are strongly continuous in $t$.
%     \item \textbf{Vacuum Invariance:} $\mathcal{D}_{A\subset B;\omega}(z)|\Omega\rangle = |\Omega\rangle$ for all $z$ in the strip.
%     \item \textbf{Subspace Preservation:} $e_B \mathcal{D}_{A\subset B;\omega}(z) e_B = \mathcal{D}_{A\subset B;\omega}(z)$.
%     \item \textbf{Generalized Cocycle Identity:} On the real line, the flow satisfies the open-system cocycle relation:
%     \begin{equation}
%         \mathcal{D}_{A\subset B;\omega}(t+s) = \mathcal{D}_{A\subset B;\omega}(t) \sigma^{(t)}_{\omega;B}\big(\mathcal{D}_{A\subset B;\omega}(s)\big) + \mathcal{M}(t,s)
%     \end{equation}
%     where $\mathcal{M}(t,s) = e_B u_t e_B^\perp \sigma^{(t)}_{\omega;B}(u_s) e_B$ is the memory/leakage term through $e_B^\perp = 1 - e_B$, generating a non-Markovian Nakajima-Zwanzig evolution.
%     \item \textbf{Dissipative Generator:} Its strong derivative at $t=0$ yields the relative entropy generator on the code subspace:
%     \begin{equation}
%         i \frac{d}{dt} \mathcal{D}_{A\subset B;\omega}(t) \bigg|_{t=0} = e_B (K_{\omega;A} - K_{\omega;B}) e_B
%     \end{equation}
%     \item \textbf{Takesaki Equivalence:} $\mathcal{D}_{A\subset B;\omega}(t) = e_B$ for all $t \in \mathbb{R}$ (and $\mathcal{M}(t,s) = 0$) if and only if there exists a state-preserving conditional expectation $\mathcal{E}: \mA \to \mB$. 
% \end{enumerate}
% \end{theorem}


% % \begin{theorem}
% %     The subsystem flow $D_{A\subset B;\omega}(t)$ is an operator-valued function that captures the residual open-system dynamics of the subsystem and satisfies the following properties: 
% % \begin{enumerate}
% % \item $D(t)$ is a strongly continuous family of contractions on the subsystem Hilbert space $e_B\mH$bounded
% %  and extends from the real line $t$ to the strip
% % $\mathcal{S}(0,-\tfrac{1}{2}) := \{z\in\mathbb{C}: 0\geq\operatorname{Im}(z)\geq-\tfrac{1}{2}\}$ 
% %   \item $\mathcal{D}_{A\subset B}(t)$ and
% %         $\mathcal{D}_{A\subset B}(t-\tfrac{i}{2})$ are strongly
% %         continuous in $t$;

% %   \item $\mathcal{D}_{A\subset B;\omega}(z)$ is bounded analytic on
% %         $\mathcal{S}(0,-\tfrac{1}{2})$;

% %   \item $\mathcal{D}_{A\subset B;\omega}(t)|\Omega\rangle = |\Omega\rangle$ for all
% %         $t\in\mathbb{R}$;

% %   \item \textit{Cocycle condition}: $\mathcal{D}_{A\subset B;\omega}(t+s)
% %         = \sigma^{(t)}_{\omega;A}(\mathcal{D}_{A\subset B;\omega}(s))
% %           \,\mathcal{D}_{A\subset B;\omega}(t)$,
% %         where $\sigma^{(t)}_{\omega;A}$ is the modular automorphism of $\mA$
% %         associated with $|\Omega\rangle$;

% %   \item $\mathcal{D}_{A\subset B;\omega}(t-\tfrac{i}{2})
% %         = e_B J_\mA\,\mathcal{D}_{A\subset B;\omega}(t)\,J_\mB e_B$.
% %         In particular, $\mathcal{D}_{A\subset B;\omega}(-\tfrac{i}{2})=e_B J_\mA J_\mB e_B$.
% %       %  This unitary is the canonical endomorphism and the Petz map for the inclusion;
% %        % it generates the Jones tower
% %         %\begin{equation}
% %          % \mathcal{M}\supset\mathcal{N}
% %         %   \supset\mathcal{D}^*\mathcal{M}\mathcal{D}
% %         %   \supset\mathcal{D}^*\mathcal{N}\mathcal{D}\supset\cdots,
% %         %   \qquad \mathcal{D}:=\mathcal{D}_{A\subset B;\omega}(-\tfrac{i}{2});
% %         %   \label{eq:Jones-tower}
% %         % \end{equation}

% % \item  For complex values of the arguments, one finds
% % \begin{align}
% % D(t+i/2)^* e_B J_\mA e_B D(t)=D(t)^*e_B J_\mA e_B D(t+i/2)    
% % \end{align}
% %  and is independent of t.

% % \item For all $t\in \mathbb{R}$ we have $\text{Ad}( D(t)D(i/2)^*)\mA\subset \mA$.

% % \end{enumerate}

% % \end{theorem}

% \subsection{Quantum ergodic hierarchy and subsystem flows}

% The modular flow of a state $\omega$ of $\mA$ might not be reducible on subalgebra $\mB$, meaning that $ \mB_t\subset \mA$ need not be the same $\mB$. However, as $\mB_t$ explores $\mA$ it might not get everywhere. In that case, we have a proper {\it entire history subalgebra}:
% \begin{align}
%     \vee_{t\in\mathbb{R}}\alpha_t(\mB)=\mB_{hist}
% \end{align}
% It is also possible that all $\mA_t$ keep a permanent memory of a set of operators. In that case, we have a proper {\it permanent memory subalgebra}
% \begin{align}
%    \wedge_{t\in\mathbb{R}}\alpha_t(\mB)=\mB_{mem} \ .
% \end{align}
% When these proper subalgebras exist, they are by construction reducible with respect to the dynamics.





% \begin{definition}[Generating inclusion]
% If dynamics gets everywhere in $\mA$ we say the inclusion is {\it generating} meaning 
% \begin{align}\label{geneating}
%     \vee_{t\in \mathbb{R}}\mB_t=\mA\ .
% \end{align}
% \end{definition}

% \begin{definition}
%     If the dynamics keeps no permanent memory in $\mB$ we say the inclusion is {\it forgetful} meaning that 
%     \begin{align}
%         \wedge_{t\in\mathbb{R}}\mB_t=\mathbb{C}1\ .
%     \end{align}
% \end{definition}

% The relative commutant of the total history of $\mathcal{B}$ in $\mathcal{A}$ is exactly the permanent memory of $\mathcal{B} $'s relative commutant:
% \begin{align}
% (\mathcal{B}_{hist})^c &= \left( \bigvee_{t \in \mathbb{R}} \alpha_t(\mathcal{B}) \right)' \cap \mathcal{A}= \left( \bigcap_{t \in \mathbb{R}} \alpha_t(\mathcal{B}') \right) \cap \mathcal{A}  \nn\\
% &= \bigcap_{t \in \mathbb{R}} \alpha_t(\mathcal{B}' \cap \mathcal{A})= \bigcap_{t \in \mathbb{R}} \alpha_t(\mathcal{B}^c) = (\mathcal{B}^c)_{mem}
% \end{align}
% where we have used the fact that the commutant of the supremum is the intersection of the commutants, and the fact that $\alpha_t$ is an automorphism of $\mathcal{A}$.

% The physical interpretation is as follows:
% If an observer has access to the entire dynamical history of the subsystem $\mathcal{B}$, the only local observables in $\mathcal{A}$ that remain hidden from them (i.e., that commute with their entire history) are those that constitute the permanent memory of the localized environment $\mathcal{B}^c$. Furthermore, if $\mathcal{A}$ is a type III factor and the inclusion is generating ($\mathcal{B}_{hist} = \mathcal{A}$), then $(\mathcal{B}_{hist})^c = \mathcal{A}' \cap \mathcal{A} = \mathbb{C}\mathbf{1}$. This mathematically forces the relative commutant's modular flow to be forgetful.

% On the other hand, the relative commutant of the memory algebra satisfies
% \begin{align}
% (\mathcal{B}_{mem})^c = \left( \bigcap_{t \in \mathbb{R}} \alpha_t(\mathcal{B}) \right)' \cap \mathcal{A}
% \end{align}
% Compare this to the total history of the relative commutant, which is $(\mathcal{B}^c)_{hist} = \bigvee_{t} \alpha_t(\mathcal{B}' \cap \mathcal{A})$. In the lattice of von Neumann algebras, the generated supremum of intersections is generally a subalgebra of the intersection with the supremum. Because $\alpha_t(\mathcal{B}' \cap \mathcal{A}) \subset \alpha_t(\mathcal{B}') \cap \mathcal{A}$, we have:
% \begin{align}
% (\mathcal{B}^c)_{hist} \subseteq (\mathcal{B}_{mem})^c    
% \end{align}
% The physical interpretation is as follows: The total accumulated history of the local environment $(\mathcal{B}^c)_{hist}$ is contained within the space of operators that commute with the permanent memory of $\mathcal{B}$. The operators that exist in $(\mathcal{B}_{mem})^c$ but not in $(\mathcal{B}^c)_{hist}$ represent non-local correlations that commute with the subsystem's absolute invariants but cannot be dynamically generated strictly from within the local environment $\mathcal{B}^c$ alone.

% \paragraph{Quantum error correction and conjugate inclusions.}
% In the operator-algebraic formulation of quantum error correction (QEC), the relative commutant $\mathcal{B}^c = \mathcal{B}' \cap \mathcal{A}$ serves as the complementary (or conjugate) subsystem to the primary inclusion $\mathcal{B} \subset \mathcal{A}$. 

% Assuming the existence of a state-preserving normal conditional expectation $\mathcal{E}: \mathcal{A} \to \mathcal{B}$, this map acts as the primary quantum erasure channel. It represents an observer restricted to the logical subsystem $\mathcal{B}$, having traced out the microscopic degrees of freedom in $\mathcal{A}$. The complementary noise channel, which tracks the information leaked into the environment, is given by the expectation onto the relative commutant, $\mathcal{E}^c: \mathcal{A} \to \mathcal{B}^c$. 

% The true open-system dynamics experienced by the restricted observer is described by a one-parameter family of completely positive, unital (CPU) maps $T_t = \mathcal{E} \circ \alpha_t|_\mathcal{B}$. Because the global modular flow $\alpha_t$ typically leaks operators out of $\mathcal{B}$, the map $T_t$ is generally non-unitary and non-multiplicative.

% Within this framework, the memory algebra $\mathcal{B}_{mem} = \bigcap_{t \in \mathbb{R}} \alpha_t(\mathcal{B})$ has a precise QEC interpretation: it is exactly the intersection of the multiplicative domains of the maps $T_t$ for all $t \in \mathbb{R}$. In the language of QEC, $\mathcal{B}_{mem}$ is the \textit{decoherence-free subalgebra} (or noiseless subsystem). It constitutes the maximal von Neumann subalgebra on which the dissipative dynamical channel $T_t$ acts purely reversibly and unitarily. Quantum information encoded in $\mathcal{B}_{mem}$ suffers zero decoherence or environmental leakage over infinite time.

% Conversely, the history algebra $\mathcal{B}_{hist} = \bigvee_{t \in \mathbb{R}} \alpha_t(\mathcal{B})$ acts as the \textit{dynamical support} or information hull of the code. This generated hull carries three rigorous physical interpretations:
% \begin{enumerate}
%     \item \textbf{Minimal reversible closure:} It is the smallest globally invariant von Neumann subalgebra containing the logical code. To perfectly invert the open-system flow and simulate the full unitary dynamics without information loss, $\mathcal{B}_{hist}$ defines the exact minimal algebra of observables an observer must be able to access.
%     \item \textbf{Reach of the noise:} It defines the absolute dynamical spread of the localized operators. Any operator in $\mathcal{A}$ that lies in the relative commutant of $\mathcal{B}_{hist}$ is completely orthogonal to the logical code space and can never be generated or reached by the subsystem dynamics.
%     \item \textbf{Complementary erasure:} From the perspective of the localized environment ($\mathcal{B}^c$), the dynamical expansion of the primary subsystem corresponds directly to the active erasure of the environment. 
% \end{enumerate}

% The conjugate to the primary subsystem dynamics is the evolution of the relative commutant under the exact same global modular flow, $\alpha_t(\mathcal{B}^c)$. This represents the dynamical evolution of the complementary noise channel. The relationship between these two flows is completely governed by a structural duality between history and memory:
% \begin{equation}
%     (\mathcal{B}_{hist})^c = (\mathcal{B}^c)_{mem}
% \end{equation}
% This dictates that the operators in $\mathcal{A}$ permanently hidden from an observer with access to the total history of $\mathcal{B}$ are precisely the absolute, non-dissipative invariants of the conjugate environment $\mathcal{B}^c$. Furthermore, taking the history of the environment yields the strict inclusion:
% \begin{equation}
%     (\mathcal{B}^c)_{hist} \subseteq (\mathcal{B}_{mem})^c
% \end{equation}
% The operators populating the gap $(\mathcal{B}_{mem})^c \setminus (\mathcal{B}^c)_{hist}$ represent non-local boundary correlations or emergent edge modes. They are valid observables in $\mathcal{A}$ that commute with the permanent memory of the subsystem, yet cannot be dynamically reconstructed strictly from within the localized environment $\mathcal{B}^c$ without access to the global commutant $\mathcal{A}'$.

% \paragraph{Decoupling theorem:} This history/memory duality perfectly reproduces the algebraic decoupling theorem of QEC. The decoupling theorem states that a quantum code is exactly correctable against a noise channel if and only if the complementary channel completely erases the logical information. Algebraically, if the primary subsystem's flow is fully generating ($\mathcal{B}_{hist} = \mathcal{A}$), the relative commutant duality explicitly forces $(\mathcal{B}^c)_{mem} = \mathbb{C}\mathbf{1}$. Thus, the exact dynamical recovery of the full algebra by the logical subsystem mathematically forces the conjugate environment to undergo complete asymptotic thermalization (a trivial tail). Information cannot be permanently preserved in the local environment if it is dynamically recovered by the primary code.

% \paragraph{The ergodic hierarchy and subsystem flows.}
% The structural duality between history and memory provides a natural operator-algebraic framework to classify the irreversibility of the dynamics, situating it within the quantum ergodic hierarchy. For a W*-dynamical system $(\mathcal{A}, \omega, \alpha_t)$, the global flow can be categorized by its degree of chaos and information loss:
% \begin{enumerate}
%     \item \textbf{Ergodicity:} The system is ergodic if there are no non-trivial global constants of motion. Algebraically, the fixed-point algebra of the flow is trivial: $\{A \in \mathcal{A} \mid \alpha_t(A) = A \ \forall t\} = \mathbb{C}\mathbf{1}$. While time averages converge to state averages, quantum information is not necessarily lost or scrambled, merely delocalized over time.
%     \item \textbf{Mixing:} A system is strongly mixing in the future if all dynamical correlations decay asymptotically: $\lim_{t \to \infty} \omega(A \alpha_t(B)) = \omega(A)\omega(B)$ for all $A, B \in \mathcal{A}$. It is strongly mixing in the past if the same holds for the limit $t\to -\infty$. While mixing guarantees that any localized excitation eventually spreads evenly across the system (meaning arbitrary states become locally indistinguishable from the vacuum state), it does not mandate a strictly monotonic flow of information.
%     \item \textbf{Quantum $K$-systems:} The pinnacle of the ergodic hierarchy is the noncommutative $K$-system, which mathematically enforces a strict, irreversible coarse-graining. A system is a $K$-system if there exists a distinguished subalgebra $\mathcal{B} \subset \mathcal{A}$ such that the flow is monotonically nested ($\alpha_t(\mathcal{B}) \subset \mathcal{B}$ for $t \ge 0$), fully generating ($\mathcal{B}_{hist} = \mathcal{A}$), and asymptotically forgetful ($\mathcal{B}_{mem} = \mathbb{C}\mathbf{1}$).
% \end{enumerate}

% The inclusion $\mathcal{B} \subset \mathcal{A}$ in a $K$-system perfectly realizes the QEC duality discussed above. The nesting condition ensures that the completely positive maps $T_t = \mathcal{E} \circ \alpha_t|_{\mathcal{B}}$ form a strictly contractive dynamical semigroup (a Markovian macroscopic flow) in the forward time direction. The generating condition ($\mathcal{B}_{hist} = \mathcal{A}$) guarantees that the subsystem flow eventually explores the entire global algebra. 

% Most crucially, the trivial tail condition ($\mathcal{B}_{mem} = \mathbb{C}\mathbf{1}$) ensures that the multiplicative domain of the long-time flow vanishes. Consequently, in a true quantum $K$-system, the relative commutant duality $(\mathcal{B}_{hist})^c = (\mathcal{B}^c)_{mem}$ physically dictates that the complementary environment $\mathcal{B}^c$ must undergo complete asymptotic thermalization. The strict irreversibility of a $K$-system is thus mathematically synonymous with the continuous, unrecoverable leakage of logical information into the conjugate environmental channels, leading to absolute memory loss at $t \to \infty$.

% \subsection{Characteristic flows}

% \paragraph{Common cyclic and separating vector:} Clearly, the action of subsystem flow on general operators in $\mA$ is non-unitary. It becomes unitary when the vector $\ket{\Omega}$ is a common cyclic and separating vector. In that case, the code subspace is the entire Hilbert space, i.e., $e_B=1$, and the standard cocycle identity holds; i.e., $\mM(t,s)=0$. In the AQFT literature, such subsystem flows are called {\it characteristic functions}. In fact, the inclusion has a common cyclic and separating vector if and only if the subsystem flow $D(t)$ satisfies properties $1\to 5$ with $e_B=1$, and $\mM(t,s)=0$ (See theorem ... in \cite{} for a proof). 

% We are now ready to introduce the characteristic function of the inclusion $\mathcal{B}\subset\mathcal{A}$\cite{revolutionizing}:
% \begin{definition}
%     Given a pair of von Neumann algebras $\mathcal{B}\subset\mathcal{A}$ and a common cyclic and separating vector $\ket{\Omega}\in\mathcal{H}$, the characteristic function is defined as:
%     \begin{equation}
%         \mathcal{D}_{\mathcal{A}, \mathcal{B}}(t) := \Delta_\mathcal{A}^{it}\Delta_\mathcal{B}^{-it}
%     \end{equation}
% \end{definition}
% We have already seen that the characteristic function $\mathcal{D}_{\mathcal{A}, \mathcal{B}}(t)$ can be extended to the imaginary values: $\mathcal{D}_{\mathcal{A}, \mathcal{B}}(-i\alpha)$ where $0\leq \alpha\leq 1/2$. The crucial fact is that the characteristic function can be extended to the entire strip $S(0, -1/2) := \{z\in\mathbb{C}: 0\geq \Im{z}\geq -1/2\}$ \cite{revolutionizing}:
% \begin{lemma}
%     Given a pair of von Neumann algebra $\mathcal{B}\subset\mathcal{A}$ with a common cyclic and separating vector $\ket{\Omega}$, then the characteristic function $\mathcal{D}_{\mathcal{A}, \mathcal{B}}(t)$ can be extended to the strip $\mathcal{S}(0, -1/2)$. The extension is a bounded analytic function on $\mathcal{S}(0,-1/2)$.

%     Moreover, on the boundary $\partial\mathcal{S}(0,-1/2) = \{z\in\mathbb{C}:\Im{z} = 0\}\sqcup\{z\in\mathbb{C}:\Im{z} = -1/2\}$, the characteristic function is strongly continuous.
% \end{lemma}
% For a full discussion of the characteristic function, we refer the reader to \cite{revolutionizing}. Here we list the properties of the characteristic function \cite{revolutionizing}:
% \begin{enumerate}
%     \item $\mathcal{D}_{\mathcal{A},\mathcal{B}}(t)$ and $\mathcal{D}_{\mathcal{A}, \mathcal{B}}(t - \frac{i}{2})$ are unitary and strongly continuous in $t$;
%     \item $\mathcal{D}_{\mathcal{A},\mathcal{B}}(z)$ is bounded analytic on the strip $\mathcal{S}(0,-\frac{1}{2})$;
%     \item $\mathcal{D}_{\mathcal{A},\mathcal{B}}(t)\ket{\Omega} = \ket{\Omega}$;
%     \item The cocycle condition holds: $\mathcal{D}_{\mathcal{A},\mathcal{B}}(t+s) = \sigma_t^{\ket{\Omega}}(\mathcal{D}_{\mathcal{A},\mathcal{B}}(s))\mathcal{D}_{\mathcal{A},\mathcal{B}}(t)$ where $\sigma_t^{\ket{\Omega}}$ is the modular automorphism on $\mathcal{A}$ associated with the state $\ket{\Omega}$;
%     \item $\mathcal{D}_{\mathcal{A},\mathcal{B}}(t - \frac{i}{2}) = J_\mathcal{A}\mathcal{D}_{\mathcal{A},\mathcal{B}}(t)J_\mathcal{B}$. In particular, $\mathcal{D}_{\mathcal{A},\mathcal{B}}(-\frac{i}{2}) = J_\mathcal{A}J_\mathcal{B}$. This unitary operator is crucial in subfactor theory creating a chain of algebras:
%     \begin{equation}
%         \mathcal{A}\supset\mathcal{B}\supset\mathcal{D}_{\mathcal{A},\mathcal{B}}(-\frac{i}{2})^*\mathcal{A}\mathcal{D}_{\mathcal{A},\mathcal{B}}(-\frac{i}{2})\supset\mathcal{D}_{\mathcal{A},\mathcal{B}}(-\frac{i}{2})^*\mathcal{B}\mathcal{D}_{\mathcal{A},\mathcal{B}}(-\frac{i}{2})\supset...
%     \end{equation}
%     \item For all $t\in\mathbb{R}$, $\mathcal{D}_{\mathcal{A},\mathcal{B}}(t)\mathcal{D}_{\mathcal{A},\mathcal{B}}(-\frac{i}{2})^*\mathcal{A}\mathcal{D}_{\mathcal{A},\mathcal{B}}(-\frac{i}{2})\mathcal{D}_{\mathcal{A},\mathcal{B}}(t)^*\subset\mathcal{A}$ holds.
% \end{enumerate}
% Any operator-valued function associated with a pair of von Neumann algebras that satisfy the above-listed properties is called a characteristic function \cite{revolutionizing}. The following theorem verifies the intuition that characteristic function fully captures the structure of inclusions of von Neumann algebras with common cyclic and separating vectors \cite{revolutionizing}:
% \begin{theorem}
%     Given a von Neumann algebra $\mathcal{A}$ and a cyclic and separating vector $\ket{\Omega}$, then each subalgebra $\mathcal{B}\subset\mathcal{A}$ such that $\ket{\Omega}$ is cyclic for $\mathcal{B}$ is uniquely associated with a characteristic function.
% \end{theorem}
% Again for details, we refer the readers to \cite{revolutionizing}.

% From the point of view of ergodic theory, one should note that an inclusion of algebras $\mB\subset \mA$ with $\omega$ a common cyclic and separating vector is generating. Such inclusions can occur in infinite-dimensional algebras of type $I$, $II$, or type $III$. However, as we mentioned earlier, since we want the modular flow of states to be ergodic, we focus on type III$_1$ subfactors with a common cyclic and separating vector. Such examples naturally appear in QFT as subsystem inclusions.

% In infinite dimensions, the inclusion of algebras can be more exotic:
% \begin{enumerate}
%     \item {\bf Singular inclusions:} An non-trivial inclusion $\mA\subset \mB$ is singular if the relative commutant is trivial. Such inclusions cannot occur in type $I$ algebras. A physical example of singular inclusions in QFT is the inclusion of the type $III_1$ algebra of wedges with a null translation in vacuum QFT of an interacting theory in $d>2$ spacetime dimensions.
    
%     \item {\bf No uncorrelated states of $B$ and $R$:} We could also have setups where there is a non-trivial relative commutant $\mA_R$ and $\mA_R\vee \mA_B=\mA_A$; however, there is no tensor product factorization $\mA_A=\mA_B\otimes \mA_R$. The surprising statement is that for these inclusions there cannot exist any normal conditional expectations $\mE:\mA_A\to \mA_B$, even though there is a non-trivial relative commutant. This adds the following nuance to our previous statement:
%     {\it If there exists at least one normal conditional expectation $\mE:\mA_A\to \mA_B$ for subfactors $\mA_B\subset \mA_A$ then there is a one-to-one correspondence between such conditional expectations and the states of the relative commutant.} Then, here is how the logic works: if there exists a normal conditional expectation, it must preserve an uncorrelated state  $\omega_R\otimes \omega_B$, but this cannot exist because there is no tensor product decomposition \footnote{Even though there are no conditional expectations, we still say such inclusions have infinite index.}.
%     Such inclusions appear naturally, where $B$ and $R$ do not touch. Type I and II$_1$ factors always admit normal conditional expectations; however, inclusions with no conditional expectations are commonplace in QFT. They occurs in every time the regions associated to $B$ and $R$ touch. 

%     \item {\bf No uncorrelated states of $B$ and $A'$:} Given the inclusion $\mA_B\subset \mA_A$, we can always consider $\mA_B\subset \mA_B\vee \mA_{A'}$. Now, if there is no tensor factorization between $B$ and $R$, it is possible that the same holds for the second inclusion. This occurs when there is no tensor product between $B$ and $A'$ either. When this occurs, we say the inclusion $\mA_B\subset \mA_A$ does not {\it split}. This occurs when $B$ touches both $R$ and $A'$. Finally, there is the case where $B$ touches $R$, but not $A'$ (split case). 
% \end{enumerate}


% The examples we study in this work include cases in QFT where all cases above occur.
 

% \NL{Note that irreducible dynamics implies that any non-trivial inclusion is generating \footnote{To see this, note that the dynamics is reducible with respect to $\vee_{t\in \mathbb{R}}\mB_t$.}.} \NL{What is the ergodic theory interpretation of irreducibility?}

% \NL{Add paragraphs one for each notion of index of an inclusion, normal and conormal inclusions}


% We illustrate the features above by two instructive examples:

% \paragraph{Araki-Powers inclusions (non-ergodic inclusion of type III$_1$ algebras):}

% Let $\mA\subset \mB$ be a \emph{finite-index} inclusion of type~$III_1$
% factors, meaning that the relative commutant $\mathcal{R}=\mA'\cap \mB$
% is a finite-dimensional matrix algebra $M_n(\mathbb{C})$.  By the Jones--Kosaki index
% theory, the inclusion has index $[\mathcal{A}:\mathcal{B}]=n_R^2<\infty$.  Since $\mathcal{R}$
% is finite-dimensional, it is type~$I$, so the inclusion is normal (conormal) in the sense
% that $\mathcal{B}\vee\mathcal{R}=\mathcal{A}$.  By Lemma~\ref{lem:typeIII-standard} and the
% correspondence~\eqref{eq:CE-state-correspondence}, normal conditional expectations
% $\mE:\mA\to \mB$ exist and are parametrised by states on $M_n(\mathbb{C})$.
% The coarse-graining flow $U_t=\sigma_t^{\omega_\mathcal{A}}\circ(\sigma_t^{\omega_\mathcal{B}})^{-1}$
% acts non-trivially only on $\mathcal{R}=M_n(\mathbb{C})$ and preserves it; it therefore
% does \emph{not} satisfy {\it generating property} in (\ref{geneating}) and
% is \emph{not mixing}.  This inclusion is non-ergodic, because even though the algebras are III$_1$ factors, and the action of $\Delta_{\omega;A}^{it}$ on $\mB$ does not fix any operators, there are operators in the relative commutant $R$ that are in the centralizer of $\omega_A$: namely the eigenvectors of the density matrix $\omega_R$.

% \paragraph{Local algebras of QFT (no normal conditional expectations):}
% In a local conformal net $\{\mathcal{A}(I)\}$ on the circle, let $I_1,I_2$ be disjoint
% open intervals with closures sharing exactly one endpoint, and set $I=\mathrm{int}(I_1\cup I_2)$.
% Haag duality and strong additivity give $\mathcal{A}(I_1)'\cap\mathcal{A}(I)=\mathcal{A}(I_2)$
% and $\mathcal{A}(I_1)\vee\mathcal{A}(I_2)=\mathcal{A}(I)$, so the inclusion
% $\mathcal{B}:=\mathcal{A}(I_1)\subset\mathcal{M}:=\mathcal{A}(I)$ is conormal with
% $\mathcal{R}=\mathcal{A}(I_2)$.  By Theorem~\ref{thm:split-no-CE} and the fact that
% adjacent touching algebras of a non-trivial conformal net never form a split pair
% (Wiesbrock~\cite{Wiesbrock1993}), there is \emph{no} normal conditional expectation
% $\mE:\mathcal{A}(I)\to\mathcal{A}(I_1)$.  If $\mE$ existed, Proposition~\ref{prop:dense-common}
% would give $\mathcal{A}(I)\cong\mathcal{A}(I_1)\otimes\mathcal{A}(I_2)$, contradicting the
% splitting obstruction.  The inclusion is therefore in the ``unexpected'' class: the index
% is infinite, no QEC is possible, and every state on $\mathcal{A}(I)$ is universally
% entangled between $\mathcal{A}(I_1)$ and $\mathcal{A}(I_2)$.


% \paragraph{General von Neumann algebras:}
% \label{sec:conormality-expectations} Here, we provide rigorous results that generalize the discussion of cases above to general von Neumann algebras.

% Let $\mA\subset\mB\subset B(\mathcal H)$ be an inclusion of
% $\sigma$-finite type $\mathrm{III}_1$ factors.  The inclusion is
% \emph{standard} with respect to $\Omega\in\mathcal H$ if $\Omega$ is cyclic
% and separating for $\mA$, $\mB$ and the relative commutant $\mB^c=\mB'\cap \mA$. It is \emph{split} if a type I factor $\mathcal F$ exists with
% $\mB\subset\mathcal F\subset\mA$, and \emph{conormal} if
% \begin{equation}
%  \mB\vee\mB^c=\mA.
% \end{equation}
% Throughout this section, a conditional expectation
% $\mE:\mathcal \mA\to\mathcal \mB$ is assumed to be \emph{faithful and normal},
% unless the contrary is explicitly stated.  Tomiyama's theorem \cite{Tomiyama1957,Tomiyama1959,Takesaki1972} gives
% \begin{equation}
%  \mE(b_1 a b_2)=b_1\mE(a) b_2,
%  \qquad b_1,b_2\in \mB,\quad a\in \mA.
% \end{equation}

% %\subsection{The split case}



% \begin{theorem}[Sufficient condition for no uncorrelated state]\label{thm:split-no-expectation}
% If $\mB$ is a type III factor and
% $\mA\subset \mB$ is split, then no normal conditional expectation
% $\mE:\mA\subset \mB$ exists.  
% \end{theorem}

% \begin{proof}
% Let $\mB\subset\mathcal F\subset\mA$ be an interpolating type I
% factor and suppose that $\mE:\mA\to \mB$ is normal.  Its restriction
% \begin{equation}
%  \mE_{\mathcal F}:=\mE|_{\mathcal F}:\mathcal F\longrightarrow\mB
% \end{equation}
% is a normal norm-one projection with range $\mathcal B$.  Tomiyama's type
% theorem states that a normally expected von Neumann subalgebra of a type I von
% Neumann algebra is type I \cite{Tomiyama1959}\cite{Takesaki1972}. This contradicts
% with $\mB$ being a type III factor.
% \end{proof}

% Note that in the theorem above, faithfulness and preservation of the
% distinguished vector state are not needed, and 
% standardness plays no role in this implication.

% \begin{lemma}\label{lem:expectation-slice}
% If $\mathcal B$ is a factor and $\mE:\mathcal A\to\mathcal B$ is a faithful
% normal conditional expectation, there is a faithful normal state $\chi_R$ on
% the relative commutant $\mB^c$ such that
% \begin{equation}
%  \mE(c)=\chi_R(c)\mathbbm 1,
%  \qquad \mE(b c)=\chi_R(c)b,
%  \qquad b\in\mathcal B,\quad c\in\mathcal R.
% \end{equation}
% \end{lemma}

% \begin{proof}
% For $b\in\mathcal B$ and $c\in\mathcal R$,
% \begin{equation}
%  b \mE(c)=\mE(b c)=\mE(c b)=\mE(c)b.
% \end{equation}
% Thus $\mE(c)\in Z(\mathcal B)=\mathbb R\mathbbm 1$, so
% $\mE(c)=\chi_R(c)\mathbbm 1$ for a normal state $\chi$.  If $c\geq0$ and
% $\chi_R(c)=0$, faithfulness of $\mE$ implies $c=0$; hence $\chi_R$ is faithful.  The
% second identity is the bimodule property.
% \end{proof}

% For a faithful normal state $\varphi$ on $\mathcal B$, the state
% $\varphi_0:=\varphi\circ \mE$ is faithful and normal and satisfies
% \begin{equation}
%  \varphi_0(b c)=\varphi_0(b)\chi_R(c).
% \end{equation}

% \begin{proposition}[conormality and tensor product]\label{prop:conormal-factorization}
% Under the hypotheses of Lemma~\ref{lem:expectation-slice}, multiplication
% extends to a normal $*$-isomorphism
% \begin{equation}
%  m:\mathcal B\overline\otimes\mathcal R
%  \overset{\cong}{\longrightarrow}\mathcal B\vee\mathcal R,
%  \qquad m(b\otimes c)=b c.
% \end{equation}
% If the inclusion is conormal, then
% \begin{equation}\label{eq:conormal-factorization}
% \mathcal A\cong
%  \mathcal B\overline\otimes \mB^c .
% \end{equation}
% \end{proposition}

% \begin{proof}
% Multiplication is a $*$-homomorphism on $\mathcal B\odot\mathcal R$ because
% the two algebras commute.  Moreover,
% \begin{equation}
%  \varphi_0\bigl(m(b\otimes c)\bigr)
%  =(\varphi\otimes\chi_R)(b\otimes c).
% \end{equation}
% The GNS representation of the faithful normal product state
% $\varphi\otimes\chi_R$ is the spatial tensor-product representation.  Hence
% multiplication extends normally to $\mathcal B\overline\otimes\mathcal R$.
% Faithfulness of the product state gives injectivity, while surjectivity follows
% from the definition of $\mathcal B\vee\mathcal R$.  Conormality then yields
% \eqref{eq:conormal-factorization}; compare
% \cite{Faulkner2020,FaulknerSperanza2024}.
% \end{proof}

% We define the Pimsner-Popa index of a conditional expectation to be the smallest positive number $\lambda\geq 1$ sucht that for all positive operators $x_+\in \mA$ we have 
% \begin{align}
%     \mE(x_+)\geq \lambda^{-1}x_+
% \end{align}
% Then, the index of an inclusion is defined to be 
% \begin{align}
%     [\mA:\mB]_0=\inf_\mE \text{Ind}(\mE)\ .
% \end{align}
% The index is a measure of the relative size of $\mA$ to $\mB$. When the relative commutant is infinite dimensional, or there are no normal conditional expectations, we say the index is infinite. Finite index inclusions always have normal conditional expectations.

% % If the relative commutant has a finite trace (as in type $I_d$ and $II_1$

% % There is no contradiction at this abstract level.  Tensor products of type
% % $\mathrm{III}_1$ factors can again be type $\mathrm{III}_1$; for example,
% % $R_\infty\overline\otimes R_\infty\cong R_\infty$.  The contradiction below
% % is specifically geometric: it uses the realization of $\mathcal N$ and its
% % relative commutant as local algebras of adjacent AQFT regions.

% % \subsection{The AQFT obstruction for adjacent regions}

% % Let $\mathcal A$ be a nontrivial local conformal net, and let $I_1,I_2$ be
% % disjoint intervals whose closures meet at one endpoint.  If $I$ is the interior
% % of $\overline{I_1}\cup\overline{I_2}$, Haag duality and strong additivity give
% % \begin{equation}\label{eq:adjacent-geometry}
% %  \mathcal A(I_1)'\cap\mathcal A(I)=\mathcal A(I_2),
% %  \qquad
% %  \mathcal A(I_1)\vee\mathcal A(I_2)=\mathcal A(I).
% % \end{equation}
% % Thus $\mathcal N:=\mathcal A(I_1)\subset\mathcal M:=\mathcal A(I)$ is
% % conormal and $\mathcal N'\cap\mathcal M=\mathcal A(I_2)$.

% % \begin{theorem}\label{thm:adjacent-no-expectation}
% % In this adjacent-region configuration, no faithful normal conditional
% % expectation
% % \begin{equation}
% %  E:\mathcal A(I)\longrightarrow\mathcal A(I_1)
% % \end{equation}
% % exists.
% % \end{theorem}

% % \begin{proof}
% % If $E$ existed, Proposition~\ref{prop:conormal-factorization} and
% % \eqref{eq:adjacent-geometry} would give
% % \begin{equation}\label{eq:forbidden-adjacent-factorization}
% %  \mathcal A(I)\cong
% %  \mathcal A(I_1)\overline\otimes\mathcal A(I_2),
% %  \qquad a_1\otimes a_2\longmapsto a_1a_2.
% % \end{equation}
% % For commuting factors, spatiality of the multiplication map is equivalent to
% % the split property \cite{DoplicherLongo1984}.  Hence the adjacent algebras
% % $\mathcal A(I_1)$ and $\mathcal A(I_2)$ would form a split pair.

% % This is impossible for a nontrivial chiral conformal net.
% % Reference~\cite[Proposition~1.6.12]{GiorgettiThesis} proves explicitly that
% % local algebras of disjoint touching intervals never satisfy the split property
% % unless the net is trivial.  Its first proof reduces this to
% % \cite[Lemma~13]{Wiesbrock1993}, which shows that a standard inclusion of type
% % $\mathrm{III}_1$ factors with type $\mathrm{III}_1$ relative commutant cannot
% % be both half-sided modular and split.  Its second proof uses the canonical
% % intermediate type I factor of
% % \cite[Theorems~2.1 and 4.1]{DoplicherLongo1984} and geometric modular action.
% % Thus \eqref{eq:forbidden-adjacent-factorization} contradicts the
% % adjacent-region structure.
% % \end{proof}

% % The logical chain is
% % \begin{equation}
% % \begin{split}
% %  &\text{faithful normal expectation}+\text{conormality}
% %  \Longrightarrow
% %  \mathcal M\cong\mathcal N\overline\otimes
% %  (\mathcal N'\cap\mathcal M)\\
% %  &\Longrightarrow
% %  \text{$(\mathcal N,\mathcal N'\cap\mathcal M)$ is a split commuting pair}\\
% %  &\Longrightarrow
% %  \text{contradiction for adjacent AQFT regions}.
% % \end{split}
% % \end{equation}
% % The last implication is an AQFT input, not a type-classification argument.
% % Rehren describes the split property as maximal decoupling for intervals with
% % disjoint closures and strong additivity as maximal interaction for touching
% % intervals \cite{Rehren2015}.  The failure of split decoupling when chiral
% % intervals touch is also stated in \cite{LongoRehren2004}.



% \subsection{Geometry from modular K-systems}
% \label{subsec:geometry-K-systems}
 
% In this Section, we have introduced the quantum ergodic hierarchy
% for a general quantum dynamical system ($W^*$-dynamical system) $(\mA,\omega,\alpha_t)$: a subalgebra
% $\mB\subset\mA$ defines a K-system when its forward orbit is nested
% ($\alpha_t(\mB)\subset\mB$ for $t\ge0$), generating ($\mB_{\rm his}=\mA$), and
% forgetful ($\mB_{\rm mem}=\mathbb{C}\id$). The classical precursor of this
% construction, recalled in Appendix~A.3, already anticipated the quantum theory: the
% one-sided inclusion $\alpha_t(\mB_+)\subseteq\mB_+$, $t\ge0$, of~(A.35)--(A.36) was
% built from an \emph{externally specified} measure-preserving flow $\alpha_t$, with no
% relation to any state-dependent construction. In the quantum setting, however, there is
% a distinguished choice of dynamics attached to every state: the modular flow itself. The
% purpose of this subsection is to specialize the K-system construction to this
% distinguished choice, and to show that the resulting \emph{modular} K-systems are far
% more rigid than generic K-systems: the single nesting condition already forces the
% generating and forgetful conditions, forces maximal chaos with a universal rate, and,
% when strengthened to a standard inclusion, forces an emergent continuous conformal
% symmetry. These results are due to Borchers~\cite{Borchers2000} and
% Wiesbrock~\cite{Wiesbrock1993}; their formulation in the language of quantum ergodic
% theory and quantum K-systems that we follow here is due to~\cite{OFLLM2024}, with the
% sharpened geometric statement in the last part of this subsection due
% to~\cite{LLMO2025}.

% \subsubsection*{Half-sided modular inclusions}
 
% \begin{definition}[Half-sided modular inclusion]\label{def:HSMI}
% Let $\mB\subset\mA$ be a proper inclusion of von Neumann factors with a common cyclic and
% separating vector $|\Omega\rangle$. We say the inclusion is a \emph{future half-sided
% modular inclusion} (HSMI$+$), or that $\mB$ is a \emph{modular future subalgebra} of
% $\mA$, if
% \begin{equation}
%   \forall t\ge0:\qquad \Delta_{\omega;\mA}^{-it}\,\mB\,\Delta_{\omega;\mA}^{it}\subset\mB.
%   \label{eq:HSMI}
% \end{equation}
% HSMI$-$ (a \emph{modular past subalgebra}) is the same condition with $t\le0$.
% \end{definition}
 
% \noindent
% The sign in~\eqref{eq:HSMI} is conventional but not arbitrary: it is fixed so that
% modular time runs in the same direction as physical time in the K-system sense
% of Section~\ref{subsec:quantum-ergodic-hierarchy}. Concretely, setting
% $\alpha_t := \sigma_{-t}^{\omega;\mA}$ identifies~\eqref{eq:HSMI} exactly with the
% abstract nesting condition $\alpha_t(\mB)\subset\mB$, $t\ge0$, required of a quantum
% K-system; a modular K-system is then the special case of
% Section~\ref{subsec:quantum-ergodic-hierarchy}'s general construction in which the
% external dynamics $\alpha_t$ coincides with (minus) the algebra's own modular flow. This
% is the precise quantum sharpening of the classical anticipation in~(A.35): there, the
% one-sided inclusion was imposed on an arbitrary $\alpha_t$; here, it is a rigid algebraic
% condition on the state $\omega$ and the subalgebra $\mB$ alone, with no external input.
 
% \begin{theorem}[HSMI theorem; Borchers, Wiesbrock~{\cite{Borchers2000,Wiesbrock1993}}, see also~{\cite[Thm.~16]{OFLLM2024}}]
% \label{thm:HSMI}
% Suppose $\mB\subset\mA$ is HSMI$\pm$. Then the operator
% \begin{equation}
%   G := \frac{1}{2\pi}\bigl(K_{\omega;\mA} - K_{\omega;\mB}\bigr)
%   \label{eq:HSMI-generator}
% \end{equation}
% is positive (resp. negative) semi-definite for HSMI$+$ (resp. HSMI$-$), and generates a
% unitary one-parameter group $U(s) := e^{isG}$ satisfying, for all $s,t\in\mathbb{R}$,
% \begin{align}
%   \frac{1}{2\pi}\,[\,K_{\omega;\mA},\,G\,] &= -iG, \label{eq:HSMI-commutator}\\
%   \Delta_{\omega;\mA}^{-it}\,U(s)\,\Delta_{\omega;\mA}^{it}
%   = \Delta_{\omega;\mB}^{-it}\,U(s)\,\Delta_{\omega;\mB}^{it} &= U(e^{2\pi t}s),
%   \label{eq:HSMI-chaos}\\
%   J_{\omega;\mA}\,U(s)\,J_{\omega;\mA} = J_{\omega;\mB}\,U(s)\,J_{\omega;\mB} &= U(-s),
%   \label{eq:HSMI-detailed-balance}\\
%   \mB &= U(1)\,\mA\,U(-1). \label{eq:HSMI-translation}
% \end{align}
% \end{theorem}
 
% \noindent
% Equation~\eqref{eq:HSMI-chaos} is the \emph{maximal modular chaos} relation: viewed under
% conjugation by modular time, the one-parameter group $U(s)$ dilates at the universal rate
% $s\mapsto e^{2\pi t}s$. We will see in Section~3.3 that this is precisely the rate that
% saturates the modular chaos bound obeyed by the matrix elements of the subsystem flow;
% the HSMI theorem is the abstract, state-independent statement of why that bound is
% saturated whenever the inclusion is a half-sided modular one. Equation
% \eqref{eq:HSMI-detailed-balance} is the quantum detailed balance condition: it identifies
% $G$ as odd under the modular conjugations of both $\mA$ and $\mB$, the operator-algebraic
% counterpart of the classical detailed-balance condition of Boltzmann for a dynamical
% matrix $T_{ij}$ satisfying $e^{-\beta E_i}T_{ij}=e^{-\beta E_j}T_{ji}$.
 
% The following proposition identifies $G$ with the generator of the subsystem flow: 
% \begin{proposition}\label{prop:G-is-subsystem-flow-generator}
% If $\mB\subset\mA$ is HSMI$+$, the common cyclic and separating vector gives
% $e_\mB=\id$, and the generator~\eqref{eq:HSMI-generator} coincides, up to the
% normalization $2\pi$, with the generator of the subsystem flow of
% Definition~2.1 evaluated at $t=0$:
% \begin{equation}
%   2\pi G = i\,\frac{d}{dt}D_{\mA\subset\mB;\omega}(t)\Big|_{t=0}
%          = K_{\omega;\mA}-K_{\omega;\mB}.
%   \label{eq:G-equals-subsystem-generator}
% \end{equation}
% \end{proposition}
% \begin{proof}
% Since $|\Omega\rangle$ is cyclic for $\mB$, $\mB|\Omega\rangle$ is dense in $\mH$ and the
% Jones projection of Definition~2.1 is $e_\mB=\id$. The claim is then immediate from
% comparing~\eqref{eq:HSMI-generator} with~(2.29).
% \end{proof}
 
% \noindent
% Proposition~\ref{prop:G-is-subsystem-flow-generator} says that the half-sided modular
% inclusion condition is exactly the statement that the subsystem-flow generator of
% Section~\ref{subsec:quantum-ergodic-hierarchy} is of definite sign. In general, this
% generator measures the residual, non-unitary leakage of the coarse-grained description;
% Theorem~\ref{thm:HSMI} shows that when the inclusion is half-sided modular, this leakage
% organizes itself into a genuine positive-energy unitary representation of the translation
% group $\mathbb{R}$, rather than a generic completely positive semigroup.
 
% \begin{theorem}[Half-sided translations; Borchers~{\cite{Borchers2000}}, see also~{\cite[Thm.~18]{OFLLM2024}}]
% \label{thm:half-sided-translations}
% Let $U(t)=e^{iGt}$ be a unitary flow with $U(t)|\Omega\rangle=|\Omega\rangle$, and
% suppose $\mB_t := U(t)\mB U(t)^\dagger\subset\mB$ for all $t>0$. Then the following are
% equivalent:
% \begin{enumerate}
%   \item The generator is positive: $G\ge0$.
%   \item Borchers' commutation relations hold: $\Delta_{\omega;\mA}^{-is}U(t)\Delta_{\omega;\mA}^{is}=U(e^{2\pi s}t)$.
%   \item $\mB_t$ is a modular future subalgebra of $\mB$ for every $t$.
% \end{enumerate}
% \end{theorem}
 
% \noindent
% Theorem~\ref{thm:half-sided-translations} is the converse direction to
% Theorem~\ref{thm:HSMI}: it says that \emph{half-sided translations} (a positive-energy
% unitary flow nesting a subalgebra) and \emph{half-sided modular inclusions} are two
% descriptions of the same phenomenon. This equivalence is the operator-algebraic avatar
% of the physical expectation that a positive-energy flow and a thermal (KMS/modular) flow
% of the same system are compatible only through a very restrictive algebraic relation —
% precisely Borchers' commutation relations~\eqref{eq:HSMI-chaos}.
 
% % ======================================================================
% \subsubsection*{Modular K-systems are automatically maximally chaotic:}
 
% We can now return to the ergodic hierarchy of Section~\ref{subsec:quantum-ergodic-hierarchy}
% and show that it collapses for modular dynamics: unlike a generic K-system, for which
% nesting, generating, and forgetfulness are three logically independent conditions on
% $\mB\subset\mA$, a half-sided modular inclusion automatically satisfies all three.
 
% \begin{corollary}[Maximal modular chaos $\equiv$ modular K-system; {\cite[Cor.~17]{OFLLM2024}}]
% \label{cor:maximal-chaos-K-system}
% Let $\mB\subset\mA$ be HSMI$+$ with respect to a faithful normal state $\omega$ on the
% type~$III_1$ factor $\mA$. Then, with $\alpha_t:=\sigma_{-t}^{\omega;\mA}$, the inclusion
% is a quantum K-system in the sense of Section~\ref{subsec:quantum-ergodic-hierarchy}:
% $\mB_{\rm his}=\mA$ and $\mB_{\rm mem}=\mathbb{C}\id$.
% \end{corollary}
 
% \begin{proof}[Proof sketch]
% By~\eqref{eq:HSMI-translation}, $\alpha_{-1}(\mB)=U(1)\mB U(-1)^{-1}\cdots$; more directly,
% iterating~\eqref{eq:HSMI-chaos} shows that the family $\mB_t=\alpha_t(\mB)$ dilates
% exponentially under modular conjugation at rate $2\pi$, so for any operator
% $a\in\mA$ supported away from the fixed point of this dilation there is a $t$ with
% $a\in\mB_t$; this gives $\bigvee_t\mB_t=\mA$. For the tail, suppose
% $a\in\bigcap_{t\ge0}\mB_t$. Equation~\eqref{eq:HSMI-commutator} shows that $K_{\omega;\mA}$
% scales $G$ by $e^{2\pi t}$ under $t$-translation, so the spectral support of any operator
% fixed under arbitrarily large dilations of a positive generator collapses to the
% spectral projector at the origin, forcing $a$ to be a multiple of the identity because
% $\mA$ is a factor. See~\cite{OFLLM2024} for the complete argument.
% \end{proof}
 
% \begin{remark}
% Corollary~\ref{cor:maximal-chaos-K-system} is the operator-algebraic reason that the
% inclusions of wedges and double cones studied in Section~3 are automatically both
% generating and forgetful: once a subregion inclusion is recognized as a half-sided
% modular inclusion of the vacuum modular flow, these two ergodic properties are not
% additional physical input but algebraic consequences of positivity of the generator
% $G$. The maximal rate $\lambda=2\pi$ in~\eqref{eq:HSMI-chaos} is exactly the bound that
% the matrix elements of the characteristic function are shown to saturate in
% Section~3.3; the HSMI theorem is the reason a modular K-system saturates the modular
% chaos bound identically, rather than merely satisfying it.
% \end{remark}
 
% % ======================================================================
% \subsubsection*{Standard inclusions and the emergent conformal group}
 
% The HSMI theorem already produces two of the three generators of $\mathfrak{sl}(2,\mathbb{R})$:
% the modular Hamiltonian $K_{\omega;\mA}$ itself generates dilations, and $G$ generates
% positive-energy translations. The remaining question is whether a third generator closes
% these into the full conformal algebra of the line, $\psl$. This requires strengthening
% HSMI to a \emph{standard inclusion} in the sense already introduced in
% Section~\ref{subsec:quantum-ergodic-hierarchy}: $\Omega$ is cyclic and separating not
% only for $\mA$ and $\mB$ but also for the relative commutant $\mB^c=\mB'\cap\mA$.
 
% \begin{theorem}[Standard inclusion implies $\PSL$; Wiesbrock~{\cite{Wiesbrock1993}}, see also~{\cite[Thm.~29, $U=\id$]{LLMO2025}}]
% \label{thm:wiesbrock-PSL}
% Let $\mB\subset\mA$ be a standard HSMI$+$ inclusion with common cyclic and separating
% vector $|\Omega\rangle$ for $\mA,\mB,\mB^c$. Define
% \begin{equation}
%   \widehat{\mB} := \mB \vee J_{\omega;\mA}\,\mB\,J_{\omega;\mA},
%   \label{eq:B-hat}
% \end{equation}
% and
% \begin{equation}
%   G_0 := \frac{1}{2\pi}\log\Delta_{\omega;\mA}, \qquad
%   G_1 := \frac{1}{2\pi}\bigl(\log\Delta_{\omega;\mB}-\log\Delta_{\omega;\mA}\bigr), \qquad
%   G_2 := \frac{1}{2\pi}\bigl(\log\Delta_{\omega;\widehat{\mB}}-\log\Delta_{\omega;\mA}\bigr).
%   \label{eq:three-generators}
% \end{equation}
% (Note that $G_1=G$, the HSMI generator of~\eqref{eq:HSMI-generator}.) Then
% \begin{equation}
%   [G_0,G_1]=iG_1,\qquad [G_0,G_2]=-iG_2,\qquad [G_1,G_2]=-2iG_0,
%   \label{eq:psl2R-algebra}
% \end{equation}
% and $G_0,G_1,G_2$ exponentiate to a unitary representation of the universal cover
% $\PSLt$ of the conformal group of the line, with inversion map
% \begin{equation}
%   I = J_{\omega;\mA}\,J_{\omega;\widehat\mB}.
%   \label{eq:inversion}
% \end{equation}
% \end{theorem}
 
% \begin{proof}[Proof sketch]
% By Theorem~\ref{thm:HSMI}, $\mB\subset\mA$ HSMI$+$ already furnishes $G_0,G_1$ with
% $[G_0,G_1]=iG_1$, this being~\eqref{eq:HSMI-commutator}. Standardness of the inclusion
% (cyclicity and separation of $\Omega$ for the relative commutant $\mB^c$) lets one apply
% Theorem~\ref{thm:HSMI} a second time to the derived inclusion $\mB^c\subset\mA$: since
% $\mB\subset\mA$ is HSMI$+$, the relation~\eqref{eq:HSMI} applied to $\mB'$ shows
% $\mB^c=\mB'\cap\mA$ is HSMI$-$, and one checks using~\eqref{eq:HSMI-detailed-balance} that
% $\mB^c\subset\mB'$ is in turn HSMI$+$. This is precisely the triple of conditions
% required by the general (``twisted'') version of this theorem at trivial twist $U=\id$
% in~\cite{LLMO2025}; the twisted construction there reduces at $U=\id$ to the untwisted
% statement proved originally by Wiesbrock~\cite{Wiesbrock1993}. The relative-commutant
% generator built from $\mB^c\subset\mA$ is then shown, by an explicit computation using the
% chain of relations~\eqref{eq:HSMI-chaos}--\eqref{eq:HSMI-translation} applied
% at the special point $s=\log 2$, to coincide with $G_2$ as defined
% from~\eqref{eq:B-hat}; this computation is the content of~\cite[eqs.~(6.31)--(6.40)]{LLMO2025}
% and we do not reproduce it here. Given~\eqref{eq:three-generators}, the commutator
% $[G_1,G_2]=-2iG_0$ is a direct consequence of the modular commutation relations of
% $\mA$, $\mB$, and $\widehat\mB$ applied to one another. Since $\psl$ is the unique
% three-dimensional non-abelian Lie algebra with this commutator structure (the Lorentz
% algebra $\mathfrak{so}(1,2)$ in the basis of a boost and two null translations), the
% operators $G_0,G_1,G_2$ exponentiate to a representation of its simply connected group
% $\PSLt$; invariance of $|\Omega\rangle$ under $G_0,G_1$ together with the
% $J$-conjugation relations fixes the inversion element to be~\eqref{eq:inversion}. See
% \cite[\S6.2]{LLMO2025} for the complete argument.
% \end{proof}
 
% \begin{remark}
% At trivial twist, $\widehat\mB=\mB\vee J_{\omega;\mA}\mB J_{\omega;\mA}$ is invariant under
% $J_{\omega;\mA}$, so $J_{\omega;\mA}$ and $J_{\omega;\widehat\mB}$ commute and the
% inversion~\eqref{eq:inversion} squares to the identity, as required of the inversion
% element of $\PSL$. This is the statement, isolated already by Wiesbrock, that a single
% standard half-sided modular inclusion is sufficient data to generate the full conformal
% group of the line: no additional input (such as a second, independently specified
% inclusion) is required, because the second translation generator $G_2$ is \emph{derived}
% from $\mB$ and its modular mirror image~\eqref{eq:B-hat}, rather than posited
% separately. A more general ``twisted'' version of Theorem~\ref{thm:wiesbrock-PSL}, in
% which the mirror image is taken with respect to an additional unitary $U$ commuting with
% $\Delta_{\omega;\mA}^{is}$, is developed in~\cite{LLMO2025} and is needed when the naive
% inclusion $\mB\subset\mA$ fails to be standard — the generic situation for conformal
% field theories with non-integer scaling dimension.
% \end{remark}
 
% Theorem~\ref{thm:wiesbrock-PSL} is the precise sense in which \emph{geometry emerges from
% ergodicity}: the data of a single standard half-sided modular inclusion — purely
% algebraic information about how one subalgebra sits inside another, with no a priori
% reference to points, distances, or causal structure — is already sufficient to
% reconstruct an honest representation of the one-dimensional conformal group. In
% Section~3 we will see this representation made concrete: the inclusion of wedges and of
% balls inside a vacuum QFT or CFT furnishes standard HSMI inclusions whose associated
% $\PSL$ representation is literally the geometric action of boosts, null translations,
% and special conformal transformations on the boundary of the region, and the
% characteristic function of Definition~2.4 is a matrix element of this very
% representation. The exponential dependence on $e^{2\pi z}$ that will appear throughout
% the wedge and ball computations of Sections~3.1 and~3.2 is not a coincidence of the
% free-field examples considered there: it is the universal signature, guaranteed by
% Theorem~\ref{thm:wiesbrock-PSL}, of any standard half-sided modular inclusion whatsoever.
 









% %\section*{Coarse-Graining, Conditional Expectation, and Modular Invariance}

% % A useful way to compare classical ergodic theory with operator-algebraic dynamics is to formulate both in terms of algebras of observables. 



% % \NL{Araki-Woods lattice system} as an example of QECC: correspondence between the states of the relative commutant and the conditional algebra.




\section{Geometric subregion flows: wedges in QFTs and balls in CFTs}



When there is geometry, there is a net of inclusions (coarse-grainings). We call such inclusions {\it geometric inclusions}. In the special case of inclusion of wedges in vacuum QFT, and the inclusion of balls in vacuum CFT, the flow associated with a geometric inclusion acts geometrically. However, in more general states of QFTs, or for more general geometric inclusions, the flow is non-local and non-geometric. Hence, we refer to it as {\it characteristic flows}.

%\subsection{Characteristic operators for inclusion of wedges and balls}



 %\NL{Add theorems for inclusion of balls and wedges}


\paragraph{Matrix elements of characteristic functions}

List of results: 1. matrix element = function of proper distance, 2. poles only for zero proper distance, 3. non-zero proper distance in the strip = poles only on the bry, 4. two pairs of poles on the bry of the stripe (Wedge), 5. branch cut for general $\Delta$, 6. 4 pairs of poles on the bry (Ball)

Quick mention about 1. the tower ends after the first step, 2. $J_AJ_B$ thing, 3. future light cone?
\subsection{Inclusion of Wedges}


The conformal symmetry allows us to take $B = W(a>0) \subset A = W(0)$ without loss of generality, where $W(a)$ denotes the Rindler wedge with the vertex point $(t=0,x=a)$. For a 1 + 1d CFT, the characteristic operator of the inclusion of wedges is explicitly computed in \cite{}. 
\begin{theorem}[Characteristic Operator of Inclusion of Wedges]
    In a 1+1d CFT with the wedges $W(a>0) \subset W(0)$, the characteristic operator $\Delta_{W(0)}^{-iz}\Delta_{W(a)}^{iz}$ is 
    \begin{align}
        \Delta_{W(0)}^{-iz}\Delta_{W(a)}^{iz} = \exp\left( ia\lb\sinh(2\pi z)P_0 + (\cosh(2\pi z)-1) P_1\rb\right),
    \end{align}
    where $P_\mu$ is the generator of translation in $x^{\mu}$-direction.
\end{theorem} 
For simplicity, we take $O$ to be a scalar primary field with conformal dimension $2h$. We may choose to normalize $f(z)$ with $\sqrt{\braket{\mO(X)|\Delta_A^{1/2}\mO(X)}\braket{\mO(Y)|\Delta_B^{1/2}\mO(Y)}}$. The modulus of the function is bounded by $1$ but strictly less than $1$ with this normalization, in general. We will omit the normalization factor unless it is important. For $X = (t_1,x_1)$ and $Y = (t_2,x_2+a)$ and real $z$, we have 
\begin{gather}
    f(z) =\frac{C}{(\Delta s^2+i\ep\Delta t)^{2h}} = \frac{C}{(\Delta s^2)^{2h}}\\ %=\frac{C}{[\Delta s^2-i\epsilon(t_1+t_2-a\sinh(2\pi z))]^{2h}},\nonumber\\
    \Delta t = t_1+t_2-a\sinh(2\pi z), \quad \Delta x =x_1+x_2 + a\cosh(2\pi z), \quad \Delta s^2 = -\Delta t^2 + \Delta x^2, %\\
    %\Delta s^2 = (x_1+x_2)^2 - (t_1+t_2)^2 + a^2 + 2a((x_1+x_2)\cosh(2\pi z) +(t_1+t_2)\sinh(2\pi z))
\end{gather}
where $C$ is a normalization factor. The $i\ep$ term can be dropped because the transformed $X$ and $Y$ are spacelike separated.


Without loss of generality, we consider the matrix element of the characteristic function $f(z)=\bra{O(X)}\Delta_A^{1/2+iz}\Delta_B^{-iz}\ket{O(Y)}$ of a CFT in 1+1d with inclusion of wedges, where $O(X),O(Y)$ are point-like fields with $Y\in B \subset A\ni X$. The general properties of the characteristic function imply that the function $f(z)$ is analytic in the strip Im$(z)\in (-1/4,1/4)$ and continuous and bounded in the closure. For self-adjoint $O(x), O(y)$, we also have $f(z) = \overline{f(\overline{z})}$ for $z$ in the stripe, and thus $f(z)\in\mathbb{R}$ for $z\in\mathbb{R}$. 


We observe the following properties of the analytic continuation of the function. All singularities are located at $\Delta s^2 = 0$. The first four singularities are located at Re$(z) = \pm\frac{1}{2\pi}\log(\frac{x_1\mp t_1+x_2\mp t_2}{a})$ with $|\text{Im}(z)| = 1/2$. The function $f(z)$ is analytic in a larger stripe Im$(z)\in (-1/2,1/2)$. The function has no zeros in the strip and decays to zero in the limit of Re$(z) \to \pm\infty$. The function is continuous and real in the interval between $-\frac{1}{2\pi}\log(\frac{x_1- t_1+x_2- t_2}{a})$ and $\frac{1}{2\pi}\log(\frac{x_1+ t_1+x_2+ t_2}{a})$ with $|\text{Im}(z)| = 1/2$. The function being real and continuous in these intervals allows us to further extend the analytic continuation to outside the strip by the Schwarz reflection principle. The principal value and the branch cuts of the extended function $f(z)$ can be chosen such that it is periodic with $f(z+i) = f(z)$ away from the branch cut, and the singularities are located at Re$(z) = \pm\frac{1}{2\pi}\log(\frac{x_1\mp t_1+x_2\mp t_2}{a})$ with Im$(z) \in \frac{1}{2}+\mathbb{Z}$. For $2h \in \mathbb{Z}$, the singularities become poles, and $f(z)$ is a meromorphic function on the complex plane. For $2h \notin \mathbb{Z}$, these are branch points. We can choose the branch cut to start from each branch point toward $\pm\infty$ horizontally. The discontinuity across the branch cut is a phase difference $e^{i2\pi h}$. The phase of $\Delta s^2$ as a function of $z$ is plotted in Fig.\ref{fig:deltaS2_z}. $\Delta s^2$ is a analytic function of $z$. The branch cut structure of $f(z) = C(\Delta s^2)^{-2h}$ is simply related to the phase of $\Delta s^2$ around the zeros. To summarize the properties, we have:
\begin{figure}[tp]
    \centering
    \begin{subfigure}{.3\textwidth}
        \includegraphics[width=1\linewidth]{images/delta_s2_plot_ordered.pdf}
        \label{fig:placeholder1}
        %\mathcal{A}ption{a}
    \end{subfigure}
    \begin{subfigure}{.3\textwidth}
        \includegraphics[width=1\linewidth]{images/delta_s2_plot_coincide.pdf}
        \label{fig:placeholder2}
        %\mathcal{A}ption{b}
    \end{subfigure}
    \begin{subfigure}{.3\textwidth}
        \includegraphics[width=1\linewidth]{images/delta_s2_plot_reversed.pdf}
        \label{fig:placeholder3}
        %\mathcal{A}ption{c}
    \end{subfigure}     
    \caption{The phase of $\Delta s^2$ as a function of $z$. The left figure is the case with $-\log(x_1- t_1+x_2- t_2) < \log(x_1+ t_1+x_2+ t_2)$. The middle one is the case with the branch points coinciding. The right one is the case with $-\log(x_1- t_1+x_2- t_2) > \log(x_1+ t_1+x_2+ t_2)$. To illustrate the results, we have picked $a = 1, x_1 = 0.1, t_1 = 0, t_2 = 0, 2h = 1.5$. The left one has $x_2 = 3$. The middle one has $x_1+x_2 = a$. The right one has $x_2 = 1/3$.}
    \label{fig:deltaS2_z}
\end{figure}

\begin{lemma} [Wedge-inside-wedge inclusion in CFT]
    For wedge-inside-wedge inclusion, the matrix element of the characteristic operator $f(z) = \bra{O(X)}\Delta_A^{1/2+iz}\Delta_B^{-iz}\ket{O(Y)}$, with self-adjoint scalar field $O(X), O(Y)$ and $Y\in B \subset A \ni X$, has the following properties:
    \begin{enumerate}
        \item There is an analytic continuation of $f(z)$ on the strip Im$(z)\in (-1/2,1/2)$. If the conformal dimension is not an integer, there are a pair of branch points on $\text{Im}(z)=\pm1/2$ each. If the conformal dimension is an integer, there are no branch cuts, and the branch points become poles.
        \item $f(z)$ is singular only when the transformed point-like field has zero proper distance.
        \item $f(z)$ has no zero. Zero is achieved only when the transformed point-like field has an infinite proper distance. The function decays exponentially to zero for large $|\text{Re}(z)|$.
        \item The function is real on Im$(z) = 0$, and therefore $f(\overline{z}) = \overline{f(z)}$.
        \item The function is real and continuous in the interval between the pair of singularities. If the pairs of singularities do not coincide, $f(z)$ can be extended to the whole complex plane, except for the singularities and possibly the branch cuts that repeat on $\text{Im}(z) \in \mathbb{Z} +1/2$. There is a choice of branch cuts where the function is periodic with period $i$, i.e. $f(z+i)=f(z)$. 
    \end{enumerate}
\end{lemma}

Although we only consider the point-like field $O(X),O(Y)$ in this setup, some features remain if we consider the smeared field. The smearing function may remove the singularities, but it cannot remove the branch cut. The multi-valued nature remains even with smeared fields. Also, several features remain valid in QFT and in higher dimensions. In general QFT, the two-point function of two self-adjoint point-like fields is still a function of the proper distance $\Delta s^2$. This is because the theory has Lorentz symmetry, and the modular flows of wedges remain geometric in QFT. The two-point function has a singularity only if $\Delta s^2 = 0$ in general. The function can only achieve zero with infinite proper distance. The function is real if the two points are spacelike separated at some $z$ due to causality. With these general features of QFT, most of the above properties and the singularity structure remain the same in general QFT, except for the details about the singularities and potentially the branch cuts. 

\subsection{Inclusion of Ball inside Wedge}

\NL{Add theorem for the characteristic function for balls in CFT}

In this subsection, we consider the setup with ball-inside-wedge inclusion. There is a distinction compared to the case of the inclusion of wedges. The inclusion of wedges has a common end point at infinity, so that the inclusion is not split. The inclusion of balls is the same as the inclusion of a ball inside a wedge. There is a conformal transformation that maps the larger ball to a wedge. In \cite{}, the characteristic operator of inclusion of balls in computed for 1+1d CFT.
\begin{theorem}[Characteristic Operator of Inclusion of Concentric Double Cones]
    In a 1+1d CFT with concentric doable cones $B_{u+s} \subset B_u$, where $B_s$ is the double cone centered at the origin with radius $e^{-2\pi s}$, the characteristic operator $\Delta_{B_u}^{-iz}\Delta_{B_{u+s}}^{iz}$ is 
    \begin{align}
        \Delta_{B_u}^{-iz}\Delta_{B_{u+s}}^{iz} = \exp\left( -\beta\left(\frac{1}{2\pi \cosh(\pi s)}\partial_uK_{B_{u+s/2}}-\tanh(\pi z)K_F\right) \right),
    \end{align}
    where 
    \begin{gather}
        \beta = -\frac{i \cos^{-1}(A)}{2\pi \sqrt{1-A^2}}\sinh(2\pi s)\sinh(2\pi z), \quad
        A = 1-2\sinh^2(\pi s)\sinh^2(\pi z),
    \end{gather}
    $K_{B_s} = -\ln\Delta_{B_s}$ is the modular Hamiltonian of the algebra of $B_s$, and $K_F$ is the modular flow of the algebra of the future light cone $\{x^0>|x^1|\}$.
\end{theorem}
Instead of using the expression for the inclusion of concentric double cones, we consider the setup with a double cone $B$ included in the right Rindler wedge $A$ in 1+1d CFT. The conformal transformation that maps the inclusion of double cones to the inclusion of a ball inside a wedge preserves the vacuum. The inclusion of concentric double cones is thus a special case of the inclusion of a double cone inside a wedge with general relative position. The modular flow of a double cone $B$ with the two tips $(t_c\pm r,x_c)$ is
\begin{align}
    \Delta_B^{-is} O(t,x) \Delta_B^{is} = \gamma^{2h} O(t_s,x_s),
\end{align}
where
\begin{gather}
    %t_s = \frac{2r^2(t-t_c) \cosh(2\pi s) + (r^2-d^2)r \sinh(2\pi s)}{r^2+d^2 + (r^2-d^2)\cosh(2\pi s) +2r(t-t_c)\sinh(2\pi s)} + t_c,\\
    t_s = \gamma \left[(t-t_c)\cosh(2\pi s) + \frac{r^2-d^2}{2r}\sinh(2\pi s)\right] + t_c,\\
    %x_s = \frac{2r^2(x-x_c)}{r^2+d^2 + (r^2-d^2)\cosh(2\pi s) +2r(t-t_c)\sinh(2\pi s)} + x_c,\\
    x_s = \gamma (x-x_c) + x_c,\\
    \gamma = \frac{2r^2}{r^2+d^2 + (r^2-d^2)\cosh(2\pi s) +2r(t-t_c)\sinh(2\pi s)}, \\
    d^2 = -(t-t_c)^2 + (x-x_c)^2.
\end{gather}
For $(t,x) \in B$ and $s\in \mathbb{R}$ , $\gamma>0$ and the point $(t_s,x_s)$ remains in $B$ . For $(t,x) \notin B$, $\gamma$ could be non-positive for some real $s$. This corresponds to the point-like field crossing infinity and reaching the other patch of Minkowski space. There will be a phase from $\gamma^{2h}$ when $2h$ is non-integer. For computing the matrix element of the characteristic operator, we only need the case where $(t,x) \in B$.

The matrix element of the characteristic operator depends on the proper distance and the conformal factor $\gamma$. We consider $f(z)=\bra{O(X)}\Delta_A^{1/2+iz}\Delta_B^{-iz}\ket{O(Y)}$ with $Y\in B \subset A\ni X$. For $X = (t_1,x_1), Y = (t_2+t_c,x_2+x_c)$, and real $z$, we have
\begin{gather}
    %\gamma &= \frac{2r^2}{r^2+d^2 + (r^2-d^2)\cosh(2\pi z) +2rt_2\sinh(2\pi z)} \\
    \gamma= \frac{4r^2 e^{2\pi z}}{r^2-d^2+2rt_2} \frac{1}{(e^{2\pi z} + p_+)(e^{2\pi z} + p_-)},\\
    %\Delta t &= -t_1 \cosh(2\pi z) - x_1 \sinh(2\pi z) - \gamma\left[t_2\cosh(2\pi z) + \frac{r^2-d^2}{2r}\sinh(2\pi z)\right] - t_c, \\
    %\Delta x &=   -t_1 \sinh(2\pi z) - x_1 \cosh(2\pi z) - \gamma x_2 - x_c,\\
    \Delta x + \Delta t = -(x_1+t_1)e^{2\pi z} - (x_c+t_c) - r\frac{e^{2\pi z}-p_-}{e^{2\pi z}+p_-},\\
    \Delta x - \Delta t = - (x_1-t_1)e^{-2\pi z} - (x_c-t_c) - r\frac{e^{2\pi z}-p_+}{e^{2\pi z}+p_+},\\
    %f(z) &= \gamma^{h} \frac{C}{[\Delta s^2-i\epsilon(t_1 \cosh(2\pi z) + x_1 \sinh(2\pi z) + \gamma (t_2\cosh(2\pi s) + \frac{r^2-d^2}{2r}\sinh(2\pi s))+t_c)]^{h}}.
    f(z) = C \left(\frac{\gamma}{-\Delta t^2 + \Delta x^2} \right)^{2h},
\end{gather}
where
\begin{gather}
    d^2 = -t_2^2 + x_2^2, \quad p_\pm = \frac{r^2+d^2\pm 2rx_2}{r^2-d^2+2rt_2} > 0,\\
    p_+p_- = \frac{r^2-d^2-2rt_2}{r^2-d^2+2rt_2}.
\end{gather}
We can then consider the analytic continuation of $f(z)$ with complex $z$. Similar to the previous case, the singularities are located at $\Delta s^2 = 0$. The above expression may suggest that $\gamma \to \infty$ is another possible singularity. However, $\gamma \to \infty$ is achieved only if the point $Y$ is mapped to the infinity. The infinite proper distance ensures that $f(z)$ is finite even if $\gamma$ diverges. For $\Delta s^2/\gamma = 0$, we have
\begin{gather}
    (x_1+t_1)e^{4\pi z} +  (x_c + t_c + r + p_-x_1+ p_-t_1)e^{2\pi z}+ (x_c+t_c-r) p_-  = 0,\\
    \text{or } \quad (x_1-t_1)e^{-4\pi z}+ \left(x_c - t_c - r+\frac{x_1-t_1}{p_+}\right)e^{-2\pi z} + \frac{x_c-t_c + r}{p_+} = 0.
\end{gather}

The general properties are summarized in the following lemma.
\begin{lemma} [Ball-inside-wedge inclusion in CFT]
    For ball-inside-wedge inclusion $B\subset A$, the matrix element of the characteristic operator $f(z) = \bra{O(X)}\Delta_A^{1/2+iz}\Delta_B^{-iz}\ket{O(Y)}$, with self-adjoint scalar field $O(X), O(Y)$ and $Y\in B \subset A \ni X$, has the following properties:
    \begin{enumerate}
        \item There is an analytic continuation of $f(z)$ on the strip Im$(z)\in (-1/2,1/2)$. If the conformal dimension is not an integer, there are two pairs of branch points on $\text{Im}(z)=\pm1/2$ each. If the conformal dimension is an integer, there are no branch cuts, and the branch points become poles.
        \item $f(z)$ is singular only when the transformed point-like fields have zero proper distance.
        \item $f(z)$ has no zero. Zero is achieved only when $|\text{Re}(z)| \to \infty$. The function decays exponentially to zero for large $|\text{Re}(z)|$.
        \item The function is real on Im$(z) = 0$, and therefore $f(\overline{z}) = \overline{f(z)}$.
        \item There is a interval between two singularities for which the function is real and continuous, if the singularities do not coincide. In that case, $f(z)$ can be extended to the whole complex plane, except for the singularities and possibly the branch cuts that repeat on $\text{Im}(z) \in \mathbb{Z} +1/2$. There is a choice of branch cuts where the function is periodic with period $i$, i.e. $f(z+i)=f(z)$. 
    \end{enumerate}
\end{lemma}
One main difference compared to the wedge-inside-wedge case is that there are more singularities and, in general, more branch cuts on $\text{Im}(z) = \pm 1/2$. In this case, $\Delta x \pm \Delta t = 0$ gives two quadratic equations of $e^{\pm 2\pi z}$. In the wedge-inside-wedge case, $\Delta x \pm \Delta t = 0$ gives a linear equation of $e^{\pm2\pi z}$. Another difference is that these properties are more restricted to CFT only. In a QFT, the modular flow of a ball-shaped region is not geometrical in general. We cannot directly carry the proper distance argument to QFT. Because of this limitation, we consider the approach of conformal perturbation method in the next section.

\NL{Study this in modular frequency space}

\subsection{Modular Chaos Bound}
We first review the modular chaos bound condition studied in \cite{}\footnote{Different conventions maybe used but they are simply related by a shift in $z$, a sign flip, or a complex conjugate.}. We consider, for $B \subset A$,
\begin{eqnarray}
    f(z)&&=\frac{\bra{O_1}\Delta_A^{1/2+iz}\Delta_B^{-iz}\ket{O_2}}{\sqrt{\bra{O_1}\Delta_A^{1/2}\ket{O_1} \bra{O_2}\Delta_B^{1/2}\ket{O_2}}}\\
    &&=\frac{\bra{\Delta_A^{1/4}O_1} \Delta_A^{1/4+iz}\Delta_B^{-1/4-iz}\ket{\Delta_B^{1/4} O_2}}{\sqrt{\bra{O_1}\Delta_A^{1/2}\ket{O_1} \bra{O_2}\Delta_B^{1/2}\ket{O_2}}}\nn\\
    &&=1+i \mathcal{M}(z)
\end{eqnarray}
Both $\ket{\Delta_A^{1/4}O_1}$ and $\ket{\Delta_B^{1/4} O_2}$ have a finite norm even if $O_1, O_2$ are a point-like operator. $f(z)$ is bounded analytic on the strip Im$(z) \in (-1/4,1/4)$ and continuous on the closure. This property follows from $\Delta_A^{iz}\Delta_B^{-iz}$ being bounded analytic on the strip Im$(z) \in (-1/2,0)$ and strongly continuous on the closure. 
First, we find that 
\begin{eqnarray}
    1-\text{Im}(\mathcal{M}(z))=\text{Re}(f(z))\leq |f(z)|\leq 1
\end{eqnarray}
which implies that 
\begin{eqnarray}
    \text{Im}\mathcal{M}(z)\geq 0\ .
\end{eqnarray}
The change of variable $E=ie^{-2\pi z}$ sends the strip to the upper half-plane $\mathbb{H}^+$. $\mM(E)$ is an analytic map from the upper half-plane to itself. By writing $E=v+iw$, the Schwarz-Pick lemma gives
\begin{eqnarray}
    \frac{|w \p_w \mM(v+iw)|}{\text{Im} \mM(v+iw)}\leq 1\ . 
\end{eqnarray}
Since we have $|\p_w \text{Im}\mM(v+i\omega)|\leq |\p_w \mM(v+i\omega)|$, we find
\begin{eqnarray}
    |w\p_w \log \text{Im} \mM(v+iw)| \stackrel{(1)}{\leq} \frac{|w \p_w \mM(v+iw)|}{\text{Im} \mM(v+iw)} \stackrel{(2)}{\leq} 1\ .
\end{eqnarray}
This is the modular chaos bound (the exact version). To saturate this bound, the inequality $(2)$ fixes $\mM(E)=\frac{a E+b}{c E+d}$ to be a $PSL(2,\mathbb{R})$ transformation. Next, the saturation of the first inequality $(1)$ gives
\begin{eqnarray}
    \p_w \text{Re} \mM(v+iw)=0, \quad \forall v+iw\in \mathbb{H}^+ .
\end{eqnarray}
This is satisfied if and only $\mM(E)$ takes the form
\begin{eqnarray}
    \mM(E)=a E+b, \quad a>0,\ b\in\mathbb{R}.
\end{eqnarray}
However, this saturation is very restrictive and has little application. $f(z) = 1+ib - a e^{-2\pi z}$ is not bounded in the strip. Therefore, there is no $f(z)$ saturates this bounded and the inequality is not tight, i.e. $|w\p_w \log \Im \mM(v+iw)| < 1$. The bound at best gives a supremum rather than a maximum.
\begin{lemma}
    Given a function $\mM(E)$ that is analytic in the upper half-plane ($\text{Im} E > 0$). If $f(E) = 1+i\mM(E)$ is bounded by 1 in the upper half-plane, then $\mM(E)$ cannot saturate the modular chaos bound, i.e.
    \begin{align}
        |w\p_w \log \text{Im} \mM(v+iw)| < 1, \quad \text{for}\ E=v+iw.
    \end{align}
\end{lemma}
There is another way to see why we need an approximate version of this bound. We consider $\mM = -i(f-1)$, which has image within the unit disk centered at $i$ in $\mathbb{C}$. The inequality $(2)$ is saturated if and only if the map is a M\"obius transformation with real coefficients. The M\"obius transformations are a bijective function that maps the upper half-plane to itself. The image of $\mM$ is by definition a strict subset of the upper half-plane, whereas the domain is the entire upper half-plane. The inequality can never exactly saturate. Instead, we shall consider an asymptotic saturation with a suitable limit/approximation. The usual modular chaos bound is proposed in the following way\cite{}. Consider $\mM = i\varepsilon p e^{\lambda z} + O(\varepsilon^2)$ for small $\varepsilon \ll 1$, $\varepsilon e^{\lambda z} \ll 1$, and $p>0$. The modular chaos bound gives $\lambda \leq 2\pi$. This bound is only meaningful in the proper range of time scales. This bound controls the behaviour of early time chaos. 

In both CFT examples we considered, $f(z)$ is a function of $e^{2\pi z}$. Thus, the expansion of the properly shifted and normalized $f(z)$ will take the form $1-\varepsilon p e^{2 \pi z} + O(\varepsilon^2)$. Thus, the modular chaos bound is saturated in these CFT examples.

% \subsection{Modular chaos bound in GFF}

% We then consider the modular chaos bound in the 0+1d GFF example. We consider $B = (-1,1)$ and $B(s) = (-e^{-2\pi s}, e^{-2\pi s})$. We define
% \begin{eqnarray}
%     f_{s,x}(u)&&=\frac{\bra{\phi(x)}\Delta_B^{1/2+iu}\Delta_{B(s)}^{-iu}\ket{\phi(x)}}{\sqrt{\braket{\phi(x)|\Delta_B^{1/2}\phi(x)}\braket{\phi(x)|\Delta_{B(s)}^{1/2}\phi(x)}}}\\
%     &&=\frac{\bra{\Delta_A^{1/4}\phi(x)} T_{0,s}(-i/4+u)\ket{\Delta_B^{1/4} \phi(x)}}{\sqrt{\braket{\phi(x)|\Delta_B^{1/2}\phi(x)}\braket{\phi(x)|\Delta_{B(s)}^{1/2}\phi(x)}}}\nn\\
%     &&=1+i \mathcal{M}_{s,x}(u),
% \end{eqnarray}
% for $x \in B(s)\subset B$. By substituting $E = ie^{-2\pi u} \in {\mathbb{H}_+}$ for $\Im u \in (-\pi/4,\pi/4)$, we have
% \begin{align}
%     the.long.complicated.expression.here.finish.it.later
% \end{align}
% This is not a M\"obius transformation, which is expected. In the $\pi s \ll 1$ limit, we have
% \begin{align}
%     f_{s,x}(E) &= 1 - i\pi s h\lb \frac{1-x}{1+x}\frac{1}{E} - \frac{1+x}{1-x}E \rb + O((\pi s)^2),\\
%     f_{s,x}(u) &= 1 - \pi s h\lb \frac{1-x}{1+x}e^{2\pi u} + \frac{1+x}{1-x}e^{-2\pi u} \rb + O((\pi s)^2).
% \end{align}
% We then consider the late time behaviour in a double-scale limit, i.e. $e^{2\pi u} \gg1$ while keeping $1\gg\pi se^{2\pi u}$. The previous expression gives
% \begin{align} 
%     f_{s,x}(u) \approx 1 - \pi s h \frac{1-x}{1+x}e^{2\pi u},\\
%     \mM_{s,x}(u)\approx i\pi s h \frac{x-1}{x+1}e^{2\pi u}.
% \end{align}
% This is the saturation of the modular chaos bound asymptotically. It gives a similar result if we consider $e^{-2\pi u} \gg1$ instead. In other words, $\mM_{s,x}(u) \sim ise^{2\pi |u|}$ for $e^{2\pi |u|} \gg 1\gg\pi se^{2\pi |u|}$.



\section{Subregion flows in Conformal Perturbation Theory}\label{section:conformalperturbation}


%\NL{should we call them geometric or local inclusions or not?}

4.1. Setup of perturbation of modular flow, 4.2. cocycle flow, 4.3. 1st order correction of characteristic function, 4.4 analytic property 

A short paragraph about general QFT, move on to conformal perturbation theory, %mention that only 2- and 3-pt functions are universal, best is first order correction if without CFT details 
\subsection{Modular Flow with Conformal Perturbation Method}

\NL{Without loss of generality and for simplicity of presentation for the remainder of this section, we focus on 1+1 D QFT}

To study the characteristic function in QFT,  we first compute the modular flow perturbatively. We apply the conformal perturbation method to compute the modular flow of a ball-shaped region. We follow the method in \cite{,} and consider the Lagrangian perturbed with a relevant deformation 
\begin{align}
    \mathcal{L}_\text{CFT} \to \mathcal{L}_\text{CFT} - \lambda O(x), 
\end{align}
where $O(x)$ is a conformal primary in the flat space coordinate $x$ and $\lambda$ is a small parameter. The modular flow $e^{isK_\lambda}$ of an algebra $\mathcal{A}$ associated to a ball-shaped region $I$ (restricted to the algebra) can be computed perturbatively
\begin{gather}
    \delta_\lambda(e^{isK_\lambda}) = \int_0^1 dw e^{iswK_0} is\delta_\lambda K_\lambda e^{is(1-w)K_0} , \\
    \delta_\lambda K_\lambda= \int_{-\infty + i\theta}^{\infty + i\theta} \frac{2\pi du}{4\sinh^2(\pi u)} e^{iuK_0}\rho^{-1}\delta\rho e^{-iuK_0}, \quad\text{for } 0<\theta<1,\\
    \rho^{-1}\delta\rho = -\lambda\int_I d^{D-1}X\int_0^1 d\tau J(\tau,X) \Omega^{\Delta}(-i\tau,X) e^{\tau K_0} O(X) e^{-\tau K_0} ,
\end{gather}
where $J$ is the Jacobian determinant between the geometrically transformed coordinate  on $e^{\tau K_0} O(X) e^{-\tau K_0}$ and $(X, \tau)$, $\Omega$ is the conformal factor that depends on the geometric action of the flow $e^{is K_0}$, $\Delta$ is the conformal dimension of $O$\footnote{We would consider $D/2 < \Delta < D$ for the perturbation method to work \cite{}.}, and $D$ is the spacetime dimension. $K_\lambda$ here is not the full modular Hamiltonian. It is the one-sided modular Hamiltonian that acts only on the algebra. By substituting $\rho^{-1}\delta_\lambda\rho$ into the expression of $\delta_\lambda K_\lambda$ and changing the order of integration, we can choose $\theta = \tau$ to cancel the Euclidean part and obtain a  purely Lorentzian action \cite{}
\begin{align}
    \delta_\lambda K_\lambda=& -\lambda\int_I dX\int_0^1 d\tau \int_{-\infty}^{\infty} du \, \frac{2\pi J(\tau,X)}{4\sinh^2(\pi (u+i\tau))}   \Omega^{\Delta}(-i\tau,X) e^{iuK_0} O(X) e^{-iuK_0} \\
    =& \pi \lambda\int_I dX \int_{-\infty}^{\infty} du \, h(u,X) e^{iuK_0} O(X) e^{-iuK_0}, \label{eq:deltaK}
\end{align}
where we define
\begin{align}
    h(u,X) \equiv -\int_0^1\, d\tau \frac{J(\tau,X)\Omega^{\Delta}(-i\tau,X)}{2\sinh^2(\pi (u+i\tau))} 
\end{align}
Without loss of generality and for simplicity of the presentation, we focus on 1+1 D QFT. For simplicity, we consider $I$ to be the interval $(-1,1)$ at $t = 0$. The unit ball $B$ is the double cone corresponding to $I$. We can obtain the result for a double cone in general position by applying the Lorentz transformation on the unit ball $B$. We can put $r$ back for ball-shaped region with radius $r$. We can also generalize the result to ball-shaped regions in higher dimension. %In this setup, we have 
% \begin{gather}
%     e^{isK_0} O(t,x) e^{-isK_0} = \tilde{\Omega}^{-\Delta} O(t_s,x_s), \\
%     \tilde{t}_s = \tilde{\Omega}^{-1}\left[t \cosh(2\pi s) + \frac{(1+t^2-x^2)}{2}\sinh(2\pi s) \right],\\
%     \tilde{x}_s = \tilde{\Omega}^{-1}x, \\
%    \tilde {\Omega}^{-1} = \frac{2}{1-t^2+x^2+(1+t^2-x^2)\cosh(2\pi s) + 2t \sinh(2\pi s)}.
% \end{gather}
In this set up, for $O(t=0,x)$, we have
\begin{gather}
    e^{isK_0} O(t,x) e^{-isK_0} = \Omega^{-\Delta} O(t_s,x_s), \\
    t_s = {\Omega}^{-1}\left[\frac{(1-x^2)}{2}\sinh(2\pi s) \right], \quad x_s = {\Omega}^{-1}x, \\
    {\Omega}(s,x) =  \frac{1+x^2+(1-x^2)\cosh(2\pi s)}{2} = 1+(1-x^2)\sinh^2(\pi s).
    %{\Omega}(s,x) = \tilde{\Omega}|_{t=0}= \frac{1+x^2+(1-x^2)\cosh(2\pi s)}{2} = 1+(1-x^2)\sinh^2(\pi s).
\end{gather}
The conformal factor $\Omega$ is $\gamma^{-1}$ in previous section with $r = 1, t =t_c=x_c=0$. The conformal factor $\Omega$ and the Jacobian determinant $J$ in Euclidean evolution are
\begin{gather}
    \Omega(-i\tau,x)= \frac{1+x^2+(1-x^2)\cos(2\pi \tau)}{2} =1-(1-x^2)\sin^2(\pi \tau),\\
    J(\tau,x) = \det \left.\begin{pmatrix}
        \partial_st_{s} & \partial_s x_{s}\\  \partial_X t_{s}& \partial_X x_{s}
    \end{pmatrix}\right|_{s=-i\tau} =\frac{\pi (1-x^2)}{\Omega(\tau,x)^2},\\
    \implies h(u,X) = -\pi(1-X^2)\int_0^1\, d\tau \frac{\Omega^{\Delta-2}(-i\tau,X)}{2\sinh^2(\pi (u+i\tau))}.
\end{gather}
%{(Maybe write it as something like convolution, and the delta function like fattened, also try numeric/plot to see. Also see $x^2=1$ is special, maybe like interpolation between x, or some coarse graining)}
Notice that the $d\tau$ integral of $h(u,X)$ can be done explicitly if there is no $\Omega^{\Delta-2}(-i\tau,X)$, with
\begin{align}
    -\int_0^1\, \frac{ \pi d\tau }{2\sinh^2(\pi (u+i\tau))} = \delta(u).
\end{align}
A similar calculation appears in the calculation for the Rindler wedge case. The modular flow of the Rindler wedge in a CFT is a boost, which gives a trivial conformal factor. This delta function combined with the CFT modular Hamiltonian gives the new boost operator in the QFT \cite{}. With the delta function, Eq. (\ref{eq:deltaK}) becomes an integral of $X$ only. This term is the local contribution.  If $\Delta = 2$, the extra factor $\Omega^{\Delta-2}(-i\tau,X)$ becomes trivial and we obtain the delta function as in the Rindler wedge case. The intuition is that the modular flow remains geometric because the deformed CFT is just another CFT, though $\Delta = 2$ is a marginal perturbation and this may not be a suitable application of the perturbation method. For these reasons, we separate the divergent part from the integral
\begin{gather}
    h(u,X) =  (1-X^2)\left[\delta(u) + h_{reg}(u,X)\right], \\
    h_{reg}(u,X) \equiv -\pi\int_0^1\, d\tau \frac{ (\Omega^{\Delta-2}(-i\tau,X) -1)}{2\sinh^2(\pi (u+i\tau))} .
\end{gather}
By combining all the expressions, we have
\begin{align}
    u_{s,\lambda} =&\; e^{-isK_0}e^{isK_\lambda} \approx 1 + e^{-isK_0}\delta e^{isK_\lambda}\\
    =&\; 1 + i\pi s\lambda\int_0^1 dw  \int_{-1}^1 dX \int_{-\infty}^{\infty} du \, h(u,X) e^{i(u+sw-s)K_0} O(X) e^{-i(u+sw-s)K_0} \\
    =&\; 1 + i\pi \lambda\int_0^s dw  \int_{-1}^1 dX \int_{-\infty}^{\infty} du \, h(u+w,X) e^{iuK_0} O(X) e^{-iuK_0} \\
    =&\; 1 + i\pi \lambda\int_0^s dw  \int_{-1}^1 dX \int_{-\infty}^{\infty} du \, (1-X^2)\left[\delta(u+w) + h_{reg}(u+w,X)\right] e^{iuK_0} O(X) e^{-iuK_0}\\
    =&\; 1 + i\pi \lambda\int_{-s}^0 dw  \int_{-1}^1 dX \, (1-X^2)e^{iwK_0} O(X) e^{-iwK_0} \nonumber\\
    &+ {i\pi \lambda}\int_0^s dw  \int_{-1}^1 dX \int_{-\infty}^{\infty} du \, (1-X^2)h_{reg}(u+s,X) \Omega^{-\Delta}(u,X)O(t_u,x_u).
\end{align}
Thus, we have
\begin{gather}
    u_{s,\lambda} = 1 + i\pi \lambda \int_{-1}^1 dX \int_{-\infty}^{\infty} du \, (1-X^2) Q(u,X;s) \Omega^{-\Delta}(u,X) O(t_u,x_u) +O(\lambda^2),\\
    Q(u,X;s) = Q_{\text{loc}}(u,X;s) + Q_{\text{nonloc}}(u,X;s),\\
    Q_{\text{loc}}(u,X;s) =  \textbf{1}_{(-s,0)}(u),
\end{gather}
and 
\begin{align}
    Q_{\text{nonloc}}(u,X;s) \equiv&\;\int_0^s dw \; h_{reg}(u+w,X) \nonumber\\
    =&\; -\pi\int_0^s dw\int_0^1 \, d\tau  \;\frac{\left(1 - (1-X^2)\sin^2(\pi \tau)\right)^{\Delta-2} -1}{2\sinh^2(\pi (u+w+i\tau))}.
\end{align}
%{(Emphasis the local, non local, emphasis the role of $w, x, u$, rewrite the first term to look more local, combine $sdw$, maybe even bring the range $\int_0^s$ to the $1_{[-s,0]}$, bring the local term here, etc.)} 
$\textbf{1}_{(-s,0)}(x)$ is the indicator function of $(-s,0)$\footnote{If $s < 0$, we set $\textbf{1}_{(-s,0)}(x) = - \textbf{1}_{(0,-s)}(x)$.}. $Q_{\text{loc}}$ is the contribution of the delta function, which is the local contribution term to the correction of the modular Hamiltonian. $Q(u,X;s)$ should be interpreted in terms of distribution. If $\Delta > 3/2$, we can swap the order of the integral of $dw$ and $d\tau$. The divergences at $u = 0, \tau = 0 , 1$ are regulated by separating $Q_{\text{loc}}$. The divergence at $X=0, \tau = 1/2$ is mild enough so that the integrand is absolutely integrable. This allows us to change the order of integration. 
\begin{align}
    Q_{\text{nonloc}}(u,X;s) =&\; \int_0^1 \, d\tau \; \left[\left(1 - (1-X^2)\sin^2(\pi \tau)\right)^{\Delta-2} -1\right] \nonumber\\
    &\times\frac{\coth(\pi(u+s+i\tau)) - \coth(\pi(u+i\tau))}{2}.
\end{align}
The integral of $X$ and $u$ is the integration over the double cone region. The $Q_{\text{loc}}$ term is simply a indicator function. The contribution is the local perturbation from zero modular evolution to $-s$. The $Q_{\text{nonloc}}$ term has an extra integral of the Euclidean modular time $\tau$. Although the result is purely Lorentzian, the Euclidean modular evolution contributes to the weight function. For general $\Delta-2 < 0$, the integral is hard to compute.
%\KL{(play around $s,w,u$ to see if symmetrize the expression, maybe write the $s$ as some indicator function)}


\paragraph{Almost-marginal relevant operators*:} 
We consider small $\epsilon =2-\Delta\ll 1$ and expand $Q_{\text{nonloc}}(u,X;s)$ up to the leading order term. We have the following expansion\footnote{One may worry that the first order approximation is only valid for $\epsilon \ln\Omega \ll 1$ but $\Omega$ can be arbitrarily large when $X\approx 0$ and $\tau \approx 1/2$. We can split the $\tau$ integral into $(\frac{1}{2}-\delta, \frac{1}{2}+\delta)$ and the rest, with $1 \gg \pi\delta > e^{-\frac{1}{2\epsilon}}$. For $X\approx0$, $\Omega^{\Delta-2} \approx (X^2+\pi^2\delta^2)^{\Delta-2}$ such that the integral in $(\frac{1}{2}-\delta, \frac{1}{2}+\delta)$ is $O(e^{-\frac{1}{2\epsilon}})$. The width of the interval is also $O(e^{-\frac{1}{2\epsilon}})$. This is also how we see the integrand of $Q_{\text{nonloc}}$ being absolutely integrable for $\Delta > 1$. The orthogonality of cosine functions we use later also remains approximately true because of the same reason.}
\begin{align}
    \Omega^{\Delta-2}(-i\tau,X) - 1 =&\; - \epsilon \ln \left(\frac{1+X^2+(1-X^2)\cos(2\pi \tau)}{2}\right) + O(\epsilon^2)\\
    =&\; - \epsilon \ln \left[\frac{1+X^2}{2}\left(1+\frac{1-X^2}{1+X^2}\cos(2\pi \tau)\right)\right] + O(\epsilon^2). %\\
    %= &\; \epsilon \sum_{n=0}^{\infty} \left( \right)\cos(2 \pi n \tau) + O(\epsilon^2) 
\end{align}
% We use the identity ($r\in (-1,1), \tau\in[0,1]$)
% \begin{align}
%     \ln(1+r^2 + 2r\cos(2\pi\tau)) =&\; \ln(1+re^{i2\pi\tau}) + \ln (1+re^{-i2\pi\tau})\\
%     =&\; -\sum_{n= 1}^{\infty} \frac{(-r)^n e^{i2\pi n\tau}}{n} -\sum_{n= 1}^{\infty} \frac{(-r)^n e^{-i2\pi n\tau}}{n} \\
%     =&\; -\sum_{n \neq 0} \frac{(-r)^{|n|}e^{i2\pi n\tau}}{|n|} .
% \end{align}
We can expand $\Omega^{\Delta-2}(-i\tau,X)$ as a Fourier series $\Omega^{\Delta-2}(-i\tau,X) = \sum_{n} a_n e^{i2\pi n \tau}.$ We have\footnote{We use the identity $\ln(1+r^2 + 2r\cos(2\pi\tau)) = \ln\left[(1+re^{i2\pi\tau})(1+re^{-i2\pi\tau})\right] = -\sum_{n \neq 0} \frac{1}{|n|}(-r)^{|n|}e^{i2\pi n\tau}$ for $r\in (-1,1), \tau\in[0,1]$ and substituting $r = \frac{1-|X|}{1+|X|}$.}
\begin{gather}
    \Omega^{\Delta-2}(-i\tau,X) - 1 = \epsilon \left(\sum_{n} a_n^{(1)}e^{i2 \pi n \tau} \right) + O(\epsilon^2),\\
    a_0^{(1)} =  -2\ln \frac{1+|X|}{2}, \quad a_{n \neq 0}^{(1)} = \frac{1}{|n|} \left(-\frac{1-|X|}{1+|X|}\right)^{|n|}.
\end{gather}
We can also expand the $\coth$ function as
\begin{align}
    \coth(\pi (u+i\tau)) =& \frac{e^{\pi (u+i\tau)}+e^{-\pi (u+i\tau)}}{e^{\pi (u+i\tau)}-e^{-\pi (u+i\tau)}} \nonumber\\
    % =& \begin{cases}
    %     (1+e^{-2\pi (u+i\tau)}) \sum_{m=0}^{\infty} e^{-2\pi m (u+i\tau)}, & \text{for } u > 0,\\
    %     -(1+e^{2\pi (u+i\tau)}) \sum_{m=0}^{\infty} e^{2\pi m (u+i\tau)}, & \text{for } u < 0,
    % \end{cases}\\
    =& \begin{cases}
        1+\sum_{m=1}^{\infty} 2e^{-2\pi m (u+i\tau)}, & \text{for } u > 0,\\
        -(1+\sum_{m=1}^{\infty} 2e^{2\pi m (u+i\tau)}), & \text{for } u < 0,
    \end{cases}\\
    =& \;\text{sign}(u) \left(1+ \sum_{m \neq 0 } \theta_{\text{step}}(-mu)2 e^{2\pi m u} e^{i2\pi m \tau}\right)
\end{align}
Thus, we have
\begin{align}
    &\int_0^1 \, d\tau \; \left[\Omega^{\Delta-2}(-i\tau,X) -1\right] \coth(\pi(u+i\tau)) \nonumber\\
    =&\; \text{sign}(u)2\epsilon\left[-\ln\frac{1+|X|}{2}+\sum_{n\neq0} \theta_{\text{step}}(nu) e^{-2\pi n u} \frac{1}{|n|}\left(-\frac{1-|X|}{1+|X|}\right)^{|n|} \right] +O(\epsilon^2)\\
    =&\; -\text{sign}(u)2\epsilon \ln\left(\frac{1+|X| +(1-|X|)e^{-2\pi|u|}}{2}\right) +O(\epsilon^2).
\end{align}
The integral becomes
\begin{align}
    Q_{\text{nonloc}}(u,X;s) =&\; Q^{(1)}_{\text{nonloc}}(u,X) + O(\epsilon^2)\\
    Q^{(1)}_{\text{nonloc}}(u,X;s) =&\; -\epsilon\left\{ \text{sign}(u+s)\ln\left(\frac{1+|X| +(1-|X|)e^{-2\pi|u+s|}}{2}\right) \right. \nonumber\\
    & \left.- \text{sign}(u)\ln\left(\frac{1+|X| +(1-|X|)e^{-2\pi|u|}}{2}\right) \right\}.
\end{align}
When we evaluate the integral of $u$, we can break it into $\int_{-\infty}^{\text{min}(-s,0)}$, $\int_{\text{min}(-s,0)}^{\text{max}(-s,0)}$, and $\int_{\text{max}(-s,0)}^{\infty}$. For $u<\text{min}(-s,0)$, 
\begin{gather}
    Q_{\text{nonloc}}(u,X;s) = \epsilon\ln\left( \frac{1+|X| +(1-|X|)e^{2\pi(u+s)}}{1+|X| +(1-|X|)e^{2\pi u}}\right)
\end{gather}
For $\text{max}(-s,0)<u$,
\begin{align}    
    Q_{\text{nonloc}}(u,X;s) &= \epsilon\ln\left( \frac{1+|X| +(1-|X|)e^{-2\pi u}}{1+|X| +(1-|X|)e^{-2\pi(u+s)}}\right) \\
    &= \epsilon\ln\left( \frac{1-|X| + (1+|X|)e^{2\pi u} }{1-|X| + (1+|X|)e^{2\pi(u+s)} }\right) +\epsilon2\pi s.
\end{align}
For $\text{min}(-s,0)<u<\text{max}(-s,0)$, we have
\begin{align}
    \text{sign}(-s)\epsilon\ln\left(({1+|X| +(1-|X|)e^{-2\pi|u+s|}})({1+|X| +(1-|X|)e^{-2\pi|u|}})\right).
\end{align}
Notice that $Q^{(1)}_{\text{loc}}(u,X;s)$ has support only for $\text{min}(-s,0)<u<\text{max}(-s,0)$.
%{(See if the symmetrization helps to rewrite the log - log function)}
% Thus, the cocycle is 
% \begin{align}
%     u_{s,\lambda}=&\; 1 + i\pi s\lambda\int_0^1 dw  \int_{-1}^1 dX \, (1-X^2)e^{-iswK_0} O(X) e^{iswK_0} \nonumber\\
%     &+ i\pi s\lambda\int_0^1 dw  \int_{-1}^1 dX \int_{-\infty}^{\infty} du \, h_{reg}(u+sw,X) \Omega^{-\Delta}(u,X)O(t_u,x_u)\\
%     =&\; 1 + i\pi s\lambda \int_{-1}^1 dX \int_{-\infty}^{\infty} du \, Q(u,X) \Omega^{-\Delta}(u,X) O(t_u,x_u) +O(\lambda^2),
% \end{align}
% where 
% \begin{gather}
%     Q(u,X) =\; Q_{\text{nonloc}}(u,X) + Q_{\text{loc}}(u,X), \\
%     Q_{\text{loc}}(u,X) \equiv \; (1-X^2)\frac{1}{s} \textbf{1}_{(-s,0)}(u),
% \end{gather}
% and $\textbf{1}_{A}(x)$ is the indicator function of $A$. In the case where the perturbed QFT is another CFT ($\Delta = 2$), the non-local term $Q_{\text{nonloc}}$ vanishes. The local term $Q_\text{loc}$ combined with the original modular Hamiltonian gives the same geometric flow in the new CFT. 
To summarize the result, we have the following lemma.
\begin{lemma} [Perturbed Modular Flow of the Unit Ball]
    In 1+1D QFT with Lagrangian
    \begin{align}
        \mathcal{L}_\text{QFT} = \mathcal{L}_\text{CFT} - \lambda O(x),  
    \end{align}
    the modular flow (restricted to the algebra) of the unit ball centered at the origin is:
    \begin{gather}
        e^{is K_\lambda } = e^{is K_0} \left(1 + i\pi \lambda \int_{-1}^1 dX \int_{-\infty}^{\infty} du \, (1-X^2) Q(u,X;s) \Omega^{-\Delta}(u,X) O(t_u,x_u) +O(\lambda^2)\right),
    \end{gather}
    where
    \begin{gather}
        Q(u,X;s) =\; Q_{\text{loc}}(u,X;s) + Q_{\text{nonloc}}(u,X;s), \quad Q_{\text{loc}}(u,X) = \textbf{1}_{(-s,0)}(u),
    \end{gather}
    and 
    \begin{align}
        Q_{\text{nonloc}}(u,X;s) =\; -\pi\int_0^s dw\int_0^1 \, d\tau  \;\frac{\left(1 - (1-X^2)\sin^2(\pi \tau)\right)^{\Delta-2} -1}{2\sinh^2(\pi (u+w+i\tau))}
    \end{align}
    For $3/2<\Delta <2$, we have
    \begin{align}
        Q_{\text{nonloc}}(u,X;s) =&\; \int_0^1 \, d\tau \; \left[\left(1 - (1-X^2)\sin^2(\pi \tau)\right)^{\Delta-2} -1\right] \nonumber\\
        &\times\frac{\coth(\pi(u+s+i\tau)) - \coth(\pi(u+i\tau))}{2}.
    \end{align}
    For almost-marginal relevant operators with small $\ep = 2- \Delta \ll 1$, we have
    \begin{align}
        Q_{\text{nonloc}}(u,X;s) =&\; Q^{(1)}_{\text{nonloc}}(u,X) + O(\epsilon^2)\\
        Q^{(1)}_{\text{nonloc}}(u,X;s) =&\; -\epsilon\left\{ \text{sign}(u+s)\ln\left(\frac{1+|X| +(1-|X|)e^{-2\pi|u+s|}}{2}\right) \right. \nonumber\\
        & \left.- \text{sign}(u)\ln\left(\frac{1+|X| +(1-|X|)e^{-2\pi|u|}}{2}\right) \right\}.
    \end{align}
\end{lemma}
%\KL{(Pull the $1-X^2$, the $\frac{1+|X| +(1-|X|)e^{...}}{1+|X| +...}$ take a familiar form like modular flow, push the $s$ inside $w$? or somewhere. It is like 4 regims of the sign(u), sign(u+s), like the adding two modular flow, and how the vector field align)}

% ========== DON'T USE THE REST HERE ===========

% The action of the modular flow is
% \begin{align}
%     &e^{-isK_0}e^{isK_\lambda} \phi(t,x) e^{-isK_\lambda}e^{isK_0} \nonumber\\
%     \approx& \, \phi(t,x) + e^{-isK_0}\left[\delta_\lambda(e^{isK_\lambda}) \phi(t,x)e^{-isK_0} + e^{isK_0}\phi(t,x) \delta_\lambda(e^{-isK_\lambda}) \right]e^{isK_0},\\
%     %=& \, \phi(t,x) + e^{-isK_0}\left[\int_0^1 dw e^{iswK_0} is\delta_\lambda K_\lambda e^{-iswK_0} e^{isK_0}\phi(t,x)e^{-isK_0} - e^{isK_0}\phi(t,x)e^{-isK_0} \int_0^1 dw e^{is(1-w)K_0} is\delta_\lambda K_\lambda e^{-is(1-w)K_0} \right]e^{isK_0},\\
%     =& \, \phi(t,x) + is\int_0^1 dw \left[e^{-iswK_0} \delta_\lambda K_\lambda e^{iswK_0} ,\phi(t,x)\right], \\
%     =& \, \phi(t,x) + i\pi^2 \lambda s \int_0^1 dw\int_{-1}^1 dX \int_{-\infty}^{\infty} du  \,h(u+sw,X)\left[  e^{iuK_0} O(t=0,x=X) e^{-iuK_0} ,\phi(t,x)\right],\\
%     =& \, \phi(t,x) + i\pi^2 \lambda s \int_0^1 dw\int_{-1}^1 dX \int_{-\infty}^{\infty} du  \,h(u+sw,X) \Omega^{\Delta}(u,X) \left[ O(t_u,x_u) ,\phi(t,x)\right],\\
%     %=& \, \phi(t,x) + i\pi^2 \lambda s \int_{-1}^1 dX \int_{-\infty}^{\infty} du \, \tilde{h}(u,X;s)\left[  e^{iuK_0} O(0,X) e^{-iuK_0} ,\phi(t,x)\right],.
% \end{align}
% We recall the expression of $h(u+sw,X)$. We can perform the $\tau$-integral if $\Delta = 2$
% \begin{align}
%     \tilde{h}(u,X;s)=& \, \int_0^1 dw\, h(u+sw,X),\\
%     =& \, \int_0^1 dw\int_0^1 d\tau \, \frac{\pi^2}{\sinh^2(\pi (u+sw+i\tau))} \Omega^{\Delta - D}(\tau,X), \\
%     =& \, \int_0^1 d\tau \, \frac{1}{s}\Omega^{\Delta - D}(\tau,X) \left[ \coth(\pi (u+s+i\tau)) - \coth(\pi (u+i\tau))\right].
% \end{align}

% We estimate the difference between the two flows by
% \begin{align}
%     \lVert(e^{isK_\lambda} - e^{isK_0}) A\Omega_\lambda \rVert^2 = \lVert(e^{-isK_0}e^{isK_\lambda} - 1) A\Omega_\lambda \rVert^2 = 2 \text{Re}\braket{\Omega_\lambda|A(1-u_{s,\lambda})A|\Omega_\lambda},
% \end{align}
% where $A=A^\dagger \in \mathcal{A}$, $u_{s,\lambda} = e^{-isK_0}e^{isK_\lambda} \approx 1 + e^{-isK_0}\delta e^{isK_\lambda}$, and $\Omega_\lambda$ is the cyclic and separating QFT vacuum vector. In the first order perturbation, we have
% \begin{align}
%     u_{s,\lambda} &\approx 1 + is\lambda\pi^2 \int_{-1}^1 dx \int_{-\infty}^{\infty} du \left[\int_0^1 dw \, h(u+s-sw,X)\right] e^{iuK_0} O(0,x)  e^{-iuK_0}\\
%     &=1+ is\lambda\pi^2 \int_{-1}^1 dx \int_{-\infty}^{\infty} du \left[\int_0^1 dw \, h(u+s-sw,X) \tilde{\Omega}^{-\Delta}(u,x) \right] O(t_u,x_u) .
% \end{align}
% We then consider the case where $A = \int_{B}dX f(X) O(X)$, $\lVert(e^{isK_\lambda} - e^{isK_0}) A\Omega_\lambda \rVert^2$ becomes a 3-point function \KL{(The is also wrong, the leading order is a 4 pt function agian. The reason is that the Re will kill this term as it is purely imaginary)}
% \begin{align}
%     &\lVert(e^{isK_\lambda} - e^{isK_0}) A\Omega_\lambda \rVert^2\\ 
%     \approx&  -2 \text{Re } is\lambda\pi^2 \int_{B}dX_1\int_{B}dX_2\int_{-1}^1 dx \int_{-\infty}^{\infty} du \left[\int_0^1 dw \, h(u+s-sw,x) \tilde{\Omega}^{-\Delta}(u,x) f(X_1)f(X_2) \right] \nonumber\\
%     &\times \braket{\Omega_0|O(X_1) O(t_u,x_u) O(X_2)|\Omega_0}
% \end{align}

% \KL{Comment: Maybe should use $\Omega_0$ instead so that we can compute the 2/3-pt function.}

% Another quantity to estimate the difference is
% \begin{align}
%     &\lVert(e^{isK_\lambda}Ae^{-isK_\lambda} - e^{isK_0} A e^{-isK_0})\Omega_0 \rVert^2 \\
%     =&\, \lVert [A, u_{s,\lambda}^\dagger]\Omega_0 \rVert^2 \\
%     =&\, \braket{\Omega_0|A^2|\Omega_0} + \braket{\Omega_0|u_{s,\lambda} A^2 u_{s,\lambda}^\dagger|\Omega_0} - 2\text{Re }\braket{\Omega_0|A u_{s,\lambda} A u_{s,\lambda}^\dagger|\Omega_0}\\
%     \to &\braket{\Omega_0|[A^2, \lambda\delta u^\dagger]|\Omega_0} - 2\text{Re }\braket{\Omega_0|A[A, \lambda \delta u^\dagger]|\Omega_0}
% \end{align}
% \KL{The leading order} \KL{(The is also wrong, the leading order is a 4 pt function agian. The reason is that the Re will kill this term as it is purely imaginary)}
% Although they are 4-point functions, the first-order corrections are still 3-point functions, the same as the previous quantity. Since we are computing a commutator between $A$ and $u_{s,\lambda}^\dagger$, we can project out the spacelike separating part of $u_{s,\lambda}^\dagger$ relative to $A$. 

% \KL{Reminder: Check the order of integration issue, compare with the previous paper/interactive picture calculation, do the small s expansion first/order of scales, etc.}

% \KL{The leading order term of norm should be in $O(\lambda)$, thus the norm square is $O(\lambda^2)$. With two $A$ operators, the two power of $\lambda$ require two $O$ insertion, which results a 4-pt function in general.}

% \KL{=============== END =====================}



\subsection{First order Correction of Characteristic Operator}
With the perturbed modular flow, we can compute the characteristic operator with the conformal perturbation method. We compute $f(z)= \;_\lambda\bra{O_1(x)}\Delta_{A,\lambda}^{1/2+iz}\Delta_{B,\lambda}^{-iz}\ket{O_2(y)}_\lambda$ in a QFT parametrized by $\lambda$. If both $A$ and $B$ are wedge, the same result follows from the CFT case because the modular flows remain a boost. We consider the inclusion $B\subset A$, where $B$ is a unit double cone centered at the origin and $A$ is the wedge with the end point $(0,-x_w)$, $x_w >0$. We have a similar setup where $y\in B$, $x\in A$, and $O=O^\dagger$. The von Neumann algebras associated to the same region but with different $\lambda$ are in general different. We assume that the operators remain in both the original algebras in the CFT and the new algebras of the QFT. We denote the primary operator of the perturbation to be $O_0$. We expand the function perturbatively as
\begin{align}
    f_\lambda(z) =&\;_\lambda\bra{O_1(x)}\Delta_{A,\lambda}^{1/2+iz}\Delta_{B,\lambda}^{-iz}\ket{O_2(y)}_\lambda \\
    =&\; _0\bra{O_1(x)}\Delta_{A,0}^{1/2+iz}\Delta_{B,0}^{-iz}\ket{O_2(y)}_0 + \lambda f^{(1)}(z) + O(\lambda^2).
\end{align}
There are three contributions to the first order correction. We have
\begin{gather}
    u_{s,\lambda} = e^{-isK_0}e^{isK_\lambda} = 1 + \lambda u_{s}^{(1)} +O(\lambda^2),\\
    u_{s}^{(1)} = i\pi  \int_{-1}^1 dX \int_{-\infty}^{\infty} du \, (1-X^2) Q(u,X;s) \Omega^{-\Delta}(u,X) O_0(t_u,x_u)\\
    \ket{\Omega_\lambda} = \ket{\Omega_0} + \lambda\ket{\Omega^{(1)}} + O(\lambda^2),\\
    \ket{\Omega^{(1)}} = \int_\mathbb{R} dX \int_{-\infty}^{0} d\tau O_0(-i\tau,X) \ket{\Omega_0}.
\end{gather}
We can rewrite $f_\lambda(z)$ in terms of $u_{s,\lambda}$ and $\Omega_\lambda$. We have
\begin{align}
    f_\lambda(z) =&\;_\lambda\bra{O(x)}\Delta_{A,\lambda}^{1/2+iz}\Delta_{B,\lambda}^{-iz}\ket{O(y)}_\lambda \\
    =&\;\bra{\Omega_\lambda}O(x) \Delta_{A,\lambda}^{1/2+iz}\Delta_{B,\lambda}^{-iz}O(y) \ket{\Omega_\lambda}\\
    =& \; \bra{ \Omega_\lambda} \alpha_{A}^{-z}\left(O(x_J)\right) \left( \Delta_{B,0}^{-iz}\Delta_{B,0}^{+iz} \Delta_{B,\lambda}^{-iz}O(y) \Delta_{B,\lambda}^{+iz}\Delta_{B,0}^{-iz}\Delta_{B,0}^{+iz}\right)\ket{\Omega_\lambda}\\
    =& \; \bra{\Omega_\lambda}\alpha_{A}^{-z}\left(O(x_J)\right) \Delta_{B,0}^{-iz}\left(u_{s,\lambda}O(y) u_{z,\lambda}^\dagger\right) \Delta_{B,0}^{+iz}\ket{\Omega_\lambda},
\end{align}
where $x_J$ is the RT conjugation of $x$ and $\alpha_{A}^{z}$ ($\alpha_{B}^{z}$) is the geometric modular flow action of the region $A$ ($B$). Therefore, the three correction terms are
\begin{align}
    f^{(1)}(z) =&\; f^{(1)}_{\text{flow}}(z) + f^{(1)}_{\text{bra}}(z) + f^{(1)}_{\text{ket}}(z),\\
    f^{(1)}_{\text{flow}}(z) \equiv&\; \bra{\Omega_0} \alpha_{A}^{-z}\left(O(x_J)\right) \alpha_{B}^{-z}\left(\left[u_{z}^{(1)},O(y)\right]\right) \ket{\Omega_0}, \\
    =&\; i\pi  \int_{-1}^1 dX \int_{-\infty}^{\infty} du \, (1-X^2)Q(u,X;s) \Omega^{-\Delta}(u,X) \nonumber\\
    &\times\bra{\Omega_0} \alpha_{B}^{z} \circ \alpha_{A}^{-z}\left(O_1(x_J)\right) \left[O_0(t_u,X_u),O_2(y)\right] \ket{\Omega_0} \\
    f^{(1)}_{\text{bra}}(z) \equiv&\; \int_\mathbb{R} dX \int_{-\infty}^{0} d\tau \; \bra{\Omega_0} O_0(-i\tau,X)^\dagger  \alpha_{A}^{-z}\left(O_1(x_J)\right) \alpha_{B}^{-z}\left(O_2(y)\right) \ket{\Omega_0},\\
    f^{(1)}_{\text{ket}}(z) \equiv&\; \int_\mathbb{R} dX \int_{-\infty}^{0} d\tau \; \bra{\Omega_0} \alpha_{A}^{-z}\left(O_1(x_J)\right) \alpha_{B}^{-z}\left(O_2(y)\right) O_0(-i\tau,X) \ket{\Omega_0}.
\end{align}
All the first order correction terms are 3-point function evaluated in the original CFT, which is universal. The three-point function in CFT is \cite{}
\begin{align}
    \braket{O_1(x_1) O_2(x_2) O_3(x_3)} = \frac{C_{123}}{(x_{12}^2)^{(\Delta_1+\Delta_2-\Delta_3)/2}(x_{23}^2)^{(\Delta_2+\Delta_3-\Delta_1)/2}(x_{31}^2)^{(\Delta_3+\Delta_1-\Delta_2)/2}},    
\end{align}
where $C_{123}$ is some constant, $\Delta_i$ is the conformal dimension of $O_i$, and $x_{ij}^2$ is the squared proper distance with $i\epsilon$ prescription
\begin{gather}
    x_{ij}^2 = -(t_i-t_j-i(\epsilon_i-\epsilon_j))^2 + (X_i-X_j)^2, \quad \epsilon_1 > \epsilon_2 > \epsilon_3 \to 0^+.
\end{gather}
Since we are evaluating the modular flow in the original CFT, all the actions are geometrical. Schematically, we write the correction terms as
\begin{align}
    f^{(1)}_{\text{flow}}(z) =&\; i\pi  \int_{-1}^1 dX \int_{-\infty}^{\infty} du \, (1-X^2)Q(u,X;z) \Omega^{-\Delta}(u,X) \;\gamma^{\Delta_1} \nonumber\\
    &\times\left(\frac{C_{102}}{(x_{10}^2)^{\Delta_{10}}(x_{02}^2)^{\Delta_{02}}(x_{21}^2)^{\Delta_{21}}} -\frac{C_{120}}{(x_{12}^2)^{\Delta_{12}}(x_{20}^2)^{\Delta_{20}}(x_{01}^2)^{\Delta_{01}}}\right) \\
    f^{(1)}_{\text{bra}}(z) \equiv&\; \int_\mathbb{R} dX \int_{-\infty}^{0} d\tau \;\gamma^{\Delta_2} \frac{C_{012}}{(x_{01}^2)^{\Delta_{01}}(x_{12}^2)^{\Delta_{12}}(x_{20}^2)^{\Delta_{20}}},\\
    f^{(1)}_{\text{ket}}(z) \equiv&\; \int_\mathbb{R} dX \int_{-\infty}^{0} d\tau \; \gamma^{\Delta_2}\frac{C_{120}}{(x_{12}^2)^{\Delta_{12}}(x_{20}^2)^{\Delta_{20}}(x_{01}^2)^{\Delta_{01}}},
\end{align}
where $\gamma^{\Delta_i}$ is the conformal factor from the conformal transformation acting on $O_i$, $\Delta_{ij} = \frac{1}{2}(\Delta_i+\Delta_j - \Delta_k$ where $k$ is the rest of the index. We use $x_{ij}$ to schematically denote the distance between the transformed $O_i$ and $O_j$ with the $i\epsilon$ prescription mentioned above. For a clearer presentation, we keep all $z$ dependence and other parameters implicit for $x_{ij}$ and $\gamma$. 

We first consider $f^{(1)}_{\text{flow}}(z)$. Notice that we have pulled all the modular flow actions to $O_1$ so that $[O_0,O_2]$ is independent of $z$. This term only has contributions from $O_0$ that is not spacelike separated from $O_2$. By pushing $O_2$ towards the edge of the double cone, the region in the double cone that is not spacelike separated from $O_2$ shrinks. This seems to suggest a connection to how the modular flow may become approximately geometric, in a similar spirit to the Fredenhagen's statement \cite{fredenhagen}. 
%\subsection{Analytic Properties of Characteristic Function in QFT}

Another observation is that all three-point functions can be analytically extended as long as the proper distance between any two points is non-zero. The three-point function has pole/branch points only if one of the $x_{ij}^2$ becomes zero. For all three correction terms, $x_{12}^2$ contributes poles to the first order correction in the same way as to the zeroth order term, except the power is different. $f^{(0)}(z) \sim x_{12}^{\Delta}$ while $f^{(1)}(z) \sim x_{12}^{\Delta_{12}}$. For $x_{10}^2 $ or $ x_{20}^2$ being zero, we will have $z$ as a function of the integration variable $X$ and $u$. It is easy to see that it is possible to have $x_{10}^2 $ or $ x_{20}^2$ being zero even inside the strip Im$(z) \in (-1/4,1/4)$. However, we also have to consider how the integration kernel may smear out the singularity. In addition, we shall use smeared operators $O_1(g_1)$ and $O_2(g_2)$ with some well-behaved smearing function $g_1,g_2$ instead of point-like fields $O_1(x), O_2(y)$. These may further regulate the behaviour around $x_{10}^2 = 0$ and $x_{20}^2=0$. Even without the exact details of the smearing process, the analytic property of the characteristic operator ensures that the function is analytic in the strip Im$(z) \in (-1/4,1/4)$. We leave the question of how these seeming singularities are smoothed out to future study. 

% For a fixed $O_1$ (or $O_2$) with an $O_0$ that is null-like separated, 
% we consider the smearing of $O_0$ along $du$. We have
% \begin{align}
%       \int_{u_0-\epsilon}^{u_0+\epsilon} du \, Q(u,X;z) \Omega^{-\Delta}(u,X) \frac{1}{(x_{10}^2)^{\Delta_{10}}}\sim \int_{u_0-\epsilon}^{u_0+\epsilon} du \frac{C}{((u-u_0)^2\pm i\epsilon')^{\Delta_{01}}},
% \end{align}
% where $u_0$ is when $O_0(t_u,x_u)$ becomes null-like separated from $O_1$ with a given $X$, and $C$ is some constant. 





\section{Nonperturbative bounds on subregion flow of balls in QFT}\label{section:approx}
%1. As non-perturbative result, Reformulation of Fredenhagen, instead of smaller region, larger ball approx wedge 2. Wedge in Wedge known, ball in wedge new result, cutoff, 3. Ball inside ball? Ball 1 in wedge + Ball 2 in wedge? Two states

In this section, we consider a situation where characteristic functions of certain inclusions can be well-approximated by characteristic functions of wedge inclusions. The result is derived directly from the well-known result of Fredenhagen \cite{fredenhagen}.

Consider a net of algebras of local observables on $\{\mathcal{A}(\mathcal{O})\}_{\mathcal{O}\subset\mathbb{R}^{d,1}}$. Assume this net satisfies the Haag-Kastler axioms. Let $\ket{\Omega}$ be the vacuum vector of this net. By the Reeh-Schlieder theorem, the vacuum vector is a common cyclic and separating vector for all algebras in this net. In addition, by the axiom, this net of algebras is covariant under the Poincaré group and the group action is strongly continuous.

Consider two right wedges:
\begin{equation}
    W_0 := \{x\in\mathbb{R}^{d,1}: |x_0| < x_1\}
\end{equation}
\begin{equation}
    W_{r} := \{x\in\mathbb{R}^{d,1}: |x_0 | < x_1 - r \}
\end{equation}
The two algebras of local observables associated with the two wedges are related by the action of right translation:
\begin{equation}
    Ad_{T_r}(\mathcal{A}(W_0)) = \mathcal{A}(W_r)
\end{equation}
where $T_r$ is right translation by $r$ and $Ad_{T_r}$ is the adjoint action. Let $\{\Delta_{W_0}^{it}\}$ be the modular automorphism of $\mathcal{A}(W_0)$ under the vacuum state and let $\{\Delta_{W_r}^{it}\}$ be the modular automorphism of $\mathcal{A}(W_r)$ under the same vacuum state. The two modular automorphisms are related by:
\begin{equation}
    T_r\Delta_{W_0}^{it}T_r^* = T_r\Delta_{W_0}^{it}T_{-r} = \Delta^{it}_{W_r}
\end{equation}
where $T_r^* = T_{-r}$ is left translation by $r$.

It is well-known \cite{bisognanowichmann}that the modular flow $\Delta_{W_0}^{it}$ associated with a wedge is geometric. In particular, this means that For our purpose, we need to consider several double cones in the right wedge $W_0$. All double cones intersect the spatial slice $\{x\in\mathbb{R}^{d,1}: x_0 = 0\}$ and the intersections are balls. To specify such double cones, it suffices to specify the intersecting balls: 
\begin{equation}
    B_{c,r}\subset W_0
\end{equation}
where $B_{c,r}$ denotes a generic double cone in the right wedge. The intersecting ball $B_{c,r}\mathcal{A}p\{x\in\mathbb{R}^{d,1}:x_0 = 0\}$ is centered at $x_1 = c\text{ , }x_i = 0$ (for $i\neq 0,1$) and has a radius $r$. Notice that for the double cone to be included in the right wedge, $0 < r \leq c$. In case where $r = c$, the vertex of the right wedge $W_0$ is in the closure of the double cone. To simplify notations, we collect the following families of examples and introduce specialized notations:
\begin{enumerate}
    \item $B_R:=B_{R, R}\subset W_0$ where the intersecting ball $B_R\mathcal{A}p \{x\in\mathbb{R}^{d,1}: x_0 = 0\}$ is centered at $x_1 = R, x_i = 0$ (for $i\neq 0, 1$) and has a radius $R$;
    \item $B_{e^{-2\pi s}R} := B_{e^{-2\pi s}R, R}\subset W_0$ where the intersecting balls $B_{e^{-2\pi s}R}\mathcal{A}p\{x\in\mathbb{R}^{d,1}:x_0 =0\}$ is centered at $x_1 = R, x_i = 0$ (for $i\neq 0, 1$) and has a radius $e^{-2\pi s}R$. Here $s$ is a positive real parameter. These are the double cones that appear in Fredenhagen's original discussions \cite{fredenhagen}. Notice that the vertex of the right wedge $W_0$ is no longer contained in the closure of $B_{e^{-2\pi s}R}$. Instead, we consider the shifted right wedge $W_{r(s)}:= \{x\in\mathbb{R}^{d,1}: x_1 - r(s) > |x_0|\}$ where $r(s) := R - e^{-2\pi s}R$. The vertex of the wedge $W_{r(s)}$ is in the closure of the double cone $B_{e^{-2\pi s}R}$;
    \item $B_{e^{2\pi s}R}:= B_{e^{2\pi s}R, e^{2\pi s}R}\subset W_0$ where the center of the intersecting ball is moving away from the origin as the radius of the ball increases so that the vertex of $W_0$ is in the closure of the double cone. These are dilated double cone and as $s\rightarrow \infty$, these double cones would approach the right wedge $W_0$.
\end{enumerate}
In addition to these double cones in the right Rindler wedge, we would also need to consider translated double cones contained in translated wedges. We introduce the following notations for these translated wedges:
\begin{equation}
    B^s_{c,r} := T_s(B_{c,r})\subset W_s  = T_s(W_0)
\end{equation}
where $T_s$ is the right translation by $s$ along the $x_1$-direction. The intersection of the translated double cone $B^s_{r,c}$ and the spatial slice $\{x\in\mathbb{R}^{d,1}:x_0 = 0\}$ is a ball centered at $x_1 = c+ s\text{ , }x_i = 0$ (for $i\neq 0, 1$) with radius $r$.
\subsection{Review of Fredenhagen's Results}
In this subsection, we briefly review Fredenhagen's results in \cite{fredenhagen} and modify the results to suit our discussion. First recall the following main theorem:
\begin{theorem}
    Let $W_0$ be the right wedge with vertex at the origin and let $B_R\subset W_0$ be a double cone. For any $f\in L^1(\mathbb{R})$ and any constant $s > 0$, there exists a constant $c_f(s) > 0$ such that:
    \begin{equation}
        ||\int_\mathbb{R} dt f(t)(\Delta_{B_R}^{-it} - \Delta_{W_0}^{-it})A\ket{\Omega}||^2 \leq c_f(s)(||A\ket{\Omega}||^2 + ||A^*\ket{\Omega}||^2)
    \end{equation}
    where $A$ is any operator in the local algebra $\mathcal{A}(B_{e^{-2\pi s}R})$. The constant $c_f(s)$ approaches $0$ as $s\rightarrow \infty$.
\end{theorem}
Heuristically, Fredenhagen's result states that, for operators localized in an exponentially small region close to the vertex of the wedge $W_0$, the averaged modular evolution (smeared out by the function $f$ in time) can be well approximated by the (geometric) modular evolution of the wedge. The theorem verifies the intuition that the error of this approximation should become smaller as the support of the local operators shrinks (i.e. $s\rightarrow\infty$).

If the smearing function $f$ has finite support (symmetric around the origin), we would expect the approximation to hold for longer periods of time as the support of the local operators shrinks. This intuition can be qualitatively verified by studying the constant $c_f(s)$ for the case where $f$ is Gaussian. We would turn to the detailed calculation in Subsection \ref{subsection:gaussiansmear}.

Back to the proof of Fredenhagen's main result, it is instructive to recall the following key technical result \cite{fredenhagen}. Notice that the following result holds for any two generic von Neumann algebras with a common cyclic and separating vector:
\begin{proposition}
    Let $\mathcal{B}\subset \mathcal{A}$ be two von Neumann algebras with a common cyclic and separating vector $\ket{\Omega}$. For all $\mathcal{O}\in\mathcal{B}$ such that for all $|t| < t_0$ (where $t_0 > 0$ is a positive constant), we have:
    \begin{equation}
        \Delta_\mathcal{A}^{it}\mathcal{O}\Delta_\mathcal{A}^{-it}\in\mathcal{B}
    \end{equation}
    where $\Delta_\mathcal{A}$ is the modular operator associated with $\ket{\Omega}$ on $\mathcal{A}$, then the following resolvent estimate holds:
    \begin{equation}
        ||\lb(1 + \Delta_\mathcal{A})^{-1} - (1 + \Delta_\mathcal{B}^{-1})\rb\mathcal{O}\ket{\Omega}|| \leq \frac{2}{\pi}\tanh^{-1}(e^{-\pi t_0})(||\mathcal{O}\ket{\Omega}||^2 + ||\mathcal{O}^*\ket{\Omega}||^2)^{1/2}
    \end{equation}
\end{proposition}
In Fredenhagen's original theorem, the approximation is tighter as the support of local operators shrinks to zero. For our discussion, the size of the support is not a free parameter. Rather, we need to consider the opposite scenario where the support of the operator is fixed but we are embedding the operator in larger local algebras. More precisely, we want to show the following corollary to the Proposition above:
\begin{corollary}
    Consider two double cones inside the right Rindler wedge: $B_R\subset B_{e^{2\pi s}R}\subset W_0$. For any operator $\mathcal{O}\in\mathcal{A}(B)$, the following resolvent estimate holds:
    \begin{equation}
        ||\lb(1 + \Delta_{B_{e^{2\pi s}R}})^{-1} - (1 + \Delta_{W_0})^{-1}\rb\mathcal{O}\ket{\Omega}||\leq \frac{2}{\pi}\tanh^{-1}(e^{-\pi s})\lb||\mathcal{O}\ket{\Omega}||^2 + ||\mathcal{O}^*\ket{\Omega}||^2\rb^{1/2}
    \end{equation}
    where $\ket{\Omega}$ is a common cyclic and separating vector. $\Delta_{B_{e^{2\pi s}R}}$ is the modular operator of the dilated double cone $e^{2\pi s}B$ and $\Delta_{W_0}$ is the modular operator of the right Rindler wedge.
\end{corollary}
\begin{proof}
    Using the fact that the modular flow of the right Rindler wedge is the boost, it is elementary to observe that for all time $|t|\leq s$ and for all $\mathcal{O}\in\mathcal{A}(B)$, we have:
    \begin{equation}
        \Delta_{W_0}^{it}\mathcal{O}\Delta^{-it}_{W_0} \in \mathcal{A}(B_{e^{2\pi s}R})
    \end{equation}
    The statement follows directly from the Proposition.
\end{proof}
From this Corollary and using the same proof as Fredenhagen \cite{fredenhagen}, we immediately arrive at the following theorem:
\begin{theorem}\label{theorem:modifiedfredenhagen}
    Using the same set-up as the previous Corollary, for any operator $\mathcal{O}\in \mathcal{A}(B)$ and for any $f\in L^1(\mathbb{R})$ and for any constant $s > 0$, there exists a constant $c_f(s) > 0$ such that:
    \begin{equation}
        ||\int_\mathbb{R} dt f(t)\lb\Delta^{-it}_{B_{e^{2\pi s}R}} - \Delta^{-it}_{W_0}\rb\mathcal{O}\ket{\Omega}||^2 \leq c_f(s)\lb||\mathcal{O}\ket{\Omega}||^2 + ||\mathcal{O}^*\ket{\Omega}||^2\rb
    \end{equation}
    In addition, $\lim_{s\rightarrow \infty}c_f(s) = 0$
\end{theorem}
\begin{proof}
    Notice that we have:
    \begin{equation}
        \lim_{s\rightarrow\infty}\tanh^{-1}(e^{-\pi s}) = 0
    \end{equation}
    Then the theorem follows directly from Fredenhagen's proof and the previous Corollary.
\end{proof}
Notice that this modified theorem states that the averaged modular flow  of a dilating double cone is close to that of the right wedge. The approximation is in the $L^2$-sense as specified in the Theorem above. This result matches the intuition since the dilating double cone approaches the right wedge.


\subsection{Warm-up calculation in 1+1d CFT}
We consider the setup in a 1+1d CFT with $R = e^{2\pi s}$ and compute an asymptotic upper bound of $||\int_\mathbb{R} dt f(t)(\Delta_{B_R}^{-it} - \Delta_{W_0}^{-it})A\ket{\Omega}||^2$ with $f_{\sigma,\mu}(t)$ the normalized Gaussian distribution. For simplicity, We choose self-adjoint $A = \mathcal{O}(g) = \int dtdx \;g(t,x) \mathcal{O}(t,x)$ with a well-behave smearing function $g(t,x)$ sported in $B_{r}$. From the expression in section \ref{}, we have 
\begin{gather}
    \Delta_{B_R}^{-iz} \mathcal{O}(g) \Delta_{B_R}^{iz} = \mathcal{O}(\tilde{g}),\\
    \tilde{g}(X) = \left|\text{det}\frac{d\alpha_{B_R}^{-z}(X)}{dX}\right|\gamma^{2h}(\alpha_{B_R}^{-z}(X)) g(\alpha_{B_R}^{-z}(X)),
\end{gather}
where we denote $(t,x)$ as $X$ and $\alpha_{B_R}^{z}$ is the geometric action of $\text{Ad}_{ \Delta_{B_R}^{-iz}}$. The parameters of the action of $\text{Ad}_{\Delta_{B_R}^{-iz}}$ are
\begin{gather}
    t_c = 0, \qquad x_c  = R, \qquad 
    |-t^2 + (x-r)^2| < r\\
    d^2 = -t^2 +(x-R)^2 = -t^2 + (x-r)^2 +(R-r)^2 - 2(x-r)(R-r).
\end{gather}
For large $R=re^{2\pi s}\gg r$, we have the approximation
\begin{align}
    d^2 &= R^2\left[1 -2x/R +O({(r/R)}^2)\right],\\
    \gamma(t,x)&= \left(1 + \frac{x}{R}(\cosh(2\pi z) -1)+\frac{t}{R}\sinh (2\pi z) + O((r/R)^2)\right)^{-1}\nonumber\\
    &= 1 - \frac{x}{R}(\cosh(2\pi z) -1)-\frac{t}{R}\sinh (2\pi z) + O((r/R)^2),\\
    t_z &= t\cosh (2\pi z) + x\sinh (2\pi z) + r O(re^{4\pi|z|}/R), \quad \text{for } \;re^{4\pi|z|}\ll R,\\
    x_z &= t\sinh (2\pi z) + x\cosh (2\pi z) + r O(re^{2\pi|z|}/R), \quad \text{for } \;re^{2\pi|z|}\ll R.
\end{align}
For $re^{4\pi|z|}\ll R$, We have the approximation of $g(X)$ as
\begin{gather}
    g(\alpha_{B_R}^{-z}(X)) = g(\alpha_{W}^{-z}(X)) + O(re^{4\pi|z|}/R),\\
    \tilde{g}(X) = g(\alpha_{W}^{-z}(X)) + O(re^{4\pi|z|}/R).
\end{gather}
Notice that $O(re^{4\pi|z|}/R)$ is in the form of
\begin{align}
    O(re^{4\pi|z|}/R) = \sum_{n=-2}^2\frac{r}{R} e^{2n\pi z} \; a_n(X),
\end{align}
for some functions $a_n(X)$. This expression comes from the function $\sinh(2\pi z)$ and $\cosh(2\pi z)$ in $\gamma$, $t_z$, and $x_z$.

We break $||\int_\mathbb{R} dz f(t)(\Delta_{B_R}^{-iz} - \Delta_{W_0}^{-iz})A\ket{\Omega}||$ into
\begin{align}
    &||\int_\mathbb{R} dz f_{\sigma,\mu}(z)(\Delta_{B_R}^{-iz} - \Delta_{W_0}^{-iz})A\ket{\Omega}|| \nn\\
    &=\left|\left|\int_{|z|<s_0} dz f_{\sigma,\mu}(z) (\Delta_{B_R}^{-iz} - \Delta_{W_0}^{-iz})A\ket{\Omega}+\int_{|z|\geq s_0} dz f_{\sigma,\mu}(z)(\Delta_{B_R}^{-iz} - \Delta_{W_0}^{-iz})A\ket{\Omega}\right|\right|\\
    & \leq \left|\left|\int_{|z|<s_0} dz f_{\sigma,\mu}(z) (\Delta_{B_R}^{-iz} - \Delta_{W_0}^{-iz})A\ket{\Omega}\right|\right| + \left|\left| \int_{|z|\geq s_0} dz f_{\sigma,\mu}(z)(\Delta_{B_R}^{-iz} - \Delta_{W_0}^{-iz})A\ket{\Omega}\right|\right|,
\end{align}
for some $s_0$ such that $s_0\ll s$. The second term has is bounded by
\begin{align}
    \left|\left|\int_{|z|\geq s_0} dz f_{\sigma,\mu}(z)(\Delta_{B_R}^{-iz} - \Delta_{W_0}^{-iz})A\ket{\Omega}\right|\right| &\leq \int_{|z|\geq s_0} dz |f_{\sigma,\mu}(z)|\;\left|\left|(\Delta_{B_R}^{-iz} - \Delta_{W_0}^{-iz})A\ket{\Omega}\right|\right|\\
    &\leq 2||A\ket{\Omega}|| \int_{|z|\geq s_0} dz |f_{\sigma,\mu}(z)| .
\end{align}
This integral decays like a rescaled complementary error function of $s_0$. For the first term, we have $|z|< s_0\ll s$ and we can use the approximation for $re^{4\pi|z|}\ll R$. We have
\begin{align}
    &\left|\left|\int_{|z|<s_0} dz f_{\sigma,\mu}(z) (\Delta_{B_R}^{-iz} - \Delta_{W_0}^{-iz})A\ket{\Omega}\right|\right| \leq \int_{|z|<s_0} dz |f_{\sigma,\mu}(z)| \left|\left|(\Delta_{B_R}^{-iz} - \Delta_{W_0}^{-iz})A\ket{\Omega}\right|\right| \\
    &= \int_{|z|<s_0} dz |f_{\sigma,\mu}(z)| \sqrt{2||A\ket{\Omega}||^2 - 2\text{Re}\bra{\Omega}A \Delta_{W_0}^{+iz}\Delta_{B_R}^{-iz} A\ket{\Omega}}.
\end{align}
From the approximation we have
\begin{align}
    \text{Ad}_{\Delta_{W_0}^{+iz}\Delta_{B_R}^{-iz}} A &= \text{Ad}_{\Delta_{W_0}^{+iz}\Delta_{B_R}^{-iz}} \mathcal{O}(g(X))\\
    & = \text{Ad}_{\Delta_{W_0}^{+iz}} \left(\mathcal{O}(g(\alpha_{W}^{-z}(X))) + \mathcal{O}(O(re^{4\pi|z|}/R))\right)\\
    & = \mathcal{O}(g(X)) + \mathcal{O}(O(re^{4\pi|z|}/R)).
\end{align}
Thus, we can write 
\begin{align}
    2||A\ket{\Omega}||^2 - 2\bra{\Omega}A \Delta_{W_0}^{+iz}\Delta_{B_R}^{-iz} A\ket{\Omega} = \frac{r}{R} \sum_{n=-2}^2e^{2n\pi z} \bra{\Omega} \mathcal{O}(g) \tilde{\mathcal{O}}_n\ket{\Omega},
\end{align}
for some $\tilde{\mathcal{O}}_n$ that depends on the explicit form of $g(X)$. This gives 
\begin{align}
    \left|\left|\int_{|z|<s_0} dz f_{\sigma,\mu}(z) (\Delta_{B_R}^{-iz} - \Delta_{W_0}^{-iz})A\ket{\Omega}\right|\right| 
    \leq \sqrt{\frac{r}{R}}\int_{\mathbb{R}} dz |f_{\sigma,\mu}(z)| \sqrt{ \sum_{n=-2}^2e^{2n\pi z} \bra{\Omega} \mathcal{O}(g) \tilde{\mathcal{O}}_n\ket{\Omega}}.
\end{align}
$f_{\sigma,\mu}(z)$ is a Gaussian function so that the integral is a finite constant. $\sqrt{r/R}$ decays as $e^{-\pi s}$. Therefore, we have the bound for large $s$
\begin{align}
    \left|\left|\int_\mathbb{R} dz f_{\sigma,\mu}(z)(\Delta_{B_R}^{-iz} - \Delta_{W_0}^{-iz})A\ket{\Omega}\right|\right| \lesssim  Ke^{-\pi s} + 2||A\ket{\Omega}|| \int_{|z|\geq s_0} dz |f_{\sigma,\mu}(z)|,
\end{align}
for some constant $K$. We can take $s_0/s = \epsilon\ll 1$. The integral $\int_{|z|\geq \epsilon s} dz |f_{\sigma,\mu}(z)|$ decays much faster than $e^{-\pi s}$. Thus, the asymptotic behaviour of the difference is
\begin{align}
    \left|\left|\int_\mathbb{R} dz f_{\sigma,\mu}(z)(\Delta_{B_R}^{-iz} - \Delta_{W_0}^{-iz})A\ket{\Omega}\right|\right|^2 \lesssim  K'e^{-2\pi s},
\end{align}
for some constant $K'$.

The calculation for the characteristic operators is similar in this case. Without loss of generality, we compare the characteristic operator of the inclusion $W_0\subset W_{-a}$ and that of $B_R\subset W_{-a}$, with $a>0$. We have
\begin{align}
    &\left|\left|\int_\mathbb{R} dz f_{\sigma,\mu}(z)(\mathcal{D}_{W_{-a},W_0}(z) - \mathcal{D}_{W_{-a},B_R}(z))A\ket{\Omega}\right|\right| = \left|\left|\int_\mathbb{R} dz f_{\sigma,\mu}(z)\Delta_{W_{-a}}^{-iz}(\Delta_{B_R}^{iz} - \Delta_{W_0}^{iz})A\ket{\Omega}\right|\right| \\
    & \leq \left|\left|\int_{|z|<s_0} dz f_{\sigma,\mu}(z)\Delta_{W_{-a}}^{-iz}(\Delta_{B_R}^{iz} - \Delta_{W_0}^{iz})A\ket{\Omega}\right|\right| + \left|\left|\int_{|z|\geq s_0} dz f_{\sigma,\mu}(z)\Delta_{W_{-a}}^{-iz}(\Delta_{B_R}^{iz} - \Delta_{W_0}^{iz})A\ket{\Omega}\right|\right|\\
    &\leq \int_{|z|<s_0} dz |f_{\sigma,\mu}(z)|\;\left|\left|(\Delta_{B_R}^{iz} - \Delta_{W_0}^{iz})A\ket{\Omega}\right|\right| + 2 \left|\left|A\ket{\Omega}\right|\right|\int_{|z|\geq s_0} dz |f_{\sigma,\mu}(z)|.
\end{align}
This is the same expression as in the case of modular flow, but with $z\to-z$ in the first term because of the sign convention. Thus, the difference in the characteristic operator has a similar large $s$ asymptotic behaviour
\begin{align}
    \left|\left|\int_\mathbb{R} dz f_{\sigma,\mu}(z)(\mathcal{D}_{W_{-a},W_0}(z) - \mathcal{D}_{W_{-a},B_R}(z))A\ket{\Omega}\right|\right|^2 \lesssim K''e^{-2\pi s}.
\end{align}
We observe that the asymptotic behaviour decays much faster than the general bound. This is expected as the general bound uses neither the details of the QFT nor the extra symmetry in a CFT. In addition, it is straightforward to derive the bound for the characteristic operator case in this example. The reason is that we use the triangle inequality to bound the norm of an integral with an integral of the norm. In general QFT, $\mathcal{A}_{B_{R}} \subset \mathcal{A}_{B_{R'}}$ for $R<R'$ and $(\cup_R \mathcal{A}_{B_{R}})'' = \mathcal{A}_{W_0}$ implies $\text{s-}\lim_{R\to\infty} \Delta^{-iz}_{B_{R}} = \Delta^{-iz}_{W_0}$ in the strong operator topology. We have
\begin{align}
    \text{s-}\lim_{R\to \infty}\left|\left|(\Delta_{B_{R}}^{-iz} - \Delta_{W_0}^{-iz})A\ket{\Omega}\right|\right| = 0
\end{align}
for all $z$. It is a question of whether we have s-$\lim_{R\to\infty}\int_{|z|<s_0} dz |f_{\sigma,\mu}(z)|\;\left|\left|(\Delta_{B_R}^{iz} - \Delta_{W_0}^{iz})A\ket{\Omega}\right|\right| = 0$ as well. In other words, we need $\left|\left|(\Delta_{B_{R}}^{-iz} - \Delta_{W_0}^{-iz})A\ket{\Omega}\right|\right| \to 0$ uniformly for the range of $z$. In this CFT example, the integral of the norm indeed converges to zero. The calculation for the characteristic operator case is thus reduced to the same calculation as in the modular flow case. The geometric action of the modular flows allows us to bound all $\left|\left|(\Delta_{B_{R}}^{-iz} - \Delta_{W_0}^{-iz})A\ket{\Omega}\right|\right|$ in the truncated range of $z$.




\subsection{Approximation of Characteristic Functions - Part I. Double Cone Inside Wedge}\label{subsection:coneinsidewedge}
One of the main insights of algebraic quantum field theory is that the information of a field theory is encoded in the inclusions of local algebras. Hence, we are naturally more interested in the characteristic function associated with inclusions of a pair of local algebras. While Fredenhagen's results \cite{fredenhagen} are concerned with approximations of modular operators, we will focus on approximations of characteristic functions. Recall the definition of characteristic function \cite{revolutionizing}:
\begin{definition}
    Given a pair of von Neumann algebra $\mathcal{B}\subset\mathcal{A}$ and a common cyclic and separating vector $\ket{\Omega}\in\mathcal{H}$, the characteristic function is defined as:
    \begin{equation}
        \mathcal{D}_{\mathcal{A},\mathcal{B}}(t) := \Delta_\mathcal{A}^{-it}\Delta_\mathcal{B}^{it}
    \end{equation}
\end{definition}
It is well-known that the characteristic function can be analytically continued to the horizontal strip $S(0, \frac{1}{2}) := \{z\in\mathbb{C}:0\leq \Im{z} \leq \frac{1}{2}\}$ in the complex plane. In addition, on the boundary of the horizontal strip, the characteristic function is unitary and strongly continuous. Characteristic function is suitable to study the inclusion structure of nets of local algebras \cite{revolutionizing}. However, to calculate characteristic function is not straight-forward. Given a net of local algebras satisfying the Haag-Kastler axioms, one can calculate the characteristic functions of wedge inclusions (i.e. a wedge included in another wedge). This is because the modular flow associated with a wedge is geometric \cite{bisognanowichmann}. This raises the question if we can approximate other characteristic functions using those of wedge inclusions. In this subsection and the next, we explore these approximations. In this subsection, we study characteristic function of a double cone included in a wedge. In the next subsection, we study characteristic function of a double cone included in another double cone. Geometrically, as a double cone grows in size, for an observer localized in a bounded region in spacetime, the double cone would eventually be indistinguishable from a wedge. This is the underlying phyiscal intuition behind Fredenhagen's approximation of modular operators \cite{fredenhagen}\footnote{More precisely, it is the physical intuition behind our reformulation of Fredenhagen's original result.}. The same intuition can help us obtain approximation results for characteristic functions.

In this subsection, our main theorem concerns the approximation of two characteristic functions:
\begin{enumerate}
    \item One characteristic function is associated with the inclusion of a shifted wedge $W_r := T_r(W_0)$ in the standard right Rindler wedge $W_0$:
    \begin{equation}
        \mathcal{D}_{W_0, W_r}(t) = \Delta_{W_0}^{-it}\Delta_{W_r}^{it}
    \end{equation}
    Because the modular operators of the two wedges are both geometric, this characteristic function can be calculated explicitly;
    \item The other characteristic function is associated with the inclusion of a (one-parameter family of) shifted double cone $B^r_{e^{2\pi s }R}:=B^r_{e^{2\pi s}R, e^{2\pi s}R} = T_r(B_{e^{2\pi s}R, e^{2\pi s}R})$ in the standard right wedge $W_0$:
    \begin{equation}
        \mathcal{D}_{W_0, B^r_{e^{2\pi s}R}}(t) = \Delta_{W_0}^{-it}\Delta_{B^r_{e^{2\pi s}R}}^{it}
    \end{equation}
\end{enumerate}
Intuitively, as $s\rightarrow\infty$, the double cone $B^r_{e^{2\pi s}R}$ approaches the shifted right wedge $W_r$ and it is reasonable to expect $\mathcal{D}_{W_0,B^r_{e^{2\pi s}R}}$ to approach $\mathcal{D}_{W_0, W_r}$. Our main theorem formalizes this intuition.

Before stating the theorem, we first introduce some notation. Given a modular operator $\Delta$, we denote $E(\mu_0)$ be the spectral projection of the corresponding modular Hamiltonian $\log\Delta$:
\begin{equation}
    E(\mu_0) := \text{Proj}(-\infty, \mu_0]
\end{equation}
We will need two-sided spectral cut-off:
\begin{equation}
    F(\mu_0) := E(\mu_0) - E(-\mu_0) = \text{Proj}(-\mu_0, \mu_0)
\end{equation}
where $\mu_0 > 0$. Notice that $F(0) = 0$. 
\begin{theorem}\label{theorem:mainresult}
    Consider a net of local algebras $\{\mathcal{A}(\mathcal{O})\}_{\mathcal{O}\in\mathbb{R}^{d,1}}$ that satisfies the Haag-Kastler axioms. Let $W_0$ be the right Rindler wedge and let $W_r$ be a right-translated Rindler wedge shifted to the right by $r > 0$ along the $x_1$ direction. Consider a one-parameter family of the shifted double cones inside the shifted wedge $\{B^r_{e^{2\pi s}R} \subset W_r\}_{s > 0}$. 
    
    In addition, consider the following modular operators associated with the local algebras:
    \begin{enumerate}
        \item Let $\Delta_{W_0}$ be the modular operator of $\mathcal{A}(W_0)$;
        \item Let $\Delta_{W_r}$ be the modular operator of $\mathcal{A}(W_r)$;
        \item Let $\Delta_{B^r_{e^{2\pi s}R}}$ be the modular operator $\mathcal{A}(B^r_{e^{2\pi s}R})$
    \end{enumerate}
    Finally consider the spectral decomposition of the modular flow of the right Rindler wedge:
    \begin{equation}
        \Delta_{W_0}^{it} = \int_{-\infty}^\infty e^{i\mu t} dE(\mu)
    \end{equation}
    Then for any spectral cut-off $\mu_0 > 0$\footnote{For simplicity, the spectral cut-off $\mu_0 > 0$ restricts the spectrum from $[-\mu_0, \mu_0]$.}, for any Schwartz function $f(t)\in \mathcal{D}(\mathbb{R})$, for any operator $A\in \mathcal{A}(B^r_R)$ and for any constant $s > 0$, there exists a constant $b_f(\mu_0, s) > 0$ such that:
    \begin{equation}
        ||F(\mu_0)\lb\int_\mathbb{R} dtf(t)\lb\mathcal{D}_{W_0, B^r_{e^{2\pi s}R}}(t) - \mathcal{D}_{W_0, W_r}(t)\rb A\ket{\Omega}\rb||^2 \leq b_f(\mu_0, s)\lb||A\ket{\Omega}||^2 + ||A^*\ket{\Omega}||^2\rb
    \end{equation}
    In addition, for any spectral cut-off, we have a well-defined limit:
    \begin{equation}
        \lim_{s\rightarrow \infty}||F(\mu_0)\lb\int_\mathbb{R} dt f(t)\lb\mathcal{D}_{W_0, B^r_{e^{2\pi s}R}}(t) - \mathcal{D}_{W_0, W_r}(t)\rb A\ket{\Omega}\rb||^2 = 0
    \end{equation}
\end{theorem}
When the smearing function $f$ is Gaussian, the constant $b_f(\mu_0, s)$ can be calculated. The qualitative behavior of this constant matches the intuition (c.f. Subsection \ref{subsection:gaussiansmear}). In particular, as $s\rightarrow\infty$, the approximation error decreases for any fixed spectral cut-off. 

The proof of the theorem is straight-forward. We first consider the following lemma which is an adaptation of Fredenhagen's original result \cite{fredenhagen} coupled with the covariant translation group action:
\begin{lemma}\label{lemma:extensionfredenhagen}
    Using the same notation as the Theorem above. Specifically, consider the double cone $B^r_{e^{2\pi s}R}$ contained inside the wedge $W_r$. Then for any $f\in L^1(\mathbb{R})$ and any constant $s > 0$, there exists a constant $c_f'(s) > 0$ such that:
    \begin{equation}
        ||\int_\mathbb{R} dt f(t) (\Delta^{-it}_{B^r_{e^{2\pi s}R}} - \Delta_{W_r}^{-it})A\ket{\Omega}||^2 \leq c'_f(s)(||A\ket{\Omega}||^2 + ||A^*\ket{\Omega}||^2)
    \end{equation}
    where $A$ is any operator in the local algebra $\mathcal{A}(B_R)$. The constant $c'_f(s)$ approaches $0$ as $s\rightarrow \infty$.
\end{lemma}
\begin{proof}
    Let $T_r$ be the right translation. Then we have:
    \begin{equation}
        T_r^*\lb B^r_{e^{2\pi s}R}\rb = B_{e^{2\pi s}R}\subset W_0
    \end{equation}
    \begin{equation}
        T_r^*\lb W_r \rb = W_0
    \end{equation}
    Then the modular automorphisms for the local algebras $\mathcal{A}\lb B^r_{e^{2\pi s}R}\rb$ and $\mathcal{A}\lb B_{e^{2\pi s}R} \rb$ are related by conjugation of the translation operator:
    \begin{equation}
        \Delta^{it}_{B^r_{e^{2\pi s}R}} = T_r\Delta^{it}_{B_{e^{2\pi s}R}} T_r^*
    \end{equation}
    Similarly, the modular automorphisms for the local algebras $\mathcal{A}\lb W_r\rb$ and $\mathcal{A}\lb W_0\rb$ are related by conjugation of the translation operator:
    \begin{equation}
        \Delta^{it}_{W_r} = T_r\Delta^{it}_{W_0}T_r^*
    \end{equation}
    Using these relations, we can rewrite the norm of the smeared modular time evolution of $A\ket{\Omega}$ as follows:
    \begin{align}
        \begin{split}
            &||\int_\mathbb{R} dt f(t)\lb\Delta^{-it}_{B^r_{e^{2\pi s}R}} - \Delta^{-it}_{W_r}\rb A\ket{\Omega}|| = ||\int_\mathbb{R} dt f(t)T_r\lb\Delta^{-it}_{B_{e^{2\pi s}R}} - \Delta^{-it}_{W_0}\rb T^*_r A\ket{\Omega}||
            \\&
            =||T_r\int_\mathbb{R} dt f(t) \lb\Delta^{-it}_{B_{e^{2\pi s}R}}-\Delta^{-it}_{W_0}\rb Ad_{T_r^*}\lb A\rb\ket{\Omega}||
            \\&=||\int_\mathbb{R} dt f(t) \lb \Delta^{-it}_{B_{e^{2\pi s}R}} - \Delta^{-it}_{W_0}\rb Ad_{T^*_r}\lb A\rb\ket{\Omega}||
        \end{split}
    \end{align}
    where we have used the fact that the translation action is unitary and hence does not change the norm. By Theorem \ref{theorem:modifiedfredenhagen}, we have the estimates:
    \begin{align}
        \begin{split}
            &||\int_\mathbb{R} dt f(t) \lb \Delta^{-it}_{B^r_{e^{2\pi s}R}}- \Delta^{-it}_{W_r} \rb A\ket{\Omega}|| = ||\int_\mathbb{R} dt f(t) \lb \Delta^{-it}_{B_{e^{2\pi s}R}} - \Delta^{-it}_{W_0} \rb Ad_{T_r^*}\lb A\rb \ket{\Omega}||
            \\&\leq c^{1/2}_f(s)\lb||Ad_{T_r^*}\lb A \rb\ket{\Omega}||^2 + ||Ad_{T^*_r}\lb A^*\rb\ket{\Omega}||^2\rb^{1/2}
            \\&
            =c_f^{1/2}(s)\lb ||A\ket{\Omega}||^2 + ||A^*\ket{\Omega}||^2 \rb^{1/2}
        \end{split}
    \end{align}
    where the last equation follows from the fact that translation action is unitary and the vacuum vector is invariant under translation:
    \begin{equation}
        ||Ad_{T_r^*}\lb A \rb\ket{\Omega}|| = ||T^*_rAT_r\ket{\Omega}|| = ||T_r^*A\ket{\Omega}|| = ||A\ket{\Omega}||\qedhere
    \end{equation}
\end{proof}
Using the observation above, we are now ready to prove Theorem \ref{theorem:mainresult}:
\begin{proof}
    Using the spectral decomposition we can rewrite the norm as follows:
    \begin{align}
        \begin{split}
            &||F(\mu_0)\lb\int_\mathbb{R} dt f(t)\lb \mathcal{D}_{W_0, B^r_{e^{2\pi s}R}}(t) - \mathcal{D}_{W_0, W_r}(t) \rb A\ket{\Omega}\rb||
            \\&= ||\int_\mathbb{R} dtf(t)\int_{-\mu_0}^{\mu_0}dE(\mu) e^{-i\mu t}\lb\Delta_{B^r_{e^{2\pi s}R}}^{it} - \Delta_{W_r}^{it}\rb A\ket{\Omega}||
            \\&
            =\sup_{\substack{B\in\mathcal{A}(W_0)\\s.t.||B\ket{\Omega}||\leq 1}}|\langle B\Omega\text{ , }\int_\mathbb{R}dt f(t)\int_{-\mu_0}^{\mu_0}dE(\mu)e^{-i\mu t}\lb\Delta_{B^r_{e^{2\pi s}R}}^{it} - \Delta_{W_r}^{it}\rb A\Omega\rangle|
            \\&
            =\sup_{\substack{B\in\mathcal{A}(W_0)\\ s.t.||B\ket{\Omega}||\leq 1}}|\int_\mathbb{R}dt f(t)\int_{-\mu_0}^{\mu_0}e^{-i\mu t}d\bra{E(\mu) B\Omega}\text{ , }\lb\Delta_{B^r_{e^{2\pi s}R}}^{it} - \Delta_{W_r}^{it}\rb A\Omega\rangle|
        \end{split}
    \end{align}
\NL{Add a physical argument to justify that we can exchange the order of $dt$ integral with $dE(\mu)$. Are there choice of $f(t)$ that make it work?}
    
    To simplify notation, we denote the $(t,\mu)$-dependent measurable function as:
    \begin{equation}
        g_t(\mu) := \langle E(\mu) B\Omega\text{ , }\lb\Delta^{it}_{B^r_{e^{2\pi s}R}} - \Delta^{it}_{W_r}\rb A\Omega\rangle
    \end{equation}
    By integration by parts, we have:
    \begin{align}
        \begin{split}
            &||F(\mu_0)\lb\int_\mathbb{R} dt f(t)\lb\mathcal{D}_{W_0, B^r_{e^{2\pi s}R}}(t) - \mathcal{D}_{W_0, W_r}(t)\rb A\ket{\Omega}\rb||
            \\&
            =\sup_{\substack{B\in\mathcal{A}(W_0)\\s.t.||B\ket{\Omega}||\leq 1}}|\int_\mathbb{R}dt f(t)\lb e^{-i\mu t}g_t(\mu)|_{-\mu_0}^{\mu_0} - \int_{-\mu_0}^{\mu_0}g_t(\mu)\lb-ite^{-i\mu t}\rb d\mu \rb|
            \\&
            =\sup_{\substack{B\in\mathcal{A}(W_0)\\s.t.||B\ket{\Omega}||\leq 1}}|\int_\mathbb{R} dt f(t) \lb e^{-i\mu_0 t}g_t(\mu_0)- e^{i\mu_0 t}g_t(-\mu_0)\rb + i\int_\mathbb{R} dt \int_{-\mu_0}^{\mu_0} d\mu \lb tf(t)\rb e^{-i\mu t} g_t(\mu)|
        \end{split}
    \end{align}
    For the first term, we have the following estimate:
    \begin{align}
        \begin{split}
            &|\int_\mathbb{R} dt f(t)e^{-i\mu_0 t}\langle E(\mu_0) B\Omega\text{ , }\lb\Delta^{it}_{B^r_{e^{2\pi s}R}} - \Delta^{it}_{W_r}\rb A\Omega\rangle|
            \\&
            =|\langle E(\mu_0)B\Omega\text{ , }\int_\mathbb{R}dt f(t)e^{-i\mu_0 t}\lb\Delta^{it}_{B^r_{e^{2\pi s}R}} - \Delta^{it}_{W_r}\rb A\Omega\rangle|
            \\&
            \leq ||B\ket{\Omega}||\cdot||\int_\mathbb{R} dt f(t) e^{-i\mu_0 t}\lb\Delta^{it}_{B^r_{e^{2\pi s}R}} - \Delta^{it}_{W_r}\rb A\ket{\Omega}||
        \end{split}
    \end{align}
    Similarly, for the second term, we have the following estimate:
    \begin{align}
        \begin{split}
            &|\int_\mathbb{R} dt f(t) e^{i\mu_0 t}\langle E(-\mu_0)B\Omega, \lb \Delta^{it}_{B^r_{e^{2\pi s}R}} - \Delta^{it}_{W_r} \rb A \Omega\rangle|
            \\&
            = |\langle E(-\mu_0)B\Omega, \int_\mathbb{R} dt f(t) e^{i\mu_0 t}\lb\Delta^{it}_{B^r_{e^{2\pi s}R}} - \Delta^{it}_{W_r}\rb A \Omega \rangle|
            \\&
            \leq ||B\ket{\Omega}||\cdot||\int_\mathbb{R} dt f(t) e^{i\mu_0 t}\lb \Delta^{it}_{B^r_{e^{2\pi s}R}} - \Delta^{it}_{W_r} \rb A\ket{\Omega}||
        \end{split}
    \end{align}
    Therefore we have:
    \begin{align}
        \begin{split}
            &|\int_\mathbb{R} dt f(t)\lb e^{-i\mu_0 t}g_t(\mu_0) - e^{i\mu_0 t}g_t(-\mu_0) \rb| \leq |\int_\mathbb{R} dt f(t) e^{-i\mu_0 t}g_t(\mu_0)| + |\int_\mathbb{R} dt f(t) e^{i\mu_0 t}g_t(-\mu_0)|
            \\&
            \leq ||B\ket{\Omega}||\cdot \lb||\int_\mathbb{R} dt f(t) e^{-i\mu_0 t}\lb\Delta^{it}_{B^r_{e^{2\pi s}R}} - \Delta^{it}_{W_r}\rb A\ket{\Omega}|| + ||\int_\mathbb{R} dt f(t) e^{i\mu_0 t}\lb\Delta^{it}_{B^r_{e^{2\pi s}R}} - \Delta^{it}_{W_r}\rb A\ket{\Omega}||\rb
        \end{split}
    \end{align}
    Since the functions $f_{\mu_0}(t) := e^{-i\mu_0 t}f(t)$ and $f_{-\mu_0}(t)$ are integrable, by Lemma \ref{lemma:extensionfredenhagen} we can provide the following upper bound:
    \begin{align}
        \begin{split}
            &\sup_{\substack{B\in\mathcal{A}(W_0)\\s.t.||B\ket{\Omega}||\leq 1}}|\int_\mathbb{R} dt f(t) e^{-i\mu_0 t}\langle E(\mu_0)B\Omega\text{ , }\lb\Delta^{it}_{B^r_{e^{2\pi s}R}} - \Delta^{it}_{W_r}\rb A\Omega\rangle \\& - e^{i\mu_0 t}\langle E(-\mu_0)B\Omega, \lb\Delta^{it}_{B^r_{e^{2\pi s}R}} - \Delta^{it}_{W_r}\rb A\Omega\rangle|
            \\&
            \leq \sup_{\substack{B\in\mathcal{A}(W_0)\\s.t.||B\ket{\Omega}||\leq 1}}||B\ket{\Omega}||\cdot\lb \lb c_{f_{\mu_0}}^{1/2}(s) + c_{f_{-\mu_0}}^{1/2}(s)\rb\lb||A\ket{\Omega}||^2 + ||A^*\ket{\Omega}||^2\rb^{1/2}\rb
            \\&
            \leq \lb c_{f_{\mu_0}}^{1/2}(s) + c_{f_{-\mu_0}}^{1/2}(s)\rb\lb||A\ket{\Omega}||^2 + ||A^*\ket{\Omega}||^2\rb^{1/2}
        \end{split}
    \end{align}
    For the second term, by Fubini's theorem we have the following estimate:
    \begin{align}
        \begin{split}
            &\sup_{\substack{B\in\mathcal{A}(W_0)\\s.t.||B\ket{\Omega}||\leq 1}}|i\int_\mathbb{R} dt \int_{-\mu_0}^{\mu_0}d\mu \lb tf(t)\rb e^{-i\mu t}\langle E(\mu)B\Omega\text{ , }\lb\Delta^{it}_{B^r_{e^{2\pi s}R}} - \Delta^{it}_{W_r}\rb A\Omega \rangle|
            \\&
            \leq \sup_{\substack{B\in\mathcal{A}(W_0)\\s.t.||B\ket{\Omega}||\leq 1}}\int_{-\mu_0}^{\mu_0}d\mu\int_\mathbb{R} dt |\lb tf(t)\rb e^{-i\mu t}\langle E(\mu)B\Omega\text{ , }\lb \Delta^{it}_{B^r_{e^{2\pi s}R}} - \Delta^{it}_{W_r} \rb A\Omega \rangle|
            \\&
            \leq \sup_{\substack{B\in\mathcal{A}(W_0)\\s.t.||B\ket{\Omega}||\leq 1}}\int_{-\mu_0}^{\mu_0}d\mu ||E(\mu) B\ket{\Omega}||\cdot||\int_\mathbb{R} dt\lb tf(t) \rb e^{-i\mu t}\lb\Delta^{it}_{B^r_{e^{2\pi s}R}}- \Delta^{it}_{W_r}\rb A\ket{\Omega}||
        \end{split}
    \end{align}
    Again since the function $tf_\mu(t):= \lb tf(t) \rb e^{-i\mu t}$ is integrable, by Lemma \ref{lemma:extensionfredenhagen} we can provide the following upper bound:
    \begin{align}
        \begin{split}
            &\sup_{\substack{B\in\mathcal{A}(W_0)\\s.t.||B\ket{\Omega}||\leq 1}}|i\int_\mathbb{R} dt \int_{-\mu_0}^{\mu_0}d\mu \lb tf(t)\rb e^{-i\mu t}\langle E(\mu)B\Omega\text{ , }\lb\Delta^{it}_{B^r_{e^{2\pi s}R}} - \Delta^{it}_{W_r}\rb A\Omega \rangle|
            \\&
            \leq \sup_{\substack{B\in\mathcal{A}(W_0)\\s.t.||B\ket{\Omega}||\leq 1}}\int_{-\mu_0}^{\mu_0}d\mu ||E(\mu)B\ket{\Omega}||\cdot\lb c^{1/2}_{t f_\mu}(s)\lb||A\ket{\Omega}|| + ||A^*\ket{\Omega}||^2\rb^{1/2}\rb
            \\&
            \leq \int_{-\mu_0}^{\mu_0}d\mu \lb c^{1/2}_{tf_\mu}(s)\lb||A\ket{\Omega}||^2 + ||A^*\ket{\Omega}||^2\rb^{1/2}\rb 
            \\&
            \leq 2\mu_0\lb\sup_{-\mu_0\leq \mu\leq \mu_0}c^{1/2}_{tf_\mu}(s)\rb\cdot\lb||A\ket{\Omega}||^2 + ||A^*\ket{\Omega}||^2\rb^{1/2}
        \end{split}
    \end{align}
    Combining these two upper bounds, we have the estimates:
    \begin{align}
        \begin{split}
            &||F(\mu_0)\lb\int_\mathbb{R} dt f(t)\lb\mathcal{D}_{W_0, B^r_{e^{2\pi s}R}}(t) - \mathcal{D}_{W_0, W_r}(t)\rb A\ket{\Omega}\rb||
            \\&
            \leq \lb c^{1/2}_{f_{\mu_0}}(s) + c^{1/2}_{f_{-\mu_0}}(s) + 2\mu_0\sup_{-\mu_0\leq \mu\leq \mu_0}c^{1/2}_{tf_\mu}(s) \rb\cdot\lb||A\ket{\Omega}||^2 + ||A^*\ket{\Omega}||^2\rb^{1/2}
        \end{split}
    \end{align}
    By Lemma \ref{lemma:shiftedfredenhagenconstant} below (which requires the assumption that $tf(t)$ is integrable), the constant $b_f(\mu_0, s) := \lb c^{1/2}_{f_{\mu_0}}(s) + c^{1/2}_{f_{-\mu_0}}(s) + 2\mu_0\sup_{0\leq \mu\leq \mu_0}c^{1/2}_{tf_\mu}(s)\rb$ has a well-defined limit as $s\rightarrow\infty$:
    \begin{equation}
        \lim_{s\rightarrow\infty}b_f(\mu_0, s) = 0
    \end{equation}
    This concludes the proof of our main theorem. 
\end{proof}
\subsection{Approximation of Characteristic Functions - Part II. Double Cone Inside Double Cone}\label{subsection:coneinsidecone}
This subsection is a continuation of the previous one. Here we study characteristic function of double cone inclusions. More precisely, using the same notation as the previous subsection, we consider the following geometric set-up. Let $W_0$ be the standard right Rindler wedge, $W_r$ be the right translated wedge, $B_{e^{2\pi(u+s)}R}\subset W_0$ be a double cone where $u, s \geq 0$ are two parameters controlling the size of the double cone, and $B^r_{e^{2\pi s}R}\subset W_r$ be a right translated double cone whose size is controlled by a single parameter $s\geq 0$. Note the parameter $s$ simultaneously control the size of both double cones.
\begin{equation}
    \begin{tikzcd}[column sep=huge, row sep=large]
        B^r_s:=B^r_{e^{2\pi s}R} \arrow[r, "\subset"] \arrow[ d, "\subset"] & B_{s,u}:=B_{e^{2\pi(u+s)}R} \arrow[d, "\subset"] 
        \\
        W_r \arrow[r, "\subset"]                 & W_0
    \end{tikzcd}    
\end{equation}
We adjust the two parameters so that this commutative diagram always holds. This requires the following relation holds:
\begin{equation}
    e^{2\pi(u + s)}R > r + e^{2\pi s}R
\end{equation}
For simplicity, we use the following simplified notation for the rest of this subsection:
\begin{equation}
    B^r_s:=B^r_{e^{2\pi s}R}\text{ , }B_{s,u}:=B_{e^{2\pi (u + s)}R}
\end{equation}
Analogous to the previous subsection, we are interested in the following two characteristic functions:
\begin{enumerate}
    \item One characteristic function associated with the double cone inclusion:
    \begin{equation}
        \mathcal{D}_{B_{s,u}, B^r_s} = \Delta_{B_{s,u}}^{-it}\Delta_{B^r_s}^{it}
    \end{equation}
    \item The other characteristic function is associated with the standard wedge inclusion:
    \begin{equation}
        \mathcal{D}_{W_0,W_r}(t) = \Delta_{W_0}^{-it}\Delta_{W_r}^{it}
    \end{equation}
\end{enumerate}
Intuitively, as $s,u\rightarrow\infty$ such that $B^r_s\subset B_{s,u}$, the double cone $B^r_s$ approaches $W_r$ and the double cone $B_{s,u}$ approaches $W_0$. Under this limit, the modular operator $\Delta_{B^r_s}$ approaches the modular operator $\Delta_{W_r}$ in the strong resolvent sense\cite{takesaki2} and the modular operator $\Delta_{B_{s,u}}$ approaches the modular operator $\Delta_{W_0}$ in the strong resolvent sense. Therefore, it is reasonable to conjecture that the characteristic function $\mathcal{D}_{B^{s,u}, B^r_s}(t)$ approaches $\mathcal{D}_{W_0,W_r}(t)$. Following Theorem \ref{theorem:mainresult} from the previous subsection, we prove an analogous result:
\begin{theorem}\label{theorem:ballinsideball}
    Consider a net of local algebras $\{\mathcal{A}(\mathcal{O})\}_{\mathcal{O}\subset\mathbb{R}^{d,1}}$ that satisfies the Haag-Kastler axioms. Let $W_0$ be the right Rindler wedge and let $W_r$ be a right translated Rindler wedge shifted to the right by $r > 0$ along the $x_1$ direction. Consider a one-parameter family of shifted double cones inside the shifted wedge $\{B^r_s\subset W_r\}_{s > 0}$. In addition, consider a two-parameter family of dilated double cones inside the right Rindler wedge $\{B_{s,u}\subset W_0\text{ s.t.  }B^r_s\subset B_{s,u}\}_{s > 0, e^{2\pi(s+u)}R > r + e^{2\pi s}R}$. 

    In addition, consider the following modular operators associated with the local algebras:
    \begin{enumerate}
        \item Let $\Delta_{W_0}$ be the modular operator of $\mathcal{A}(W_0)$;
        \item Let $\Delta_{W_r}$ be the modular operator of $\mathcal{A}(W_r)$;
        \item Let $\Delta_{B^r_s}$ be the modular operator of $\mathcal{A}(B^r_s)$;
        \item Let $\Delta_{B_{s,u}}$ be the modular operator of $\mathcal{A}(B_{s,u})$
    \end{enumerate}
    Finally consider the spectral decomposition of the modular flow of the right Rindler wedge:
    \begin{equation}
        \Delta^{it}_{W_0} = \int_0^\infty e^{i\mu t}dE(\mu)
    \end{equation}
    Then for any spectral cut-off $\mu_0 > 0$, for any essentially bounded measurable function with bounded support $f(t)\in L^\infty_b(\mathbb{R})$, for any operator $A\in\mathcal{A}(B^r_0)$, we have a well-defined limit:
    \begin{align}
        \begin{split}
            \lim_{s\rightarrow\infty}\lim_{u\rightarrow\infty}||E(\mu_0)\lb\int_\mathbb{R} dt f(t)\lb \mathcal{D}_{B_{s,u}, B^r_s}(t) - \mathcal{D}_{W_0,W_r}(t) \rb A\ket{\Omega}\rb||^2 = 0
        \end{split}
    \end{align}
\end{theorem}
The proof of this theorem is relatively straight-forward by using an intermediate characteristic function associated with the inclusion: $B_{s,u}\subset W_0$. Notice that Theorem \ref{theorem:mainresult} already shows that the following limit is well-defined:
\begin{equation}
    \lim_{s\rightarrow\infty}\lim_{u\rightarrow\infty}||E(\mu_0)\lb\int_\mathbb{R} dt f(t)\lb \mathcal{D}_{W_0, B^r_s}(t) - \mathcal{D}_{W_0, W_r}(t) \rb A\ket{\Omega}\rb||^2 = 0
\end{equation}
Therefore, we are actually left with proving the following limit holds:
\begin{equation}
    \lim_{s\rightarrow\infty}\lim_{u\rightarrow\infty}||E(\mu_0)\lb\int_\mathbb{R} dt f(t)\lb \mathcal{D}_{B_{s,u}, B^r_s}(t) - \mathcal{D}_{W_0, B^r_s}(t) \rb A\ket{\Omega}\rb||^2 = 0
\end{equation}
\begin{lemma}
    For every fixed $s > 0$, the following limit holds:
    \begin{equation}
        \lim_{u\rightarrow\infty}||\int_\mathbb{R} dt f(t)\lb\mathcal{D}_{B_{s,u}, B^r_s}(t) - \mathcal{D}_{W_0, B^r_s}(t)\rb A\ket{\Omega}|| = 0
    \end{equation}
\end{lemma}
\begin{proof}
    As $u\rightarrow\infty$, the modular operator $\Delta_{B_{s,u}}\rightarrow \Delta_{W_0}$ in the strong resolvent sense\cite{takesaki2}. Equivalently, the modular automorphism $\Delta_{B_{s,u}}^{it}\rightarrow\Delta_{W_0}^{it}$ strongly pointwise and strongly uniformly on any bounded interval in time. Using this, we have the following:
    \begin{align}
        \begin{split}
            ||\int_\mathbb{R}& dt f(t) \lb\mathcal{D}_{B_{s,u},B^r_s}(t) - \mathcal{D}_{W_0, B^r_s}(t)\rb A\ket{\Omega}|| \leq ||f||_1\sup_{t\in\text{supp}(f)}||\lb\mathcal{D}_{B_{s,u},B^r_s}(t) - \mathcal{D}_{W_0, B^r_s}(t)\rb A\ket{\Omega}||
            \\
            &\leq ||f||_1\sup_{t\in\text{supp}(f)}\sup_{B\ket{\Omega}\text{s.t.}||B\ket{\Omega}||\leq 1}|\bra{\Omega}B, \lb\mathcal{D}_{B_{s,u}, B^r_s}(t) - \mathcal{D}_{W_0, B^r_s}(t)\rb A\ket{\Omega}|
            \\&
            =||f||_1\sup_{t\in\text{supp}(f)}\sup_{B\ket{\Omega}\text{s.t.}||B\ket{\Omega}||\leq 1}|\bra{\Omega}B\lb\Delta^{it}_{B_{s,u}}- \Delta^{it}_{W_0}\rb, \Delta^{it}_{B_s^r}A\ket{\Omega}|
            \\&
            \leq||f||_1\cdot||A||\cdot\sup_{B\ket{\Omega}\text{s.t.}||B\ket{\Omega}||\leq 1}\sup_{t\in\text{supp}(f)}||\lb\Delta^{it}_{B_{s,u}} - \Delta^{it}_{W_0}\rb B\ket{\Omega}||
        \end{split}
    \end{align}
    By uniform convergence on bounded interval, we have:
    \begin{equation}
        \lim_{u\rightarrow\infty}\sup_{t\in\text{supp}(f)}||\lb\Delta^{it}_{B_{s,u}} - \Delta^{it}_{W_0}\rb B\ket{\Omega}|| = 0
    \end{equation}
    Hence we have the desired limit.\qedhere
\end{proof}
Combining the lemma and Theorem \ref{theorem:mainresult}, Theorem \ref{theorem:ballinsideball} follows.

%\subsection{Optical theorem}
%Unitarity of $S$ implies the optical theorem 
%\begin{eqnarray}
%i (M^\dagger- M)=M^\dagger M\ .
%\end{eqnarray}


%\subsection{Dispersion relations}

%The dispersion relations follow from the Titchmarsh theorem; see Appendix \ref{app:titchmarsh}. 

%\subsection{Regge limit in Scattering}
%In defining the scattering amplitude, we take the $z\to \infty$ limit so that $\mathcal{M}$ is only a function of the kinematics of the incoming and outgoing asymptotic states. In particular, for $2\to 2$ scattering it is a function of Mandelstam variables $s$ and $t$\footnote{We warn the readers not to confuse $t$ with time $\Re(z)$ in our modular scattering matrix.}. 

%In $2\to 2$ scattering of same particles with incoming momenta $p,k$ and outgoing momenta $p',k'$ the Mandelstam variables are $s=-(p+k)^2$ and $t=-(k'-k)^2$. In the center of mass coordinates $\vec{p}=-\vec{k}$, $s=4(m^2+\vec{p}^2)$ has the interpretation of the center of mass energy squared and positive and $t=-2\vec{p}^2(1-\cos\theta)\leq 0$ has the interpretation of the momentum transfer squared and neagtive \footnote{Note that we are using the  metric signature $(-,+,\cdots , +)$. For the mostly negative signature, often we defines $s=(p+k)^2$ and $t=(k'-k)$ so that in the physical regime $s>0$ and $t<0$.}. Similarly, we have the variable $u=-(p'-k)^2$ which in the center of mass frame is $u=-2\vec{p}^2(1+\cos\theta)<0$. At very high energies $|s|\gg 1$ we can ignore the masses and obtain relations
%\begin{eqnarray}
%    t\simeq -s \sin^2(\theta/2), \qquad u\simeq -s \cos^2(\theta/2)\ .
%\end{eqnarray}
%The forward scattering limit is the limit of $\theta\simeq 0$. The analyticity often discussed in dispersion relations is in the ``forward scattering limit" of large $s$ and fixed $t$. The assumption is that the scattering amplitude is analytic in the upper half-plane of $s$ and we can write $\mathcal{M}_{phys}(s,t)=\lim_{\ep\to 0}\mathcal{M}(s+i\ep, t)$. Here, the idea can be summarized as follows: the physical scattering amplitude is the boundary value of an analytic function.



%In the Regge limit $s\to \infty$ and $t$ fixed, the scattering amplitude is expected to behave as 
%\begin{eqnarray}
%    \mM(s,t)\sim s^{\alpha(t)}
%\end{eqnarray}
%where $\alpha(t)$ are called Regge trajectories.
%In Regge theory, one views the $2\to 2$ scattering matrix as a function of energy ($E\sim \sqrt{s}$) and scattering angle $\theta$ expanded in angular momentum basis
%\begin{eqnarray}
%    \mM(E,x=\cos\theta)=\sum_{l=0}^\infty(2l+1)P_l(x)\mM_l(E) 
%\end{eqnarray}
%where $P_l(x)$ are Legendre polynomials. 
%Regge's insight was to study $\mM_l(E)$ in terms of a continuous and complex variable angular momentum variable $l$. He established that the scattering amplitude is analytic in the complex $z$ plane, except for potential cuts aong the real axis.


%\subsection{Conformal Regge limit}

%{\bf Inclusion of algebras:} Another Regge limit discussed in the literature is the Regge limit defined in \cite{chandrasekaran2022scattering} (see eq 1.11, 6.11 and 6.12). The correlator of interest is \footnote{Note the definition above differ from the notation of \cite{chandrasekaran2022scattering} by $z\to z/(2\pi)$.}
%\begin{eqnarray}
%    \mathcal{F}(t_L,t_R,z)&&=\braket{\mO_L(t_L)|\Delta_{LI}^{-iz}\Delta_L^{iz}|\mO_R(t_R)}
%\end{eqnarray}
%where $\mO_L\in \mA_L$ and $\mO_R\in \mA_{L'}=\mA_R$, and the operators are evolved with modular flows
%\begin{eqnarray}
%    &&\mO_L(2\pi t)\equiv \rho_L^{it}\mO_L \rho_L^{-it}=\Delta_L^{it}\mO_L \Delta_L^{-it}\nn\\
%    &&\mO_R(2\pi t)\equiv \rho_R^{it}\mO_R \rho_R^{-it}=\Delta_L^{-it}\mO_R \Delta_L^{it}\ .
%\end{eqnarray}
%In this geometric setup, $L$ and $I$ are disjoint, and we have the causal complements $L'=R$ and $(LI)'=r$. Since $\Delta_A^{iz}=\Delta_{A'}^{-iz}$ we can rewrite the measure $(\ref{})$ in terms of the inclusion $r\subset R=r I$, where now, the relative commutant is $I$; i.e. $R=r I$, but there is no split between $r$ and $I$:
%\begin{eqnarray}
%    \mathcal{F}(2\pi t_L,2\pi t_R,z)&&=\braket{\mO_L| \Delta_L^{-i t_L}\Delta_{LI}^{-iz}\Delta_L^{i(z-t_R)}|\mO_R}\nn\\
%    &&=\braket{\mO_L|\Delta_R^{it_L}\Delta_r^{iz}\Delta_R^{-i(z-t_R)}|\mO_R}\nn\\
%    &&=\braket{\mO_L|\Delta_R^{i(t_L+z/2)} S(-z/2)\Delta_R^{-i(z/2-t_R)}|\mO_R}
%\end{eqnarray}
%where $T_{r\subset R}(z)$ and $S_{r\subset R}(z)$ are the characteristic function and modular scattering matrix of the inclusion $r\subset R$, respectively:
%\begin{eqnarray}
%    &&T_{r\subset R}(z)=\Delta_R^{iz}\Delta_r^{-iz}\nn\\
%    &&S_{r\subset R}(z)=T(z)T(-z)^{-1}\ .
%\end{eqnarray}
% The authors in \cite{chandrasekaran2022scattering} call the limit of large $t_L-t_R$ the Regge limit. We consider two particular limits of the above expression:
%\begin{eqnarray}\label{cases}
%    &&\mathcal{F}(t-z/2,-t+z/2,z)=\braket{\mO_L|\Delta_R^{it}S(-z/2)\Delta_R^{-it}|\mO_R}=\braket{\mO_L(t)|S(-z/2)|\mO_R(-t)}\nn\\
%    &&\mathcal{F}(t,z-t,z)=\braket{\mO_L|\Delta_R^{it}\Delta_r^{iz}\Delta_R^{-it}|\mO_R}=\braket{\mO_L(t)|\Delta_r^{iz}|\mO_R(-t)}\ .
%\end{eqnarray}
%The first case computes various matrix elements of the modular scattering matrix, whereas the second case computes matrix elements of the modular flow $\Delta_r^{iz}$. This is simply a consequence of 
%\begin{eqnarray}
%    S(-z/2)=\Delta^{-iz/2}_R \Delta_r^{iz}\Delta_R^{-iz/2}
%\end{eqnarray}
%which can be interpreted as looking at matrix elements of $\Delta_r^{iz}$ in a basis where we are evolving both times $t_L$ and $t_R$ upwards, stretching the wormhole! Note that in both cases in (\ref{cases}) $t_L-t_R=2t-z$ and the Regge limit is the large $t$ for fixed $z$. \NL{Maybe this can be interpreted as exploring the action of $S(-z/2)$ and $\Delta_r^{iz}$ from the point of view of operators that are boosted to the area near the entangling surface of $R$. See figure 2 in \cite{chandrasekaran2022scattering}.}



% In other words, this correlator can be written as 
% \begin{eqnarray}
%     \mathcal{F}(t_L, t_R, z)&&=\braket{\mO_L| \Delta_L^{-i t_L}\Delta_{LI}^{-iz}\Delta_L^{i(z+t_R)}|\mO_R}\nn\\
%     &&=\braket{\mO_L|\Delta_L^{-i(z+t_L/\pi)/2}T(z/2)^{-1}T(-z/2)\Delta_L^{i(z-t_R/\pi)/2}|\mO_R}
% \end{eqnarray}\ .





%we consider a one-parameter family of (inner or outer) automorphisms of the smaller algebra $\alpha_t:\mA\to \mA$ realized as $U(t)$ and consider 
%\begin{eqnarray}
%    U(t)^\dagger T_{A\subset B}(z)U(t)\ .
%\end{eqnarray}
%The Regge limit is defined to be the limit where the relative size of the inclusion $s\ll 1$ and $t$ is large such that $s e^{2\pi t}\ll 1$. 

%Two such cases of such Regge limits are discussed in the literature:
%In \cite{chandrasekaran2022scattering} (see eq 1.11, 6.11 and 6.12), the authors consider the case where $U(t)=\Delta_A^{it}$ and $\mB=\mA \otimes \mA_I$. This is a standard split inclusion, and we can take the size of the relative commutant $|I|$ to be small.

%In this case, what we obtain is
%\begin{eqnarray}
% U(t)^\dagger T_{A\subset B}(z)U(t)=\Delta_A^{-it}\Delta_B^{iz}\Delta_A^{-iz}\Delta_A^{it}\ .
%\end{eqnarray}
%Equation 6.12 of \cite{chandrasekaran2022scattering} establishes that this correlator can be viewed as a regulator 




%\begin{eqnarray}
%\braket{\Phi_L(0)| T_{A\subset B}(z)|\Phi_R(T)}
%\end{eqnarray}
%where the relative size of the inclusion is small $s\ll 1$, and we look at large $T\gg 1$ such that $s e^{2\pi T}\ll 1$. 



%{\bf Inclusion of algebras (standard splits):} Consider the inclusion of algebras $\mA\subset \mB$ that is standard; the relative commutant $\mB\mathcal{A}p \mA'$ is cyclic. Then, we can consider the inclusion $\mA\in \mA\otimes \mB'$ that is a standard split. We consider the characteristic function 
%\begin{eqnarray}
%    T_{A\subset A\otimes B'}(z)=\Delta_{A\otimes B'}^{iz}\Delta_A^{-iz}\ .
%\end{eqnarray}
%The relative size of the inclusion is $|B'|$. The Regge limit is defined to be the limit where $|B'|\ll 1$, and $z\gg 1$ such that $|B'| e^{2\pi z}\ll 1$.





%To see the physics of this limit, consider the 



%\subsection{Modular asymptotic states}

%The S-matrix 

\section{Discussion}


%\section{Characteristic function as a regulator, double cone, and Modular quasi-normal modes}




% \section{1D Lattice examples}


% Consider the one-dimensional quantum XY-model with the Hamiltonian
% \begin{eqnarray}
%     H--\frac{1}{2}\sum_i\lb \frac{1+\gamma}{2}X_i X_{i+1}+\frac{1-\gamma}{2}Y_i Y_{i+1}+\lambda Z_i\rb\ .
% \end{eqnarray}
% The cases $\gamma=0$ and $\gamma=1$ are called the XX model and the Ising model, respectively. This model can be solved using the Jordan-Wigner transform 
% that sends it to a second-order action in fermionic modes.

% We start with the case $\lambda=1$, the transverse Ising model. For $\lambda<1$ the system is in the disordered phase, and for $\lambda>1$ it is ordered. We define $k=\min(\lambda,\lambda^{-1})$ and $k'=\sqrt{1-k^2}$.

% The modular Hamiltonian of a half-line  $W_0=\{i| i\geq 0\}$ can be written explicitly as
% \begin{eqnarray}
% &&\lambda<1:\quad \hat{K}_{W(0)}=-2I(k')\sum_{j=0}^\infty \lb j X_jX_{j+1}+k (j-1/2) Z_j\rb\nn\\
% &&\lambda>1:\quad \hat{K}_{W(0)}=-2I(k')\sum_{j=0}^\infty \lb k j X_jX_{j+1}+ (j-1/2) Z_j\rb\nn
% \end{eqnarray}
% where $I(k')$ is a complete elliptic integral of the first kind. We would like to compute the characteristic function $\Delta_B^{1/4}\Delta_A^{-1/4}$ for half-spaces $A=W(a>0)\subset W(0)=B$. We can avoid the calculation of the full modular operator, by thinking in terms of the characteristic function as the following super-operator acting on the operators in the algebra:
% \begin{eqnarray}
% &&\Delta_B^{1/4}\Delta_A^{-1/4}a\ket{\Omega}=\mathcal{D}_{i/4}(a)\ket{\Omega}\nn\\
% &&\mathcal{D}_{i/4}(a)=e^{K_B/4}e^{-K_A/4} a e^{K_A/4}e^{-K_B/4}\ .
% \end{eqnarray}
% Define the operator $\tilde{K}=K/(2I(k'))$, then for $\lambda<1$ we have
% \begin{eqnarray}
% &&[\tilde{K},X_l]=-k(l-1/2)i Y_l\nn\\
% &&[\tilde{K},Z_l]=i (l Y_lX_{l+1}+(l-1)X_{l-1}Y_l)\nn\\
% &&[\tilde{K},Y_l]=-\ .
% \end{eqnarray}
% Use Jordan-Wigner transform to write it in terms of the fermions

%Similarly, the modular Hamiltonians of the commutants are
%\begin{eqnarray}
%&&\lambda<1:\quad \hat{K}_{W'(0)}=2I(k')\sum_{j=0}^{-\infty} \lb j X_jX_{j-1}+k (j+1/2) Z_j\rb\nn\\
%&&\lambda>1:\quad \hat{K}_{W'(0)}=2I(k')\sum_{j=0}^{-\infty} \lb k j X_jX_{j-1}+ (j+1/2) Z_j\rb\nn\ .
%\end{eqnarray}
%
%
% We first write down the full modular Hamiltonians for $\lambda<1$ 
%\begin{eqnarray}
%K_{W(0)}&&=\hat{K}_{W(0)}-\hat{K}_{W'(0)}=\nn\\
%&&=-2I(k')\sum_{j=\infty}^\infty \lb j X_jX_{j+\text{sgn}(j)}+k (j-\text{sgn}(j))/2) Z_j\rb\nn
%\end{eqnarray}
%and for $\lambda>1$ 
%\begin{eqnarray}
%K_{W(0)}&&=\hat{K}_{W(0)}-\hat{K}_{W'(0)}=\nn\\
%&&=-2I(k')\sum_{j=\infty}^\infty \lb k j X_jX_{j+\text{sgn}(j)}+ (j-\text{sgn}(j))/2) Z_j\rb\nn
%\end{eqnarray}
% 
%\NL{Finish the calculation; how about the XY-model?}

% As the second example, consider a 1d harmonic chain
% \begin{eqnarray}
% H=\sum_i \lb \frac{1}{2m}p_i^2+\frac{m\omega^2}{2}q_i^2+\frac{\kappa}{2}(q_i-q_{i+1})^2\rb
% \end{eqnarray}
% It is a discretization of massive scalar field in 1 dimensions. The modular Hamiltonian of a subregion $A$ of $L$ sites is
% \begin{eqnarray}
% &&\hat{K}_A=\frac{1}{2}\vec{r}^T\cdot {\bf K}_A \cdot \vec{r},\qquad \vec{r}=\begin{pmatrix}&\vec{q}\\& \vec{p}\end{pmatrix}
% \end{eqnarray}
% The matrix ${\bf K}_A$ is real, symmetric, and positive definite so that $\hat{K}_A$ is hermitian. Consider the two-point functions $Q_{ij}=\braket{q_i q_j}$ and $P_{ij}=\braket{p_i p_j}$. Since ${\bf Q}$ and ${\bf P}$ are symmetric we define ${\bf C}_A=\sqrt{{\bf Q}_A{\bf P}_A}$ and have $\sqrt{{\bf Q}_A{\bf P}_A}^T=\sqrt{{\bf P}_A{\bf Q}_A}$. Then,
% \begin{eqnarray}
% &&\hat{K}_A=({\bf P}_A\oplus {\bf Q}_A)\lb h({\bf C}_A)\oplus h({\bf C}^T_A)\rb\nn\\
% &&h(y)=\frac{1}{y}\log\lb \frac{y+1/2}{y-1/2}\rb
% \end{eqnarray}
% Therefore, we can split $\hat{K}_A$ into two terms
% \begin{eqnarray}
% &&{\bf K}_A=\hat{K}^{(M)}_A+\hat{K}^{(N)}_A\nn\\
% &&{\bf K}^{(M)}_A=L\sum_{ij}\frac{M_{ij}}{L}q_i q_j,\qquad {\bf K}_A^{(N)}=L\sum_{ij}\frac{N_{ij}}{L}\pi_i \pi_j,
% \end{eqnarray}
% where we have multiplied and dividing by a factor of $L$ to compare to the continuum limit. Note that in the continuum limit which motivates the existence of the following limits
% \begin{eqnarray}
% &&\mu_k(x_k)=\lim_{L\to \infty}\frac{M_{i,i+k}}{L},\qquad \nu_k(x_k)=\lim_{L\to \infty}\frac{N_{i,i+k}}{L}\nn\\
% &&x_k=\frac{1}{L}\lb i+k/2\rb\ .
% \end{eqnarray}

% The calculation of full modular Hamiltonian is somewhat 

% Take the transverse Ising model and compute the characteristic functions for half-space inclusions.

% Take massless free bosons and fermions and compute the characteristic function for an interval



\NL{Kubo-Ando mean for inclusions in QFT Rindler wedges, and }

\NL{Zeros of characteristic function, and a chaos bound}

%\section{Characteristic function as a regulator, double cone, and Modular quasi-normal modes}




% \section{1D Lattice examples}


% Consider the one-dimensional quantum XY-model with the Hamiltonian
% \begin{eqnarray}
%     H--\frac{1}{2}\sum_i\lb \frac{1+\gamma}{2}X_i X_{i+1}+\frac{1-\gamma}{2}Y_i Y_{i+1}+\lambda Z_i\rb\ .
% \end{eqnarray}
% The cases $\gamma=0$ and $\gamma=1$ are called the XX model and the Ising model, respectively. This model can be solved using the Jordan-Wigner transform 
% that sends it to a second-order action in fermionic modes.

% We start with the case $\lambda=1$, the transverse Ising model. For $\lambda<1$ the system is in the disordered phase, and for $\lambda>1$ it is ordered. We define $k=\min(\lambda,\lambda^{-1})$ and $k'=\sqrt{1-k^2}$.

% The modular Hamiltonian of a half-line  $W_0=\{i| i\geq 0\}$ can be written explicitly as
% \begin{eqnarray}
% &&\lambda<1:\quad \hat{K}_{W(0)}=-2I(k')\sum_{j=0}^\infty \lb j X_jX_{j+1}+k (j-1/2) Z_j\rb\nn\\
% &&\lambda>1:\quad \hat{K}_{W(0)}=-2I(k')\sum_{j=0}^\infty \lb k j X_jX_{j+1}+ (j-1/2) Z_j\rb\nn
% \end{eqnarray}
% where $I(k')$ is a complete elliptic integral of the first kind. We would like to compute the characteristic function $\Delta_B^{1/4}\Delta_A^{-1/4}$ for half-spaces $A=W(a>0)\subset W(0)=B$. We can avoid the calculation of the full modular operator, by thinking in terms of the characteristic function as the following super-operator acting on the operators in the algebra:
% \begin{eqnarray}
% &&\Delta_B^{1/4}\Delta_A^{-1/4}a\ket{\Omega}=\mathcal{D}_{i/4}(a)\ket{\Omega}\nn\\
% &&\mathcal{D}_{i/4}(a)=e^{K_B/4}e^{-K_A/4} a e^{K_A/4}e^{-K_B/4}\ .
% \end{eqnarray}
% Define the operator $\tilde{K}=K/(2I(k'))$, then for $\lambda<1$ we have
% \begin{eqnarray}
% &&[\tilde{K},X_l]=-k(l-1/2)i Y_l\nn\\
% &&[\tilde{K},Z_l]=i (l Y_lX_{l+1}+(l-1)X_{l-1}Y_l)\nn\\
% &&[\tilde{K},Y_l]=-\ .
% \end{eqnarray}
% Use Jordan-Wigner transform to write it in terms of the fermions

%Similarly, the modular Hamiltonians of the commutants are
%\begin{eqnarray}
%&&\lambda<1:\quad \hat{K}_{W'(0)}=2I(k')\sum_{j=0}^{-\infty} \lb j X_jX_{j-1}+k (j+1/2) Z_j\rb\nn\\
%&&\lambda>1:\quad \hat{K}_{W'(0)}=2I(k')\sum_{j=0}^{-\infty} \lb k j X_jX_{j-1}+ (j+1/2) Z_j\rb\nn\ .
%\end{eqnarray}
%
%
% We first write down the full modular Hamiltonians for $\lambda<1$ 
%\begin{eqnarray}
%K_{W(0)}&&=\hat{K}_{W(0)}-\hat{K}_{W'(0)}=\nn\\
%&&=-2I(k')\sum_{j=\infty}^\infty \lb j X_jX_{j+\text{sgn}(j)}+k (j-\text{sgn}(j))/2) Z_j\rb\nn
%\end{eqnarray}
%and for $\lambda>1$ 
%\begin{eqnarray}
%K_{W(0)}&&=\hat{K}_{W(0)}-\hat{K}_{W'(0)}=\nn\\
%&&=-2I(k')\sum_{j=\infty}^\infty \lb k j X_jX_{j+\text{sgn}(j)}+ (j-\text{sgn}(j))/2) Z_j\rb\nn
%\end{eqnarray}
% 
%\NL{Finish the calculation; how about the XY-model?}

% As the second example, consider a 1d harmonic chain
% \begin{eqnarray}
% H=\sum_i \lb \frac{1}{2m}p_i^2+\frac{m\omega^2}{2}q_i^2+\frac{\kappa}{2}(q_i-q_{i+1})^2\rb
% \end{eqnarray}
% It is a discretization of massive scalar field in 1 dimensions. The modular Hamiltonian of a subregion $A$ of $L$ sites is
% \begin{eqnarray}
% &&\hat{K}_A=\frac{1}{2}\vec{r}^T\cdot {\bf K}_A \cdot \vec{r},\qquad \vec{r}=\begin{pmatrix}&\vec{q}\\& \vec{p}\end{pmatrix}
% \end{eqnarray}
% The matrix ${\bf K}_A$ is real, symmetric, and positive definite so that $\hat{K}_A$ is hermitian. Consider the two-point functions $Q_{ij}=\braket{q_i q_j}$ and $P_{ij}=\braket{p_i p_j}$. Since ${\bf Q}$ and ${\bf P}$ are symmetric we define ${\bf C}_A=\sqrt{{\bf Q}_A{\bf P}_A}$ and have $\sqrt{{\bf Q}_A{\bf P}_A}^T=\sqrt{{\bf P}_A{\bf Q}_A}$. Then,
% \begin{eqnarray}
% &&\hat{K}_A=({\bf P}_A\oplus {\bf Q}_A)\lb h({\bf C}_A)\oplus h({\bf C}^T_A)\rb\nn\\
% &&h(y)=\frac{1}{y}\log\lb \frac{y+1/2}{y-1/2}\rb
% \end{eqnarray}
% Therefore, we can split $\hat{K}_A$ into two terms
% \begin{eqnarray}
% &&{\bf K}_A=\hat{K}^{(M)}_A+\hat{K}^{(N)}_A\nn\\
% &&{\bf K}^{(M)}_A=L\sum_{ij}\frac{M_{ij}}{L}q_i q_j,\qquad {\bf K}_A^{(N)}=L\sum_{ij}\frac{N_{ij}}{L}\pi_i \pi_j,
% \end{eqnarray}
% where we have multiplied and dividing by a factor of $L$ to compare to the continuum limit. Note that in the continuum limit which motivates the existence of the following limits
% \begin{eqnarray}
% &&\mu_k(x_k)=\lim_{L\to \infty}\frac{M_{i,i+k}}{L},\qquad \nu_k(x_k)=\lim_{L\to \infty}\frac{N_{i,i+k}}{L}\nn\\
% &&x_k=\frac{1}{L}\lb i+k/2\rb\ .
% \end{eqnarray}

% The calculation of full modular Hamiltonian is somewhat 

% Take the transverse Ising model and compute the characteristic functions for half-space inclusions.

% Take massless free bosons and fermions and compute the characteristic function for an interval
%\subsection{older stuff}
\NL{No zeros and recovery maps}

\NL{Characteristic Functions and Noncommutative $L_p$-Space}




% \section{CFT Hamiltonian from a ball}

% %Emphasize $T(z)T(-z)^{-1}$.

% Consider the vacuum state of a CFT in $\mathbb{R}^{d-1,1}$ conformally embedded inside the Lorentzian cylinder $\mathbb{S}^{d-1}\times \mathbb{R}_T$ with $T$ the global time on the cylinder; see figure \ref{}. The universal cover of the conformal group $\widetilde{SO(d+1,2)}$ acts on the Lorentzian cylinder. Consider concentric ball-shaped regions $B(s)$ centered at the $x^\mu=0$ and radius $e^{-2\pi s}$. 
% \begin{theorem}
%     For any real $t$ and $s>0$ we have the inclusion of ball-shaped region $B(t+s)\subset B(t)$ we have the characteristic function for $\Im (z)\in (0,1/2)$:
%     \begin{eqnarray}
%         &&\log T_{t,s}(z)=-(\beta_BK_B+\beta_H H+\beta_F K_F)\nn\\
%         &&\beta_F=\beta\cosh(\pi s)\coth(\pi z)\nn\\
%         &&\beta_B=\beta \sinh(\pi (2t+s))\nn\\
%         &&\beta_H=\beta\cosh(\pi (2t+s))\nn\\
%         &&\beta=\frac{2i\cosh^{-1}(A)}{\sqrt{\sinh^2(\pi z)\sinh^2(\pi s)-1}}\nn\\
%         &&A=1-2\sinh^2(\pi s)\sinh^2(\pi z)
%     \end{eqnarray}
% where $B=B(0)$ is a unit ball, $s=|B(t+s)|/|B(t)|$ is the relative size of the inclusion and $|B|$ is the linear size of $B$.
% \end{theorem}
% \begin{corollary}
% In this case, at small $s$ the characteristic function is
% \begin{eqnarray}
% T_{t,s}(z)&&=1-2\pi i s \lb  \frac{\sinh(\pi z)}{2\pi}\p_t K_{B(t)}+\cosh(\pi z)K_F\rb \nn\\
% &&+O(s^2)\nn\\
%  &&\frac{1}{2\pi}\p_t K_{B(t)}=\sinh(2\pi t)K_B+\cosh(2\pi t) H\nn\\
%     &&K_{B(t)}=\cosh(2\pi t)K_B+\sinh(2\pi t) H\ .
% \end{eqnarray}
% The modular S-matrix in this limit is
% \begin{eqnarray}
% S_t(z)=1-2 i s \sinh(\pi z)\p_t K_{B(t)}+O(s^2)\ .
% \end{eqnarray}
% Further expanding for large $|z|$ we obtain the scattering amplitude operator
% \begin{eqnarray}
% \tilde{M}_t=-2 s e^{\pi z}\p_t K_{B(t)}
% \end{eqnarray}
% which once again saturates the chaos bound.
% \end{corollary}


% A local observer on a unit ball $B$ has access to $K_B$. 


% Therefore, if we want to reconstruct the global Hamiltonian $H$ using only local data $K_B$ we need to remove the $K_F$ term. Under $z\to -z$ the coefficient $\beta_F$ is invariant whereas $\beta_B$ and $\beta_H$ are invariant. 

% We are particularly interested in the limit $s\ll 1$ and $s e^{\pi \Re z}\ll 1$. We define $T_{t,s}(z)=1+i M_{t,s}(z)$ so that in this limit we find
% \begin{eqnarray}
% %\beta&&\sim 4\pi i s \sinh(\pi z)\nn\\
%     &&M_{t,s}(z)=\nn\\
%     &&-4\pi s\lb \frac{\sinh(\pi z)}{2\pi}\p_t K_{B(t)}+\cosh(\pi z)K_F\rb+O(s^2)\nn\\
%     &&\frac{1}{2\pi}\p_t K_{B(t)}=\sinh(2\pi t)K_B+\cosh(2\pi t) H\nn\\
%     &&K_{B(t)}=\cosh(2\pi t)K_B+\sinh(2\pi t) H\ .
%     \end{eqnarray}
% Note that $M_{t,s}(z)$ is an operator-valued function, but its matrix elements are analogous to the scattering amplitude defined in scattering theory.

% The expansion above makes it clear that if we compute $M_{t,s}(z)$ 

% The connection to scattering theory suggests writing $z=x+iy$ and taking a further limit of large $x$ to find
% \begin{eqnarray}
%     M=-4\pi s e^{\pi z}\lb\sinh(2\pi t)K_B+\cosh(2\pi t)K_B+K_F\rb\nn
% \end{eqnarray}


% ----

% \begin{enumerate}
%     \item   The positive operator below is independent of $s$:
%     \begin{eqnarray}
%     \tilde{H}_{t}&&\equiv \frac{1}{4\pi s}T_{t,s}(-i/4)T_{t,s}(-i/4)^\dagger \nn\\
%     &&=\cosh(2\pi t)H+\sinh(2\pi t)K_B\nn\\
%     &&=\frac{1}{2\pi}\p_t K_{B(t)}\Big|_{t=0}=-R\p_R K_{B(R)}\ .
%     \end{eqnarray}
%     \item The boundary value at $z=-i/2$ is
%     \begin{eqnarray}
%         T_{t,s}(-i/2)=4\pi s (\tanh(\pi s)K_B+H)\ .
%     \end{eqnarray}
%     \item We can construct both the global Hamiltonian $J_{(-1)0}$ and the generator of rotations on the Lorentzian cylinder in terms of the local modular Hamiltonian
%        \begin{eqnarray}
%     \begin{pmatrix}
%     J_{(-1)0}\\
%         J_{0(d+1)}
%     \end{pmatrix}
%     &&=\begin{pmatrix}
%         &\cosh(2\pi t) & \sinh(2\pi t)  \\
%          &\sinh(2\pi t)& \cosh(2\pi t)
%     \end{pmatrix}
%     \begin{pmatrix}
%         & \tilde{H}_{t} \\
%         & -K_{B(t)}
%     \end{pmatrix}\nn\\
%     &&=\frac{-1}{2R}\begin{pmatrix}
%         &1+R^2 & 1-R^2  \\
%          &1-R^2& 1+R^2
%     \end{pmatrix}
%     \begin{pmatrix}
%         & R\p_R K_{B(R)}\\
%         & K_{B(R)}
%     \end{pmatrix}\ .
% \end{eqnarray}
%     \end{enumerate}




% \begin{eqnarray}
%     T_{t,s}(x)T_{t,s}(-x)^{-1}=\mathbb{I}+2\pi s\sinh(2\pi x)\begin{pmatrix}
%         &0 & -1  \\
%          &1& 0
%     \end{pmatrix}+O(s^2)\nn\ .
% \end{eqnarray}
% This means that $T_{t,s}(x) T_{t,s}(-x)^{-1}$ generates an infinitesimal step of real-time evolution $e^{i s H}$. Interestingly, this is independent of $t$.


% Our result generalize almost trivially to Conformally flat spacetimes. The results generalize to conformally flat spacetimes including double cones in de Sitter and Anti-de Sitter space \cite{frob2023modular}. In curved spacetime, ...

% Note that nuclearity implies that $e^{-\ep \tilde{H}}$ is trace class for all $\ep>0$, which in turn implies that the spectrum of $\tilde{H}$ is discrete \footnote{Note that in Minkowski space QFT, $e^{-\beta H}$ is not trace class, and the spectrum of Hamiltonian contains cuts corresponding to the continuum of multi-particle states.}. 

% Next, consider the eternal black hole in AdS in the limit of $G_N\to 0$. The boundary theory is a GFF with time-band algebras. Assuming spherical symmetry, we can view time-band algebras as time-interval algebras of $0+1$ GFF with an infinite tower of fields corresponding to the spherical harmonics. Then, the assumption of nuclearity implies that the split property, which in turn implies that the emergent type III$_1$ algebras of GFF are hyperfinite. This is related to the conjecture of Sam and Hong \cite{}.








\section*{Acknowledgements}
We thank Keiichiro Furuya, Mudassir Moosa, Shoy Ouseph, Laimei Nie, 

% Acknowledgments
%\begin{acknowledgments}
%\end{acknowledgments}

\appendix

%\begin{widetext}
% \section{Proof of Proposition \ref{prop:density}}\label{app:proof}
% In this appendix, we prove Proposition \ref{prop:density}. To prove this proposition, recall that we need to answer two questions:
% \begin{enumerate}
%     \item if there exists a common cyclic and separating vector;
%     \item if there exists a dense subset of common cyclic and separating vectors.
% \end{enumerate}
% The existence question is answered by the following observation due to Dixmier and Marechal\cite{dixmiermarechal}:
% \begin{lemma}
%     For a sequence of von Neumann algebras $(A_1,A_2,...)$ acting on the same Hilbert space $\mathcal{H}$, if each von Neumann algebra has a cyclic and separating vector, then there exists a common cyclic and separating vector for all $A_i$'s. 
% \end{lemma}
% This lemma directly shows that $\mathcal{C}_\mathcal{A}\mathcal{A}p\mathcal{C}_\mathcal{B} \neq \emptyset$. The density question explores whether such common cyclic and separating vectors are hard to find. The answer to this question follows from two simple observations:
% \begin{lemma}
%     For any von Neumann algebra $\mathcal{A}$ in its standard form acting on a Hilbert space $\mathcal{H}$, the set of cyclic vectors is dense in $\mathcal{H}$ and the set of separating vectors is dense in $\mathcal{H}$.
% \end{lemma}
% \begin{proof}
%     Since $\mathcal{A}$ is in its standard form, there exists a cyclic vector $\ket{\xi}\in\mathcal{H}$. For any invertible operator $\mathcal{O}\in\mathcal{A}$, the vector $\mathcal{O}\ket{\xi}\in\mathcal{H}$ is also a cyclic vector:
%     \begin{equation}
%         \overline{A\mathcal{O}\ket{\xi}} = \overline{A\mathcal{O}^{-1}\mathcal{O}\ket{\xi}} = \overline{A\ket{\xi}} = \mathcal{H}
%     \end{equation}
    
%     Finally by an observation due to Dixmier and Marechal\cite{dixmiermarechal}, any element in $\mathcal{A}$ is a strong-operator limit of a sequence of invertible elements in $\mathcal{A}$. Therefore the set $\{\mathcal{O}\ket{\xi}: \mathcal{O}\in \text{Inv}(\mathcal{A})\}$ is dense in $\mathcal{H}$ where $\text{Inv}(\mathcal{A})$ is the set of invertible operators in $\mathcal{A}$.

%     Since $\mathcal{A}$ is in its standard form, its commutant $\mathcal{A}'$ acting on $\mathcal{H}$ is also in its standard form. Hence, the set of cyclic vectors of $\mathcal{A}'$ is dense in $\mathcal{H}$. Thus, the set of separating vectors of $\mathcal{A}$ is dense in $\mathcal{H}$.
% \end{proof}
% The next observation is again due to Dixmier and Marachel\cite{dixmiermarechal}. The statement applies to general von Neumann algebras.
% \begin{lemma}
%     For any von Neumann algebra $A$ acting on a Hilbert space $\mathcal{H}$ the set of cyclic vector forms a $G_\delta$ subset of $\mathcal{H}$ and the set of separating vector forms a $G_\delta$ subset of $\mathcal{H}$
% \end{lemma}
% The proof of this lemma is elementary and short. So we include it here. Recall a $G_\delta$ set of a topological space is the intersection of countably infinite open subsets. 
% \begin{proof}
%     An empty set is a $G_\delta$ set by definition.

%     Suppose the set of cyclic vectors of $A$ is non-empty, i.e. there exists $\ket{\xi_0}\in \mathcal{H}$ such that $\overline{A\ket{\xi_0}}= \mathcal{H}$. Then for any operator $\mathcal{O}\in A$, we consider the countably infinite collection of open sets:
%     \begin{equation}
%         U(\mathcal{O},n):=\{\ket{\xi}\in\mathcal{H}:||\mathcal{O}\ket{\xi} - \ket{\xi_0}|| < \frac{1}{n}\}
%     \end{equation}
%     If $\ket{\xi}\in\mathcal{H}$ is a cyclic vector of $A$, then for all $n > 0$ there exists an operator $\mathcal{O}_n$ such that $||\mathcal{O}_n\ket{\xi} - \ket{\xi_0}|| < \frac{1}{n}$. Hence we have:
%     \begin{equation}
%         \ket{\xi}\in \mathcal{A}p_{n > 0}\cup_{\mathcal{O}\in A}U(\mathcal{O}, n)
%     \end{equation}
%     Note that because $\cup_{\mathcal{O}\in A} U(\mathcal{O}, n)$ is open, the intersection $\mathcal{A}p_{n > 0}\cup_{\mathcal{O}\in A}U(\mathcal{O}, n)$ is a $G_\delta$ set. 
    
%     Conversely, any vector in this intersection is a cyclic vector of $A$. This is because the reference vector $\ket{\xi_0}$ is a cyclic vector. Thus for an arbitrary vector $\ket{\xi}\in\mathcal{H}$ and for any $\epsilon  > 0$, there exists an operator $\mathcal{O}\in A$ such that:
%     \begin{equation}
%         ||\mathcal{O}\ket{\xi_0} - \ket{\xi}|| < \frac{\epsilon}{2}
%     \end{equation}
%     For any vector $\ket{\xi'}$ in the aforementioned $G_\delta$ set, there exists another operator $\mathcal{O}'$ such that:
%     \begin{equation}
%         ||\mathcal{O}'\ket{\xi'
%         } - \ket{\xi_0}|| < \frac{\epsilon}{2 ||\mathcal{O}||}
%     \end{equation}
%     Then we have:
%     \begin{equation}
%         ||\mathcal{O}\mathcal{O}'\ket{\xi'} - \ket{\xi}||\leq||\mathcal{O}||\cdot||\mathcal{O}\ket{\xi'} - \ket{\xi_0}|| + ||\mathcal{O}\ket{\xi_0} - \ket{\xi}|| <\epsilon
%     \end{equation}
%     Hence $\overline{\mathcal{A}\ket{\xi'}} = \mathcal{H}$. 

%     Thus it follows that the set of cyclic vectors is exactly the $G_\delta$ set $\mathcal{A}p_{n >0}\cup_{\mathcal{O}\in A}U(\mathcal{O},n)$. Apply the same argument to the commutant $A'$ and since the set of separating vectors of $A$ is the set of cyclic vectors of $A'$, we arrive at the conclusion.
% \end{proof}
% Based on these two observations, the density of common cyclic and separating vectors follows from the Baire category theorem. 
% \begin{corollary}
%     For two type $III$ von Neumann algebras $\mathcal{A}, \mathcal{B}$ acting on a common separable Hilbert space $\mathcal{H}$, then the set of common cyclic and separating vectors is a dense $G_\delta$ subset of $\mathcal{H}$.
% \end{corollary}
% \begin{proof}
%     Following the previous two lemmas, for the algebra $\mathcal{A}$, the set of cyclic vectors is a dense $G_\delta$ subset of $\mathcal{H}$ and the set of separating vectors is also a dense $G_\delta$ subset of $\mathcal{H}$. By Baire category theorem, the intersection of these two sets is also a dense $G_\delta$ subset of $\mathcal{H}$, i.e. $\mathcal{C}_\mathcal{A}$ is a dense $G_\delta$ subset of $\mathcal{H}$. Similarly, $\mathcal{C}_\mathcal{B}$ is a dense $G_\delta$ subset of $\mathcal{H}$.

%     Finally, by another application of Baire category theorem, the intersection $\mathcal{C}_\mathcal{A}\mathcal{A}p\mathcal{C}_\mathcal{B}$ is a dense $G_\delta$ subset of $\mathcal{H}$.
% \end{proof}

%\subsection{$L^2$-Nuclearity Condition}\label{subsection:l2nuclearity}
%We are now ready to present the definition of $L^2$-nuclearity.
%\begin{definition}
%    Given a pair of von Neumann algebras $\mathcal{B}\subset\mathcal{A}$ with a common cyclic and separating vector $\ket{\Omega}\in\mathcal{H}$, $\mathcal{D}_{\mathcal{A}, \mathcal{B}}(-i/4) = \Delta_\mathcal{A}^{1/4}\Delta_\mathcal{B}^{-1/4}$ is a well-defined map between the standard Hilbert spaces:
%    \begin{equation}
%        \Delta_\mathcal{A}^{1/4}\Delta_\mathcal{B}^{-1/4}: L_2(\mathcal{B})\rightarrow L_2(\mathcal{A})
%    \end{equation} 
%    If $\Delta_\mathcal{A}^{1/4}\Delta_\mathcal{B}^{-1/4}$ is a nuclear operator, then we say the triplet $(\mathcal{B}\subset\mathcal{A}, \ket{\Omega})$ satisfies the \textit{$L^2$-nuclearity condition}.
%\end{definition}
%For a map between two separable Hilbert spaces, it is nuclear if and only if it is trace-class. For maps between general Banach spaces, the notion of a nuclear operator generalizes the notion of a trace-class operator \cite{nuclearandsumming}.

%For algebraic quantum field theories, one of the most important consequences of $L^2$-nuclearity is the following observation \cite{buchholz2007nuclearity}:
%\begin{theorem}
%    Given a pair of von Neumann algebras $\mathcal{B}\subset\mathcal{A}$ with a common cyclic and separating vector $\ket{\Omega}$, if it satisfies $L^2$-nuclearity condition, then the inclusion is split. More specifically, we have isomorphism:
%    \begin{equation}
%        \mathcal{B}'\vee\mathcal{A}\cong\mathcal{B}'\otimes\mathcal{A}
%    \end{equation}
%\end{theorem}
%For other important consequences of $L^2$-nuclearity, we refer the readers to \cite{buchholz2007nuclearity}. 

%In the context of algebraic quantum field theory, $L^2$-nuclearity is usually proved using representation theory. For example, in the context of massless Lorentzian conformal field theories of dimension $d + 1$, the global conformal group $SO(d + 1, 2)$ admits discrete series representations (i.e. positive energy representations). In this case, any operator representing an element in the maximal compact sub Lie algebra has discrete spectrum that is bounded from below. This provides ample examples of nuclear operators. The calculations in the main text reflect this fact.

%\section{Modular Scattering matrix}

%In scattering theory, we start with an interacting Hamiltonian that splits as 
%\begin{eqnarray}
%H=H_0+V
%\end{eqnarray}
%where $H_0$ is some reference (free) Hamiltonian with multi-particle states $\ket{\Phi_\alpha}$ and eigen-energies $E_\alpha$:
%\begin{eqnarray}
%H_0\ket{\Phi_\alpha}=E_\alpha \ket{\Phi_\alpha}\ .
%\end{eqnarray}
%The asymptotic states of the interacting theory are defined as two sets of eigenstates of the interacting Hamiltonian $H$ with the same eigenvalues $E_\alpha$ as those of $H_0$:
%\begin{eqnarray}
%H\ket{\Psi_\alpha^\pm}=E_\alpha \ket{\Psi_\alpha^\pm}
%\end{eqnarray}
%with the property that coarse-grained multi-particle states $\ket{\Phi_\alpha^\pm}$ behave like free multi-particle eigenstates $\ket{\Phi_\alpha}$. To make this coarse-graining rigorous, we consider states of the interacting theory that belong to a microcanonical ensemble of small width $\Delta E$ with a smooth wave-packet $g(\alpha)$ and evolve them for long times $t\gg 1/(\Delta E)$. Then, the requirement is that 
%\begin{eqnarray}
%\|\int e^{-i E_\alpha t}g_\alpha (\ket{\Psi^\pm_\alpha}-\ket{\Phi_\alpha})\|
%\end{eqnarray}
%can be made arbitrarily small for large $t$. The intuition is that while the free multi-particle wave-packet $\ket{\Phi_\alpha}$ can never be an eigenket of the interacting theory, as we separate excitations by a large amount, the interaction falls off so fast that they can be approximated arbitrarily well by suitably coarse-grained eigenkets $\ket{\Psi_\alpha^\pm}$. The asymptotic states can be defined as 
%\begin{eqnarray}
%&&\ket{\Psi_\alpha^\pm}=\Omega(\mp \infty)\ket{\Phi_\alpha}\nn\\
%&&\Omega(t)=e^{i Ht}e^{-i H_0t}\ .
%\end{eqnarray}
%Of course, this is only meaningful after coarse-graining, and not as an exact equality. The S-matrix is defined to be the overlap of the in and out asymptotic states
%\begin{eqnarray}
%\braket{\Phi_\beta|S \Phi_\alpha}&&=\braket{\Psi_\beta^+|\Psi_\alpha^-}=\lim_{t\to \infty}\braket{\Phi_\beta|\Omega(t)^\dagger \Omega(-t)|\Phi_\alpha}\ .
%\end{eqnarray}
%In other words, we have
%\begin{eqnarray}
%S=\lim_{t\to \infty} e^{i H_0t} e^{-2i Ht}e^{i H_0t}\ .
%\end{eqnarray} 

%In direct analogy, we define a modular S-matrix for any pair of modular Hamiltonians $K$ and $K_0$ as 
%\begin{eqnarray}
%S=\lim_{t\to \infty} e^{i K_0t} e^{-2i Kt} e^{i K_0t}\ .
%\end{eqnarray}
%There are two obvious choices: 
%\begin{enumerate}
%\item {\bf Two algebras and one state:} For a general pair of algebras $\mA$ and $\mB$ and a single state $\ket{\Omega}$, we define the modular S-matrix as 
%\begin{eqnarray}
%S_{A,B}(z)&&=e^{2\pi i z K_B }e^{-4\pi i z K_A}e^{2\pi i z K_B}\nn\\
%&&=\Delta_B^{-iz}\Delta_A^{2iz}\Delta_B^{-iz}\ .
%\end{eqnarray}
%In the special case, $\mA\subset \mB$ the operator $T(z)$ and admits an analytic in the complex $z$-domain $\Im(z)\in (0,1/2)$: 
%\begin{eqnarray}
%    S_{A\subset B}(z)=T(z)T(-z)^{-1}\ .
%\end{eqnarray}
% In particular, for imaginary values of $z=i\alpha$, the modular S-matrix is the norm of the characteristic function
% \begin{eqnarray}
%     S(i\alpha)=|\Delta^{\alpha}\Delta^{-\alpha}|\ .
% \end{eqnarray}
%In the remainder of this work, we focus on the case of inclusion, because this analyticity allows us to prove a modular chaos bound \NL{cite Maldacena et al, Simon, Faulkner}. In the special case of a common cyclic and separating vector for the inclusion of algebras, $T(z)$ is a characteristic function which involves more conditions than just analyticity such as unitarity on the other boundary (see appendix \ref{}).

%\item {\bf One algebra and two states:} For a pair of states $\psi$ and $\phi$ of the algebra $\mA$ we define
%\begin{eqnarray}
%S_{\psi,\phi}(z)=e^{i z K_\phi} e^{-2i z K_\psi}e^{i z K_\phi}=u_{\psi|\phi}(z)u_{\psi|\phi}(-z)^{-1}
%\end{eqnarray}
%where $u_{\psi|\phi}$ is a cycle which is analytic in the complex $z$-domain $\Im(z)\in (0,1/2)$. In the special case, $\phi$ majorizes $\psi$ we obtain further properties on the boundaries.

%\NL{cite de Boer et al.} \NL{Is there a chaos bound here?}
%\end{enumerate}

%When the modular S-matrix admits an analytic extension to the complex $z$ strip, for imaginary values $z=i\alpha$, the modular S-matrix is the norm of the characteristic function or the cocycle:
%\begin{eqnarray}
%&&S_{A\subset B}(z)=|T(z)|^2\nn\\
%&&S_{\psi,\phi}(z)=|u_{\psi|\phi}(i\alpha)|^2\ .
%\end{eqnarray}

%The modular S-matrix expressions can be viewed as an operator algebraic regulator for modular Hamiltonian. In the previous section, we regulated by the operator using the intuition of Euclidean path-integrals and the geometric action of $P$ and $H$ we introduced the regulator 
%\begin{eqnarray}
%&&\Delta^{1/2+it}\to e^{i \ep Ht} \Delta^{1/2+it}e^{i \ep Ht}\nn\\
%&&\Delta^{1/2+it}\to e^{i \ep Pt}\Delta^{1/2+it} e^{-i \ep Pt}\ .
%\end{eqnarray}
%Our modular $S-matrix$ can be interpreted as a regulator procedure in the absence of operators $P$ and $H$. However, in the absence of Euclidean path-integral picture, it is not apriori clear that the modular S-matrix has discrete spectrum. The assumption of nuclearity is the operator algebraic avatar of our intuition of locality coming from the Euclidean path-integrals. In the case of inclusion of algebras $\mA\subset \mB$, $L^2$-nuclearity posits that for purely the analytic continuation of the modular S-matrix to purely imaginary $z=i\alpha$ is trace-class for all $\alpha\in (0,1/2)$:
%\begin{eqnarray}
%&&\forall \alpha\in (0,1/2):\quad \tr(Z(i\alpha)) <\infty\ .
%\end{eqnarray}
%In analogy with scattering theory, we define the scattering amplitude operator $M$ using
%\begin{eqnarray}
%S=1+i M\ .
%\end{eqnarray}
%Consider a sequence of interpolating algebras $\mA\subset \mA_{n-1}\subset \mA_n\subset \mB$ such that $\lim_{n\to -\infty}\mA_n=\mA$ and $\lim_{n\to \infty}\mA_n=\mB$. Then,  we have
%\begin{eqnarray}
%&&\lim_{n\to \infty} K_{\mA_n}=K_B\nn\\
%&&\lim_{n\to -\infty} K_{\mA_n}=K_A\ .
%\end{eqnarray}
%For the characteristic function we find
%\begin{eqnarray}
%&&T_{A\subset B}=T_{A_n\subset B}T_{A\subset A_n}\nn\\
%&&\lim_{n\to \infty}T_{A_n\subset B}=1\nn\\
%&&\lim_{n\to -\infty}T_{A\subset A_n}=1\ .
%\end{eqnarray}
%Therefore, we can intuit that the relative size of the inclusion $A\subset B$ controls the size of the regulator, and we should keep it small.   
%If the inclusion of algebras has finite index, one can use it as a measure of relative size of an inclusion. In our application to QFT, we will be interested in inclusion of type III$_1$ algebras with infinite index that split. In these cases, we will use the space volume of regions to define a measure of relative size. For instance, in the inclusion of Rindler wedges, we need to take the proper distance between entangling surfaces to be small, whereas for the inclusion of balls, we take the conformal cross ratio as a measure of relative size.

%We often refer to the relative size of inclusion as the cut-off parameter $\ep$. We first expand the S-martix in small $\ep$ to obtain non-trivial scattering:
%\begin{eqnarray}
%S(t)=1+\ep i \tilde{M}(t)+O(\ep^2)
%\end{eqnarray}
%Then, take the large $t$ limit. Analyticity of $S(t)$ results in a modular chaos bound 
%\begin{eqnarray}
%\lim_{t\to \infty}\frac{\log |M(t)|}{2\pi}\leq 1\ .
%\end{eqnarray}
%Therefore, we will have to work in the approximation $\ep\ll 1$ and $1\gg 1$ but $\ep e^{2\pi t}\ll 1$. In the standard chaos bound of \cite{}, the gravitational constant $G_N$ or $\hbar$ plays the role of $\ep$, and the interpretation is that we are quantizing a classical theory corresponding to $G_N=0$ and $\hbar=0$. In contrast, our definition of chaos is genuinely quantum, and there is no clear classical limit.\NL{make a milder statement}


% \section{Subsystem flows in classical ergodic theory}\label{app:classical}

% %\subsection{Coarse-gainings and Classical net of subregion algebras}

% We begin by reviewing the measure-theoretic definition of a classical dynamical system in subsection \ref{subsec:classical-definitions}, the subspaces and subsystems of observables in subsection \ref{subsec:classical-region-subsystem}, and their role in classical ergodic theory in subsection \ref{subsec:classical-ergodic-hierarchy}. 

% In subsection \ref{subsec:classical-state-flow}, we argue that when the configuration space $X$ is equipped with a Poisson bracket (as suggested by the phase space intuition from physics) {\it the Poisson flow of a state} is a classical avatar of modular dynamics associated to every state. Similar to quantum modular flow, the Poisson flow of a state generates an automorphism of the algebra leaving the state invariant, and satisfies a classical avatar of the KMS condition \footnote{Recall that in the quantum case, these two properties uniquely fix the modular flow of a state \cite{}.}. However, crucially for us, as opposed to the quantum modular Hamiltonian, the Poisson Hamiltonian of a state does not satisfy a monotonicity condition under coarse-graining as in the quantum case \footnote{Recall that in the quantum case, this monotonicity originates from the fact that the Tomita operator of the algebra is an extension of its restriction to the subalgebra. In the classical case, there is no analog of the Tomita operator.}. This monotonicity is at the heart of non-commutative results on the reconstruction of geometry from a net of subsystems, such as half-sided modular inclusions and intersections as we review in Section \ref{}. 

% For each coarse-graining we define a notion of a coarse-graining flow. In the case of subspaces, the {\it subspace flow} has the interpretation of a generating function for entropy fluctuations. However, the focus of this work is {\it subsystem flows} which are richer.
% Classical many-body systems and field theory come equipped with a net of subsystems, and the quantum generalization of subsystem flows plays a crucial role in the emergence of geometry from a net of abstract subsystems. 

% \subsection{Reveiw of classical dynamical system}\label{subsec:classical-definitions}
% % We will see that it is the failure of conditional expectation to intertwine the microscopic phase action with its coarse-grained counterpart that is of interest to us.  The same relative density that measures this failure generates a cocycle whose expectation is exactly the Kullback--Leibler divergence, and this divergence is exactly the entropy increase produced by coarse-graining.  These results provide a rigorous kinematic account of information loss, but they do not by themselves derive an irreversible physical time evolution.

% \paragraph{Algebra of observables:} In classical mechanics, one often begins with a phase space that we take to be an orientable smooth manifold \(X\) that admits a strictly positive volume measure. Points \(x\in X\) represent microstates of the system. But from the algebraic viewpoint, the primary object is not the set of points but the commutative von Neumann algebra of observables (complex measurable functions)
% \[
% \mA=L^\infty(X,\Sigma,\mu),
% \]
% where \((X,\Sigma,\mu)\) is a standard probability space. When $X$ is a smooth phase space of finite volume, we take $\mu$ to be the normalized volume measure\footnote{$\Sigma$ is the completion of the Borel $\sigma$-algebra with respect to $\mu$. Elements of $L^\infty(X,\Sigma,\mu)$ are equivalence classes of essentially bounded measurable functions that agree $\mu$-almost everywhere.}.


% \paragraph{States:} The normalized measure $\mu$ above is interpreted as a distinguished reference state\footnote{A $\sigma$-finite volume measure of infinite total volume instead defines a distinguished weight. We return to this extension in later sections.}. General normal states correspond to probability measures $\nu$ on $X$ that are absolutely continuous with respect to $\mu$ (denoted $\nu\ll \mu$), meaning that for every measurable set $E\subseteq X$, if $\mu(E)=0$ then $\nu(E)=0$.
% There exists a $\mu$-measurable function $\rho_\nu:X\to [0,\infty)$ such that $\nu(E)=\int_E \rho_\nu d\mu$ \footnote{In mathematical terminology, $\rho_\nu$ is called a Radon-Nikodym derivative.}. In other words, every normal state can be viewed as a probability density function $\rho_\nu$ on $X$, with expectation values of an observable $f$ in this state captured by integrals on $X$: 
% \begin{align}
% \braket{f}_\nu=\int_{x\in X}f(x) d\nu(x)=\int_{x\in X} \rho_\nu(x) f(x) d\mu(x)\ .
% \end{align}

% \paragraph{GNS Hilbert space:} The GNS Hilbert space of the reference state $\mu$ is
% \[
%   \mH=L^2(X,\Sigma,\mu),
%   \qquad
%   \mA=L^\infty(X,\Sigma,\mu).
% \]
% Observables act by multiplication, and the reference state $\mu$ is represented by the constant unit vector $\ket{1}$. A general normal state with density $\rho$ is represented by
% \[
%   \ket{\sqrt{\rho}}\in\mH,
%   \qquad
%   \|\ket{\sqrt{\rho}}\|_2^2=\int_X\rho\,\dd\mu=1.
% \]
% The GNS construction builds a connection between observables and states.


% \paragraph{Dynamics:} Dynamics will be described by a measurable one-parameter group of invertible, measure-preserving transformations $\Phi_t:X\to X$. General dynamics need not be smooth. In physics, we are often interested in smooth flows described by a one-parameter family of volume-preserving diffeomorphisms of $X$. Taking a $t$-derivative of how a point $x\in X$ moves under this flow defines a vector field $V$, which preserves the volume measure if and only if $\mathcal{L}_V\mu=0$; i.e., it is divergence-free with respect to $\mu$. 




% \paragraph{Coarse-grainings:} Classical coarse-graining retains a subalgebra of observables and discards the information needed to distinguish states that agree on this subalgebra. Two useful constructions arise from a subregion of the space of microstates and from a subsystem of observed variables. Both can be formulated in measure theory, before introducing a Hamiltonian or a Poisson bracket. We first describe these constructions and then ask when the retained observables evolve autonomously.

% In the measure-theoretic setting above, a coarse description is specified by a sub-$\sigma$-algebra $\Sigma_0\subseteq\Sigma$, with associated unital von Neumann subalgebra
% \begin{align}
%     \mB=L^\infty(X,\Sigma_0,\mu)\subseteq\mA\ .
% \end{align}
% The reference measure $\mu$ determines the coarse-graining map
% \begin{align}
%     E_{\mB}(f)=\mathbb E_\mu[f\mid\Sigma_0]\ .
% \end{align}
% This is the unique $\mu$-preserving normal conditional expectation onto $\mB$. On $L^2(X,\mu)$ it acts as the orthogonal projection onto $L^2(X,\Sigma_0,\mu)$. A general normal state $\omega_\rho(f)=\int_X\rho f\,d\mu$ is coarse-grained to $\omega_0=\omega_\rho\circ E_{\mB}$, whose density is $\rho_0=E_{\mB}(\rho)$, using the $L^1$ extension of conditional expectation. The two states agree on $\mB$. Thus the subalgebra specifies the retained information, while the chosen expectation specifies how a state on this subalgebra is extended back to the full algebra.


% \paragraph{Entropy of a coarse-graining:} Consider a $\mu$-preserving coarse-graining $E$, a normal state with density $\rho$, and its coarse-grained density $\rho_0=E(\rho)$. Then, the relative entropy (KL-divergence) of these two states associates an entropy to this coarse-graining:
% \begin{align}
%     S(\rho\|\rho_0)=\int_X \rho \log(\rho/\rho_0)d\mu=H(\rho_0)-H(\rho)
% \end{align}
% where $H(\rho)=-\int_X\rho\log\rho\,d\mu$ is the entropy relative to the chosen reference measure, and the equality assumes that the two entropies are finite. The nonnegativity of relative entropy implies that coarse-graining cannot decrease this entropy. This compares a state with its coarse-grained version; it does not by itself imply monotonic entropy growth under time evolution.

% \paragraph{Generating function for entropy:} For faithful densities $\rho,\rho_0>0$ almost everywhere, the multiplicative relative phase 
% \[
% u_t
% =
% \left(\frac{\rho}{\rho_0}\right)^{it}
% \]
% is the generating function for coarse-graining entropy because 
% \begin{align}
% &-i\left.\frac{\dd}{\dd t}u_t\right|_{t=0}
% =
% \log\frac{\rho}{\rho_0},\nn\\
% &S(\rho\|\rho_0)=-i \int_X \rho (\partial_t u_t\big|_{t=0})d\mu\ .
% \end{align}
% Here the differentiation under the integral is justified when $\rho\log(\rho/\rho_0)$ is integrable. More generally, for every faithful density $\rho$ we can define the strongly continuous unitary multiplication group
% \[
%   U_t^\rho:=\Mop_{\rho^{it}}
%   =\exp\!\bigl(it\Mop_{\log\rho}\bigr),
%   \qquad
%   (U_t^\rho\psi)(x)=e^{it\log\rho(x)}\psi(x).
% \]
% Similarly, the coarse-grained density defines
% \[
%   U_t^{\rho_0}:=\Mop_{\rho_0^{it}}.
% \]
% The parameter $t$ is an auxiliary phase parameter. Nothing in this construction identifies it with laboratory time. In particular, $U_t^\rho$ does not transport points of $X$; it only rotates the phase of a Hilbert space vector pointwise\footnote{For faithful normal states on the same abelian algebra, $u_t=(\rho/\rho_0)^{it}$ is the Connes Radon-Nikodym cocycle. Relative modular operators remain nontrivial, even though the modular automorphism group of each state is trivial.}.

% \subsection{Classical subspace (single-particle) and subsystem (multi-particle) coarse-grainings}
% \label{subsec:classical-region-subsystem}

% Subspace and subsystem coarse-grainings organize incomplete information differently. An example of a subspace coarse-graining resolves microstates within a selected region of the phase space and identifies those outside it. A microstate is either inside the coarse-graining region or outside it. The set of microstates adds up in this case; that is why, in anticipation of the generalization to Hilbert space, we refer to them as {\it subspaces} so that addition is a direct sum. As we will see, there is a net of subspace algebras on the phase space of a single particle. As an example of coarse-graining to a subspace, consider the reduced state to the part of the phase space of a pendulum where the angle $\theta$ is in some limited range. 

% Next, consider a microstate described using several (potentially independent) degrees of freedom. In this case, the set of microstates multiplies, and in anticipation of generalization to Hilbert spaces, we refer to this case as {\it subsystems}, so that multiplication is tensor product. As a motivating example, consider two particles with the phase space $(q_1,p_1,q_2,p_2)$, and the restriction to the first particle as a coarse-graining on the second particle degree of freedom. 

% Note that since subspaces could have the same dimension as the entire phase space, they can be discussed even in a single-particle setup. However, subsystems necessarily live on lower-dimensional submanifolds. This intuition naturally suggests that a net of subsystems describes many-system systems or field theory. 


% % Evolution transports the resolved region in the first case and the fibers of indistinguishable microstates in the second. Their invariance has correspondingly different meanings: an invariant region defines a conserved event, while an invariant subsystem defines an autonomous factor dynamics. Ergodicity excludes nontrivial conserved events but permits nontrivial autonomous factors.

% \paragraph{Subspace coarse-grainings:}
% For a measurable region $O\subset X$, define the algebra generated by observables that are constant outside of $O$:
% \begin{align}
%     \mA_O=\{f\in\mA:\ f=f \cdot 1_O +c 1_{O^c}\text{ for some }c\in\mathbb C\}
% \end{align}
% where $1_O$ is the indicator function on $O$, and $(f\cdot 1_O)(x)=f(x)1_O(x)$ is point by point multiplication.
% This algebra resolves the microstate completely within $O$ and retains only the event ``outside $O$'' on its complement. Assuming $0<\mu(O)<1$, its conditional expectation is
% \begin{align}
%     E_O(f)=1_O f+
%     1_{O^c}\mathbb{E}_{O^c}(f)
% \end{align}
% where 
% \begin{align}
%     \mathbb{E}_{O^c}(f)=\frac{1}{\mu(1_{O^c})}\int_{O^c} f d\mu
% \end{align}
% is the average of $f$ on $O^c$.
% For $O_1\subseteq O_2$, we have $\mA_{O_1}\subseteq\mA_{O_2}$, so this construction gives a net of coarse observable algebras. Since a microstate $x\in X$ is either in $O$ or $O^c$ we have 
% \begin{align}
%     \mA_X=1_O\mA\oplus 1_{O^c}\mA
%     \simeq L^\infty(O,\mu|_O)\oplus L^\infty(O^c,\mu|_{O^c})\ .
% \end{align}
% Our unital subspace algebras have the form $\mA_O=1_O\mA\oplus\mathbb C1_{O^c}$, so $\mA_O$ and $\mA_{O^c}$ are not themselves the two direct summands. In fact,
% \begin{align}
%     \mA_O\cap\mA_{O^c}=\operatorname{span}\{1_O,1_{O^c}\}\ .
% \end{align}

% \paragraph{Subsystem coarse-grainings:}
% Suppose instead that the probability space has a product decomposition
% \begin{align}
%     (X,\Sigma,\mu)=(Y\times Z,\Sigma_Y\otimes\Sigma_Z,\mu_Y\otimes\mu_Z)\ .
% \end{align}
% A microstate now specifies both an observed variable $y$ and a hidden variable $z$. The observable algebra factorizes as
% \begin{align}
%     \mA\simeq\mA_Y\,\bar\otimes\,\mA_Z,
%     \qquad \mA_Y=L^\infty(Y,\mu_Y),\quad
%     \mA_Z=L^\infty(Z,\mu_Z)\ .
% \end{align}
% Retaining the subsystem $Y$ means retaining $\mB_Y=\mA_Y\otimes 1$. The corresponding conditional expectation averages over the hidden variables:
% \begin{align}
%     (\mathcal{E}_Y f)(y,z)=\int_Z f(y,z')\,d\mu_Z(z')\ .
% \end{align}
% This operation forgets the hidden variables at every microstate. The tensor factors describe simultaneously specified degrees of freedom. Although the reference measure is a product, a general state with density $\rho(y,z)$ can have correlations between them. 

% More generally, a measurable map $\pi:X\to Y$ in a fiber-bundle specifies coarse observables through
% \begin{align}
%     &\mB_Y=\mA_Y\circ \pi\subset \mA,
%     \qquad \mu_Y=\mu\circ \pi^{-1}\nn\\
%     &\mA_Y=L^\infty(Y,\mu_Y)
% \end{align}
% where the pre-image $\pi^{-1}$ is defined as
% \begin{align}
% \pi^{-1}(Y)=\{x\in X|\pi(x)\in Y\}
% \end{align}
% In other words, in mathematical notation, $\mB_\pi=\pi^*\mA_Y$ is the pullback of $\mA_Y$, and $\mu_Y=\pi_*\mu$ is the push-forward of the measure $\mu$ on $X$.
% % Under the standard Borel assumptions, disintegration gives conditional probability measures $\mu_y$ concentrated on the fibers $\pi^{-1}(y)$, for $\mu_Y$-almost every $y$. 
% The conditional expectation is
% \begin{align}
%     (\mathcal{E}_\pi f)(x)=\int_{\pi^{-1}(\pi(x))}f(x')\,d\mu_{\pi(x)}(x')\ .
% \end{align}
% The fiber measures may depend on $y$, and no global product decomposition is required. 

% Note that both subspace and subsystem conditional expectations are special cases of a common measure-theoretic notion of coarse-graining.

% \paragraph{Dynamically closed and invariant coarse-grainings:}
% Let $\Phi_t:X\to X$ be a measurable one-parameter group of invertible, measure-preserving transformations. It evolves observables according to
% \begin{align}
%     \alpha_t(f)=f\circ\Phi_{-t}\ .
% \end{align}
% For any coarse algebra $\mB\subseteq\mA$, we can follow the family $\mB_t=\alpha_t(\mB)$. We call the coarse description dynamically closed when
% \begin{align}\label{eq:classical-coarse-invariance}
%     \alpha_t(\mB)=\mB,\qquad t\in\mathbb R\ .
% \end{align}
% For subspace algebras $\mA_O$ as we see below, invariance of $\mA_O$ violates ergodicity.
% For the subsystem $Y$ in a product space $Y\times Z$, dynamical closure means that the $Y$-component of $\Phi_t(y,z)$ depends only on $y$ \footnote{Note that this does not follow from preservation of the product measure.}. For a coarse-graining specified by $\pi$, this means that there is a measure-preserving evolution $\Psi_t$ on the coarse space such that
% \begin{align}\label{eq:classical-factor-intertwining}
%     \pi\circ\Phi_t=\Psi_t\circ\pi\ ,
% \end{align}
% understood modulo null sets. Evolving the microstate and then retaining its coarse variables gives the same result as evolving the coarse variables themselves. In ergodic theory, such a coarse description is called a {\it factor system} in ergodic theory, and we say the coarse-grained variables evolve autonomously. Intuitively, this means that the future of the macroscopic (or observed) variables can be determined from the present macroscopic state, without the need to know the hidden microscopic details. In other words, there is an effective theory that describes the coarse-grained variables.



% % For a subregion, evolution transports its associated algebra:
% % \begin{align}
% %     \alpha_t(\mA_O)=\mA_{\Phi_t(O)}\ .
% % \end{align}
% % An invariant region therefore defines an invariant coarse algebra. For nonatomic $\mu$ and $0<\mu(O)<1$, the converse holds as well: the single atom representing the unresolved complement identifies $O^c$ from $\mA_O$. Ergodicity excludes nontrivial invariant regions, but it does not exclude nontrivial factor systems. Invariance of a subalgebra permits its observables to evolve into one another; they need not be individually fixed.

% %  For example, on the two-torus with uniform product measure, the flow
% % \begin{align}
% %     \Phi_t(y,z)=\bigl(y+t\sin(2\pi z)\!\!\pmod 1,\ z\bigr)
% % \end{align}
% % preserves the measure, but the future observed coordinate depends on the hidden coordinate. Hence $\mB_Y$ is not invariant. An autonomous observed subsystem is also weaker than a product dynamics: the hidden variables may evolve in a way that depends on the observed variables.

% \paragraph{Effective theory (Compatibility of coarse-graining and evolution):}
% For the measure-preserving dynamics above, invariance of $\mB$ is equivalent to
% \begin{align}\label{eq:classical-expectation-equivariance}
%     \mE_{\pi}\circ\alpha_t=\alpha_t\circ \mE_{\pi},
%     \qquad t\in\mathbb R\ .
% \end{align}
% Indeed, when $\mB_Y$ is invariant, $\alpha_t\circ \mE_{\pi}\circ\alpha_{-t}$ is another $\mu$-preserving conditional expectation onto $\mB_Y$, so uniqueness gives the equality. Conversely, this equality ensures that $\alpha_t$ maps the range of $\mE_{\pi}$ onto itself. The physical interpretation of this is macroscopic Predictability: The future state of the observed, macroscopic variables can be completely determined using only the present macroscopic state. This compatibility guarantees the existence of an effective Theory that governs the dynamics of the coarse-grained variables.


% \paragraph{Correlation functions of observables and subalgebras:}
% For complex functions $f,g\in\mA$, define the correlation function and its connected part as
% \begin{align}
%     &C_{f,g}(t)=\mu\bigl(f^*\alpha_t(g)\bigr)\nn\\
%     &C^{con}_{f,g}(t)
%     :=\mu\bigl(f^*\alpha_t(g)\bigr)-\mu(f^*)\mu(g),
%     \qquad \mu(f)=\int_X f\,d\mu\ .
% \end{align}
% Ergodicity is equivalent to the vanishing of the time-averaged connected correlation for every pair of observables:
% \begin{align}
%     \lim_{T\to\infty}\frac{1}{T}\int_0^T C_{f,g}(t)\,dt=0\ .
% \end{align}
% A stronger condition, {\it weak mixing}, requires
% \begin{align}
%     \lim_{T\to\infty}\frac{1}{T}\int_0^T |C_{f,g}(t)|^2\,dt=0\ .
% \end{align}
% The flow is {\it mixing} if $C_{f,g}(t)\to0$ as $t\to+\infty$. These conditions distinguish cancellation after time averaging from actual loss of correlation at late times. Mixing does not specify a decay rate and does not imply that there are no autonomous subsystems.

% \paragraph{Correlations between subspaces:} Of course, the definition of correlation above changes if we multiply the function by a constant. To account this and define a measure of correlations between two subalgebras, we define
% \begin{align}
%     C(O_1,O_2)\equiv \sup_{\text{supp}(f_i)\in O_i,|f_i|_\infty\leq 1}|C(f_1,f_2)|
% \end{align}
% where $|f|_\infty=\sup_{x\in X}|f(x)|$. It follows from the definition above that 
% \begin{align}
%     C(O_1,O_2)\leq \int_{O_1\cap O_2} 1=\mu(O_1\cap O_2)=C(1_{O_1},1_{O_2})
% \end{align}
% Therefore, we have
% \begin{align}
%     C(O_1,O_2)=C(1_{O_1},1_{O_2})\ .
% \end{align}
% For connected correlators, we find 
% \begin{align}
%     \alpha(O_1,O_2)\equiv \sup_{\text{supp}(f_i)\in O_i,|f_i|_\infty\leq 1}|C^{con}(f_1,f_2)|=\sup_{A\subset O_1,B\subset O_2}C^{con}(1_{A},1_{B})\ .
% \end{align}
% This quantity is central to the ergodic theory of single-particle systems.
% For indicator functions associated with subspaces, we have
% \begin{align}
%     C_{1_{O_1},1_{O_2}}(t)
%     =\mu\bigl(O_1\cap\Phi_t(O_2)\bigr)-\mu(O_1)\mu(O_2)\ .
% \end{align}
% which is a measure of correlations between a subspace and its time-evolved one.

% \paragraph{Correlations between subsystems:} To evaluate these suprema for subsystems, we consider factorized subsystems $Y$ and $Z$ that make up the total phase space $X = Y \times Z$ with a joint measure $\mu_{YZ}$.The test functions are $f \in \mathcal{A}_Y$ (depending only on $y$) and $g \in \mathcal{A}_Z$ (depending only on $z$), subject to the bounds $\Vert{}f\Vert{}_\infty \le 1$ and $\Vert{}g\Vert{}_\infty \le 1$. Then, the result of this optimization gives
% \begin{align}
%     &\sup_{f_1\in \mA_Y, f_2\in \mA_Z, |f_i|_\infty\leq 1}|C(f_1,f_2)|=1\nn\\
%      \alpha(Y,Z)&\equiv \sup_{f_1\in \mA_Y, f_2\in \mA_Z, |f_i|_\infty\leq 1}|C^{con}(f_1,f_2)|=\nn\\
%      &=\sup_{A\subset Y,B\subset Z}|\mu_{Y\times Z}(A\times B)-\mu_Y(A)\otimes \mu_Z(B)|
% \end{align}
% Note that $\alpha(Y,Z)$ is bounded above by trace distance
% \begin{align}
%     \alpha(Y,Z)\leq \sup_{E\subset Y\times Z}|\mu_{Y\times Z}(E)-(\mu_Y\otimes \mu_Z)(E)|=\|\mu_{Y\times Z}-\mu_Y\otimes \mu_Z\|_1
% \end{align}
% Both $\alpha(Y,Z)$ and trace distance vanish if and only if $\mu_{Y\otimes Z}=\mu_Y\otimes \mu_Z$


% \paragraph{Subspace coarse-grainings of a single non-relativistic particle:}


% \paragraph{Subspace coarse-grainings of a single free relativistic particle:}
% Consider a free relativistic particle moving on the Minkowski space $\mathbb{R}^{d,1}$, and the set of closed causally-complete regions $D$. For instance, $\{x^\in\mathbb{R}^{d,1}: x^1\geq x^0\}$ is a closed Rindler wedge. Points on the phase space (microstates) correspond to trajectories of this particle, which either intersect $D$ that we denote by $O_D$, or its causal complement $D'$ and its corresponding points on the phase space $O_{D'}$. If we ignore the trajectories that passes through the intersection $D\cap D'$ and belong to both, we see that the set of microstates add up; therefore it is natural to expect a subspace coarse-graining. The set of classical operators (measurable functions on the phase space) that are constant on $O_{D'}$ form a subalgebra that we denote by $\mA_{O_D}$. This is a subspace coarse-graining in analogy with our discussions above.

% \paragraph{Subsystem coarse-grainings of a many-body lattice model:} 

% \paragraph{Subsystem coarse-grainings of a free relativistic field theory:} Classical field theory provides a natural setting for a net of subsystem algebras with infinitely many degrees of freedom. For an unconstrained real scalar field on a Cauchy surface $\Sigma$, a phase-space point is a pair of initial data $(\phi,\pi)$. An open spatial region $U\subset\Sigma$ selects the unital polynomial algebra
% \begin{align}
%     \mathcal B_U
%     :=\operatorname{Alg}_{\mathbb R}\{1,\phi(f),\pi(g):f,g\in C_c^\infty(U)\},
%     \qquad
%     \{\phi(f),\pi(g)\}=\int_\Sigma fg\,d\Sigma\ .
% \end{align}
% Here $\phi(f)=\int_\Sigma f\phi\,d\Sigma$, and similarly for $\pi(g)$; the other brackets between generators vanish. Thus $\mathcal B_U$ is a Poisson subalgebra, and algebras associated with disjoint spatial regions Poisson commute. Retaining $\mathcal B_U$ resolves the field data in $U$ while discarding the exterior data. Although labeled by a spatial region, this is a subsystem coarse-graining: $U$ selects degrees of freedom, whereas a region of phase space selects entire field configurations.



% \subsection{Ergodic hierarchy, subspaces and subalgebras}
% \label{subsec:classical-ergodic-hierarchy}

% The ergodic hierarchy describes how much information about a state remains accessible at late times. In the algebraic language, it can be studied both through correlations of observables and through the evolution of coarse observable algebras. For our purposes, the relevant distinction is between conserved events, autonomous subsystems, and subalgebras that are preserved only in one time direction.

% \paragraph{Ergodicity, conserved subspaces and conserved subsystems:}
% Let $(\mA,\mu,\alpha_t)$ be the measure-preserving dynamical system introduced above. The flow is {\it ergodic} if every measurable set invariant under all $\Phi_t$ has measure zero or one. Equivalently, its fixed-point algebra is trivial:
% \begin{align}
%     \mA^\alpha
%     :=\{f\in\mA:\alpha_t(f)=f\text{ for all }t\in\mathbb R\}
%     =\mathbb C1\ .
% \end{align}
% Ergodicity forbids invariant subspace algebras. If for some $O$ we have 
% \begin{align}
%     \alpha_t(\mA_O)=\mA_{\Phi_t(O)}
% \end{align}
% then $1_O\in\mA^\alpha$ (it is a {\it conserved event}), which contradicts ergodicity. However, ergodicity does not exclude an invariant subsystem $\mB\subset\mA$: the condition $\alpha_t(\mB)=\mB$ permits observables in $\mB$ to evolve into one another. Such a subsystem is an autonomous factor of the dynamics (conserved subsystem). Taking two identical independent copies of a classical dynamical system always results in autonomous factors; however, a standard result in ergodic theory implies that if the original system is weakly mixing, the combined dynamical system is ergodic \footnote{As a simple example, consider two identical independent copies of a discrete-time dynamical system that is the doubling map on a unit circle $X=[0,1)$ with uniform Lebesgue measure $\mu$: $\Phi(x)=2x \mod 1$.}.

% \paragraph{Information recovered by subalgebra evolution:}
% Given a coarse algebra $\mB$, the family $\mB_t=\alpha_t(\mB)$ describes the observables accessible as this coarse description is transported by the dynamics. If $\mB_t=\mB$, evolution remains within the original description. Otherwise, the transported observables can depend on variables that were initially hidden. The algebra
% \begin{align}
%     \mA_{\mB}:=\bigvee_{t\in\mathbb R}\alpha_t(\mB)
% \end{align}
% contains the information accessible from all time translates of the coarse observables. Here $\bigvee$ denotes the generated von Neumann algebra. The condition $\mA_{\mB}=\mA$ says that these translates generate the full observable algebra. This is a property of the chosen coarse description, distinct from ergodicity of the dynamics.

% %{\paragraph Correlations among subalgebras:} We can also define 





% \paragraph{Half-sided subalgebra translations and the Kolmogorov property:}
% Of particular interest are subalgebras satisfying
% \begin{align}\label{eq:classical-one-sided-inclusion}
%     \alpha_t(\mB_+)\subseteq\mB_+,\qquad t\geq0\ .
% \end{align}
% The family $\mB_t=\alpha_t(\mB_+)$ is then decreasing: $\mB_{t+s}\subseteq\mB_t$ for $s\geq0$. To interpret this inclusion, one can take $\mB_+$ to be generated by observations on a half-infinite time interval, with the orientation chosen so that positive $t$ removes the earliest part of the record.

% The Kolmogorov, or $K$, property can be formulated through the existence of such a one-sided algebra whose translates generate the full algebra and whose tail is trivial:
% \begin{align}\label{eq:classical-K-algebra}
%     \bigvee_{t\in\mathbb R}\mB_t=\mA,
%     \qquad
%     \bigcap_{t\geq0}\mB_t=\mathbb C1\ .
% \end{align}
% The first condition says that the complete record resolves all observables. The second says that no nonconstant observable remains accessible arbitrarily far into the retained half of the record. The analogous conditions with integer time describe a $K$-automorphism. They provide an inclusion-based formulation of loss of memory and motivate the future and past algebras used in quantum ergodic theory~\cite{ouseph2024local}.

% \paragraph{Example: a Bernoulli shift:}
% Consider independent binary variables $x_n$, indexed by $n\in\mathbb Z$, with the uniform product measure on $X=\{0,1\}^{\mathbb Z}$. Let $\alpha(x_n)=x_{n+1}$, and let $\mB_+$ be the von Neumann algebra generated by $x_n$ with $n\geq0$. Then
% \begin{align}
%     \alpha^k(\mB_+)
%     =\operatorname{vN}\{x_n:n\geq k\},
%     \qquad
%     \alpha(\mB_+)\subsetneq\mB_+\ .
% \end{align}
% All integer translates generate $\mA$, and the tail algebra is trivial by the independence of the coordinates and the Kolmogorov zero-one law. At each step the retained description forgets one more coordinate, although the shift on the full system is invertible and preserves the measure. This is a subsystem coarse-graining: the retained and discarded coordinates give tensor factors of the observable algebra.

% This example also distinguishes one-sided subsystem inclusions from shrinking phase-space regions. In a finite measure space, $\Phi_t(O)\subseteq O$ and measure preservation imply $\mu(O\setminus\Phi_t(O))=0$. Thus a region cannot shrink strictly modulo null sets under a measure-preserving flow. In the nonatomic setting of Section~\ref{subsec:classical-region-subsystem}, the same obstruction applies to proper one-sided inclusions of the algebras $\mA_O$. It does not apply to the subsystem algebras of the Bernoulli example.

% \paragraph{Coarse-graining along a decreasing family:}
% Each $\mB_t$ has a $\mu$-preserving conditional expectation $\mathcal{E}_t:\mA\to\mB_t$. For $s\geq t$, the tower property gives
% \begin{align}
%     \mathcal{E}_s\mathcal{E}_t=\mathcal{E}_s\ .
% \end{align}
% For a fixed state with density $\rho$, its coarse-grained densities $\rho_t=\mathcal{E}_t(\rho)$ therefore satisfy
% \begin{align}
%     H(\rho_s)\geq H(\rho_t),
%     \qquad
%     S(\rho\|\rho_s)\geq S(\rho\|\rho_t),
%     \qquad s\geq t,
% \end{align}
% whenever the entropies are finite. Here the loss of information follows from the nesting of the retained algebras. It does not require irreversible microscopic dynamics. For a trivial tail, the reverse martingale theorem further gives $E_t(f)\to\mu(f)1$ in $L^2$ for every $f\in L^2(X,\mu)$.

% % \paragraph{Towards state-dependent subalgebra flows:}
% % The structures above depend on a specified dynamics $\alpha_t$. Measure theory alone does not assign a nontrivial flow to a state, and all the coarse algebras considered here admit $\mu$-preserving conditional expectations. In the quantum setting, modular theory supplies a distinguished dynamics, and compatibility with that dynamics controls the existence of state-preserving expectations. The one-sided relation \eqref{eq:classical-one-sided-inclusion} anticipates the inclusion structure of a half-sided modular flow, without by itself supplying its quantum modular properties.

% % To make a classical comparison with state-generated modular evolution, we next introduce a Poisson bracket and the Hamiltonian associated with $-\log\rho$. The ergodic properties of this dynamics must be assessed on its appropriate invariant sectors: on a symplectic phase space it conserves $\rho$ pointwise, so a nonconstant smooth density already supplies nontrivial conserved events on the full phase space.

% \subsection{Classical avatar of modular flow in Hamiltonian systems}
% \label{subsec:classical-state-flow}

% So far, the state and the dynamics have been independent inputs. A probability density assigns weights to microstates, but measure theory alone does not turn these weights into a motion of the microstates. A Poisson bracket supplies the missing structure: it converts a function into a Hamiltonian vector field. Applied to the logarithm of a density, it gives a distinguished state-dependent dynamics. This construction will let us compare the dynamics associated with a state and with its coarse description.

% \paragraph{Poisson brackets and Hamiltonian motion:}
% A Poisson structure on a smooth manifold $X$ is an antisymmetric bivector $\Pi$ whose bracket
% \begin{align}
%     \{f,g\}=\Pi(df,dg),\qquad f,g\in C^\infty(X),
% \end{align}
% satisfies the Jacobi identity. The bracket obeys the Leibniz rule in each argument, so fixing a smooth real Hamiltonian $H$ defines a vector field. We use the convention
% \begin{align}\label{eq:classical-Hamiltonian-convention}
%     X_H(f)&=\{f,H\},\qquad
%     \alpha_t^H(f)=f\circ\Phi_{-t}^H,\nn\\
%     \delta_H(f)&:=\left.\partial_t\alpha_t^H(f)\right|_{t=0}=\{H,f\},
% \end{align}
% where $\Phi_t^H$ is the flow of $X_H$. Thus the generator on observables has the opposite sign from the vector field transporting points. The Jacobi identity ensures that Hamiltonian evolution preserves the Poisson bracket. In canonical coordinates, $\{q_i,p_j\}=\delta_{ij}$ and the point trajectories satisfy the usual Hamilton equations.

% \paragraph{Casimirs and symplectic leaves:}
% A Casimir is a function $c$ with $\{c,f\}=0$ for every smooth $f$. It is conserved by every Hamiltonian, rather than by one chosen dynamics. The points accessible from a given microstate by successive Hamiltonian motions form a symplectic leaf. Every Casimir is constant on each leaf, although global Casimirs need not distinguish all leaves. A leaf is therefore a classical analogue of an irreducible sector for Hamiltonian motion. This analogy concerns the Poisson bracket: the pointwise algebra $L^\infty(X,\mu)$ remains abelian, and restricting it to a leaf does not turn it into a nonabelian factor.

% \paragraph{Liouville measure and the state Hamiltonian:}
% For the explicit constructions below, take $X$ to be a connected compact symplectic manifold without boundary, with symplectic form $\omega$ and normalized Liouville measure
% \begin{align}
%     d\mu=\frac{\omega^n/n!}{\int_X\omega^n/n!},\qquad \dim X=2n.
% \end{align}
% Hamiltonian flows preserve this measure. Equivalently,
% \begin{align}\label{statepreserve}
%     \int_X\{f,g\}\,d\mu=0
% \end{align}
% for smooth $f,g$. Compactness makes smooth Hamiltonian flows complete and avoids boundary terms. On a noncompact phase space these properties, together with the required integrability, must be checked separately.

% Let $\rho$ be a smooth strictly positive probability density relative to $\mu$. Define its {\it state Hamiltonian} and the associated evolution by
% \begin{align}\label{eq:classical-state-Hamiltonian}
%     K_\rho=-\log\rho,\qquad
%     \Phi_t^\rho:=\Phi_t^{K_\rho},\qquad
%     \alpha_t^\rho:=\alpha_t^{K_\rho}.
% \end{align}
% The infinitesimal evolution of an observable is
% \begin{align}\label{eq:classical-state-derivation}
%     \delta_\rho(f)=\{K_\rho,f\}=\{f,\log\rho\}.
% \end{align}
% The flow moves along level sets of the density, since $X_{K_\rho}(\rho)=0$. It preserves both $\mu$ and the state $\omega_\rho(f)=\int_X\rho f\,d\mu$. For a Gibbs density $\rho=Z^{-1}e^{-\beta H}$, we have $K_\rho=\beta H+\log Z$ and $\Phi_t^\rho=\Phi_{\beta t}^H$. Thus the state recovers the Hamiltonian motion, with a rescaled parameter; for a general density there is no prior identification of this parameter with physical time.

% \paragraph{The classical KMS identity:}
% The logarithm is distinguished by an integration-by-parts identity. Using \eqref{statepreserve},
% \begin{align}
%     0=\int_X\{g,f\rho\}\,d\mu
%     =-\omega_\rho(\{f,g\})+
%     \omega_\rho\bigl(f\{g,\log\rho\}\bigr).
% \end{align}
% Consequently,
% \begin{align}\label{KMSclassical}
%     \omega_\rho(\{f,g\})=\omega_\rho(f\delta_\rho(g)).
% \end{align}
% For a Gibbs state this reads $\omega_\beta(\{f,g\})=\beta\omega_\beta(f\delta_H(g))$. This is the classical, or static, KMS condition, written with the time orientation in \eqref{eq:classical-Hamiltonian-convention}~\cite{vanDeVen2022}. It is stronger than stationarity: for a given smooth positive density, testing against all smooth $f$ determines $\delta_\rho(g)$ for every $g$. Conversely, on a connected symplectic manifold, imposing this identity for a specified Hamiltonian dynamics determines a smooth positive density to be Gibbs, up to normalization. The identity also motivates formulations on algebras of local observables in infinite systems, where existence of the dynamics and equilibrium states requires additional hypotheses; a global Gibbs density need not exist.

% \paragraph{Poisson modular flow:}
% There is a geometric formulation of the same construction. For a smooth positive measure $\nu$, define its Poisson modular vector field by
% \begin{align}
%     \mathcal M_\nu(f):=\operatorname{div}_\nu X_f.
% \end{align}
% This measures the failure of Hamiltonian motions to preserve $\nu$~\cite{Weinstein1997}. For $d\nu=\rho\,d\mu$ and Liouville $\mu$, the divergence formula gives
% \begin{align}
%     \mathcal M_\nu(f)=X_f(\log\rho)
%     =\{f,-\log\rho\}=X_{K_\rho}(f).
% \end{align}
% Thus $\Phi_t^\rho$ is precisely this Poisson modular flow; its action on observables uses the inverse pullback, so $\delta_\rho=-\mathcal M_\nu$. On a general Poisson manifold and for an arbitrary smooth reference measure, the formula instead reads $\mathcal M_{\rho\mu}=\mathcal M_\mu+X_{-\log\rho}$. A reference measure invariant under every Hamiltonian need not exist. Our symplectic setting removes this additional obstruction.

% The construction uses the Poisson structure in an essential way. It is distinct from the multiplication group $U_t^\rho$ introduced above and from the trivial modular automorphism group of the abelian von Neumann algebra. It also supplies no automatic ergodic or entropy-production property. A nonconstant $\rho$ is itself a conserved observable, so the flow cannot be ergodic on the full probability space. Questions about ergodicity and mixing must be posed on invariant sectors with suitable invariant measures. The comparison of full and coarse state Hamiltonians, however, is meaningful without assuming ergodicity.

% \subsection{Subsystem Poisson flows in classical many-body systems and field theory}
% \label{subsec:classical-subalgebra-Hamiltonian}

% A coarse description determines a reduced state and hence, when the coarse space is Poisson, its own state-generated dynamics. We can then ask whether the reduced dynamics agrees with what the full dynamics does to the observed variables. Comparing the two motions gives a classical model of the subalgebra flows studied in this paper. The comparison separates evolution already visible within the retained algebra from evolution involving the discarded variables.

% \paragraph{Poisson subsystems:} We say a subalgebra $\mB_Y\subset \mB$ is a Poisson subsystem if for all $b_1,b_2\in \mB_Y$, we have $\{b_1,b_2\}\in \mB_Y$. Geometrically, this means that $\pi:X\to Y$ is a smooth surjective Poisson submersion:
% \begin{align}\label{eq:classical-Poisson-subsystem}
%     \{\pi^*f,\pi^*g\}_X=\pi^*\{f,g\}_Y,
%     \qquad f,g\in C^\infty(Y).
% \end{align}
% The smooth coarse observables $\mathcal B=\pi^*C^\infty(Y)$ are closed under the bracket. This condition allows us to speak of Hamiltonian dynamics within the coarse description.

% The simplest setting makes both the geometric and probabilistic structures explicit. Take a product $X=Y\times Z$ of connected compact symplectic manifolds, with product symplectic structure and normalized Liouville measure $\mu=\mu_Y\otimes\mu_Z$. The projection to $Y$ is Poisson, and the associated von Neumann subalgebra is $\mB_Y=L^\infty(Y,\mu_Y)\otimes1$. For a smooth positive density $\rho$, restriction of the state gives the marginal density
% \begin{align}\label{eq:classical-marginal-Hamiltonian}
%     \rho_Y(y)=\int_Z\rho(y,z)\,d\mu_Z(z),\qquad
%     K_Y=-\log\rho_Y.
% \end{align}
% The expectation $\mE_Y$ from Section~\ref{subsec:classical-region-subsystem} extends this marginal back to $X$ as $\rho_0(y,z)=\rho_Y(y)$. Here, $\rho_Y=\pi_*\rho$ is the classical analog of the reduced state on $Y$, and $\rho_0=\pi^*\rho_Y=\pi^*\pi_* \rho$ is the pullback of $\rho_Y$ to $X$, which is the coarse-grained version of $\rho$. Its Hamiltonian is $K_0=\pi^*K_Y$, whose flow is
% \begin{align}
%     \Phi_t^0(y,z)=(\Phi_t^Y(y),z).
% \end{align}
% This provides a definite lift of the reduced motion to the full phase space. For a general Poisson submersion, a specified coarse Hamiltonian can likewise be pulled back, but the marginal measure need not be Liouville and a complementary subsystem need not exist. We use the product setting below to keep the marginal, its reference measure, and its lifted Hamiltonian unambiguous.


% \NL{Include discussion of singular inclusions, if they exist}

% \paragraph{Agreement between reduced dynamics and full dynamics means no correlations:}
% Agreement means that the observed trajectory is determined by its initial coarse state and follows the flow generated by $K_Y$:
% \begin{align}\label{eq:classical-state-intertwining}
%     \pi\circ\Phi_t^\rho=\Phi_t^Y\circ\pi.
% \end{align}
% For our complete smooth flows, this is equivalent to the infinitesimal condition
% \begin{align}\label{eq:classical-state-defect}
%     \{\pi^*g,K_\rho\}_X-\pi^*\{g,K_Y\}_Y
%     =\{\pi^*g,K_\rho-K_0\}_X=0,
%     \qquad g\in C^\infty(Y).
% \end{align}
% The left side measures the part of the observed velocity missed by the reduced state Hamiltonian. A Hamiltonian invisible to all coarse observables belongs to the Poisson centralizer
% \begin{align}
%     \mathcal B^{\mathrm c}_{\mathrm P}
%     :=\{h\in C^\infty(X):\{h,b\}=0\text{ for every }b\in\mathcal B\}.
% \end{align}
% Thus agreement is precisely the condition $K_\rho-K_0\in\mathcal B^{\mathrm c}_{\mathrm P}$. This is an analogue of a relative commutant defined using Poisson brackets. The ordinary operator commutant of a classical subalgebra contains the whole abelian observable algebra and cannot express this condition.

% In the product setting, nondegeneracy of the symplectic form and connectedness of $Y$ imply that this centralizer consists of functions of $z$ alone. Hence \eqref{eq:classical-state-intertwining} is equivalent to
% \begin{align}\label{eq:classical-product-state-closure}
%     K_\rho(y,z)=K_Y(y)+K_Z(z),
%     \qquad
%     \rho(y,z)=\rho_Y(y)\rho_Z(z).
% \end{align}
% Normalization of the marginal fixes $\int_Z\rho_Z\,d\mu_Z=1$. Integrating the condition in (\ref{eq:classical-state-defect}) already implies that the Poisson flow of coarse-grained observables $\pi^*g$ with the coarse-grained state $\rho$ is the same as its Poisson flow with the microscopic state $\rho$. In other words:
% \begin{align}
%     \forall g\in \mB_Y:\qquad \Phi^\rho_t\circ \Phi_{-t}^0(\pi^*g)=\pi^*g\ .
% \end{align}
% This is the classical analog of {\it Takesaki's condition} in Section \ref{}.

% \paragraph{Comparing the two flows:}
% Define the relative transformation by first undoing the lifted coarse motion and then applying the full motion:
% \begin{align}\label{eq:classical-relative-transformation}
%     R_t=\Phi_t^\rho\circ\Phi_{-t}^0,
%     \qquad
%     \mathcal T_t(f):=f\circ R_t^{-1}
%     =\alpha_t^\rho\circ\alpha_{-t}^0(f).
% \end{align}
% For each $t$, this is a symplectic transformation and a measure-preserving automorphism, respectively. Since the lifted coarse flow preserves $\mB_Y$,
% \begin{align}\label{eq:classical-relative-algebra-flow}
%     \mathcal T_t(\mB_Y)=\alpha_t^\rho(\mB_Y).
% \end{align}
% The relative transformation therefore follows the same moving subalgebra as the full dynamics, while canceling the chosen internal coarse evolution. It fixes every coarse observable exactly when \eqref{eq:classical-state-intertwining} holds. For a product state, it can still move hidden variables (the relative commutant): $R_t(y,z)=(y,\Phi_t^Z(z))$.

% The infinitesimal generator of $R_t$ is the difference of the state Hamiltonians:
% \begin{align}\label{eq:classical-relative-Hamiltonian}
%     K_{\mathrm{rel}}&=K_\rho-K_0=-\log\frac{\rho}{\rho_0},\nn\\
%     \left.\partial_tR_t\right|_{t=0}&=X_{K_{\mathrm{rel}}},\qquad
%     \left.\partial_t\mathcal T_t(f)\right|_{t=0}=\{K_{\mathrm{rel}},f\}.
% \end{align}
% Every choice of state $\rho_Z(z)$ on the relative commutant $Z$ results in a conditional expectation $\pi_{\rho_Z}:X\to Y$ defined by 
% \begin{align}
%     \pi_{\rho_Z}(f(y,z))=\int_Z \rho_z(y)f(y,z)\ .
% \end{align}
% In fact, there is a one-to-one correspondence between the states of the relative commutant and conditional expectations $\mE_\pi:X\to Y$.


% % At finite time, $R_t$ need not be a one-parameter group: the two Hamiltonian flows need not commute. It reduces to the autonomous flow generated by $K_{\mathrm{rel}}$ when they do commute. The relative transformation and the autonomous flow of the difference Hamiltonian agree to first order, but should not be identified in general. This distinction also appears when composing the modular evolutions of two quantum algebras.

% \paragraph{Example: a uniform marginal with nonautonomous observables:}
% Take $Y$ and $Z$ to be two-dimensional tori with angular coordinates $(q,p)$ and $(Q,P)$, each of period $2\pi$, and the product symplectic form $dq\wedge dp+dQ\wedge dP$. Consider
% \begin{align}
%     \rho(q,p,Q,P)=\frac{e^{\epsilon\cos(q-Q)}}{\mathcal Z_\epsilon},
%     \qquad
%     \mathcal Z_\epsilon=\int_X e^{\epsilon\cos(q-Q)}\,d\mu,
%     \qquad \epsilon\ne0.
% \end{align}
% Translation invariance of the angular integral gives $\rho_Y=1$, so $K_0=0$ and the reduced state predicts no motion. The full state Hamiltonian is $K_\rho=-\epsilon\cos(q-Q)+\log\mathcal Z_\epsilon$, giving
% \begin{align}
%     \dot q=\dot Q=0,\qquad
%     \dot p=-\epsilon\sin(q-Q),\qquad
%     \dot P=\epsilon\sin(q-Q).
% \end{align}
% For example, the bounded coarse observable $e^{ip}$ evolves as
% \begin{align}
%     \alpha_t^\rho(e^{ip})
%     =\exp\bigl(i[p+\epsilon t\sin(q-Q)]\bigr).
% \end{align}
% Its future value depends on the hidden angle $Q$, so it leaves $\mB_Y$ for $t\ne0$. The full state and its uniform marginal are stationary, but the retained observables do not evolve autonomously. Here $R_t=\Phi_t^\rho$ directly displays the motion that the marginal state Hamiltonian misses.

% \paragraph{Relative entropy and the information discarded:}
% The same relative logarithm appeared in the phase construction at the beginning of this section. Here it generates a vector field, while its expectation gives
% \begin{align}\label{eq:classical-relative-entropy-Hamiltonian}
%     S(\rho\|\rho_0)=-\omega_\rho(K_{\mathrm{rel}})
%     =H(\rho_0)-H(\rho).
% \end{align}
% The vector field alone does not determine this number: on a connected symplectic manifold it determines its Hamiltonian only up to a constant. The normalized densities fix that constant. On a general Poisson manifold, Casimirs give a larger ambiguity.

% Information loss and failure of autonomous coarse dynamics are also distinct. For a product state, the relative transformation acts only on $Z$, yet
% \begin{align}
%     S(\rho_Y\rho_Z\|\rho_Y\otimes1)
%     =\int_Z\rho_Z\log\rho_Z\,d\mu_Z
% \end{align}
% can be positive. Coarse-graining has discarded the hidden distribution even without correlations. More generally,
% \begin{align}
%     S(\rho\|\rho_Y\otimes1)
%     =S(\rho\|\rho_Y\rho_Z)+S(\rho_Z\|1),
% \end{align}
% where $\rho_Z$ is the other marginal. The first term is the mutual information and the second measures the hidden marginal relative to the chosen reference. Neither this entropy nor the relative transformation supplies a monotone entropy-production law without an additional coarse-graining or dynamical assumption.

% \paragraph{Subregions and the passage to field theory:}
% The measurable region algebras $\mA_O$ remain valid coarse descriptions. What can fail is the smoothness needed to assign their coarse densities a Hamiltonian flow. For the region expectation $E_O$,
% \begin{align}
%     \rho_0=E_O(\rho)=\chi_O\rho+c_O\chi_{O^c},\qquad
%     c_O=\frac{1}{\mu(O^c)}\int_{O^c}\rho\,d\mu.
% \end{align}
% For an open region with smooth boundary this density is generally discontinuous there. In the interior of $O$, $K_0=K_\rho$; in the interior of the complement, $K_0$ is constant. These local statements do not specify a smooth global Hamiltonian flow across the boundary. Additional regularization or boundary conditions are needed. This is a limitation of the Hamiltonian construction, not of the measure-theoretic net.

% Finally, every classical coarse algebra still admits a state-preserving conditional expectation for a faithful probability state, even when its Poisson state flow fails to preserve that algebra. In the product example this expectation is explicitly
% \begin{align}
%     E_{\rho,Y}(f)(y)=\int_Z f(y,z)\frac{\rho(y,z)}{\rho_Y(y)}\,d\mu_Z(z).
% \end{align}
% Its existence does not imply \eqref{eq:classical-state-intertwining}. In quantum modular theory the relation changes: for a faithful normal state $\varphi$ on $\mA$, Takesaki's theorem identifies invariance of a unital von Neumann subalgebra $\mB$ under $\sigma_t^\varphi$ with the existence of a $\varphi$-preserving normal conditional expectation onto $\mB$~\cite{Takesaki1972}. Failure of invariance rules out an expectation preserving this state; it does not by itself rule out expectations preserving other states. This is the additional structure that the quantum theory contributes to the comparison of full and coarse dynamics.

% The classical construction thus motivates studying how relative state evolutions transport retained observables. In QFT a spacetime region selects local field observables, whereas a region $O\subset X$ above selects possible microstates. The connection developed in the next section is through the inclusions of observable algebras and their state-dependent flows; it does not identify these two notions of region.

% % Bibliography addition for the new citation key (insert into modularscattering.bib):
% % @misc{vanDeVen2022,
% %   author = {van de Ven, Christiaan J. F.},
% %   title = {Gibbs states and their classical limit},
% %   year = {2022},
% %   eprint = {2211.01755},
% %   archivePrefix = {arXiv},
% %   primaryClass = {math-ph},
% %   url = {https://arxiv.org/abs/2211.01755}
% % }

% Choose an ambient Poisson algebra $\mathcal A$ containing these observables and stable under a complete Hamiltonian evolution. With our convention for the pullback, the transported subsystem is
% \begin{align}
%     \alpha_t^H(F)=F\circ\Phi_{-t}^H,
%     \qquad \partial_t\alpha_t^H(F)=\{H,\alpha_t^H(F)\},
%     \qquad \mathcal B_{U,t}=\alpha_t^H(\mathcal B_U)\ .
% \end{align}
% Hamiltonian evolution preserves Poisson brackets, so every $\mathcal B_{U,t}$ remains a Poisson subalgebra. Invariance for all $t$ gives an autonomous reversible subsystem. More generally, propagation can make the transported observables depend on initially discarded data. Whether their time translates generate $\mathcal A$, or instead form a nested family, is a property of the chosen subsystem and dynamics together.

% A simple example of strict nesting is supplied by a classical chiral current $j(x)$ on the line. On rapidly decreasing smooth configurations, choose
% \begin{align}
%     \{j(x),j(y)\}=-\partial_x\delta(x-y),
%     \qquad P=\frac12\int_{\mathbb R}j(x)^2\,dx,
%     \qquad \tau_a j(x)=j(x+a)\ .
% \end{align}
% The nonnegative Hamiltonian $P$ generates $\tau_a$ through $\{P,j(x)\}=\partial_xj(x)$. Let $\mathcal B_b$ be generated by $1$ and $j(f)$ with $f\in C_c^\infty((b,\infty))$. Since $\tau_a j(f)=j(f(\,\cdot-a))$,
% \begin{align}
%     \tau_a(\mathcal B_b)=\mathcal B_{b+a}\subsetneq\mathcal B_b\quad(a>0),
%     \qquad
%     \bigcup_{b\in\mathbb R}\mathcal B_b=\mathcal A_{\mathrm{loc}},
%     \qquad
%     \bigcap_{b\geq0}\mathcal B_b=\mathbb R1\ .
% \end{align}
% Here $\mathcal A_{\mathrm{loc}}$ is the polynomial algebra of compactly smeared currents; $P$ is a global functional generating its automorphisms. The full translation group is reversible, but its positive half acts on $\mathcal B_0$ by injective, nonsurjective Poisson endomorphisms. The decreasing family progressively discards accessible field degrees of freedom.

% The dilation Hamiltonian $D=\frac12\int_{\mathbb R}xj(x)^2\,dx$ gives a further connection with half-sided modular inclusions. Its flow $\delta_s$ satisfies
% \begin{align}
%     \{D,P\}=P,
%     \qquad \delta_s(\mathcal B_b)=\mathcal B_{e^s b},
%     \qquad \delta_s\tau_a\delta_{-s}=\tau_{e^s a}\ .
% \end{align}
% Thus dilations preserve $\mathcal B_0$ and act half-sidedly on $\mathcal B_b$ for $b>0$. These are classical Hamiltonian realizations of the translation--dilation relations. Their construction uses the Poisson structure; identifying the dilation with a state-generated modular evolution requires additional input.


\section{Subsystem flows in classical ergodic theory}%
\label{app:classical}
 
We begin by reviewing the measure-theoretic definition of a classical dynamical system
in Section~\ref{subsec:classical-definitions}, the subspaces and subsystems of
observables in Section~\ref{subsec:classical-region-subsystem}, and their role in
classical ergodic theory in Section~\ref{subsec:classical-ergodic-hierarchy}.
 
In Section~\ref{subsec:classical-state-flow}, we argue that when the configuration
space $X$ is equipped with a Poisson bracket (as suggested by the phase-space intuition
from physics), the \emph{Poisson flow of a state} is a classical avatar of modular
dynamics associated to every state.  Similar to quantum modular flow, the Poisson flow
of a state generates an automorphism of the algebra leaving the state invariant, and
satisfies a classical avatar of the KMS condition.%
\footnote{Recall that in the quantum case, these two properties uniquely fix the modular
  flow of a state~\cite{}.}
However, crucially for us, as opposed to the quantum modular Hamiltonian, the Poisson
Hamiltonian of a state does not satisfy a monotonicity condition under coarse-graining
as in the quantum case.%
\footnote{Recall that in the quantum case, this monotonicity originates from the fact
  that the Tomita operator of the algebra is an extension of its restriction to the
  subalgebra.  In the classical case, there is no analog of the Tomita operator.}
This monotonicity is at the heart of non-commutative results on the reconstruction of
geometry from a net of subsystems, such as half-sided modular inclusions and
intersections as reviewed in Section~\ref{subsec:geometry-K-systems}.
 
For each coarse-graining we define a notion of a coarse-graining flow.  In the case of
subspaces, the \emph{subspace flow} has the interpretation of a generating function for
entropy fluctuations.  The focus of this work is \emph{subsystem flows}, which are
richer.  Classical many-body systems and field theory come equipped with a net of
subsystems, and the quantum generalization of subsystem flows plays a crucial role in
the emergence of geometry from a net of abstract subsystems.
 
% -------------------------------------------------------------------
\subsection{Review of classical dynamical systems}%
\label{subsec:classical-definitions}
 
\paragraph{Algebra of observables.}
In classical mechanics, one often begins with a phase space that we take to be an
orientable smooth manifold $X$ that admits a strictly positive volume measure.  Points
$x\in X$ represent microstates of the system.  From the algebraic viewpoint, the
primary object is not the set of points but the commutative von Neumann algebra of
observables (complex measurable functions)
\[
\mA=L^\infty(X,\Sigma,\mu),
\]
where $(X,\Sigma,\mu)$ is a standard probability space.  When $X$ is a smooth phase
space of finite volume, we take $\mu$ to be the normalized volume measure.%
\footnote{$\Sigma$ is the completion of the Borel $\sigma$-algebra with respect to
  $\mu$.  Elements of $L^\infty(X,\Sigma,\mu)$ are equivalence classes of essentially
  bounded measurable functions that agree $\mu$-almost everywhere.}
 
\paragraph{States.}
The normalized measure $\mu$ is interpreted as a distinguished reference state.%
\footnote{A $\sigma$-finite volume measure of infinite total volume instead defines a
  distinguished weight.  We return to this extension in later sections.}
General normal states correspond to probability measures $\nu$ on $X$ that are
absolutely continuous with respect to $\mu$ (denoted $\nu\ll \mu$): for every
measurable set $E\subseteq X$, $\mu(E)=0$ implies $\nu(E)=0$.  There exists a
$\mu$-measurable function $\rho_\nu:X\to [0,\infty)$ such that
$\nu(E)=\int_E \rho_\nu\,d\mu$.%
\footnote{In mathematical terminology, $\rho_\nu$ is called a Radon-Nikodym derivative.}
Every normal state is thus a probability density function $\rho_\nu$ on $X$, with
expectation values
\begin{align}
\langle f\rangle_\nu=\int_{x\in X}f(x)\,d\nu(x)=\int_{x\in X} \rho_\nu(x) f(x)\,d\mu(x).
\end{align}
 
\paragraph{GNS Hilbert space.}
The GNS Hilbert space of the reference state $\mu$ is
\[
  \mH=L^2(X,\Sigma,\mu), \qquad \mA=L^\infty(X,\Sigma,\mu).
\]
Observables act by multiplication, and the reference state $\mu$ is represented by the
constant unit vector $|1\rangle$.  A general normal state with density $\rho$ is
represented by
\[
  |\sqrt{\rho}\rangle\in\mH, \qquad
  \||\sqrt{\rho}\rangle\|_2^2=\int_X\rho\,d\mu=1.
\]
 
\paragraph{Dynamics.}
Dynamics is described by a measurable one-parameter group of invertible,
measure-preserving transformations $\Phi_t:X\to X$.  In physics, one is often
interested in smooth flows described by volume-preserving diffeomorphisms.  Taking a
$t$-derivative of how a point $x\in X$ moves under this flow defines a vector field
$V$, which preserves the volume measure if and only if $\mathcal{L}_V\mu=0$, i.e., it
is divergence-free with respect to $\mu$.
 
\paragraph{Coarse-grainings.}
Classical coarse-graining retains a subalgebra of observables and discards the
information needed to distinguish states that agree on this subalgebra.  Two useful
constructions arise from a subregion of the space of microstates and from a subsystem
of observed variables.  Both can be formulated in measure theory, before introducing a
Hamiltonian or a Poisson bracket.
 
In the measure-theoretic setting above, a coarse description is specified by a
sub-$\sigma$-algebra $\Sigma_0\subseteq\Sigma$, with associated unital von Neumann
subalgebra
\begin{align}
    \mB=L^\infty(X,\Sigma_0,\mu)\subseteq\mA.
\end{align}
The reference measure $\mu$ determines the coarse-graining map
\begin{align}
    E_{\mB}(f)=\mathbb{E}_\mu[f\mid\Sigma_0].
\end{align}
This is the unique $\mu$-preserving normal conditional expectation onto $\mB$.  On
$L^2(X,\mu)$ it acts as the orthogonal projection onto $L^2(X,\Sigma_0,\mu)$.  A
general normal state $\omega_\rho(f)=\int_X\rho f\,d\mu$ is coarse-grained to
$\omega_0=\omega_\rho\circ E_{\mB}$, whose density is $\rho_0=E_{\mB}(\rho)$.  The
two states agree on $\mB$.  Thus the subalgebra specifies the retained information,
while the chosen expectation specifies how a state on this subalgebra is extended back
to the full algebra.
 
\paragraph{Entropy of a coarse-graining.}
For a $\mu$-preserving coarse-graining $E$, a normal state with density $\rho$, and
its coarse-grained density $\rho_0=E(\rho)$, the relative entropy (KL-divergence)
associates an entropy to this coarse-graining:
\begin{align}
    S(\rho\|\rho_0)=\int_X \rho \log(\rho/\rho_0)\,d\mu=H(\rho_0)-H(\rho),
\end{align}
where $H(\rho)=-\int_X\rho\log\rho\,d\mu$ is the entropy relative to the reference
measure $\mu$, and the equality assumes the entropies are finite.  Non-negativity of
relative entropy implies that coarse-graining cannot decrease this entropy.  This
compares a state with its coarse-grained version; it does not by itself imply monotonic
entropy growth under time evolution.
 
\paragraph{Generating function for entropy.}
For faithful densities $\rho,\rho_0>0$ almost everywhere, the multiplicative relative
phase
\[
u_t = \left(\frac{\rho}{\rho_0}\right)^{it}
\]
is the generating function for coarse-graining entropy:
\begin{align}
&-i\left.\frac{d}{dt}u_t\right|_{t=0} = \log\frac{\rho}{\rho_0},\qquad
S(\rho\|\rho_0)=-i \int_X \rho (\partial_t u_t\big|_{t=0})\,d\mu.
\end{align}
More generally, for every faithful density $\rho$ we define the strongly continuous
unitary multiplication group
\[
  U_t^\rho:=\Mop_{\rho^{it}}
  =\exp\!\bigl(it\Mop_{\log\rho}\bigr),
  \qquad
  (U_t^\rho\psi)(x)=e^{it\log\rho(x)}\psi(x).
\]
Similarly, the coarse-grained density defines $U_t^{\rho_0}:=\Mop_{\rho_0^{it}}$.
The parameter $t$ is an auxiliary phase parameter with no identification with laboratory
time.  In particular, $U_t^\rho$ does not transport points of $X$; it only rotates the
phase of a Hilbert space vector pointwise.%
\footnote{For faithful normal states on the same abelian algebra,
  $u_t=(\rho/\rho_0)^{it}$ is the Connes Radon-Nikodym cocycle.  Relative modular
  operators remain nontrivial, even though the modular automorphism group of each state
  is trivial.}
 
% -------------------------------------------------------------------
\subsection{Subspace (single-particle) and subsystem (multi-particle) coarse-grainings}
\label{subsec:classical-region-subsystem}
 
Subspace and subsystem coarse-grainings organize incomplete information differently.
An example of a subspace coarse-graining resolves microstates within a selected region
of the phase space and identifies those outside it.  A microstate is either inside the
coarse-graining region or outside it.  The set of microstates adds up in this case;
that is why, in anticipation of the generalization to Hilbert space, we refer to them
as \emph{subspaces} so that addition is a direct sum.  As we will see, there is a net
of subspace algebras on the phase space of a single particle.
 
Next, consider a microstate described using several (potentially independent) degrees
of freedom.  In this case, the set of microstates multiplies, and in anticipation of
generalization to Hilbert spaces, we refer to this case as \emph{subsystems}, so that
multiplication is tensor product.  Note that since subspaces could have the same
dimension as the entire phase space, they can be discussed even in a single-particle
setup.  However, subsystems necessarily live on lower-dimensional submanifolds.  This
intuition naturally suggests that a net of subsystems describes many-body systems or
field theory.
 
\paragraph{Subspace coarse-grainings.}
For a measurable region $O\subset X$, define the algebra generated by observables that
are constant outside of $O$:
\begin{align}
    \mA_O=\{f\in\mA:\; f=f \cdot 1_O +c\, 1_{O^c}\text{ for some }c\in\mathbb{C}\},
\end{align}
where $1_O$ is the indicator function on $O$ and $(f\cdot 1_O)(x)=f(x)1_O(x)$.  This
algebra resolves the microstate completely within $O$ and retains only the event
``outside $O$'' on its complement.  Assuming $0<\mu(O)<1$, its conditional expectation
is
\begin{align}
    E_O(f)=1_O f + 1_{O^c}\mathbb{E}_{O^c}(f),\qquad
    \mathbb{E}_{O^c}(f)=\frac{1}{\mu(1_{O^c})}\int_{O^c} f\,d\mu.
\end{align}
For $O_1\subseteq O_2$, we have $\mA_{O_1}\subseteq\mA_{O_2}$, giving a net of coarse
observable algebras.  Since a microstate $x\in X$ is either in $O$ or $O^c$,
\begin{align}
    \mA_X=1_O\mA\oplus 1_{O^c}\mA
    \simeq L^\infty(O,\mu|_O)\oplus L^\infty(O^c,\mu|_{O^c}).
\end{align}
Our unital subspace algebras have the form $\mA_O=1_O\mA\oplus\mathbb{C}1_{O^c}$, so
$\mA_O\cap\mA_{O^c}=\operatorname{span}\{1_O,1_{O^c}\}$.
 
\paragraph{Subsystem coarse-grainings.}
Suppose instead that the probability space has a product decomposition
\begin{align}
    (X,\Sigma,\mu)=(Y\times Z,\Sigma_Y\otimes\Sigma_Z,\mu_Y\otimes\mu_Z).
\end{align}
A microstate now specifies both an observed variable $y$ and a hidden variable $z$.
The observable algebra factorizes as
$\mA\simeq\mA_Y\,\bar\otimes\,\mA_Z$, with
$\mA_Y=L^\infty(Y,\mu_Y)$ and $\mA_Z=L^\infty(Z,\mu_Z)$.  Retaining the subsystem
$Y$ means retaining $\mB_Y=\mA_Y\otimes 1$.  The corresponding conditional expectation
averages over the hidden variables:
\begin{align}
    (\mathcal{E}_Y f)(y,z)=\int_Z f(y,z')\,d\mu_Z(z').
\end{align}
More generally, a measurable map $\pi:X\to Y$ in a fiber-bundle specifies coarse
observables $\mB_\pi=\pi^*\mA_Y$ and marginal $\mu_Y=\pi_*\mu$.  The conditional
expectation is
\begin{align}
    (\mathcal{E}_\pi f)(x)=\int_{\pi^{-1}(\pi(x))}f(x')\,d\mu_{\pi(x)}(x').
\end{align}
Note that both subspace and subsystem conditional expectations are special cases of a
common measure-theoretic notion of coarse-graining.
 
\paragraph{Dynamically closed and invariant coarse-grainings.}
Let $\Phi_t:X\to X$ be a measurable one-parameter group of invertible,
measure-preserving transformations.  It evolves observables according to
$\alpha_t(f)=f\circ\Phi_{-t}$.  For any coarse algebra $\mB\subseteq\mA$, we can
follow the family $\mB_t=\alpha_t(\mB)$.  We call the coarse description dynamically
closed when
\begin{align}\label{eq:classical-coarse-invariance}
    \alpha_t(\mB)=\mB,\qquad t\in\mathbb{R}.
\end{align}
For subspace algebras $\mA_O$, invariance of $\mA_O$ violates ergodicity.  For the
subsystem $Y$ in a product space $Y\times Z$, dynamical closure means the $Y$-component
of $\Phi_t(y,z)$ depends only on $y$.%
\footnote{Note that this does not follow from preservation of the product measure.}
For a coarse-graining specified by $\pi$, this means there is a measure-preserving
evolution $\Psi_t$ on the coarse space such that
\begin{align}\label{eq:classical-factor-intertwining}
    \pi\circ\Phi_t=\Psi_t\circ\pi,
\end{align}
modulo null sets.  Such a coarse description is called a \emph{factor system}, and we
say the coarse-grained variables evolve autonomously.
 
\paragraph{Effective theory (compatibility of coarse-graining and evolution).}
For the measure-preserving dynamics above, invariance of $\mB$ is equivalent to
\begin{align}\label{eq:classical-expectation-equivariance}
    \mE_{\pi}\circ\alpha_t=\alpha_t\circ \mE_{\pi},\qquad t\in\mathbb{R}.
\end{align}
The physical interpretation is macroscopic predictability: the future state of the
observed, macroscopic variables can be completely determined using only the present
macroscopic state, guaranteeing the existence of an effective theory for the
coarse-grained variables.
 
\paragraph{Correlation functions of observables and subalgebras.}
For complex functions $f,g\in\mA$, define the correlation function and its connected
part as
\begin{align}
    &C_{f,g}(t)=\mu\bigl(f^*\alpha_t(g)\bigr),\quad
    C^{con}_{f,g}(t)
    :=\mu\bigl(f^*\alpha_t(g)\bigr)-\mu(f^*)\mu(g),\quad
    \mu(f)=\int_X f\,d\mu.
\end{align}
Ergodicity is equivalent to the vanishing of the time-averaged connected correlation:
$\lim_{T\to\infty}\frac{1}{T}\int_0^T C_{f,g}(t)\,dt=0$.  A stronger condition,
\emph{weak mixing}, requires
$\lim_{T\to\infty}\frac{1}{T}\int_0^T |C_{f,g}(t)|^2\,dt=0$.  The flow is
\emph{mixing} if $C_{f,g}(t)\to0$ as $t\to+\infty$.  These conditions distinguish
cancellation after time averaging from actual loss of correlation at late times.
 
\paragraph{Correlations between subspaces.}
\begin{align}
    C(O_1,O_2)\equiv \sup_{\operatorname{supp}(f_i)\in O_i,|f_i|_\infty\leq 1}|C(f_1,f_2)|,\qquad
    C(O_1,O_2)=C(1_{O_1},1_{O_2}).
\end{align}
For connected correlators,
$\alpha(O_1,O_2)\equiv\sup_{\operatorname{supp}(f_i)\in O_i,|f_i|_\infty\leq1}|C^{con}(f_1,f_2)|
=\sup_{A\subset O_1,B\subset O_2}C^{con}(1_A,1_B)$, and
$C_{1_{O_1},1_{O_2}}(t)=\mu(O_1\cap\Phi_t(O_2))-\mu(O_1)\mu(O_2)$.
 
\paragraph{Correlations between subsystems.}
For factorized subsystems $Y$ and $Z$ with $X=Y\times Z$ and joint measure $\mu_{YZ}$,
with test functions $f\in\mathcal{A}_Y$ and $g\in\mathcal{A}_Z$ subject to
$\|f\|_\infty\le1$, $\|g\|_\infty\le1$:
\begin{align}
     \alpha(Y,Z)&\equiv \sup_{f_1\in \mA_Y, f_2\in \mA_Z, |f_i|_\infty\leq 1}|C^{con}(f_1,f_2)|
     =\sup_{A\subset Y,B\subset Z}|\mu_{Y\times Z}(A\times B)-\mu_Y(A)\otimes \mu_Z(B)|.
\end{align}
Note that $\alpha(Y,Z)$ is bounded above by trace distance:
\begin{align}
    \alpha(Y,Z)\leq \|\mu_{Y\times Z}-\mu_Y\otimes \mu_Z\|_1.
\end{align}
Both $\alpha(Y,Z)$ and the trace distance vanish if and only if $\mu_{Y\times Z}=\mu_Y\otimes\mu_Z$.
 
\paragraph{Subspace coarse-grainings of a single non-relativistic particle.}
Consider a particle moving on a compact phase space $X=(q,p)\in\mathbb{T}^{2n}$
(position $q\in\mathbb{T}^n$ and momentum $p\in\mathbb{T}^n$), with normalized
Liouville measure $\mu$.  For a region $O\subset X$ with $0<\mu(O)<1$, the algebra
$\mA_O$ defined above resolves the microstate completely within $O$ and retains only
the event ``outside $O$'' on its complement.  For example, for a one-dimensional
harmonic oscillator on the torus, $O$ can be the half of phase space with $p>0$; the
algebra $\mA_O$ then contains all functions of $(q,p)$ that are constant for $p\le0$.
The dynamical closure condition $\alpha_t(\mA_O)=\mA_O$ requires $\Phi_t(O)=O$ modulo
null sets, which for a measure-preserving flow and nonatomic $\mu$ forces $O$ to be
essentially trivial or invariant—contradicting ergodicity.  Thus ergodicity precludes
nontrivial invariant subspace algebras, and the ergodic properties of the system must
be studied on the full phase space or on invariant level sets of constants of motion.
 
\paragraph{Subspace coarse-grainings of a single free relativistic particle.}
Consider a free relativistic particle moving on Minkowski space $\mathbb{R}^{d,1}$,
and the set of closed causally-complete regions $D$.  For instance,
$\{x\in\mathbb{R}^{d,1}:x^1\ge x^0\}$ is a closed Rindler wedge.  Points on the
phase space (microstates) correspond to trajectories of this particle, which either
intersect $D$ (denoted $O_D$) or its causal complement $D'$ (denoted $O_{D'}$).
Ignoring trajectories through $D\cap D'$, the set of microstates adds up; this is
therefore a subspace coarse-graining.  The set of classical operators (measurable
functions on the phase space) that are constant on $O_{D'}$ form a subalgebra
$\mA_{O_D}$.  This subspace coarse-graining is the classical precursor of the local
algebra $\mathcal{A}(D)$ of quantum field theory.
 
\paragraph{Subsystem coarse-grainings of a many-body lattice model.}
Consider a one-dimensional spin chain: microstates are elements $x=(x_n)_{n\in\mathbb{Z}}$
of a product space $X=\prod_{n\in\mathbb{Z}}X_n$, where each site $X_n=\{0,1,\ldots,d-1\}$
is a finite set with uniform measure, and the overall reference measure
$\mu=\bigotimes_n\mu_n$ is the product measure.  A spatial subsystem is a subset
$\Lambda\subset\mathbb{Z}$ of sites, selecting the product subalgebra
\begin{align}
    \mB_\Lambda = L^\infty\!\left(\prod_{n\in\Lambda}X_n,\mu|_\Lambda\right)\otimes 1_{X\setminus\Lambda}.
\end{align}
This is a subsystem coarse-graining: the retained degrees of freedom are the
configuration of spins on $\Lambda$, and the conditional expectation $\mE_\Lambda$
averages over the spin configuration on $\mathbb{Z}\setminus\Lambda$.  For a
half-infinite subsystem $\Lambda_+=\{n\ge0\}$, the shifted subsystem
$\alpha_k(\mB_{\Lambda_+})=\mB_{\{n\ge k\}}$ satisfies the one-sided nesting property
$\alpha_k(\mB_{\Lambda_+})\subset\mB_{\Lambda_+}$ for all integers $k\ge1$, and the
full system is generated by all shifts.  This is the Bernoulli shift structure of
Section~\ref{subsec:classical-ergodic-hierarchy}.
 
\paragraph{Subsystem coarse-grainings of a free relativistic field theory.}
Classical field theory provides a natural setting for a net of subsystem algebras with
infinitely many degrees of freedom.  For an unconstrained real scalar field on a Cauchy
surface $\Sigma$, a phase-space point is a pair of initial data $(\phi,\pi)$.  An open
spatial region $U\subset\Sigma$ selects the unital polynomial algebra
\begin{align}
    \mathcal{B}_U
    :=\operatorname{Alg}_{\mathbb{R}}\{1,\phi(f),\pi(g):f,g\in C_c^\infty(U)\},
    \qquad
    \{\phi(f),\pi(g)\}=\int_\Sigma fg\,d\Sigma.
\end{align}
Thus $\mathcal{B}_U$ is a Poisson subalgebra, and algebras associated with disjoint
spatial regions Poisson commute.  Retaining $\mathcal{B}_U$ resolves the field data in
$U$ while discarding the exterior data.  Although labeled by a spatial region, this is
a subsystem coarse-graining: $U$ selects degrees of freedom, whereas a region of phase
space selects entire field configurations.
 
% -------------------------------------------------------------------
\subsection{Ergodic hierarchy, subspaces and subalgebras}%
\label{subsec:classical-ergodic-hierarchy}
 
The ergodic hierarchy describes how much information about a state remains accessible
at late times.  In the algebraic language, it can be studied through correlations of
observables and through the evolution of coarse observable algebras.  For our purposes,
the relevant distinction is between conserved events, autonomous subsystems, and
subalgebras that are preserved only in one time direction.
 
\paragraph{Ergodicity, conserved subspaces, and conserved subsystems.}
Let $(\mA,\mu,\alpha_t)$ be the measure-preserving dynamical system introduced above.
The flow is \emph{ergodic} if every measurable set invariant under all $\Phi_t$ has
measure zero or one.  Equivalently, its fixed-point algebra is trivial:
\begin{align}
    \mA^\alpha
    :=\{f\in\mA:\alpha_t(f)=f\text{ for all }t\in\mathbb{R}\}
    =\mathbb{C}1.
\end{align}
Ergodicity forbids invariant subspace algebras: if $\alpha_t(\mA_O)=\mA_{\Phi_t(O)}$
for some $O$, then $1_O\in\mA^\alpha$ (a \emph{conserved event}), contradicting
ergodicity.  However, ergodicity does not exclude an invariant subsystem
$\mB\subset\mA$: the condition $\alpha_t(\mB)=\mB$ permits observables in $\mB$ to
evolve into one another.  Such a subsystem is an autonomous factor of the dynamics
(conserved subsystem).  Taking two identical independent copies of a classical
dynamical system always results in autonomous factors; however, a standard result in
ergodic theory implies that if the original system is weakly mixing, the combined
dynamical system is ergodic.%
\footnote{A clean example is the \emph{Chacon rotation system}: the transformation
  $T$ on $[0,1)$ defined by the Chacon substitution is weakly mixing (all eigenvalues
  of the Koopman operator $U_T:f\mapsto f\circ T$ are trivial except $\lambda=1$) but
  not strongly mixing (correlations do not decay to zero pointwise in $t$).  Two
  independent copies of a weakly mixing system are ergodic, because the tensor product
  of two weakly mixing Koopman operators has only the trivial eigenvalue.  This is the
  prototypical example showing the gap between ergodicity and mixing in the classical
  hierarchy.}
 
\paragraph{Information recovered by subalgebra evolution.}
Given a coarse algebra $\mB$, the family $\mB_t=\alpha_t(\mB)$ describes the
observables accessible as this coarse description is transported by the dynamics.  The
algebra
\begin{align}
    \mA_{\mB}:=\bigvee_{t\in\mathbb{R}}\alpha_t(\mB)
\end{align}
contains the information accessible from all time translates of the coarse observables.
The condition $\mA_{\mB}=\mA$ says that these translates generate the full observable
algebra.
 
\paragraph{Half-sided subalgebra translations and the Kolmogorov property.}
Of particular interest are subalgebras satisfying
\begin{align}\label{eq:classical-one-sided-inclusion}
    \alpha_t(\mB_+)\subseteq\mB_+,\qquad t\geq0.
\end{align}
The family $\mB_t=\alpha_t(\mB_+)$ is then decreasing: $\mB_{t+s}\subseteq\mB_t$ for
$s\geq0$.  To interpret this inclusion, one can take $\mB_+$ to be generated by
observations on a half-infinite time interval, with the orientation chosen so that
positive $t$ removes the earliest part of the record.
 
The Kolmogorov, or $K$, property can be formulated through the existence of such a
one-sided algebra whose translates generate the full algebra and whose tail is trivial:
\begin{align}\label{eq:classical-K-algebra}
    \bigvee_{t\in\mathbb{R}}\mB_t=\mA,\qquad
    \bigcap_{t\geq0}\mB_t=\mathbb{C}1.
\end{align}
The first condition says that the complete record resolves all observables.  The second
says that no nonconstant observable remains accessible arbitrarily far into the retained
half of the record.  The analogous conditions with integer time describe a
$K$-automorphism.  They provide an inclusion-based formulation of loss of memory and
motivate the future and past algebras used in quantum ergodic theory~\cite{ouseph2024local}.
 
\paragraph{Example: a Bernoulli shift.}
Consider independent binary variables $x_n$, indexed by $n\in\mathbb{Z}$, with the
uniform product measure on $X=\{0,1\}^{\mathbb{Z}}$.  Let $\alpha(x_n)=x_{n+1}$, and
let $\mB_+$ be the von Neumann algebra generated by $x_n$ with $n\geq0$.  Then
\begin{align}
    \alpha^k(\mB_+)=\operatorname{vN}\{x_n:n\geq k\},\qquad
    \alpha(\mB_+)\subsetneq\mB_+.
\end{align}
All integer translates generate $\mA$, and the tail algebra is trivial by the
independence of the coordinates and the Kolmogorov zero-one law.  At each step the
retained description forgets one more coordinate, although the shift on the full system
is invertible and preserves the measure.  This is a subsystem coarse-graining: the
retained and discarded coordinates give tensor factors of the observable algebra.
 
This example also distinguishes one-sided subsystem inclusions from shrinking
phase-space regions.  In a finite measure space, $\Phi_t(O)\subseteq O$ and measure
preservation imply $\mu(O\setminus\Phi_t(O))=0$.  Thus a region cannot shrink
strictly modulo null sets under a measure-preserving flow.  In the nonatomic setting
of Section~\ref{subsec:classical-region-subsystem}, the same obstruction applies to
proper one-sided inclusions of the algebras $\mA_O$.  It does not apply to the
subsystem algebras of the Bernoulli example.
 
\paragraph{Coarse-graining along a decreasing family.}
Each $\mB_t$ has a $\mu$-preserving conditional expectation $\mathcal{E}_t:\mA\to\mB_t$.
For $s\geq t$, the tower property gives $\mathcal{E}_s\mathcal{E}_t=\mathcal{E}_s$.
For a fixed state with density $\rho$, its coarse-grained densities $\rho_t=\mathcal{E}_t(\rho)$ satisfy
\begin{align}
    H(\rho_s)\geq H(\rho_t),\qquad
    S(\rho\|\rho_s)\geq S(\rho\|\rho_t),\qquad s\geq t,
\end{align}
whenever the entropies are finite.  The loss of information follows from the nesting of
the retained algebras.  It does not require irreversible microscopic dynamics.  For a
trivial tail, the reverse martingale theorem further gives
$\mathcal{E}_t(f)\to\mu(f)1$ in $L^2$ for every $f\in L^2(X,\mu)$.
 
% -------------------------------------------------------------------
\subsection{Classical avatar of modular flow in Hamiltonian systems}%
\label{subsec:classical-state-flow}
 
So far, the state and the dynamics have been independent inputs.  A probability density
assigns weights to microstates, but measure theory alone does not turn these weights
into a motion of the microstates.  A Poisson bracket supplies the missing structure: it
converts a function into a Hamiltonian vector field.  Applied to the logarithm of a
density, it gives a distinguished state-dependent dynamics.  This construction will let
us compare the dynamics associated with a state and with its coarse description.
 
\paragraph{Poisson brackets and Hamiltonian motion.}
A Poisson structure on a smooth manifold $X$ is an antisymmetric bivector $\Pi$ whose
bracket
\begin{align}
    \{f,g\}=\Pi(df,dg),\qquad f,g\in C^\infty(X),
\end{align}
satisfies the Jacobi identity.  The bracket obeys the Leibniz rule in each argument,
so fixing a smooth real Hamiltonian $H$ defines a vector field.  We use the convention
\begin{align}\label{eq:classical-Hamiltonian-convention}
    X_H(f)&=\{f,H\},\qquad
    \alpha_t^H(f)=f\circ\Phi_{-t}^H,\nonumber\\
    \delta_H(f)&:=\left.\partial_t\alpha_t^H(f)\right|_{t=0}=\{H,f\},
\end{align}
where $\Phi_t^H$ is the flow of $X_H$.  Thus the generator on observables has the
opposite sign from the vector field transporting points.  The Jacobi identity ensures
that Hamiltonian evolution preserves the Poisson bracket.  In canonical coordinates,
$\{q_i,p_j\}=\delta_{ij}$ and the point trajectories satisfy the usual Hamilton
equations.
 
\paragraph{Casimirs and symplectic leaves.}
A Casimir is a function $c$ with $\{c,f\}=0$ for every smooth $f$.  It is conserved
by every Hamiltonian.  The points accessible from a given microstate by successive
Hamiltonian motions form a symplectic leaf.  A leaf is a classical analogue of an
irreducible sector for Hamiltonian motion.  This analogy concerns the Poisson bracket:
the pointwise algebra $L^\infty(X,\mu)$ remains abelian, and restricting it to a leaf
does not turn it into a nonabelian factor.
 
\paragraph{Liouville measure and the state Hamiltonian.}
For the explicit constructions below, take $X$ to be a connected compact symplectic
manifold without boundary, with symplectic form $\omega$ and normalized Liouville
measure
\begin{align}
    d\mu=\frac{\omega^n/n!}{\int_X\omega^n/n!},\qquad \dim X=2n.
\end{align}
Hamiltonian flows preserve this measure.  Equivalently,
\begin{align}\label{statepreserve}
    \int_X\{f,g\}\,d\mu=0
\end{align}
for smooth $f,g$.  Let $\rho$ be a smooth strictly positive probability density
relative to $\mu$.  Define its \emph{state Hamiltonian} and the associated evolution by
\begin{align}\label{eq:classical-state-Hamiltonian}
    K_\rho=-\log\rho,\qquad
    \Phi_t^\rho:=\Phi_t^{K_\rho},\qquad
    \alpha_t^\rho:=\alpha_t^{K_\rho}.
\end{align}
The infinitesimal evolution of an observable is
\begin{align}\label{eq:classical-state-derivation}
    \delta_\rho(f)=\{K_\rho,f\}=\{f,\log\rho\}.
\end{align}
The flow moves along level sets of the density, since $X_{K_\rho}(\rho)=0$.  It
preserves both $\mu$ and the state $\omega_\rho(f)=\int_X\rho f\,d\mu$.  For a Gibbs
density $\rho=\mathcal{Z}^{-1}e^{-\beta H}$, we have $K_\rho=\beta H+\log\mathcal{Z}$
(where $\mathcal{Z}$ is the partition function) and $\Phi_t^\rho=\Phi_{\beta t}^H$.
Thus the state recovers the Hamiltonian motion, with a rescaled parameter; for a
general density there is no prior identification of this parameter with physical time.
 
\paragraph{The classical KMS identity.}
The logarithm is distinguished by an integration-by-parts identity.  Using
\eqref{statepreserve},
\begin{align}
    0=\int_X\{g,f\rho\}\,d\mu
    =-\omega_\rho(\{f,g\})+
    \omega_\rho\bigl(f\{g,\log\rho\}\bigr).
\end{align}
Consequently,
\begin{align}\label{KMSclassical}
    \omega_\rho(\{f,g\})=\omega_\rho(f\delta_\rho(g)).
\end{align}
For a Gibbs state this reads $\omega_\beta(\{f,g\})=\beta\omega_\beta(f\delta_H(g))$.
This is the classical, or static, KMS condition, written with the time orientation in
\eqref{eq:classical-Hamiltonian-convention}~\cite{vanDeVen2022}.  It is stronger than
stationarity: testing against all smooth $f$ determines $\delta_\rho(g)$ for every
$g$.  Conversely, on a connected symplectic manifold, imposing this identity for a
specified Hamiltonian dynamics determines a smooth positive density to be Gibbs, up to
normalization.
 
\paragraph{Poisson modular flow.}
For a smooth positive measure $\nu$, define its Poisson modular vector field by
$\mathcal{M}_\nu(f):=\operatorname{div}_\nu X_f$.  This measures the failure of
Hamiltonian motions to preserve $\nu$~\cite{Weinstein1997}.  For $d\nu=\rho\,d\mu$
and Liouville $\mu$,
\begin{align}
    \mathcal{M}_\nu(f)=X_f(\log\rho)=\{f,-\log\rho\}=X_{K_\rho}(f).
\end{align}
Thus $\Phi_t^\rho$ is precisely this Poisson modular flow.  The construction is
distinct from the multiplication group $U_t^\rho$ introduced above and from the trivial
modular automorphism group of the abelian von Neumann algebra.  It also supplies no
automatic ergodic or entropy-production property.  A nonconstant $\rho$ is itself a
conserved observable, so the flow cannot be ergodic on the full probability space.
Questions about ergodicity and mixing must be posed on invariant sectors with suitable
invariant measures.
 
% -------------------------------------------------------------------
\subsection{Subsystem Poisson flows in classical many-body systems and field theory}%
\label{subsec:classical-subalgebra-Hamiltonian}
 
A coarse description determines a reduced state and hence, when the coarse space is
Poisson, its own state-generated dynamics.  We can then ask whether the reduced
dynamics agrees with what the full dynamics does to the observed variables.  Comparing
the two motions gives a classical model of the subalgebra flows studied in this paper.
The comparison separates evolution already visible within the retained algebra from
evolution involving the discarded variables.
 
\paragraph{Poisson subsystems.}
We say a subalgebra $\mB_Y\subset\mA$ is a \emph{Poisson subsystem} if for all
$b_1,b_2\in\mB_Y$ we have $\{b_1,b_2\}\in\mB_Y$.  Geometrically, this means that
$\pi:X\to Y$ is a smooth surjective Poisson submersion:
\begin{align}\label{eq:classical-Poisson-subsystem}
    \{\pi^*f,\pi^*g\}_X=\pi^*\{f,g\}_Y,\qquad f,g\in C^\infty(Y).
\end{align}
The smooth coarse observables $\mathcal{B}=\pi^*C^\infty(Y)$ are closed under the
bracket.  This condition allows us to speak of Hamiltonian dynamics within the coarse
description.
 
The simplest setting makes both the geometric and probabilistic structures explicit.
Take a product $X=Y\times Z$ of connected compact symplectic manifolds, with product
symplectic structure and normalized Liouville measure $\mu=\mu_Y\otimes\mu_Z$.  The
projection to $Y$ is Poisson, and the associated von Neumann subalgebra is
$\mB_Y=L^\infty(Y,\mu_Y)\otimes1$.  For a smooth positive density $\rho$, restriction
of the state gives the marginal density
\begin{align}\label{eq:classical-marginal-Hamiltonian}
    \rho_Y(y)=\int_Z\rho(y,z)\,d\mu_Z(z),\qquad K_Y=-\log\rho_Y.
\end{align}
The expectation $\mE_Y$ from Section~\ref{subsec:classical-region-subsystem} extends
this marginal back to $X$ as $\rho_0(y,z)=\rho_Y(y)$.  Its Hamiltonian is
$K_0=\pi^*K_Y$, whose flow is $\Phi_t^0(y,z)=(\Phi_t^Y(y),z)$.
 
Singular inclusions (those where the relative commutant of $\mB_Y$ in $\mA$ is
trivial, so that $\mB_Y$ and its Poisson centralizer generate all of $\mA$) can in
principle arise on noncompact phase spaces or for degenerate Poisson structures, but
they do not occur in the product setting we study here, where the symplectic form is
non-degenerate on each factor.
 
\paragraph{Agreement between reduced dynamics and full dynamics means no correlations.}
Agreement means that the observed trajectory is determined by its initial coarse state
and follows the flow generated by $K_Y$:
\begin{align}\label{eq:classical-state-intertwining}
    \pi\circ\Phi_t^\rho=\Phi_t^Y\circ\pi.
\end{align}
For our complete smooth flows, this is equivalent to the infinitesimal condition
\begin{align}\label{eq:classical-state-defect}
    \{\pi^*g,K_\rho\}_X-\pi^*\{g,K_Y\}_Y
    =\{\pi^*g,K_\rho-K_0\}_X=0,\qquad g\in C^\infty(Y).
\end{align}
The left side measures the part of the observed velocity missed by the reduced state
Hamiltonian.  A Hamiltonian invisible to all coarse observables belongs to the Poisson
centralizer
\begin{align}
    \mathcal{B}^{\mathrm{c}}_{\mathrm{P}}
    :=\{h\in C^\infty(X):\{h,b\}=0\text{ for every }b\in\mathcal{B}\}.
\end{align}
Thus agreement is precisely the condition $K_\rho-K_0\in\mathcal{B}^{\mathrm{c}}_{\mathrm{P}}$.
This is an analogue of a relative commutant defined using Poisson brackets.  In the
product setting, nondegeneracy and connectedness of $Y$ imply this centralizer consists
of functions of $z$ alone.  Hence \eqref{eq:classical-state-intertwining} is equivalent to
\begin{align}\label{eq:classical-product-state-closure}
    K_\rho(y,z)=K_Y(y)+K_Z(z),\qquad
    \rho(y,z)=\rho_Y(y)\rho_Z(z).
\end{align}
In other words, agreement is equivalent to the state being a product state.  This is
the classical analog of Takesaki's condition in Section~\ref{subsec:quantum-coarse-2}.
 
\paragraph{Comparing the two flows.}
Define the relative transformation by first undoing the lifted coarse motion and then
applying the full motion:
\begin{align}\label{eq:classical-relative-transformation}
    R_t=\Phi_t^\rho\circ\Phi_{-t}^0,\qquad
    \mathcal{T}_t(f):=f\circ R_t^{-1}=\alpha_t^\rho\circ\alpha_{-t}^0(f).
\end{align}
For each $t$, this is a symplectic transformation and a measure-preserving
automorphism.  Since the lifted coarse flow preserves $\mB_Y$,
\begin{align}
    \mathcal{T}_t(\mB_Y)=\alpha_t^\rho(\mB_Y).
\end{align}
The relative transformation follows the same moving subalgebra as the full dynamics,
while canceling the chosen internal coarse evolution.  For a product state, it can
still move hidden variables: $R_t(y,z)=(y,\Phi_t^Z(z))$.
 
The infinitesimal generator of $R_t$ is the difference of the state Hamiltonians:
\begin{align}\label{eq:classical-relative-Hamiltonian}
    K_{\mathrm{rel}}&=K_\rho-K_0=-\log\frac{\rho}{\rho_0},\nonumber\\
    \left.\partial_tR_t\right|_{t=0}&=X_{K_{\mathrm{rel}}},\qquad
    \left.\partial_t\mathcal{T}_t(f)\right|_{t=0}=\{K_{\mathrm{rel}},f\}.
\end{align}
Every choice of state $\rho_Z(z)$ on the relative commutant $Z$ results in a
conditional expectation $\pi_{\rho_Z}:X\to Y$, and there is a one-to-one
correspondence between states of the relative commutant and conditional expectations
$\mE_\pi:X\to Y$.
 
\paragraph{Example: a uniform marginal with nonautonomous observables.}
Take $Y$ and $Z$ to be two-dimensional tori with angular coordinates $(q,p)$ and
$(Q,P)$, each of period $2\pi$, and the product symplectic form
$dq\wedge dp+dQ\wedge dP$.  Consider
\begin{align}
    \rho(q,p,Q,P)=\frac{e^{\epsilon\cos(q-Q)}}{\mathcal{Z}_\epsilon},\qquad
    \mathcal{Z}_\epsilon=\int_X e^{\epsilon\cos(q-Q)}\,d\mu,\qquad \epsilon\ne0.
\end{align}
Translation invariance gives $\rho_Y=1$, so $K_0=0$ and the reduced state predicts no
motion.  The full state Hamiltonian is $K_\rho=-\epsilon\cos(q-Q)+\log\mathcal{Z}_\epsilon$,
giving $\dot q=\dot Q=0$, $\dot p=-\epsilon\sin(q-Q)$, $\dot P=\epsilon\sin(q-Q)$.
The bounded coarse observable $e^{ip}$ evolves as
$\alpha_t^\rho(e^{ip})=\exp(i[p+\epsilon t\sin(q-Q)])$.  Its future value depends on
the hidden angle $Q$, so it leaves $\mB_Y$ for $t\ne0$.  Here $R_t=\Phi_t^\rho$
directly displays the motion that the marginal state Hamiltonian misses.
 
\paragraph{Relative entropy and the information discarded.}
\begin{align}\label{eq:classical-relative-entropy-Hamiltonian}
    S(\rho\|\rho_0)=-\omega_\rho(K_{\mathrm{rel}})=H(\rho_0)-H(\rho).
\end{align}
The vector field alone does not determine this number: on a connected symplectic
manifold it determines its Hamiltonian only up to a constant.  The normalized densities
fix that constant.  On a general Poisson manifold, Casimirs give a larger ambiguity.
 
Information loss and failure of autonomous coarse dynamics are also distinct.  For a
product state, the relative transformation acts only on $Z$, yet
\begin{align}
    S(\rho_Y\rho_Z\|\rho_Y\otimes1)
    =\int_Z\rho_Z\log\rho_Z\,d\mu_Z = -H(\rho_Z)
\end{align}
can be positive, where $-H(\rho_Z)=S(\rho_Z\|\mu_Z)$ is the relative entropy of
$\rho_Z$ with respect to the reference measure $\mu_Z$ (i.e., the negative of the
differential entropy of $\rho_Z$).  Coarse-graining has discarded the hidden
distribution even without correlations.  More generally,
\begin{align}
    S(\rho\|\rho_Y\otimes1)
    =S(\rho\|\rho_Y\rho_Z)+S(\rho_Z\|\mu_Z),
\end{align}
where $\rho_Z$ is the other marginal.  The first term is the mutual information and the
second, $S(\rho_Z\|\mu_Z)=-H(\rho_Z)$, measures the hidden marginal relative to the
reference measure $\mu_Z$.  Neither this entropy nor the relative transformation
supplies a monotone entropy-production law without an additional coarse-graining or
dynamical assumption.
 
\paragraph{Subregions and the passage to field theory.}
The measurable region algebras $\mA_O$ remain valid coarse descriptions.  What can fail
is the smoothness needed to assign their coarse densities a Hamiltonian flow.  For the
region expectation $E_O$,
\begin{align}
    \rho_0=E_O(\rho)=\chi_O\rho+c_O\chi_{O^c},\qquad
    c_O=\frac{1}{\mu(O^c)}\int_{O^c}\rho\,d\mu.
\end{align}
For an open region with smooth boundary, this density is generally discontinuous there.
In the interior of $O$, $K_0=K_\rho$; in the interior of the complement, $K_0$ is
constant.  These local statements do not specify a smooth global Hamiltonian flow
across the boundary.  Additional regularization or boundary conditions are needed.
This is a limitation of the Hamiltonian construction, not of the measure-theoretic net.
 
Finally, every classical coarse algebra still admits a state-preserving conditional
expectation for a faithful probability state, even when its Poisson state flow fails to
preserve that algebra.  In the product example this expectation is
\begin{align}
    E_{\rho,Y}(f)(y)=\int_Z f(y,z)\frac{\rho(y,z)}{\rho_Y(y)}\,d\mu_Z(z).
\end{align}
Its existence does not imply \eqref{eq:classical-state-intertwining}.  In quantum
modular theory the relation changes: for a faithful normal state $\varphi$ on $\mA$,
Takesaki's theorem identifies invariance of a unital von Neumann subalgebra $\mB$ under
$\sigma_t^\varphi$ with the existence of a $\varphi$-preserving normal conditional
expectation onto $\mB$~\cite{Takesaki1972}.  Failure of invariance rules out an
expectation preserving this state; it does not by itself rule out expectations
preserving other states.  This is the additional structure that the quantum theory
contributes to the comparison of full and coarse dynamics.
 
The classical construction thus motivates studying how relative state evolutions
transport retained observables.  In QFT a spacetime region selects local field
observables, whereas a region $O\subset X$ above selects possible microstates.  The
connection developed in Section~\ref{section:basics} is through the inclusions of
observable algebras and their state-dependent flows; it does not identify these two
notions of region.
 
Choose an ambient Poisson algebra $\mathcal{A}$ containing these observables and stable
under a complete Hamiltonian evolution.  With our convention for the pullback, the
transported subsystem is
\begin{align}
    \alpha_t^H(F)=F\circ\Phi_{-t}^H,\qquad
    \partial_t\alpha_t^H(F)=\{H,\alpha_t^H(F)\},\qquad
    \mathcal{B}_{U,t}=\alpha_t^H(\mathcal{B}_U).
\end{align}
Hamiltonian evolution preserves Poisson brackets, so every $\mathcal{B}_{U,t}$ remains
a Poisson subalgebra.  Invariance for all $t$ gives an autonomous reversible subsystem.
More generally, propagation can make the transported observables depend on initially
discarded data.
 
A simple example of strict nesting is supplied by a classical chiral current $j(x)$
on the line.  On rapidly decreasing smooth configurations, choose
\begin{align}
    \{j(x),j(y)\}=-\partial_x\delta(x-y),\qquad
    P=\frac{1}{2}\int_{\mathbb{R}}j(x)^2\,dx,\qquad
    \tau_a j(x)=j(x+a).
\end{align}
The nonnegative Hamiltonian $P$ generates $\tau_a$ through $\{P,j(x)\}=\partial_xj(x)$.
Let $\mathcal{B}_b$ be generated by $1$ and $j(f)$ with $f\in C_c^\infty((b,\infty))$.
Since $\tau_a j(f)=j(f(\,\cdot-a))$,
\begin{align}
    \tau_a(\mathcal{B}_b)=\mathcal{B}_{b+a}\subsetneq\mathcal{B}_b\quad(a>0),\qquad
    \bigcup_{b\in\mathbb{R}}\mathcal{B}_b=\mathcal{A}_{\mathrm{loc}},\qquad
    \bigcap_{b\geq0}\mathcal{B}_b=\mathbb{R}1.
\end{align}
Here $\mathcal{A}_{\mathrm{loc}}$ is the polynomial algebra of compactly smeared
currents; $P$ is a global functional generating its automorphisms.  The full
translation group is reversible, but its positive half acts on $\mathcal{B}_0$ by
injective, nonsurjective Poisson endomorphisms.  The decreasing family progressively
discards accessible field degrees of freedom.
 
The dilation Hamiltonian $D=\frac{1}{2}\int_{\mathbb{R}}xj(x)^2\,dx$ gives a further
connection with half-sided modular inclusions.  Its flow $\delta_s$ satisfies
\begin{align}
    \{D,P\}=P,\qquad
    \delta_s(\mathcal{B}_b)=\mathcal{B}_{e^s b},\qquad
    \delta_s\tau_a\delta_{-s}=\tau_{e^s a}.
\end{align}
Thus dilations preserve $\mathcal{B}_0$ and act half-sidedly on $\mathcal{B}_b$ for
$b>0$.  These are classical Hamiltonian realizations of the translation--dilation
relations.  Their construction uses the Poisson structure; identifying the dilation
with a state-generated modular evolution requires additional input.

\section{Review of quantum dynamical systems}
\label{subsec:review-quantum}
 
In this subsection, we review and fix our notation.  Readers familiar with standard
definitions can skip this subsection.
 
\paragraph{Algebra of quantum observables.}
Quantum observables are described by a \emph{von Neumann algebra} $\mathcal{A}$, a
unital $*$-subalgebra of the bounded operators $B(\mathcal{H})$ on some separable
Hilbert space $\mathcal{H}$, closed in the weak operator topology.  A von Neumann
algebra $\mathcal{A}$ is a \emph{factor} if its center
$\mathcal{Z}(\mathcal{A}) = \mathcal{A}\cap\mathcal{A}'$ consists only of scalar
multiples of the identity, where
$\mathcal{A}' = \{b\in B(\mathcal{H}) : [b,a]=0\text{ for all }a\in\mathcal{A}\}$ is
the commutant.  Physically, a factor algebra describes a system with no classical
superselection sectors.%
\footnote{As we review in Appendix~\ref{app:classical}, in classical physics the
  commutator is replaced by the Poisson bracket and the commutant $\mathcal{A}'$ is
  the quantum analog of the Poisson centralizer.  It consists of all observables
  invisible to $\mathcal{A}$.  The classical analog of a factor is a non-degenerate
  Poisson bracket that restricts to a symplectic leaf.}
 
\paragraph{Quantum states.}
A \emph{normal state} on $\mathcal{A}$ is a positive, normalized,
$\sigma$-weakly continuous linear functional
$\varphi:\mathcal{A}\to\mathbb{C}$, with $\varphi(\id)=1$.  Normal states are in
one-to-one correspondence with density matrices $\rho_\varphi\in\mathcal{A}_*^+$ via
$\varphi(a)=\Tr(\rho_\varphi\, a)$ when $\mathcal{A}$ is represented on a Hilbert
space with a trace.  In the finite-dimensional setting, these are precisely the
positive semidefinite matrices of unit trace.
 
A state $\varphi$ is \emph{faithful} if $\varphi(a^*a)=0$ implies $a=0$;
equivalently, its density matrix has no zero eigenvalues.  Faithfulness is the quantum
counterpart of the condition that the probability density $\rho_\nu$ is strictly
positive almost everywhere, and it is required for the modular theory to be
non-trivial.
 
A state $\varphi$ on $\mathcal{A}$ is \emph{tracial} if $\varphi(ab)=\varphi(ba)$ for
all $a,b\in\mathcal{A}$.  Tracial states play the role of the reference Liouville
measure in classical dynamics: every normal state is absolutely continuous with
respect to a faithful tracial state, in direct analogy with $\nu\ll\mu$ in the
classical setting.  Not every von Neumann algebra admits a faithful tracial state;
those that do are called \emph{finite} von Neumann algebras.  In the
infinite-dimensional algebras relevant to quantum field theory, one works instead with
faithful normal states.
 
\paragraph{GNS Hilbert space and modular operators.}
Given a von Neumann algebra $\mathcal{A}$ and a faithful normal state $\varphi$, the
Gelfand--Naimark--Segal (GNS) construction produces a Hilbert space
$\mathcal{H}_\varphi$, a cyclic unit vector $|\Omega\rangle\in\mathcal{H}_\varphi$
with $\varphi(a)=\langle\Omega|a|\Omega\rangle$, and a $*$-representation
$\pi:\mathcal{A}\to B(\mathcal{H}_\varphi)$.%
\footnote{The classical analog is $\mathcal{H}=L^2(X,\mu)$ with the constant function
  $|\id\rangle$ as the reference vector.}
The Tomita--Takesaki modular theory assigns to the pair $(\mathcal{A},|\Omega\rangle)$
an antilinear operator $S_\varphi$ defined on the dense domain
$\mathcal{A}|\Omega\rangle$ by
\begin{equation}
  S_\varphi\, a|\Omega\rangle = a^*|\Omega\rangle, \qquad a\in\mathcal{A}.
  \label{eq:Tomita-S}
\end{equation}
Its polar decomposition $S_\varphi = J_\varphi\,\Delta_\varphi^{1/2}$ defines the
\emph{modular conjugation} $J_\varphi$, an antiunitary involution, and the
\emph{modular operator} $\Delta_\varphi$, a positive self-adjoint operator.
 
\paragraph{Modular dynamics.}
The Tomita--Takesaki theorem~\cite{Takesaki1970} guarantees that the one-parameter
group of unitaries $\Delta_\varphi^{it}$, $t\in\mathbb{R}$, implements an
automorphism group of $\mathcal{A}$:
\begin{equation}
  \sigma_t^\varphi(a) = \Delta_\varphi^{it}\, a\, \Delta_\varphi^{-it},
  \qquad a\in\mathcal{A},\quad t\in\mathbb{R}.
  \label{eq:modular-flow}
\end{equation}
This is the \emph{modular flow} (or \emph{modular automorphism group}) of $\varphi$.
 
The modular operator $\Delta_\varphi$ and $J_\varphi$ have no direct counterparts in
classical dynamics.  They are genuinely quantum objects encoding the non-commutativity
of $\mathcal{A}$.  \textcolor{blue}{However, in Appendix~\ref{app:classical} (see especially
Sections~\ref{subsec:classical-state-flow} and~\ref{subsec:classical-subalgebra-Hamiltonian})
we argue that modular flow can be interpreted as a quantum generalization of
\emph{Poisson flow}: the GNS Hilbert space associates observables with states, and the
modular flow associates a flow with a state generated by the modular Hamiltonian
$K_{\omega;\mA}=-\log\Delta_{\omega;\mA}$.  In classical physics the Radon--Nikodym
derivative plays the role of the density matrix, and the Poisson bracket with
$-\log\rho$ generates a flow $\alpha_t^\rho$ associated to every state.}
The state $\varphi$ is invariant under its own modular flow:
$\varphi\circ\sigma_t^\varphi = \varphi$ for all $t\in\mathbb{R}$, in direct analogy
with the classical fact that $\omega_\rho$ is preserved by $\alpha_t^\rho$.
 
The modular flow satisfies the \emph{KMS condition} at inverse temperature $\beta=-1$:
for all $a,b\in\mathcal{A}$ the function
\begin{equation}
  F_{a,b}(t) = \varphi\bigl(a\,\sigma_t^\varphi(b)\bigr)
  \label{eq:KMS-function}
\end{equation}
extends to a function analytic in the strip $0<\operatorname{Im}(t)<1$ and continuous
on its closure, satisfying the boundary condition
$F_{a,b}(t+i)=\varphi(\sigma_t^\varphi(b)\,a)$.  This is the precise quantum analog
of the classical KMS identity~\eqref{KMSclassical},
which reduced to an integration-by-parts identity for the Poisson bracket.
Differentiating \eqref{eq:KMS-function} at $t=0$ recovers
\begin{equation}
  \varphi([a,b]) = \varphi(a\,\delta_\varphi(b)),
  \label{eq:KMS-infinitesimal}
\end{equation}
where $\delta_\varphi(b)=[\log\Delta_\varphi,b]$ is the derivation generating
$\sigma^\varphi$.  In the abelian case $[a,b]=0$ and this becomes
$0=\varphi(a\{b,\log\rho\})$, matching~\eqref{KMSclassical}.
 
 \section{Quantum coarse-grainings}\label{subsec:quantum-coarse-2}
 
\begin{definition}
    A \emph{quantum coarse-graining} is a von Neumann subalgebra
$\mathcal{B}\subseteq\mathcal{A}$ together with a \emph{normal conditional expectation}
$\mE_{\mathcal{B}}:\mathcal{A}\to\mathcal{B}$, a completely positive, unital,
$\mathcal{B}$-bimodular map satisfying
\begin{equation}
  \mE_{\mathcal{B}}(b_1 a b_2) = b_1\,\mE_{\mathcal{B}}(a)\,b_2,
  \qquad a\in\mathcal{A},\;b_1,b_2\in\mathcal{B}.
  \label{eq:CE-bimodular}
\end{equation}
\end{definition}

This is the quantum analog of the classical conditional expectation
$\mE_{\mathcal{B}}(f) = \mathbb{E}_\mu[f|\Sigma_0]$ reviewed in
Appendix~\ref{app:classical}.  On the GNS Hilbert space, $\mE_{\mathcal{B}}$ extends
to an orthogonal projection
$e_{\mathcal{B}}:\mathcal{H}_\varphi\to\overline{\mathcal{B}|\Omega\rangle}$, the
\emph{Jones projection} associated to the inclusion $\mathcal{B}\subset\mathcal{A}$.
It is the operator representative of the superoperator $\mE_\mB$ in the GNS Hilbert
space.
 
\paragraph{Entropy of a coarse-graining.}
Given a faithful normal state $\varphi$ on $\mathcal{A}$, a von Neumann subalgebra
$\mathcal{B}\subseteq\mathcal{A}$, and a $\varphi$-preserving conditional expectation
$\mE:\mathcal{A}\to\mathcal{B}$, the coarse-grained state is
$\varphi_0 = \varphi\circ \mE$, and the \emph{quantum relative entropy} of $\varphi$
with respect to $\varphi_0$ is
\begin{equation}
  S(\varphi\|\varphi_0)
  = \langle\Omega|\log\Delta_{\varphi|\varphi_0;\mathcal{A}} |\Omega\rangle
  =i\frac{d}{dt}\langle\Omega|\Delta_{\varphi|\varphi_0}^{it}\Omega\rangle\Big|_{t=0},
  \label{eq:relative-entropy}
\end{equation}
where $\Delta_{\varphi|\varphi_0;\mathcal{A}}$ is the relative modular operator of
$\varphi$ and $\varphi_0$ on $\mathcal{A}$:
\begin{align}
  &S_{\varphi|\varphi_0;A}a|\Omega\rangle=a^*|\Omega_0\rangle,\qquad
   \Delta_{\varphi|\varphi_0;A}=S^\dagger_{\varphi|\varphi_0;A}S_{\varphi|\varphi_0;A}
   \nonumber\\
  &\langle a\Omega|\Delta_{\varphi|\varphi_0}|b\Omega\rangle = \varphi_0(b^*a),
  \qquad a,b\in\mathcal{A},
  \label{eq:relative-modular}
\end{align}
which generalizes the classical Radon--Nikodym derivative $\rho/\rho_0$.  The quantum
relative entropy is the quantum analog of the KL-divergence:
\begin{equation}
  S(\rho\|\rho_0) = \int_X \rho\log\!\left(\frac{\rho}{\rho_0}\right)d\mu
  = H(\rho_0)-H(\rho).
  \label{eq:classical-KL}
\end{equation}
Non-negativity $S(\varphi\|\varphi_0)\geq 0$, with equality iff $\varphi=\varphi_0$ on
$\mathcal{A}$ (i.e.\ iff $\varphi$ is already supported on $\mathcal{B}$), is the
quantum analog of the monotonicity of entropy under coarse-graining: discarding
information can only increase entropy.
 
\paragraph{Different types of coarse-grainings.}
Just as in the classical theory, there are two natural classes of subalgebras that we
call subspaces and subsystems.  To build intuition, we first motivate them using finite
quantum systems (matrix algebras):
\begin{enumerate}
    \item \textbf{Subspace inclusions:} The Hilbert space splits:
    $\mH=P\mH\oplus(1-P)\mH$.  The coarse subalgebra is
    $\mA_P=\pi(\mA)=P\mA P$, with inclusion
    $\pi_*(\mA_P)=\mA_P\oplus \mathbb{C}(1-P)\subset \mA$.%
    \footnote{The inclusion is the pull-back of the projection map.}
    The inclusion has a non-trivial center $\mathbb{C}P\oplus \mathbb{C}(1-P)=Z_P\subset \pi_*(\mA_P)$.
    There is a one-to-one correspondence between conditional expectations
    $\mE_{\omega_{1-P}}:\mA\to \mA_P$ and states $\omega_P$ on $(1-P)\mH$:
    \begin{align}
    \forall a\in \mA:\qquad \mE_P(a)=Pa P\oplus \omega_{1-P}((1-P)a(1-P)).
    \end{align}
    A state $\omega$ on $\mA$ is preserved by this conditional expectation if and only
    if it has no off-diagonal elements between $P\mH$ and $(1-P)\mH$—density matrices
    that are decohered in the binary measurement basis $\{P,1-P\}$.
 
    Subspace inclusions are particularly relevant for single-particle systems or
    Hilbert spaces that do not decompose as tensor products.  For an increasing family
    $P_i\leq P_{i+1}$, we have $\mathcal{A}_{P_i}\subseteq\mathcal{A}_{P_{i+1}}$,
    paralleling the classical net $\mathcal{A}_{O_i}\subseteq\mathcal{A}_{O_{i+1}}$
    (see Appendix~\ref{app:classical}).  Since the focus of our work is many-body
    quantum systems and QFT, in this work we do not consider such inclusions further.
 
    \item \textbf{Subsystem inclusions:} The Hilbert space decomposes as a tensor
    product $\mH_A=\mH_B\otimes \mH_R$.  The coarse operators are
    $\mA_B:\mH_B\to \mH_B$, and the operators $\mA_R:\mH_R\to \mH_R$ form the
    \emph{relative commutant} of the inclusion $\mA_B\otimes 1_R\subset \mA_A$.  The
    inclusion has a trivial center; we say we have a \emph{subfactor}.  This is a
    quantum generalization of the classical setup $X=Y\times Z$
    (Appendix~\ref{app:classical}).
 
    There is a one-to-one correspondence between conditional expectations
    $\mE_{\omega_R}:\mA_A\to \mA_B$ and states of the relative commutant $\omega_R$:
    \begin{align}\label{condexp}
       \forall b\in \mA_B, r\in \mA_R:\qquad
       \mE_{\omega_R}(b\otimes r)=b\: \omega_R(r).
    \end{align}
    A state is invariant under this coarse-graining if and only if it has no
    correlations between $B$ and $R$, i.e., $\omega_A=\omega_B\otimes \omega_R$.
    This is equivalent to $I_\omega(A:R)=0$, \textcolor{blue}{or in finite dimensions, where
    $\omega_A$ is represented by a density matrix $\rho_A$,}
    \begin{align}
    \forall t\in\mathbb{R},\; b\in\mA_B:\qquad
    \textcolor{blue}{\rho_A^{it}} (b\otimes 1_R) \textcolor{blue}{\rho_A^{-it}}
    = \textcolor{blue}{\rho_B^{it}}b\textcolor{blue}{\rho_B^{-it}}\otimes 1_R.
    \label{eq:product-state-modular}
    \end{align}
    \textcolor{blue}{Here $\rho_A,\rho_B$ are density matrices; in the general GNS setting one replaces
    $\rho^{it}$ by $\Delta_{\omega;\mA}^{it}$ and $\Delta_{\omega;\mB}^{it}$.}
 
    This type of inclusion is natural in the multi-particle quantum systems we study.
    In a many-body quantum system, we can have an increasing family of subsystems,
    e.g., half-lines in an infinite 1d spin chain or wedges in QFT,
    $W_i\leq W_{i+1}$; we have $\mathcal{A}_{W_i}\subseteq\mathcal{A}_{W_{i+1}}$,
    giving a net of coarse observable subsystems.  See
    Appendix~\ref{subsec:classical-subalgebra-Hamiltonian} for a classical analog
    of subsystem nets in classical field theory.  This work focuses on such inclusions.
 
    \item \textbf{Mixed inclusions:} The most general case is
    $\mH_A=\oplus_a \mH_B^{(a)}\otimes \mH_R^{(a)}$:
    \begin{align}
    \pi_*(\mA_B)=\oplus_a \mA_B^{(a)}\otimes 1_R^{(a)}.
    \end{align}
    This has a non-trivial center $\oplus_a \mathbb{C}1_B^{(a)}\otimes \mathbb{C}1_R^{(a)}$.  The
    states invariant under the associated conditional expectation are quantum Markov
    states satisfying $I(B:R|Z)=0$:
    \begin{align}
        \omega=\oplus_a p_a \omega_B^{(a)}\otimes \omega_R^{(a)},\qquad
        S(\omega_A\|\omega_A\circ \mE)=0.
    \end{align}
\end{enumerate}
 
Since we are interested in ergodic modular dynamics, we focus on case~2 above:
subsystems or subfactors.
 

\section{Characteristic Functions of Relative Modular Operators}\label{app:relcharacter}
We take a slight digression in this subsection to briefly discuss a slight generalization of characteristic function. Instead of modular operator, we can study the \textit{characteristic function of relative modular operators}:
\begin{definition}
    Given a pair of von Neumann algebras $\mathcal{B}\subset\mathcal{A}$ and two common cyclic and separating vector $\ket{\Omega_\varphi}, \ket{\Omega_\psi}\in\mathcal{H}$, the \textit{characteristic function of relative modular operator} is defined as:
    \begin{equation}
        \mathcal{D}_{\mathcal{A}, \mathcal{B}}(t) := \Delta_\mathcal{\varphi|\psi;\mathcal{A}}^{it}\Delta_{\varphi|\psi;\mathcal{B}}^{-it}
    \end{equation}
    where $\varphi$ (and $\psi$ respectively) is the state corresponding to the vector $\ket{\Omega_\varphi}$ (and $\ket{\Omega_\psi}$ respectively) and $\Delta_{\varphi|\psi;\mathcal{A}}$ (and $\Delta_{\varphi|\psi;\mathcal{B}}$) is the relative modular operator of the two states $\varphi,\psi$ restricted to the algebra $\mathcal{A}$ (and $\mathcal{B}$ respectively).
\end{definition}
The main result of this appendix is the following continuity theorem:
\begin{theorem}
    Let $B\subset A$ be a standard inclusion. Consider two sequences of normal faithful states $\{\varphi_n\}$ and $\{\psi_n\}$ such that the corresponding representing vectors $\{\xi_{\varphi_n}\}$ and $\{\xi_{\psi_n}\}$ are common cyclic and separating vectors of $A$ and $B$. Assume both sequences converge in norm $\lim_n\varphi_n = \varphi$ and $\lim_n\psi_n = \psi$ such that the limiting states are normal faithful state and their representing vectors are also common cyclic and separating vectors for $A$ and $B$. In this case, for all complex numbers $z$ in the strip $\Re{z}\in[0, 1/2]$, we have strong convergence:
    \begin{equation}
        \lim_n \Delta_{\varphi_n|\psi_n;A}^z \Delta_{\varphi_n|\psi_n;B}^{-z} = \Delta^z_{\varphi|\psi;A}\Delta^{-z}_{\varphi|\psi; B}
    \end{equation}
\end{theorem}
Before proving the theorem, we comment on the conditions on the theorem. When $B\subset A$ is a standard inclusion, by definition the set of common cyclic and separating vectors is not empty. In this case, it turns out that this set is actually a dense $G_\delta$ set in the representing Hilbert space. Thus, sequences like $\{\varphi_n\}$ and $\{\psi_n\}$ can be easily found within this dense $G_\delta$ set. However, it is necessary to check that the norm limit $\varphi$ and $\psi$ can still be represented by common cyclic and separating vectors. A priori, it is possible to have limiting states that cannot be represented by common cyclic and separating vectors.

The proof is based on a few simple observations.
\begin{lemma}
    Under the same conditions as the theorem above, $\{\Delta_{\varphi_n|\psi_n;A}\}_n$ converges in strong resolvent sense to $\Delta_{\varphi|\psi;A}$. Same for $\{\Delta_{\varphi_n|\psi_n;B}\}$
\end{lemma}
\begin{proof}
    Let $\mathcal{H}$ be the standard representing Hilbert space. Then by definition, we have:
    \begin{equation}
        J_A\Delta^{1/2}_{\varphi|\psi;A}x\xi_{\psi} = x^*\xi_{\varphi}
    \end{equation}
    For any $n$, we have:
    \begin{equation}
        ||\Delta^{1/2}_{\varphi_n|\psi_n;A}x\xi_{\psi} - \Delta^{1/2}_{\varphi|\psi;A}x\xi_{\psi}|| \leq ||\Delta^{1/2}_{\varphi_n|\psi_n;A}x\xi_\psi - \Delta^{1/2}_{\varphi_n|\psi;A}x\xi_\psi|| + ||\Delta^{1/2}_{\varphi_n|\psi;A}x\xi_\psi - \Delta^{1/2}_{\varphi|\psi;A}x\xi_\psi|| 
    \end{equation}
    In addition, for any $n$, we have:
    \begin{equation}
        ||\Delta^{1/2}_{\varphi_n|\psi_n;A}x\xi_\psi - \Delta^{1/2}_{\varphi_n|\psi;A}x\xi_\psi|| \leq ||\Delta^{1/2}_{\varphi_n|\psi_n;A}(x\xi_\psi - x\xi_{\psi_n})|| + ||\Delta^{1/2}_{\varphi_n|\psi_n;A}x\xi_{\psi_n} - \Delta^{1/2}_{\varphi_n|\psi;A}x\xi_{\psi}||
    \end{equation}
    Since $J_A\Delta^{1/2}_{\varphi_n|\psi_n;A}x\xi_{\psi_n} = x^*\xi_{\varphi_n}$ and $J_A\Delta^{1/2}_{\varphi_n|\psi;A}x\xi_\psi = x^*\xi_{\varphi_n}$, we have:
    \begin{align}
        \begin{split}
            ||\Delta^{1/2}_{\varphi_n|\psi_n;A}x\xi_{\psi_n} - \Delta^{1/2}_{\varphi_n|\psi;A}x\xi_{\psi}|| = ||J_A\Delta^{1/2}_{\varphi_n|\psi_n;A}x\xi_{\psi_n} - J_A\Delta^{1/2}_{\varphi_n|\psi;A}x\xi_{\psi}|| = 0
        \end{split}
    \end{align}
    From Powers-Stormer inequality, we have:
    \begin{align}
        ||x\xi_\psi -x\xi_{\psi_n}||^2 \leq ||x||^2||\psi-\psi_n||
    \end{align}
    Hence, we have:
    \begin{align}
        \begin{split}
            ||\Delta^{1/2}_{\varphi_n|\psi_n;A}(x\xi_\psi -x\xi_{\psi_n})|| &= ||J_A\Delta^{1/2}_{\varphi_n|\psi_n;A}(x\xi_\psi - x\xi_{\psi_n})||
            \\&
            \leq ||x^*(\xi_\psi - \xi_{\psi_n})|| \leq ||x||||\psi - \psi_n||\rightarrow 0
        \end{split}
    \end{align}
    Finally, by modular theory we have:
    \begin{equation}
        \lim_n||\Delta^{1/2}_{\varphi_n|\psi;A}x\xi_\psi - \Delta^{1/2}_{\varphi|\psi;A}x\xi_\psi|| = 0
    \end{equation}
    Therefore, combining all these observations, we have:
    \begin{equation}
        \lim_n||\Delta^{1/2}_{\varphi_n|\psi_n;A}x\xi_\psi - \Delta^{1/2}_{\varphi|\psi;A}x\xi_\psi|| = 0
    \end{equation}
    By a theorem in \cite{takesaki2}, this means that $\{\Delta_{\varphi_n|\psi_n;A}\}$ converges to $\Delta_{\varphi|\psi;A}$ in the strong resolvent sense. 

    Because $B\subset A$ is a standard inclusion, $B$ is also in the standard form on $\mathcal{H}$. Replacing all modular conjugation from $J_A$ to $J_B$, the proof above can be used verbatim to show that $\{\Delta_{\varphi_n|\psi_n;B}\}$ converges to $\Delta_{\varphi|\psi;B}$ in the strong resolvent sense as well.
\end{proof}
\begin{corollary}
    Under the same conditions as in the theorem, $\{\Delta^{it}_{\varphi_n|\psi_n;A}\}$ converges strongly to $\Delta^{it}_{\varphi|\psi;A}$. Same for $\{\Delta^{it}_{\varphi_n|\psi_n;B}\}$.

    In addition, the product $\{\Delta^{it}_{\varphi_n|\psi_n;A}\Delta^{-it}_{\varphi_n|\psi_n;B}\}$ converges strongly to $\Delta^{it}_{\varphi|\psi;A}\Delta^{-it}_{\varphi|\psi;B}$.
\end{corollary}
\begin{proof}
    The first statement follows from the characterization of convergence in strong resolvent sense \cite{takesaki2}.

    The second statement follows from the Uniform Boundedness Principle and the first statement.
\end{proof}
Notice that in the proof above, we have used the fact that both $A$ and $B$ represented on $\mathcal{H}$ are in their standard forms. Therefore, the modular conjugations do not depend on the states under consideration. That is why we have used the notations $J_A$ and $J_B$ to emphasize that the modular conjugations only depend on the algebra not the states.

Next we recall some simple relations regarding relative modular operator:
\begin{equation}
    J_A\Delta_{\varphi|\psi;A}J_A = \Delta^{-1}_{\psi|\varphi;A}
\end{equation}
where $\varphi,\psi$ are both normal faithful states.

Next, we make a simple calculation:
\begin{align}
    \begin{split}
        J_A \Delta^{it + 1/2}_{\varphi|\psi;A}\Delta^{-it-1/2}_{\varphi|\psi;B} J_B &= \Delta^{it}_{\psi|\varphi;A} J_A\Delta^{1/2}_{\varphi|\psi;A}J_B\Delta^{1/2}_{\psi|\varphi;B}\Delta_{\psi|\varphi;B}^{-it}
        \\&=\Delta^{it}_{\psi|\varphi;A}S_{\varphi|\psi;A}S_{\psi|\varphi;B}\Delta^{-it}_{\psi|\varphi;B}
    \end{split}
\end{align}
where $S_{\varphi|\psi;A}$ and $S_{\psi|\varphi;B}$ are the relative modular operators. A simple calculation shows:
\begin{equation}
    S_{\varphi|\psi;A}S_{\psi|\varphi;B}x\xi_{\varphi} = S_{\varphi|\psi;A}x^*\xi_{\psi} = x\xi_{\varphi}
\end{equation}
Hence, we have the identity:
\begin{equation}
    J_A\Delta^{it+1/2}_{\varphi|\psi;A}\Delta^{-it-1/2}_{\varphi|\psi;B}J_B = \Delta^{it}_{\psi|\varphi;A}\Delta^{-it}_{\psi|\varphi;B}
\end{equation}
Using this identity and the previous corollary, it is straight-forward to see:
\begin{corollary}
    Under the same condition as the theorem, we have strong convergence:
    \begin{equation}
        \lim_n\Delta^{it + 1/2}_{\varphi_n|\psi_n;A}\Delta^{-it-1/2}_{\varphi_n|\psi_n;B} = \Delta^{it + 1/2}_{\varphi|\psi;A}\Delta^{-it-1/2}_{\varphi|\psi;B}
    \end{equation}
\end{corollary}
%For any pairs of vectors $\xi,\eta\in\mathcal{H}$, we consider the functions:
%\begin{equation}
%    f^{(n)}_{\xi,\eta}(z) := \bra\eta, \Delta^{z}_{\varphi_n|\psi_n;A}\Delta^{-z}_{\varphi_n|\psi_n;B} - \Delta^z_{\varphi|\psi;A}\Delta^{-z}_{\varphi|\psi;B}\ket{\xi} 
%\end{equation}
%For each given $n$, the function is analytic inside the strip $\Re{z}\in(0,1/2)$
%Consider the relative modular operator $\Delta_{\phi|\psi}$. For the special case, $\phi=\psi$ this is simply the modular operator of $\psi$. The first observation is that a sequence of states $\phi_n\to \phi$ and $\psi_n\to \psi$ and any bounded continuous function $f$ we have
%\begin{eqnarray}
%    \lim_n f(\Delta_{\phi_n|\psi_n})=f(\Delta_{\phi|\psi})\ .
%\end{eqnarray}
%In particular this implies for all $\beta>0$ we have
%\begin{eqnarray}
%    &&\lim_n \frac{1}{\beta+\Delta_{\phi_n|\psi_n}} =\frac{1}{\beta+\Delta_{\phi|\psi}}
    %\nn\\
    %&&\lim_n \frac{\Delta_{\phi_n|\psi_n}}{\beta(\beta+\Delta_{\phi_n|\psi_n})}=\frac{1}{\beta}-\lim_n\frac{1}{\beta+\Delta_{\phi_n|\psi_n}} =\frac{\Delta_{\phi_n|\psi_n}}{\beta(\beta+\Delta_{\phi|\psi})}
%\end{eqnarray}
%in operator norm topology. Note that for the sequence $(\beta+\Delta_{\phi_n|\psi_n})^{-1}$ is uniformly bounded by $1/\beta$. 

%If $A_n\to A$ and $B_n\to B$ in operator norm topology and $A_n$ are uniformly bounded meaning that $\sup_n\|A_n\|<\infty$ then we have $\lim_n A_nB_n=A B$. The proof is straightforward:
%\begin{eqnarray}
%    \lim_n \|A_nB_n-A B\|&=&\lim_n\|A_n (B_n-B)\|+\lim_n \|(A_n-A)B\|\nn\\
%    & \leq &\sup_n \|A_n\|\lim_n\|B_n-B\|+\|B\|\lim_n\|A_n-A\|=0\nn\ .
%\end{eqnarray}
%Now, consider two algebras $\mA$ and $\mB$ (need not satisfy inclusion relation), then for all $\beta,\beta'>0$ we have
%\begin{eqnarray}\label{limitmultiply}
%    \lim_n \lb\frac{1}{\beta+\Delta_{\phi_n|\psi_n;\mA}}\rb\lb \frac{1}{\beta'+\Delta_{\phi_n|\psi_n;\mB}}\rb=\lb\frac{1}{\beta+\Delta_{\phi|\psi;\mA}}\rb\lb \frac{1}{\beta'+\Delta_{\phi|\psi;\mB}}\rb
%\end{eqnarray}
%Next, we use the following expansions for $\alpha\in (0,1)$ and positive operator $\Delta$:
%\begin{eqnarray}
%    &&\Delta^{-\alpha}=\frac{\sin(\pi\alpha)}{\pi}\int_0^\infty d\beta \:\beta^{-\alpha}\lb \frac{1}{\beta+\Delta}\rb\nn\\
%    &&\Delta^{\alpha}=\frac{\sin(\pi\alpha)}{\pi}\int_0^\infty d\beta \:\beta^{\alpha}\lb \frac{1}{\beta}-\frac{1}{\beta+\Delta}\rb
%\end{eqnarray}
%As a result, we obtain
%\begin{eqnarray}
%    \Delta_{\phi_n|\psi_n;\mA}^\alpha\Delta_{\phi|\psi;\mB}^{-\alpha}=\frac{\sin^2(\pi \alpha)}{\pi^2}\int_0^\infty d\beta d\beta' \lb\frac{\beta}{\beta'}\rb^\alpha \lb \frac{1}{\beta}-\frac{1}{\beta+\Delta_{\phi_n
%|\psi_n;\mA}}\rb\lb \frac{1}{\beta'+\Delta_{\phi_n|\psi_n;\mB}}\rb 
%\end{eqnarray}
%To commute the limit and integral, we can use the dominated convergence them.
%It follows from (\ref{limitmultiply}) that the integral above converges pointwise, but apriori we do not know the integral for all $n$ is dominated by a single integrable operator. If $\mA\subset \mB$ from the monotonicity condition we know that for all $z$ in the complex strip $\Re(z)\in (0,1/2)$ we have
%\begin{eqnarray}
%    |\Delta_{\phi|\psi;\mA}^z \Delta_{\phi|\psi;\mB}^{-z}|\leq 1\ .
%\end{eqnarray}
%Therefore, for inclusion of algebra, we can use the dominated convergence theorem to find
%\begin{eqnarray}
%    \lim_n \Delta^\alpha_{\phi_n|\psi_n;\mA}\Delta_{\phi_n|\psi_n;\mB}^{-\alpha}=\Delta^\alpha_{\phi|\psi;\mA}\Delta_{\phi|\psi;\mB}^{-\alpha}\ .
%\end{eqnarray}
%\NL{Does this justify continuity?}


% Wrong, don't use it \KL{We consider $f_{ab}(z) \equiv \bra{a}\Delta_A^{1/2+iz}\Delta_B^{-iz}\ket{b} = \bra{\Delta_A^{1/4}a} T_{B\subset A}(z-i/4)\ket{\Delta_B^{1/4} b} $ with some special choice of $a,b$. $f_{ab}(z)$ is holomorphic on the stripe Im$(z) \in (-\frac{1}{4}, \frac{1}{4})$ and continuous in the closure. The characteristic function has the boundary value that, for real $t$,
% \begin{align}
%     T_{B\subset A}(t-i/2) = J_AT_{B\subset A}(t)J_B.
% \end{align}
% By choosing $a\in A$ and $b\in B \subset A$ as positive operators, we have
% \begin{align}
%     J_A \ket{\Delta_A^{1/4}a} = \ket{\Delta_A^{1/4}a},\nonumber\\
%     J_B \ket{\Delta_B^{1/4}b} = \ket{\Delta_B^{1/4}b},
% \end{align}
% because both vectors are in the corresponding natural cones. Thus, the boundary value of $f_{ab}$ satisfies, for real $t$,
% \begin{align}
%     f_{ab}(t+i/4) = f_{ab}(t-i/4).
% \end{align}
% \textbf{Mistake, check the anti-unitary steps}
% The claim is that, with this choice of $a$ and $b$, we can extend $f_{ab}(z)$ to an entire function, i.e. a holomorphic function on the whole complex plane. In fact, $f_{ab}(z)$ is a constant function. 
% \begin{lemma}
%     If $a\in A$ and $b\in B \subset A$ are both positive operators, $f_{ab}(z)$ is a constant function.
% \end{lemma}
% To prove it, we need the following result.
% \begin{lemma}
%     If $f(z)$ is holomorphic on the stripe Im$(z) \in (0,1)$ and continuous in the closure with boundary values
%     \begin{align}
%         f(t+i) = f(t)
%     \end{align}
%     for real $t$, there exists an entire function that is an extension of $f(z)$ to the whole complex plane.
% \end{lemma}
% \begin{proof}
%     The function $\tilde{g}(z) = \overline{f(\overline{z})}$ is holomorphic on the stripe Im$(z) \in (-1,0)$ and continuous in the closure with boundary values
%     \begin{align}
%         \tilde{g}(t-i) = \tilde{g}(t) = \overline{f(t)}
%     \end{align}  
%     for real $t$. We define $g(z) = \tilde{g}(z-i)$ which shifts the domain to Im$(z) \in (0,1)$. Let $F_+(z) = \frac{f+g}{2}$ and $F_-(z) = \frac{f-g}{2i}$. $F_\pm(z)$ are both holomorphic on the stripe Im$(z) \in (0,1)$ and continuous in the closure with real boundary values. By the Schwarz reflection principle, we can extend $F_\pm(z)$ to Im$(z) \in (-1,1)$ with $F_\pm(z) = \overline{F_\pm(\overline{z})}$. Since the boundary values on two sides of the stripe are real, we can shift the stripe and apply the Schwarz reflection principle repeatedly to extend $F_\pm(z)$ to the whole complex plane. Thus, $f(z) = F_+(z) + i F_-(z)$ can be extended to the whole complex plane.
% \end{proof}
% If we have $|f(z)|\leq 1$ on the stripe Im$z \in [0,1]$, then $|g(z)|$ and $|F_\pm(z)|$ are also bounded by $1$ on the stripe. The extension through the Schwarz reflection principle does not change the bound of $F_\pm(z)$. A bounded entire function must be a constant function. Therefore, $F_\pm(z)$ are constant functions, and $f(z)$ is also a constant function. This finishes the proof, as $f_{ab}(z)$ is a bounded function. }

% \section{Titchmarsh theorem}\label{app:titchmarsh}

% Titchmarsh theorems are a class of Paley-Weiner type theorems that relate the behavior of complex-valued square-integrable functions $S(t)\in L^2(\mathbb{R})$ to the analytic domain of their Fourier transform
% \begin{equation}
%     \tilde{S}(\omega) = \int_{-\infty}^{\infty} S(t) e^{i\omega t} \, dt\ .
% \end{equation}
% Theorem \ref{titchmarshUHP} gives a condition that is equivalent to $\tilde{S}(\omega)$ being analytic on the upper-half-plane. Theorem \ref{titchmarshStrip} generalizes it to the condition equivalent to the analyticity of $\tilde{S}(\omega)$ in the strip $\Im(\omega)\in (-A,B)$.

% \subsection{Analyticity in frequency upper-half-plane:}
% \begin{theorem}[Titchmarsh]\label{titchmarshUHP}
% The following three statements are mutually equivalent:

% \begin{enumerate}
%     \item \textbf{Causality:} $S(t) = 0$ for almost every $t < 0$: $S(t)=\text{Sgn}(t)S(t)$.
%     \item \textbf{Analyticity:} The function $\tilde{S}(\omega)$ is the boundary limit of a function $G(z)$ analytic in the upper half-plane $\mathbb{C}^+$ (belonging to Hardy space $H^2$).
%     \item \textbf{Hilbert Transform Pair:}
%     The real and imaginary parts of the spectrum $\tilde{S}(\omega) = U(\omega) + i V(\omega)$ are Hilbert transforms of each other:
%     \begin{align}
%         V(\omega) &= -\mathcal{H}[U](\omega) \\
%         U(\omega) &= \mathcal{H}[V](\omega)
%     \end{align}
%     where the Hilbert transform is defined as:
%     \begin{equation}
%         \mathcal{H}[f](\omega) = \frac{1}{\pi} \mathcal{P} \int_{-\infty}^{\infty} \frac{f(\omega')}{\omega' - \omega} \, d\omega'
%     \end{equation}
% \end{enumerate}
% \end{theorem}

% \begin{proof}
% \textbf{1) $\implies$ 2) (Causality to Analyticity)} \\
% Assume $S(t) = 0$ for $t < 0$. We define the analytic extension $G(z)$ for $z = \omega + i\eta$ with $\eta > 0$:
% \begin{equation}
%     G(z) = \int_{0}^{\infty} S(t) e^{i(\omega + i\eta)t} \, dt = \int_{0}^{\infty} S(t) e^{-\eta t} e^{i\omega t} \, dt
% \end{equation}
% Because $S(t) \in L^2$ and $e^{-\eta t} \le 1$ for $t > 0$, the integral converges absolutely and defines a holomorphic function in the upper half-plane. By Plancherel's theorem, it satisfies the Hardy space condition. As $\eta \to 0^+$, $G(z)$ converges to $\tilde{S}(\omega)$ in the $L^2$ norm.

% \vspace{0.5cm}
% \textbf{2) $\implies$ 3) (Analyticity to Hilbert Transform)} \\
% By Cauchy's Integral Formula for $z \in \mathbb{C}^+$:
% \begin{equation}
%     G(z) = \frac{1}{2\pi i} \int_{-\infty}^{\infty} \frac{\tilde{S}(\omega')}{\omega' - z} \, d\omega'
% \end{equation}
% Taking the limit $z \to \omega + i0^+$ using the Plemelj-Sokhotski relations:
% \begin{equation}
%     \tilde{S}(\omega) = \frac{1}{2}\tilde{S}(\omega) + \frac{1}{2\pi i} \mathcal{P} \int_{-\infty}^{\infty} \frac{\tilde{S}(\omega')}{\omega' - \omega} \, d\omega'
% \end{equation}
% Rearranging to isolate $\tilde{S}(\omega)$:
% \begin{equation}\label{contourint}
%     \tilde{S}(\omega) = \frac{1}{\pi i} \mathcal{P} \int_{-\infty}^{\infty} \frac{\tilde{S}(\omega')}{\omega' - \omega} \, d\omega' = -i \mathcal{H}[\tilde{S}](\omega)
% \end{equation}
% Substituting $\tilde{S} = U + iV$:
% \begin{equation}
%     U + iV = -i (\mathcal{H}[U] + i\mathcal{H}[V]) = \mathcal{H}[V] - i\mathcal{H}[U]
% \end{equation}
% Comparing real and imaginary parts yields the relations: $U = \mathcal{H}[V]$ and $V = -\mathcal{H}[U]$.

% \vspace{0.5cm}
% \textbf{3) $\implies$ 1) (Hilbert Transform to Causality)} \\
% From the derivation in step 2, condition (3) implies that the full complex spectrum satisfies:
% \begin{equation}
%     \tilde{S}(\omega) = -i \mathcal{H}[\tilde{S}](\omega)
% \end{equation}
% We apply the inverse Fourier transform, $\mathcal{F}^{-1}$, to both sides. We utilize the property that the Hilbert transform in the frequency domain corresponds to multiplication by the signum function in the time domain. Specifically, for the transform definition used here ($\int e^{i\omega t}$):
% \begin{equation}
%     \mathcal{F}^{-1} \big\{ -i \mathcal{H}[\tilde{S}] \big\}(t) = \text{sgn}(t) S(t)
% \end{equation}
% Substituting this into the equation:
% \begin{equation}
%     S(t) = \text{sgn}(t) S(t)
% \end{equation}
% which implies $S(t) = 0$ for almost all $t < 0$. 
% \end{proof}
% Note that the Hilbert transform can be viewed as the convolution of the function with $f(\omega)=-\frac{1}{\pi\omega}$:
% \begin{eqnarray}
%     U(\omega)=(f*V)(\omega)\ .
% \end{eqnarray}
% Since $\mathcal{F}^{-1}(\frac{-1}{\pi \omega})=\frac{i}{2\pi}\text{sgn}(t)$ we find
% \begin{eqnarray}
%     \tilde{U}(t)=i\text{sgn}(t) \tilde{V}(t)\ .
% \end{eqnarray}
% However, note that both $\tilde{U}(t)$ and $\tilde{V}(t)$ are complex, as they are defined as the inverse Fourier transform of $U(\omega)$ and $V(\omega)$. 

% \begin{corollary}
%     If $S(t)$ that satisfies the conditions of Theorem \ref{titchmarshUHP} is a real function on $\mathbb{R}$, then we have
%     \begin{eqnarray}
%         &&U(\omega)=\frac{2}{\pi}\mathcal{P}\int_0^\infty \frac{\omega' V(\omega')}{\omega^2-(\omega')^2}\nn\\
%         &&V(\omega)=\frac{-2\omega}{\pi}\mathcal{P}\int_0^\infty \frac{U(\omega')}{\omega^2-(\omega')^2}
%     \end{eqnarray}
% \end{corollary}
% \begin{proof}
%     Since $S(t)$ is real for all $t\in\mathbb{R}$, then $S(-\omega)=S(\omega)^*$ which implies that $U(\omega)$ and $V(\omega)$ are even and odd functions of $\omega$, respectively. Then, the corollary follows from multiplying the numerator and denominator of the integrand inside the Contour integral in (\ref{contourint}) by $\omega+\omega'$.
% \end{proof}



% \subsection{Analyticity in the frequency strip:}

% We would like to generalize the Titchmarsh theorem to functions that are analytic in a complex frequency strip. The key intuition is the following: If a complex function of a real parameter $t$ decays exponentially fast in time its Fourier transform is analytic on a strip, and decays exponentially fast in real frequency space. More precisely, we have the following generalized Titchmarsh theorem


% \begin{theorem}[Generalized Titchmarsh for a Strip]\label{titchmarshStrip}
% Let $S(t)$ be a complex-valued function and $\tilde{S}(\omega)$ its Fourier transform. Let $A \ge 0$ and $B \ge 0$ (with $A+B > 0$) be real constants defining a horizontal strip $\mathcal{D}$ in the complex frequency plane:
% \begin{equation}
%     \mathcal{D} = \{ \omega \in \mathbb{C} \mid -A < \text{Im}(\omega) < B \}
% \end{equation}

% The following three statements are mutually equivalent:

% \begin{enumerate}
%     \item \textbf{Two-sided exponential decay:}
%     The function $S(t)$ satisfies the asymptotic decay conditions:
%     \begin{align}
%         S(t) &\sim e^{-B|t|} \quad (t \to -\infty) \\
%         S(t) &\sim e^{-At} \quad (t \to +\infty)
%     \end{align}
%     In other words, for any $\epsilon > 0$, the weighted norms are finite:
%     \begin{equation}
%         e^{(B-\epsilon)t} S(t) \in L^2(-\infty, 0] \quad \text{and} \quad e^{-(A-\epsilon)t} S(t) \in L^2[0, \infty)\ .
%     \end{equation}

%     \item \textbf{Strip Analyticity:}
%     The function $\tilde{S}(\omega)$ is the boundary limit of a function $G(z)$ holomorphic inside the strip $-A < \text{Im}(z) < B$, satisfying the uniform Hardy condition:
%     \begin{equation}
%         \sup_{-A < y < B} \int_{-\infty}^{\infty} |G(x + iy)|^2 \, dx < \infty
%     \end{equation}

%     \item \textbf{Spectral Decomposition and shifted dispersion Relations:}
%     The spectrum $\tilde{S}(\omega)$ uniquely decomposes into:
%     \begin{equation}
%         \tilde{S}(\omega) = \tilde{S}_+(\omega) + \tilde{S}_-(\omega)
%     \end{equation}
%     where $\tilde{S}_+(\omega)$ is analytic for $\text{Im}(\omega) > -A$ and $\tilde{S}_-(\omega)$ is analytic for $\text{Im}(\omega) < B$. \\
    
%     Crucially, the real and imaginary parts of these components satisfy \textbf{Shifted Hilbert Transform} relations on the boundary lines:
    
%     \begin{itemize}
%         \item \textbf{For $\tilde{S}_+$ (on the line $\omega = \nu - iA$):} \\
%         Let $\tilde{S}_+(\nu - iA) = U_+(\nu) + i V_+(\nu)$. Then:
%         \begin{equation}
%             U_+(\nu) = \mathcal{H}[V_+](\nu), \quad V_+(\nu) = -\mathcal{H}[U_+](\nu)
%         \end{equation}
        
%         \item \textbf{For $\tilde{S}_-$ (on the line $\omega = \nu + iB$):} \\
%         Let $\tilde{S}_-(\nu + iB) = U_-(\nu) + i V_-(\nu)$. Then:
%         \begin{equation}
%             U_-(\nu) = -\mathcal{H}[V_-](\nu), \quad V_-(\nu) = \mathcal{H}[U_-](\nu)
%         \end{equation}
%     \end{itemize}
% \end{enumerate}
% \end{theorem}

% \begin{proof}
% \textbf{1) $\implies$ 2) (Decay to Analyticity)} \\
% We define the complex integral $G(z) = \int_{-\infty}^{\infty} S(t) e^{i z t} \, dt$ for $z = x + iy$.
% Splitting the integral:
% \begin{equation}
%     G(z) = \int_{-\infty}^{0} S(t) e^{(B-y)t} e^{-(B-i x)t} \, dt + \int_{0}^{\infty} S(t) e^{-(A+y)t} e^{(A+ix)t} \, dt
% \end{equation}
% Convergence requires $y < B$ (for the first term) and $y > -A$ (for the second term). Thus, $G(z)$ is analytic in $-A < y < B$. The Hardy condition follows from Plancherel's theorem applied to $S(t)e^{-yt}$.

% \vspace{0.5cm}
% \textbf{2) $\implies$ 3) (Analyticity to Decomposition)} \\
% Using Cauchy's Integral Formula on a rectangular contour inside the strip and sending the horizontal width to infinity, we obtain:
% \begin{equation}
%     \tilde{S}(z) = \frac{1}{2\pi i} \int_{-\infty - iA}^{\infty - iA} \frac{\tilde{S}(\zeta)}{\zeta - z} \, d\zeta - \frac{1}{2\pi i} \int_{-\infty + iB}^{\infty + iB} \frac{\tilde{S}(\zeta)}{\zeta - z} \, d\zeta
% \end{equation}
% We define the first integral as $\tilde{S}_+(z)$ and the second as $\tilde{S}_-(z)$. The analyticity domains and shifted Hilbert relations follow directly from the standard Titchmarsh theorem applied to the shifted functions $\tilde{S}_+(z-iA)$ and $\tilde{S}_-(z+iB)$.

% \vspace{0.5cm}
% \textbf{3) $\implies$ 1) (Decomposition to Decay)} \\
% Inverse transforming $\tilde{S}_+$ yields a function supported on $t>0$ with weight $e^{-At}$. Inverse transforming $\tilde{S}_-$ yields a function supported on $t<0$ with weight $e^{Bt}$. The sum satisfies the two-sided decay.
% \end{proof}

% \subsection*{Discussion of Limits}

% \paragraph{1. The Causal Limit ($A \to 0, B \to \infty$)}
% The strip becomes the Upper Half-Plane. The condition $e^{Bt}S(t) \in L^2$ as $B \to \infty$ forces $S(t) = 0$ for $t<0$. This recovers the standard Kramers-Kronig relations on the real axis.

% \paragraph{2. The Anti-Causal Limit ($A \to \infty, B \to 0$)}
% The strip becomes the Lower Half-Plane. The condition $e^{-At}S(t) \in L^2$ as $A \to \infty$ forces $S(t) = 0$ for $t>0$.

% \paragraph{3. The Degenerate Limit ($A \to 0, B \to 0$)}
% The strip collapses to the real axis. The region of analyticity vanishes. The theorem reduces to the Plancherel Theorem ($S \in L^2 \iff \tilde{S} \in L^2$). The real and imaginary parts of the spectrum become independent, and no dispersion relations exist.
% %\section{Coherent Perturbations of GFF}



\section{Fredenhagen's bound on modular flow of balls}\label{subsection:gaussiansmear}

In this subsection, we collect several results (including explicit calculations) of Fredenhagen's constant \cite{fredenhagen}. Recall the Fredenhagen's original estimate depends on the constant:
\begin{equation}
    c_f(s):=\inf_{\substack{g\in\mathcal{D}(\mathbb{R})\\ c > 0}}\{2||f - g||_1 + 2\int_{|u|\geq c}du|\varphi_g(u)| + \frac{2}{\pi}\tanh^{-1}(e^{-\pi s})e^{c/2}||\varphi_g||_1\}
\end{equation}
where $g\in\mathcal{D}(\mathbb{R})$ is a Schwartz function and $\varphi_g(u)$ is  a twisted Fourier transform of $g$:
\begin{equation}
    \varphi_g(u):=\frac{1}{2\pi}\int_\mathbb{R} dt g(t)\lb\frac{2}{i}\sinh(\pi t)\rb e^{itu}
\end{equation}
We first make the following observation on Fredenhagen's constant for shifted function $f_\mu(t) = e^{-i\mu t}f(t)$:
\begin{lemma}\label{lemma:shiftedfredenhagenconstant}
    For a Schwartz function $f\in \mathcal{D}(\mathbb{R})$, the following limits hold:
    \begin{equation}
        \lim_{s\rightarrow\infty}c_{f_\mu}(s) = 0
    \end{equation}
    \begin{equation}
        \lim_{s\rightarrow\infty}\sup_{-\mu_0\leq \mu\leq\mu_0}c_{tf_\mu}(s) = 0
    \end{equation}
\end{lemma}
\begin{proof}
    The proof is based on an explicit calculation of $c_{f_\mu}(s)$. Using the definition of $c_f(s)$, we have the following:
    \begin{align}
        \begin{split}
            c_{f_\mu}(s) &= \inf_{\substack{g\in\mathcal{D}(\mathbb{R})\\ c>0}}\{2||f_\mu - g|| + 2\int_{|u|\geq c}du|\varphi_g(u)| + \frac{2}{\pi}\tanh^{-1}(e^{-\pi s})e^{c/2}||\varphi_g||_1\}
            \\&
            =\inf_{\substack{g_\mu\in\mathcal{D}(\mathbb{R})\\ c > 0}}\{2||f_\mu - g_\mu||_1 + 2\int_{|u|\geq c}du|\varphi_{g_\mu}(u)| + \frac{2}{\pi}\tanh^{-1}(e^{-\pi s})e^{c/2}||\varphi_{g_\mu}||_1\}
            \\&
            =\inf_{\substack{g_\mu\in\mathcal{D}(\mathbb{R})\\ c > 0}}\{2||f-g||_1 + 2\int_{|u| \geq c}du|\varphi_g(u + \mu)| + \frac{2}{\pi}\tanh^{-1}(e^{-\pi s})e^{c/2}||\varphi_g||_1\}
            \\&
            =\inf_{\substack{g\in\mathcal{D}(\mathbb{R})\\ c > 0}}\{2||f-g||_1 + 2\int_{|u - \mu|\geq c}du|\varphi_g(u)| + \frac{2}{\pi}\tanh^{-1}(e^{-\pi s})e^{c/2}||\varphi_g||_1\}
        \end{split}
    \end{align}
    where $g_\mu(t) := e^{-i\mu t}g(t)\in\mathcal{D}(\mathbb{R})$. In addition, by definition of $\varphi_g$, we have:
    \begin{align}
        \begin{split}
            \varphi_{g_\mu}(u)&=\frac{1}{2\pi}\int_\mathbb{R} dt g_\mu(t)\lb\frac{2}{i}\sinh(\pi t)\rb e^{-itu} = \frac{1}{2\pi}\int_\mathbb{R} dt g(t)\lb\frac{2}{i}\sinh(\pi t)\rb e^{-i(\mu + u)t} = \varphi_g(u + \mu)
        \end{split}
    \end{align}
    In particular by translation invariance of the Lebesgue measure, we have $||\varphi_g||_1 = ||\varphi_{g_\mu}||_1$.

    Since the function $f$ is a Schwartz function, we have the following upper bound by taking $g = f$:
    \begin{align}
        \begin{split}
            c_{f_\mu}(s) &\leq \inf_{c > 0}\{2\int_{|u - \mu| > c}du|\varphi_f(u)| + \frac{2}{\pi}\tanh^{-1}(e^{-\pi s})e^{c/2}||\varphi_f||_1\}
        \end{split}
    \end{align}
    In addition, because $f$ is a Schwartz function, the twisted Fourier transform $\varphi_f$ is integrable. Hence, for any small constant $\epsilon > 0$, there exists $c_\epsilon > 0$ such that:
    \begin{equation}
        2\int_{|u - \mu| > c_\epsilon}du|\varphi_f(u)| < \epsilon
    \end{equation}

    Hence for any $\epsilon > 0$, we have:
    \begin{align}
        \begin{split}
            \lim_{s\rightarrow\infty}c_{f_\mu}(s) &\leq \lim_{s\rightarrow\infty}2\int_{|u-\mu|>c_\epsilon}du|\varphi_f(u)| + \frac{2}{\pi}\tanh^{-1}(e^{-\pi s})e^{c_\epsilon/2}||\varphi_f||_1
            \\&
            <\lim_{s\rightarrow\infty}2\epsilon + \frac{2}{\pi}\tanh^{-1}(e^{-\pi s})e^{c_\epsilon/2}||\varphi_f||_1 = 2\epsilon
        \end{split}
    \end{align}
    Since the bound holds for any $\epsilon > 0$, we have $\lim_{s\rightarrow\infty}c_{f_\mu}(s) = 0$.

    To calculate the limit $\lim_{s\rightarrow\infty}c_{tf_\mu}(s)$, first notice that because $f$ is a Schwartz function, $tf_\mu(t) = tf(t)e^{-i\mu t}$ is also a Schwartz function. Using the same arguments as above, we have the following upper bound:
    \begin{equation}
        \sup_{-\mu_0\leq \mu\leq\mu_0}c_{tf_\mu}(s) \leq \sup_{-\mu_0\leq \mu\leq\mu_0}\inf_{c > 0}\{2\int_{|u - \mu| > c}du|\varphi_{tf}(u)|+ \frac{2}{\pi}\tanh^{-1}(e^{-\pi s})e^{c/2}||\varphi_{tf}||_1\}
    \end{equation}
    Consider $\int_{|u-\mu| > c}du|\varphi_{tf}(u)|$ as a family of functions $\{F_\mu(c)\}_{-\mu_0\leq\mu\leq\mu_0}$. For any $\mu \in [-\mu_0, \mu_0]$, the function $F_\mu$ is Lipschitz with the same constant:
    \begin{align}
        \begin{split}
            ||\frac{d}{dc}F_\mu(c)||_\infty &= ||\frac{d}{dc}\int_{|u| > c}|\varphi_{tf}(u + \mu)|du||_\infty \leq 2||\varphi_{tf}||_\infty
        \end{split}
    \end{align}
    Hence the family of functions $\{F_\mu\}_{0\leq \mu\leq \mu_0}$ is equicontinuous. In particular, since for any $\mu\in [-\mu_0, \mu_0]$ the limit $\lim_{c\rightarrow \infty}F_\mu(c) = 0$ (by the integrability of $\varphi_{tf}$), then for any $\epsilon > 0$ there exists $c_\epsilon > 0$ such that for all $\mu \in [-\mu_0, \mu_0]$, we have $|F_\mu(c_\epsilon)| < \epsilon$. Thus for any $\epsilon > 0$, we have:
    \begin{align}
        \begin{split}
            &\sup_{-\mu_0\leq \mu\leq \mu_0}c_{tf_\mu}(s) \leq \sup_{-\mu_0\leq \mu\leq\mu_0}\{2\int_{|u - \mu| > c_\epsilon} du |\varphi_{tf}(u)| + \frac{2}{\pi}\tanh^{-1}(e^{-\pi s})e^{c_\epsilon / 2}||\varphi_{tf}||_1\}
            \\&
            < 2\epsilon + \frac{2}{\pi}\tanh^{-1}(e^{-\pi s})e^{c_\epsilon / 2}||\varphi_{tf}||_1
        \end{split}    
    \end{align}
    Taking the limit $s\rightarrow\infty$, for any $\epsilon > 0$ we have the limit:
    \begin{align}
        \begin{split}
            \lim_{s\rightarrow\infty}\sup_{-\mu_0\leq \mu\leq \mu_0}c_{tf_\mu}(s) \leq \lim_{s\rightarrow\infty}2\epsilon + \frac{2}{\pi}\tanh^{-1}(e^{-\pi s})e^{c_\epsilon / 2}||\varphi_{tf}||_1 = 2\epsilon
        \end{split}
    \end{align}
    Since the bound holds for any $\epsilon > 0$, we have the limit $\lim_{s\rightarrow\infty}\sup_{0\leq \mu\leq \mu_0}c_{tf_\mu}(s) = 0$.
\end{proof}
In the remainder of this section, we provide explicit upper bound on Fredenhagen's constant when the smearing function is Gaussian. The purpose of this calculation is to confirm several physical intuitions.
\begin{lemma}
    For Gaussian $f_{\sigma,\mu}(t) := \frac{1}{\sigma\sqrt{2\pi}}e^{-\frac{(t - \mu)^2}{2\sigma^2}}$, there exists $s_0 > 0$ and positive constants $K_3, K_4, K_5 > 0$ such that for $s > s_0$, to the first order in $s$ the following bound holds:
    \begin{equation}
        c_{f_{\sigma,\mu}}(s) \leq K_3e^{-\frac{K_4}{2}s^2} + K_5e^{-K_4 s^2}
    \end{equation}
    These positive constants depend on $\sigma,\mu$.
\end{lemma}
\begin{proof}
    Since Gaussian function is a Schwartz function, we have the following bound:
    \begin{equation}\label{equation:bound}
        c_{f_{\sigma,\mu}}(s) \leq \inf_{c > 0}\{2\int_{|u|\geq c}du |\varphi_{f_{\sigma,\mu}}(u)| + \frac{2}{\pi}\tanh^{-1}(e^{-\pi s})e^{c/2}||\varphi_{f_{\sigma,\mu}}||_1\}
    \end{equation}
    The modified Fourier transform of $f_{\sigma,\mu}$ can be calculated explicitly:
    \begin{align}
        \begin{split}
            \varphi_{f_{\sigma,\mu}}(u) &= \frac{1}{2\pi}\int_\mathbb{R} dt \lb \frac{1}{\sigma\sqrt{2\pi}}e^{-\frac{(t - \mu)^2}{2\sigma^2}} \rb\lb \frac{2}{i}\sinh(\pi t) \rb e^{-it u}
            \\&
            =\frac{1}{\sigma\sqrt{2\pi }2\pi i}\int_\mathbb{R} dt e^{-\frac{(t-  \mu)^2}{2\sigma^2} + (\pi - i u) t} - e^{-\frac{(t - \mu)^2}{2\sigma^2} - (\pi + i u) t}
            \\&
            =\frac{1}{2\pi i}\lb e^{\frac{(\pi - iu)(2\mu + \sigma^2(\pi - iu))}{2}} - e^{-\frac{(\pi + iu)(2\mu - \sigma^2(\pi + iu))}{2}}\rb
            \\&
            =\frac{1}{\pi i}e^{\frac{\sigma^2}{2}(\pi^2 - u^2) - i\mu u}\sinh(\pi \mu  - \pi i u\sigma^2)
        \end{split}
    \end{align}
    This $L^1$-norm of $\varphi_{f_{\sigma,\mu}}$ can be bounded above by:
    \begin{align}
        \begin{split}
            ||\varphi_{f_{\sigma,\mu}}||_1 &= \int_\mathbb{R} du |\varphi_{f_{\sigma,\mu}}(u)| = \int_\mathbb{R} du \frac{e^{\frac{\sigma^2}{2}(\pi^2 - u^2)}}{\pi}|\sinh(\pi\mu - \pi i u \sigma^2)|
            \\&
            \leq \int_\mathbb{R}du\frac{e^{\frac{\sigma^2}{2}(\pi^2 - u^2)}}{\pi}e^{\pi\mu}|\sin(\pi u \sigma^2)|\leq \frac{e^{\pi\mu + \frac{\sigma^2\pi^2}{2}}}{\pi}\frac{\sqrt{2\pi}}{\sigma} = \frac{e^{\pi\mu}\sqrt{2}}{\sqrt{\pi}}\frac{e^{\frac{\sigma^2\pi^2}{2}}}{\sigma}
        \end{split}
    \end{align}
    To calculate the infimum on the right hand side of Equation \ref{equation:bound}, notice that the function on the right hand side is differentiable in $c$. Hence the infimum can be calculated using elementary methods. The infimum is achieved at $c_0 > 0$ where $c_0$ is the solution to:
    \begin{equation}
        -|\varphi_{f_{\sigma,\mu}}(c)| - |\varphi_{f_{\sigma,\mu}}(-c)| + \frac{1}{\pi}\tanh^{-1}(e^{-\pi s})e^{c/2}||\varphi_{f_{\sigma,\mu}}||_1 = 0
    \end{equation}
    Notice that $|\varphi_{f_{\sigma,\mu}}(c)| = |\varphi_{f_{\sigma,\mu}}(-c)|$:
    \begin{align}
        \begin{split}
            |\varphi_{f_{\sigma,\mu}}(c)| &= \frac{e^{\frac{\sigma^2}{2}(\pi^2 - c^2)}}{\pi}|\sinh(\pi\mu - \pi i c\sigma^2)| \\&= \frac{e^{\frac{\sigma^2}{2}(\pi^2 - c^2)}}{\pi}\sqrt{\sin^2(\pi c\sigma^2)\sinh^2(\pi\mu) + \cos^2(\pi c \sigma^2)\cosh^2(\pi\mu)}= |\varphi_{f_{\sigma,\mu}}(-c)|
        \end{split}
    \end{align}
    Hence the infimum is achieved at $c_0$ where $c_0$ is the solution to:
    \begin{equation}
        |\varphi_{f_{\sigma,\mu}}(c_0)| = \frac{1}{2\pi}\tanh^{-1}(e^{-\pi s})e^{c_0/2}||\varphi_{f_{\sigma,\mu}}||_1
    \end{equation}
    We can provide a more explicit lower bound on $c_0$:
    \begin{align}
        \begin{split}
            \frac{1}{2\pi}\tanh^{-1}(e^{-\pi s})e^{c_0 / 2}||\varphi_{f_{\sigma,\mu}}||_1 &=\frac{e^{\frac{\sigma^2}{2}(\pi^2 - c_0^2)}}{\pi}\sqrt{\sin^2(\pi c_0\sigma^2)\sinh^2(\pi\mu) + \cos^2(\pi c_0 \sigma^2)\cosh^2(\pi\mu)}
            \\&
            = \frac{e^{\frac{\sigma^2}{2}(\pi^2 - c_0^2)}}{\pi}\sqrt{\cosh^2(\pi\mu) - \sin^2(\pi c_0\sigma^2)\lb\cosh^2(\pi\mu) - \sinh^2(\pi\mu)\rb}
            \\&
            =\frac{e^{\frac{\sigma^2}{2}(\pi^2 - c_0^2)}}{\pi}\sqrt{\cosh^2(\pi\mu) - \sin^2(\pi c_0\sigma^2)}
            \\&\geq\frac{e^{\frac{\sigma^2}{2}(\pi^2 - c_0^2)}}{\pi}\sqrt{\cosh^2(\pi\mu) - 1}
            =\frac{e^{\frac{\sigma^2}{2}(\pi^2 - c_0^2)}}{\pi}\sinh(\pi\mu)
        \end{split}
    \end{align}
    Therefore, we have:
    \begin{equation}
        e^{-\frac{\sigma^2}{2}c_0^2 - \frac{c_0}{2}} \leq \frac{\tanh^{-1}(e^{-\pi s})||\varphi_{f_{\sigma,\mu}}||_1 e^{-\frac{\pi^2\sigma^2}{2}}}{2\sinh(\pi\mu)}
    \end{equation}
    Since $c_0 > 0$, using this formula we can solve for a lower bound of $c_0$:
    \begin{equation}
        c_0 \geq c_0(\sigma, \mu, s) := -\frac{1}{2\sigma^2} + \sqrt{\frac{1}{4} - \frac{2}{\sigma^2}\log\lb\frac{\tanh^{-1}(e^{-\pi s})||\varphi_{f_{\sigma,\mu}}||_1 e^{-\frac{\pi^2\sigma^2}{2}}}{2\sinh(\pi\mu)}\rb}
    \end{equation}
    Notice that $c_0(\sigma,\mu, s)$ monotonely increases as $s$ increases. In particular, there exists $s_0 > 0$ such that:
    \begin{equation}
        c_0(\sigma,\mu,s_0) \geq \frac{1}{\sinh^2(\pi\mu)}
    \end{equation}
    Notice that for $u \geq c_0(\sigma,\mu,s_0)\geq \frac{1}{\sinh^2(\pi\mu)}$, the norm $|\varphi_{f_{\sigma,\mu}}|$ is monotone decreasing:
    \begin{align}
        \begin{split}
            \frac{d}{du}|\varphi_{f_{\sigma,\mu}}(u)| 
            &=\frac{d}{du}\frac{e^{\frac{\sigma^2}{2}(\pi^2 - u^2)}}{\pi}\sqrt{\cosh^2(\pi\mu) - \sin^2(\pi u\sigma^2)}
            \\&= -\sigma^2 u e^{\frac{\sigma^2}{2}(\pi^2 - u^2)}\sqrt{\cosh^2(\pi\mu) - \sin^2(\pi u \sigma^2)} - \frac{2\pi\sigma^2e^{\frac{\sigma^2}{2}(\pi^2 - u^2)}\sin(\pi u\sigma^2)\cos(\pi u \sigma^2)}{2\pi\sqrt{\cosh^2(\pi\mu) - \sin^2(\pi u \sigma^2)}}
            \\&
            =\frac{\sigma^2 e^{\frac{\sigma^2}{2}(\pi^2 - u^2)}}{\sqrt{\cosh^2(\pi\mu) - \sin^2(\pi u\sigma^2)}}\lb -u\lb \cosh^2(\pi\mu) - \sin^2(\pi u \sigma^2) \rb - \sin(\pi u \sigma^2)\cos(\pi u \sigma^2) \rb
            \\&
            \leq \frac{\sigma^2 e^{\frac{\sigma^2}{2}(\pi^2 - u^2)}}{\sqrt{\cosh^2(\pi\mu) - \sin^2(\pi u \sigma^2)}}\lb -u\lb \cosh^2(\pi\mu) - 1 \rb + 1 \rb
            \\&
            =\frac{\sigma^2 e^{\frac{\sigma^2}{2}(\pi^2 - u^2)}}{\sqrt{\cosh^2(\pi\mu) - \sin^2(\pi u \sigma^2)}}\lb  -u\sinh^2(\pi\mu) + 1\rb
        \end{split}
    \end{align}
    Hence for $u \geq c_0(\sigma,\mu,s_0) \geq \frac{1}{\sinh^2(\pi\mu)}$, $\frac{d}{du}|\varphi_{f_{\sigma,\mu}}(u)| \leq 0$ and thus $|\varphi_{f_{\sigma,\mu}}(c_0)| \geq |\varphi_{f_{\sigma,\mu}}(u)|$. And since $|\varphi_{f_{\sigma,\mu}}(u)|$ is even, for $u \leq -c_0(\sigma,\mu, s_0)\leq -\frac{1}{\sinh^2(\pi\mu)}$, the function $|\varphi_{f_{\sigma,\mu}}(u)|$ is monotonely increasing and thus $|\varphi_{f_{\sigma,\mu}}(-c_0)|\geq |\varphi_{f_{\sigma,\mu}}(u)|$. Moreover, to the leading order in $s$, $c_0(\sigma,\mu, s)$ is approximated by:
    \begin{equation}
        c_0(\sigma,\mu,s) = -\frac{1}{2\sigma^2} + \sqrt{\frac{1}{4} - \frac{2}{\sigma^2}\log\lb\frac{||\varphi_{f_{\sigma,\mu}}||_1 e^{-\frac{\pi^2\sigma^2}{2}}}{2\sinh(\pi \mu)}\rb + \frac{2\pi s }{\sigma^2}} \sim K_1 + K_2 s
    \end{equation}
    where $K_1, K_2$ are two constants which depend on $\sigma, \mu$ and $K_2 > 0$. Finally using the formula for $|\varphi_{f_{\sigma,\mu}}(c_0)|$, we can provide an upper bound:
    \begin{align}
        \begin{split}
            |\varphi_{f_{\sigma,\mu}}(c_0)| &=\frac{e^{\frac{\sigma^2}{2}(\pi^2 - c_0^2)}}{\pi}|\sinh(\pi\mu - \pi i c_0\sigma^2)| \leq \frac{e^{\pi\mu}e^{\frac{\sigma^2}{2}(\pi^2 - c_0^2)}}{\pi}
        \end{split}
    \end{align}
    Using these observations, we can calculate the following bound for all $s\geq s_0$ (i.e. $c_0(\sigma,\mu,s_0)\geq\frac{1}{\sinh^2(\pi\mu)}$):
    \begin{align}
        \begin{split}
            c_{f_{\sigma,\mu}}(s) &\leq 2\int_{|u| \geq c_0}du |\varphi_{f_{\sigma,\mu}}(u)| + \frac{2}{\pi}\tanh^{-1}(e^{-\pi s})e^{c_0 / 2}||\varphi_{f_{\sigma,\mu}}||_1
            \\&
            \leq 2\lb|\varphi_{f_{\sigma,\mu}}(c_0)|^{1/2} + |\varphi_{f_{\sigma,\mu}}(-c_0)|^{1/2}\rb\int_\mathbb{R} du |\varphi_{f_{\sigma,\mu}}(u)|^{1/2} + \frac{2}{\pi}\tanh^{-1}(e^{-\pi s })e^{c_0 / 2}||\varphi_{f_{\sigma,\mu}}||_1
            \\&
            \leq 4|\varphi_{f_{\sigma,\mu}}(c_0)|^{1/2}\frac{e^{\frac{\sigma^2\pi^2}{4}}}{\sqrt{\pi}}e^{\frac{\pi\mu}{2}}\int_\mathbb{R}du e^{-\frac{\sigma^2}{4}u^2} + \frac{2}{\pi}\tanh^{-1}(e^{-\pi s})e^{c_0/2}||\varphi_{f_{\sigma,\mu}}||_1
            \\&
            =\frac{8e^{\frac{\sigma^2\pi^2}{4}}e^{\frac{\pi\mu}{2}}}{\sigma}|\varphi_{f_{\sigma,\mu}}(c_0)|^{1/2} + 4|\varphi_{f_{\sigma,\mu}}(c_0)|
            \\&
            \leq\frac{8e^{\frac{\sigma^2\pi^2}{4}}e^{\frac{\pi\mu}{2}}}{\sigma}\frac{e^{\frac{\pi\mu}{2}}e^{\frac{\sigma^2}{4}(\pi^2 - c_0^2)}}{\sqrt{\pi}} + 4\frac{e^{\pi\mu}e^{\frac{\sigma^2}{2}(\pi^2 - c_0^2)}}{\pi}
            \\&
            \leq\frac{8e^{\frac{\sigma^2\pi^2}{4}}e^{\frac{\pi\mu}{2}}}{\sigma}\frac{e^{\frac{\pi\mu}{2}}e^{\frac{\sigma^2}{4}(\pi^2 - c_0(\sigma,\mu,s)^2)}}{\sqrt{\pi}} + 4\frac{e^{\pi\mu}e^{\frac{\sigma^2}{2}(\pi^2 - c_0(\sigma,\mu,s)^2)}}{\pi}
        \end{split}
    \end{align}
    As $s$ increases, the upper bound on $c_{f_{\sigma,\mu}}(s)$ is monotonely decreasing. And to the leading order in $s$, the upper bound can be approximated by:
    \begin{equation}
        c_{f_{\sigma,\mu}}(s) \leq K_3e^{-\frac{K_4 s^2}{2}} + K_5e^{-K_4 s^2}
    \end{equation}
    where $K_3, K_4, K_5 > 0$ are positive constants which depend on $\sigma, \mu$. 
\end{proof}
In addition, as $\sigma \rightarrow 0$, the upper bound of $c_{f_{\sigma, \mu}}$ approaches:
\begin{align}
    \begin{split}
        c_{f_{\sigma,\mu}}(s) &\leq C_1 \frac{e^{-C_2/\sigma^2}e^{C_3\sigma^2}}{\sigma}e^{-C_4s^2}
    \end{split}
\end{align}
where $C_1,C_2,C_3, C_4 > 0$ are positive constants that depend on $\mu$. The right hand side as a well-defined limit as $\sigma\rightarrow 0^+$ and the limit is $0$ (uniformly in $s$).




\bibliographystyle{unsrt}
\bibliography{modularscattering}

\end{document}

